\chapter{Chiral Koszul duality}
\label{chap:koszul-pairs}

\providecommand{\Conil}{\mathsf{Conil}}
\providecommand{\BP}{\mathrm{BP}}

\index{chiral Koszul duality|textbf}
\index{twisting morphism!chiral|textbf}
\index{Francis--Gaitsgory|textbf}

Classical Koszul duality, in the sense of
Priddy~\cite{Priddy70} and
Beilinson--Ginzburg--Soergel~\cite{BGS96}, works for graded
algebras over a field: the bar complex is a tensor coalgebra on a
finite-dimensional vector space, the cobar functor inverts it
when bar cohomology is concentrated in degree one, and the
duality exchanges quadratic algebras with quadratic coalgebras.
Three features of chiral algebras on algebraic curves defeat
this framework simultaneously. First, chiral algebras are
\emph{curved}: the bar differential on the genus expansion
satisfies $d^2_{\mathrm{fib}} = \kappa(\cA) \cdot \omega_g$
rather than $d^2 = 0$, so the classical bar-cobar adjunction,
which requires strict nilpotence, does not apply beyond genus
zero. Second, chiral algebras live over a curve~$X$, not over a
point: the tensor power $V^{\otimes n}$ is replaced by the
factorization product $\cA^{\boxtimes n}$ over the
Fulton--MacPherson compactification $\FM_n(X)$, and the
differential is a sum of residues at boundary strata governed by
OPE singularities, not a combinatorial sum of face maps. Third,
the conformal-weight components are infinite-dimensional, so
finite generation hypotheses fail and convergence of spectral
sequences becomes a separate problem.

What forces the extension is not a desire for generality but the
structure of the objects themselves. The Virasoro algebra has a
quartic OPE pole, $\cW$-algebras have poles of arbitrarily high
order, Yangians have spectral-parameter relations: none of these
admit a quadratic presentation, and the BGS theory does not see
them. The Francis--Gaitsgory adjunction
$\chirLie \dashv \chirCom$ on $\Ran(X)$~\cite{FG12} settles
Koszul duality of the Lie--Com pair at the operadic level, but
it does not determine whether a specific chiral algebra is
Koszul: it does not see the genus corrections in the bar
complex, and it does not detect the shadow obstruction tower
whose depth varies from two to infinity across the standard
families.

The bar-cobar adjunction of Theorem~A
(Theorem~\ref{thm:bar-cobar-isomorphism-main}) is the functorial pair
constructed here for the first two obstructions. It is the
adjunction that (a)~reconstructs on the Koszul locus, producing
$\Omega_X(\barB_X(\cA)) \xrightarrow{\sim} \cA$
(Theorem~\ref{thm:bar-cobar-inversion-qi}); (b)~intertwines
Verdier duality on $\Ran(X)$, connecting $\barB_X(\cA)$ with the
homotopy dual factorization algebra~$\cA^!_\infty$. Lagrangian
complementarity
$Q_g(\cA) \oplus Q_g(\cA^!) \simeq
H^*(\overline{\mathcal{M}}_g, Z(\cA))$
(Theorem~\ref{thm:quantum-complementarity-main}) uses this
bar-cobar/Verdier package together with the derived-centre pairing.
Thus $\Omega_X\barB_X(\cA)\simeq\cA$ is reconstruction, while
$\mathbb{D}_{\Ran}\barB_X(\cA)\simeq\cA^!_\infty$ is the
Verdier-Koszul route to the dual.

\smallskip\noindent\textit{Standing notation for the overview.}
A chiral twisting datum means a quadruple
$(\cA,\cC,\tau,F_\bullet)$ consisting of an augmented chiral algebra,
a conilpotent complete factorization coalgebra, a degree-$+1$
twisting morphism $\tau\colon \cC\to \cA$ satisfying the
Maurer--Cartan equation, and a compatible filtration. We write
$K_\tau^L(\cA,\cC)$ and $K_\tau^R(\cC,\cA)$ for the corresponding
twisted tensor products, and $\Theta_\cA^{\leq r}$ for the degree-$r$
truncation of the universal Maurer--Cartan element of the modular
convolution algebra. A chiral Koszul pair is such a datum on the
Koszul locus, so the twisted tensor products are acyclic and the
associated graded datum is classically Koszul.
Definitions~\ref{def:chiral-twisting-datum},
\ref{def:chiral-koszul-morphism}, and \ref{def:chiral-koszul-pair}
record the precise versions used below.

Theorem~\ref{thm:koszul-equivalences-meta} assembles twelve
tests and consequences around chiral Koszulness. Seven are
intrinsic, genuinely independent bidirectional equivalences:
$E_2$-collapse of the bar spectral sequence, vanishing of
transferred $A_\infty$ operations $m_k$ for $k \geq 3$,
acyclicity of the twisted tensor product
$K_\tau^L(\cA, \cC)$, truncation of the shadow tower
$\Theta_\cA^{\leq r}$, and three others. The
Barr--Beck--Lurie comparison is a proved consequence of the
bar-cobar counit quasi-isomorphism, listed for completeness
because it is the $(\infty,1)$-categorical formulation most
useful for applications. The factorization-homology item is the
external integration avatar only after the conditional comparison
model of Proposition~\ref{prop:bar-fh}; it is not an
unconditional route to Koszulness. One item (Hochschild duality
and concentration in degrees $\{0,1,2\}$) is a proved one-way
consequence on the Koszul locus. One (the Lagrangian eigenspace
decomposition) is conditional on perfectness of the bar-cobar
normal complex. One ($\mathcal{D}$-module purity) is a
one-directional implication. The count is exact:
$7 + 1 + 1 + 1 + 1 + 1$.

Two distinctions require emphasis. First, the four complexity
classes G/L/C/M, with shadow depths
$\{2, 3, 4, \infty\}$
(Theorem~\ref{thm:shadow-archetype-classification}), record the
degree at which the shadow obstruction tower first becomes
nontrivial; they do \emph{not} by themselves record failure of
Koszulness. The PBW-flat universal families treated here
\textup{(}Heisenberg, universal affine Kac--Moody, universal
Virasoro, universal principal $\mathcal W$-algebras, and the free
field systems on their stated lanes\textup{)} are chirally Koszul by
the PBW criterion. Simple admissible quotients, minimal models, and
quotients with vacuum null vectors require the separate
Shapovalov/Li-bar tests below and may fail. The classes separate
algebras by the complexity of their higher-genus behaviour, not by
the validity of the PBW hypothesis. Second,
Koszulness and SC formality are logically independent properties.
Koszulness is the condition that bar cohomology is concentrated
in degree one; SC formality is the condition
$m_k^{\mathrm{SC}} = 0$ for $k \geq 3$ in the Swiss-cheese
bar complex. On the PBW-flat universal surface the families above
are Koszul; only class~$G$ (Heisenberg and its relatives) is
SC-formal
(Proposition~\ref{prop:sc-formal-iff-class-g}).
Confusing the two misidentifies what the shadow tower measures.

On a point, the bar complex
$\barB(A) = T^c(s^{-1}\bar{A})$ is built from tensor powers of
the augmentation ideal; a classical twisting morphism
$\tau \colon \barB(A) \to A$ is a degree-$+1$ element in the
convolution dg~Lie algebra satisfying $d\tau + \tau \star \tau = 0$.
Its bar cohomology is the intrinsic coalgebra~$A^{\mathrm i}$; the
dual algebra~$A^!$ appears only after finite linear or Verdier duality.
On a curve~$X$, the chiral twisting morphism
(Definition~\ref{def:chiral-twisting-datum}) lives in the chiral
convolution algebra, where the star product is mediated by the
integration kernel on $X^2 \setminus \Delta$. The acyclicity of
the twisted tensor product $K_\tau^L(\cA, \cC)$, the defining
property of a chiral Koszul pair, is checked by spectral sequence
methods adapted to the PBW filtration by bar degree
(Theorem~\ref{thm:pbw-koszulness-criterion}).

Three phenomena distinguish the chiral setting from its
point-set ancestor, and all three are invisible at genus zero:
the modular characteristic $\kappa(\cA)$, which measures
curvature of the genus expansion
(Theorem~\ref{thm:modular-characteristic});
Lagrangian complementarity between $\cA$ and $\cA^!$
(Theorem~\ref{thm:quantum-complementarity-main});
and the critical discriminant $\Delta = 8\kappa S_4$, whose
vanishing controls termination of the shadow obstruction tower
(Theorem~\ref{thm:riccati-algebraicity}).
At genus zero all three collapse: $\kappa$ plays no role,
complementarity reduces to the genus-$0$ bar-cobar adjunction,
$\Delta$ is vacuous. On the formal disk, this further reduces
to classical Koszul duality; on the full genus-$0$
curve~$\mathbb{P}^1$, the Fulton--MacPherson compactifications
and Arnold relations already refine the classical picture. The
genus-zero input below is the source from which the
full modular theory is built. The standard families exemplify the four
classes: Heisenberg (class~$G$, depth~$2$, formal), affine
Kac--Moody (class~$L$, depth~$3$, single Massey product),
$\beta\gamma$ (class~$C$, depth~$4$, contact then rigid),
Virasoro and $\mathcal{W}_N$ (class~$M$, depth~$\infty$,
intrinsically non-formal). The genus-zero degree-two projection
of the universal MC element $\Theta_\cA$
(Theorem~\ref{thm:mc2-bar-intrinsic}) recovers the classical
$R$-matrix, linking the Koszul pairs of this chapter to the
Drinfeld--Kohno bridge of
Chapter~\ref{chap:koszul-pair-structure}.

The Heisenberg calculation supplies the first model of the
twisting datum. The PBW criterion turns that model into a
recognition theorem for non-quadratic algebras. A chiral Koszul
pair is the completed, finite-type, Verdier-compatible form of
that datum; Theorem~A proves bar-cobar inversion on this locus,
and the higher-genus chapter adds the curved term that appears
when the fiberwise differential no longer squares to zero.

% ================================================================
% SECTION 8.1: FROM THE HEISENBERG EXAMPLE TO GENERAL TWISTING DATA
% ================================================================


\section{From the Heisenberg example to general twisting data}
\label{sec:chiral-twisting-data}

The bar complex of $\mathcal{H}_k$ (Chapter~\ref{ch:heisenberg-frame},
Theorem~\ref{thm:frame-heisenberg-koszul-dual})
produces the coalgebra $\bar{B}(\mathcal{H}_k)
\simeq \mathrm{coLie}^{\mathrm{ch}}(V^*)$,
whose cohomology gives the intrinsic bar-dual coalgebra
$\mathcal{H}_k^{\mathrm i}$. On the finite-type Verdier lane this
coalgebra dualizes to the Koszul dual algebra
$\mathcal{H}_k^! = \mathrm{Sym}^{\mathrm{ch}}(V^*)$.
The bar-cobar counit
$\Omega_X(\bar{B}_X(\mathcal{H}_k)) \xrightarrow{\;\sim\;} \mathcal{H}_k$
is a quasi-isomorphism (recovering the original algebra, not its dual). Two structural features make this work:

\begin{enumerate}
\item The twisted tensor product
 $K_\tau^L(\mathcal{H}_k, \mathrm{coLie}^{\mathrm{ch}}(V^*))$
 is acyclic: the twisting morphism
 $\tau\colon \mathrm{coLie}^{\mathrm{ch}}(V^*) \to \mathcal{H}_k$
 solves the Maurer--Cartan equation.
\item The PBW filtration on $\mathcal{H}_k$ produces an
 associated graded that is quadratic-Koszul in the classical sense.
\end{enumerate}

These two properties (acyclicity of twisted tensor products and
classical Koszulness of the associated graded) are the
recognition criteria for chiral Koszulness. The point is not that
the chiral algebra itself is quadratic (it is not), but that the PBW
filtration by bar-length makes the leading-order structure
quadratic-Koszul (Theorem~\ref{thm:pbw-koszulness-criterion}).
These criteria apply equally to
the Virasoro algebra (quartic pole in the $TT$ OPE),
$\mathcal{W}$-algebras (poles of arbitrarily high order),
and Yangians (spectral-parameter relations), none of which
admit a quadratic presentation.

\subsection{Construction layer: twisting data}

To recognize Koszulness, we need to compare a chiral algebra with
its bar coalgebra. The bridge between them is a \emph{twisting morphism},
a degree-$+1$ map satisfying a Maurer--Cartan equation, and
to ensure convergence of the resulting spectral sequences, we need a
compatible filtration. The Koszul property is then a
\emph{recognition criterion} on these data, not a prerequisite for
their existence.

\begin{definition}[Chiral twisting datum]\label{def:chiral-twisting-datum}
\index{twisting datum!chiral|textbf}
A \emph{chiral twisting datum} on a smooth curve $X$ is a quadruple
$(\cA, \cC, \tau, F_\bullet)$ consisting of:
\begin{enumerate}
\item an augmented chiral algebra
 $\cA \in \operatorname{Fact}^{\mathrm{aug}}(X)$;
\item a conilpotent complete factorization coalgebra
 $\cC \in \operatorname{CoFact}^{\mathrm{conil,comp}}(X)$;
\item a degree $+1$ morphism
 $\tau\colon \cC \to \cA$ satisfying the Maurer--Cartan equation
 $d\tau + \tau \star \tau = 0$ in the convolution dg~Lie algebra;
\item an exhaustive, complete, bounded-below filtration $F_\bullet$ on $\cA$,
 $\cC$, $\bar{B}_X(\cA)$, and $\Omega_X(\cC)$, compatible with
 all structure maps.
\end{enumerate}
Associated to $\tau$ are the \emph{twisted tensor products}
\[
K_\tau^L(\cA,\cC) := (\cA \otimes \cC,\; d_\cA + d_\cC + d_\tau^L),
\quad
K_\tau^R(\cC,\cA) := (\cC \otimes \cA,\; d_\cC + d_\cA + d_\tau^R).
\]
\end{definition}

\begin{remark}[Explicit twisted differential]
\label{rem:explicit-twisted-differential}
\index{twisted tensor product!differential formula}
The twisting differentials $d_\tau^L$ and $d_\tau^R$ in
Definition~\ref{def:chiral-twisting-datum} are:
\begin{equation}\label{eq:chiral-twisted-differential-L}
d_\tau^L(a \otimes c)
\;=\;
\sum (-1)^{|a|}\, a \cdot \tau(c_{(1)}) \otimes c_{(2)},
\end{equation}
\begin{equation}\label{eq:chiral-twisted-differential-R}
d_\tau^R(c \otimes a)
\;=\;
\sum (-1)^{|c|}\, c_{(1)} \otimes \tau(c_{(2)}) \cdot a,
\end{equation}
where $\Delta(c) = \sum c_{(1)} \otimes c_{(2)}$ is the factorization
coproduct. The MC equation $d\tau + \tau \star \tau = 0$ is
\emph{equivalent} to $(d_\cA + d_\cC + d_\tau^L)^2 = 0$
\textup{(}cf.\ \cite[Lemma~2.1.4]{LV12}\textup{)}: the twisting
morphism condition is the nilpotence of the twisted differential.
In the chiral setting, $\tau(c_{(1)})$ is the OPE extraction:
the coproduct $\Delta(c) = \sum c_{(1)} \otimes c_{(2)}$ separates
the coalgebra element, and $\tau$ evaluates $c_{(1)}$ against the
propagator on~$\overline{C}_2(X)$ by residue at the collision
divisor
\textup{(}Proposition~\textup{\ref{prop:twisting-morphism-propagator})}.
\end{remark}

\begin{definition}[Chiral Koszul morphism]\label{def:chiral-koszul-morphism}
\index{Koszul morphism!chiral|textbf}
A chiral twisting datum $(\cA, \cC, \tau, F_\bullet)$ is
\emph{Koszul} if:
\begin{enumerate}
\item both $K_\tau^L(\cA,\cC)$ and $K_\tau^R(\cC,\cA)$ are acyclic;
\item the associated graded
 $(\operatorname{gr}\cA, \operatorname{gr}\cC, \operatorname{gr}\tau)$
 is quadratic/Koszul in the ordinary operadic sense;
\item the filtration converges strongly on $\bar{B}_X(\cA)$ and
 $\Omega_X(\cC)$.
\end{enumerate}
\end{definition}

\begin{remark}[Construction versus resolution]\label{rem:construction-vs-resolution}
For every augmented chiral algebra $\cA$, the bar
$\bar{B}_X(\cA)$ and cobar $\Omega_X(\bar{B}_X(\cA))$ exist as
constructions. But the counit
$\Omega_X(\bar{B}_X(\cA)) \to \cA$ is a quasi-isomorphism
\emph{only when $\tau$ is a Koszul morphism}
(Theorem~\ref{thm:bar-cobar-inversion-qi}).
Off the Koszul locus, the bar-cobar object persists in the
provisional coderived category
$\mathsf{D}^{\mathrm{co}}_{\mathrm{prov}}(X)$
(Appendix~\ref{sec:coderived-models}),
not in the ordinary derived category.
This is the chiral/factorization analogue of the fundamental
theorem of twisting morphisms \cite[Theorem~2.3.1]{LV12}.
\end{remark}

The following lemmas make the mapping-cone identifications
explicit, as required for Theorem~\ref{thm:fundamental-twisting-morphisms}.

\begin{lemma}[Left twisted tensor product as mapping cone; \ClaimStatusProvedHere]
\label{lem:twisted-product-cone-counit}
\index{twisted tensor product!mapping cone}
Let $(\cA, \cC, \tau, F_\bullet)$ be a chiral twisting datum.
There is a natural isomorphism of \emph{graded} factorisation
modules
\begin{equation}\label{eq:twisted-cone-iso-L}
 \Phi^L\colon K_\tau^L(\cA,\cC) \;\xrightarrow{\;\cong\;}\;
 \operatorname{Cone}(\varepsilon_\tau)[-1]
\end{equation}
where $\varepsilon_\tau\colon \Omega_X(\cC) \to \cA$ is the
canonical counit, and $\Phi^L$ intertwines the twisted
differential $d_\cA + d_\cC + d_\tau^L$ on the left with the
mapping-cone differential on the right. In particular,
$K_\tau^L$ is acyclic if and only if $\varepsilon_\tau$ is a
quasi-isomorphism.
\end{lemma}

\begin{proof}
The claim is chain-level (strict equality of differentials after
the explicit $\Phi^L$), not merely a weak equivalence. We
construct $\Phi^L$ on graded generators and verify compatibility
with the differentials.

\emph{Step 1. Graded-module isomorphism $\Phi^L$.}
By definition $\Omega_X(\cC) = T_X(s^{-1}\bar{\cC})
= \bigoplus_{n\geq 0} (s^{-1}\bar{\cC})^{\otimes_X n}$,
the free factorisation algebra on the desuspended reduced
coalgebra. The mapping cone is, by convention, the graded module
\[
 \operatorname{Cone}(\varepsilon_\tau)
 \;=\; \Omega_X(\cC)[1] \oplus \cA,
 \quad
 \operatorname{Cone}(\varepsilon_\tau)[-1]
 \;=\; \Omega_X(\cC) \oplus \cA[-1].
\]
Equivalently, $\Omega_X(\cC) = \mathbf{1} \oplus \bigoplus_{n\geq 1}
(s^{-1}\bar{\cC})^{\otimes_X n}$, with the unit summand giving the
augmentation. The tensor factor $\cA$ in $K_\tau^L = \cA
\widehat{\otimes} \cC$ is the $\cA$-module acted on by the
twisted differential, and the completed tensor
$\widehat{\otimes}$ is the pro-object limit over the
conformal-weight filtration $F_\bullet$.
Define $\Phi^L$ on graded tensors by
\begin{equation}\label{eq:Phi-L-formula}
 \Phi^L\bigl(a \otimes [c_1 | \cdots | c_n]\bigr)
 \;=\;
 a \cdot \bigl[s^{-1}\bar{c}_1 | \cdots | s^{-1}\bar{c}_n\bigr]
 \;\in\; \cA \otimes_X (s^{-1}\bar{\cC})^{\otimes_X n},
\end{equation}
where $[c_1 | \cdots | c_n] \in \cC$ is an iterated coproduct
component of~$c \in \cC$ under the shuffle/deconcatenation of
Definition~\ref{def:chiral-twisting-datum}(3), and the right-hand
side is read inside $\cA \otimes \Omega_X(\cC)$, which is the
cone $\cA \oplus \Omega_X(\cC) \cdot \cA$ under the free-module
identification. On the grading $n = 0$ summand the map is
$a \otimes 1 \mapsto a \in \cA[-1]$. The
bijection at each fixed weight level follows from the cofree
decomposition $\Delta(c) = \sum c_{(1)}\otimes c_{(2)}$ of the
chiral coproduct (Definition~\ref{def:chiral-twisting-datum}(3)),
whose deconcatenation dual is the tensor-algebra basis of
$\Omega_X(\cC)$. Finiteness of the decomposition at each
conformal-weight slice is hypothesis (H3) in the sense of
Theorem~\ref{thm:bar-cobar-platonic}; under this hypothesis
$\Phi^L$ is a graded-module isomorphism.

\emph{Step 2. Degree and sign verification.}
With desuspension $|s^{-1}v| = |v| - 1$ and cohomological sign
$|d| = +1$, we track the total grading carefully. On $K_\tau^L
= \cA\widehat{\otimes}\cC$ a pure tensor $a\otimes c$ lies in
degree $|a| + |c|$. Under the conilpotent decomposition
$c \mapsto [c_1|\cdots|c_n]$ with each $c_i$ in the reduced
coalgebra~$\bar{\cC}$, the component at bar-length $n$ lies in
degree $|a| + \sum |c_i|$. On the cone side
$\operatorname{Cone}(\varepsilon_\tau)[-1] = \Omega_X(\cC) \oplus
\cA[-1]$, the element $a\cdot[s^{-1}\bar{c}_1|\cdots|s^{-1}\bar{c}_n]$
living in the free-$\cA$-module generated by $\Omega_X(\cC)$
occupies degree $|a| + \sum(|c_i| - 1) + n = |a| + \sum |c_i|$,
where each desuspension contributes $-1$ and the cofree-tensor
coalgebra structure on $\cC$ carries the opposite $+n$ shift
(Definition~\ref{def:chiral-twisting-datum}(3); the conilpotent
coalgebra~$\cC$ is modelled as $\cC = \bigoplus_n s\cdot
(s^{-1}\bar{\cC})^{\otimes n}$ by the standard shuffle-versus-deconcatenation
duality, so $|c_i|_\cC = |c_i|_{\Omega} + 1$). Hence both sides
carry the same degree, and $\Phi^L$ is a degree-$0$ graded-module
map.

The Koszul sign rule then applies uniformly: the sign
$(-1)^{|a|}$ in $d_\tau^L(a\otimes c) = \sum(-1)^{|a|}\,
a\cdot\tau(c_{(1)})\otimes c_{(2)}$
(equation~\eqref{eq:chiral-twisted-differential-L}) matches the
Koszul sign picked up when $d_\Omega$ passes~$a$ through the
$\Omega$-factor in the cone. Specifically, on a generator
$a \otimes [c_1|\cdots|c_n]$ the mapping-cone differential equals
\[
 d^{\mathrm{cone}}\bigl(a \otimes [s^{-1}\bar{c}_1|\cdots|
 s^{-1}\bar{c}_n]\bigr)
 = d_\cA(a)\otimes[s^{-1}\bar{c}_1|\cdots]
 + (-1)^{|a|} a \otimes d_\Omega[s^{-1}\bar{c}_1|\cdots]
 - \varepsilon_\tau\text{-contribution at }n=1,
\]
where the $-\varepsilon_\tau$-contribution is the cone attaching
map. Under $\Phi^L$ this matches, term by term, $d_\cA + d_\cC +
d_\tau^L$ with $d_\cC$ acting on $\cC$ by the coproduct-twist:
the key identity is that the cobar differential on
$(s^{-1}\bar{\cC})^{\otimes n}$, which deconcatenates adjacent
pairs via the coproduct $\Delta_\cC$ and contracts by $\tau$, is
the image under $\Phi^L$ of the twisted $d_\tau^L$-differential
on $\cA \otimes \cC^{\otimes n}$. The compatibility
$\Phi^L\circ d_\tau^L = d^{\mathrm{cone}}\circ\Phi^L$ on each
bar-degree is a standard shuffle-Koszul identity
\cite[Lemma~2.1.4]{LV12}, whose chiral upgrade uses
Proposition~\ref{prop:twisting-morphism-propagator} in place of
the algebraic coaction.

\emph{Step 3. Chiral-setting verification (pole orders and
Arnold relation).}
Unlike the classical setting of
\cite[\S2.1]{LV12}, the chiral factorisation coproduct
$\Delta_\cC$ on $\cC$ is realised geometrically as residues
along the collision divisor $D_{12} \subset \overline{C}_2(X)$
via the propagator $\eta(z,w) = d\log(z - w)$
(Proposition~\ref{prop:twisting-morphism-propagator}, genus
zero; Szeg\H{o} kernel at higher genus). The twisted differential
$d_\tau^L$ on $\cA \otimes \cC$ is thus an iterated simple-pole
residue, and the cobar differential $d_\Omega$ on
$(s^{-1}\bar{\cC})^{\otimes n}$ is the Stokes boundary on
$\overline{C}_n(X)$. These two differentials differ only in which
tensor factor absorbs the residue extraction: in $d_\tau^L$ the
$\cA$-factor absorbs the OPE coefficient produced by $\tau$; in
$d_\Omega$ the $\Omega$-factor absorbs it via deconcatenation.
By the Arnold relation (Theorem~\ref{thm:arnold-relations})
\[
 \eta_{12}\wedge\eta_{23} + \eta_{23}\wedge\eta_{31}
 + \eta_{31}\wedge\eta_{12} = 0
\]
the two extractions agree modulo the Arnold ideal: a triple
residue $(z, w_1, w_2)$ with $w_1, w_2 \to z$ gives the same
coefficient whether the collision is read $w_1 = z$ first or
$w_2 = z$ first, provided the three-form
$\eta_{12}\wedge\eta_{13}$ is anti-symmetrised. Hence
$\Phi^L(d_\tau^L x) = d_\Omega(\Phi^L x) \bmod
(\mathrm{Arnold})$; since the Arnold-relation ideal is killed
on $\overline{C}_n(X)$, the equality is strict.

\emph{Step 4. Mittag-Leffler convergence of the completion.}
The chiral-setting correction to the classical LV12 proof is the
use of the \emph{completed} tensor product $\cA\widehat{\otimes}\cC$
(Definition~\ref{def:chiral-twisting-datum}), which is the
pro-object limit $\lim_n \cA/F^n \otimes \cC/F^n$ over the
filtration~$F_\bullet$. The isomorphism $\Phi^L$ of Step~1
constructed on each finite slice $F^{\leq n}$ is compatible with
the projection $F^{\leq n+1}\twoheadrightarrow F^{\leq n}$ by
naturality of the coproduct. The Mittag-Leffler condition
$\lim^1_n F^n \cA \otimes F^n \cC = 0$ holds because the
filtration is \emph{bounded below} and \emph{exhaustive}
(hypothesis~(4) of Definition~\ref{def:chiral-twisting-datum}):
every surjection $F^{\leq n+1}\twoheadrightarrow F^{\leq n}$ is
eventually an isomorphism at each fixed degree, which is the
strong Mittag-Leffler bound. Passing to the limit promotes the
finite-slice isomorphism to a chain-level isomorphism on
$\cA\widehat{\otimes}\cC$, as required; see
\cite[\S3.5]{LV12} for the bar-cobar pro-object conventions.

\emph{Step 5. Acyclicity corollary.}
Since $\Phi^L$ is a chain-level isomorphism, $H^\bullet(K_\tau^L)
\cong H^\bullet(\operatorname{Cone}(\varepsilon_\tau)[-1])$. The
long exact sequence of the mapping cone gives $K_\tau^L$ acyclic
if and only if $\varepsilon_\tau\colon \Omega_X(\cC) \to \cA$ is
a quasi-isomorphism.
\end{proof}

\begin{lemma}[Right twisted tensor product as mapping cone; \ClaimStatusProvedHere]
\label{lem:twisted-product-cone-unit}
\index{twisted tensor product!mapping cone}
Under the same hypotheses, there is a natural isomorphism of
graded factorisation modules
\[
 \Phi^R\colon K_\tau^R(\cC,\cA) \;\xrightarrow{\;\cong\;}\;
 \operatorname{Cone}(\eta_\tau)[-1]
\]
where $\eta_\tau\colon \cC \to \bar{B}_X(\cA)$ is the canonical
unit, intertwining the twisted differential $d_\cC + d_\cA +
d_\tau^R$ with the mapping-cone differential. In particular,
$K_\tau^R$ is acyclic if and only if $\eta_\tau$ is a weak
equivalence.
\end{lemma}

\begin{proof}
Dual to Lemma~\ref{lem:twisted-product-cone-counit}, with the
role of $\Omega_X(\cC)$ played by the bar coalgebra
$\bar{B}_X(\cA) = T^c_X(s^{-1}\bar{\cA})$. Explicitly, $\Phi^R$
is defined on graded tensors by
\[
 \Phi^R\bigl([c] \otimes a\bigr)
 \;=\;
 c \otimes (s^{-1}\bar{a}_{(1)} \mid \cdots \mid s^{-1}\bar{a}_{(n)})
\]
where $a = a_{(1)}\cdots a_{(n)}$ is the $\cA$-multiplication
decomposition by conformal-weight level. Steps~1--5 of the
counit-side proof transpose: the chiral coproduct on $\bar{B}_X$
is now on the bar-side, the twisted differential
$d_\tau^R(c\otimes a) = \sum(-1)^{|c|} c_{(1)}\otimes
\tau(c_{(2)})\cdot a$ of
equation~\eqref{eq:chiral-twisted-differential-R} translates
under $\Phi^R$ to the mapping-cone differential of
$\operatorname{Cone}(\eta_\tau)[-1]$, and Mittag-Leffler
convergence on $\cC\widehat{\otimes}\cA$ follows from the same
bounded-below exhaustive filtration. Acyclicity of $K_\tau^R$
thus matches weak equivalence of $\eta_\tau$.
\end{proof}

\begin{proposition}[Explicit contracting homotopy for the canonical
universal pair; \ClaimStatusProvedHere]
\label{prop:canonical-contracting-homotopy}
\index{contracting homotopy!canonical twisted tensor product}
\index{bar-cobar adjunction!explicit homotopy}
Take the canonical twisting morphism $\pi\colon \bar{B}_X(\cA)
\to \cA$ obtained by projection to bar-degree~$1$ and
desuspension
\textup{(}Proposition~\textup{\ref{prop:universal-twisting-adjunction}(i))}.
Then the left twisted tensor product $K_\pi^L(\cA,\bar{B}_X(\cA))
= \cA \widehat{\otimes} \bar{B}_X(\cA)$, equipped with the
twisted differential
$d = d_\cA + d_{\bar{B}} + d_\pi^L$, admits the explicit
contracting homotopy
\begin{equation}\label{eq:explicit-homotopy-h}
 h\colon \cA \widehat{\otimes} \bar{B}_X(\cA)
 \;\longrightarrow\; \cA \widehat{\otimes} \bar{B}_X(\cA),
 \qquad
 h\bigl(a \otimes [s^{-1}\bar{a}_1 | \cdots | s^{-1}\bar{a}_n]\bigr)
 = 1 \otimes [s^{-1}\bar{a} | s^{-1}\bar{a}_1 | \cdots |
 s^{-1}\bar{a}_n],
\end{equation}
where on the right-hand side $a \in \cA$ is reinterpreted via
its image under the projection $\cA\twoheadrightarrow \bar{\cA}
\xhookrightarrow{s^{-1}} s^{-1}\bar{\cA}$ onto the first slot of
the bar coalgebra, and $[s^{-1}\bar{a} \mid \cdots]$ denotes the
cofree-tensor insertion. The homotopy~$h$ satisfies
\[
 [d, h] \;:=\; d\circ h + h\circ d \;=\; \operatorname{id} - p,
\]
where $p\colon \cA\widehat{\otimes}\bar{B}_X(\cA)
\twoheadrightarrow \mathbf{1}\cdot\mathbf{1} \cong
\mathbf{k}$ is the augmentation onto the bar-degree-$0$
vacuum. In particular, $K_\pi^L(\cA, \bar{B}_X(\cA)) \simeq
\mathbf{k}$ is acyclic in all positive degrees.
\end{proposition}

\begin{proof}
The homotopy~$h$ inserts $a$ into the leftmost slot of the bar
cofree tensor coalgebra, leaving the remaining tensor factors
intact. This is the classical ``extra degeneracy'' homotopy of
Eilenberg--MacLane (cobar side; see \cite[Corollary~2.2.5]{LV12}
for the operadic version), adapted to the chiral setting.

\emph{Verification of $[d, h] = \operatorname{id} - p$.}
Compute on a generator $x = a \otimes
[s^{-1}\bar{a}_1|\cdots|s^{-1}\bar{a}_n]$:
\begin{enumerate}
\item $d_\cA(h x) = d_\cA(1) \otimes [\cdots] = 0$ since the
unit has trivial differential;
\item $d_{\bar{B}}(h x) = 1 \otimes d_{\bar{B}}[s^{-1}\bar{a} |
s^{-1}\bar{a}_1 | \cdots]$, which by the bar differential's
deconcatenation formula expands as a sum over adjacent pair
contractions. The term contracting the first pair $(s^{-1}\bar{a},
s^{-1}\bar{a}_1)$ produces $\pm 1\otimes[s^{-1}(\bar{a}\cdot
\bar{a}_1)|s^{-1}\bar{a}_2|\cdots]$ (Koszul sign absorbed via
desuspension), which by chiral multiplication in the bar-side
is the image under $\Phi^L$ (cf.~Lemma~\ref{lem:twisted-product-cone-counit})
of the $d_\tau^L$-term in $h\circ d_\tau^L$. All other
adjacent-pair contractions appear with opposite sign in
$h\circ d_{\bar{B}}(x)$.
\item $d_\pi^L(h x) = \sum \pm a \cdot \pi(\bar{a}_{(1)}) \otimes
[\bar{a}_{(2)}| \bar{a}_1|\cdots]$ = $\sum \pm a \cdot \bar{a}_{(1)}
\otimes [\bar{a}_{(2)}|\bar{a}_1|\cdots]$, using $\pi = s^{-1}$ on
bar-degree~$1$ (identity on primitive part).
\end{enumerate}
Summing~(1)--(3) with the corresponding terms in $h\circ d(x)
= h\circ(d_\cA + d_{\bar B} + d_\pi^L)(x)$, the telescoping of
adjacent-pair contractions leaves exactly $x$ on the
bar-degree-$\geq 1$ slice, and $\mathbf{1}\cdot\mathbf{1}$
on the bar-degree-$0$ slice, which is the
$p$-projection. Hence $[d,h] = \operatorname{id} - p$, which is
the statement of the explicit chain-homotopy witness.

\emph{Chiral correction: pole-order admissibility.}
The insertion $a \mapsto s^{-1}\bar{a}$ in the leftmost slot
respects the chiral-residue structure because $\pi\colon
\bar{B}_X(\cA)\to \cA$ is realised by the simple-pole
propagator $\eta(z,w) = d\log(z - w)$
(Proposition~\ref{prop:twisting-morphism-propagator}): the
contraction of the first pair under $d_{\bar{B}}$ is a residue
at $w = z$, which by the residue theorem cancels the
$d_\pi^L$-contribution $a\cdot\pi(\bar{a}_1)$ (same pole,
opposite sign from the Stokes boundary). This is the
\emph{chiral version} of the Eilenberg--MacLane extra-degeneracy
argument: the homotopy~$h$ is compatible with the factorisation
stack because the residue calculus on $\overline{C}_n(X)$ is
already built into the bar differential via
Proposition~\ref{prop:twisting-morphism-propagator}.

\emph{Pro-object convergence.}
Because the insertion $h$ increases bar-degree by exactly one,
each fixed bar-degree slice $F^{\leq n}(\cA\widehat{\otimes}
\bar{B}_X(\cA))$ is preserved by~$h$ up to a shift $F^{\leq n}
\to F^{\leq n+1}$. The bounded-below filtration condition in
Definition~\ref{def:chiral-twisting-datum}(4) provides the
Mittag-Leffler bound: on each conformal-weight level the sum is
finite, and the contracting homotopy commutes with the
projections defining the pro-object. Hence~$h$ descends to the
completion, delivering an explicit chain-level contraction of
$\cA\widehat{\otimes}\bar{B}_X(\cA)$ onto the unit~$\mathbf{k}$.
\end{proof}

\begin{remark}[Scope of the explicit homotopy]
\label{rem:canonical-homotopy-scope}
The homotopy~$h$ of
Proposition~\ref{prop:canonical-contracting-homotopy} is
exhibited only for the \emph{canonical} universal twisting
morphism $\pi$ (projection to bar-degree~$1$), not for a general
twisting morphism $\tau\colon \cC \to \cA$. For a general
twisting datum $(\cA,\cC,\tau,F_\bullet)$, acyclicity of
$K_\tau^L$ is equivalent to $\varepsilon_\tau$ being a
quasi-isomorphism
(Lemma~\ref{lem:twisted-product-cone-counit}); the \emph{test}
for this is that the datum be a chiral Koszul morphism in the
sense of Definition~\ref{def:chiral-koszul-morphism}, and when
the datum passes this test, an explicit contracting homotopy is
constructed by composing~$h$ with the unit $f_\tau\colon \cC
\to \bar{B}_X(\cA)$ of
Proposition~\ref{prop:bar-universal-property}. This composite
witness converges as a finite-degree chain-level contraction
precisely when the associated graded datum is classically
Koszul; for twisting data failing the Koszul-morphism recognition
criterion \textup{(}off-locus in the sense of the twisting-datum
moduli; admissible simple quotients and minimal models can lie
there\textup{)}, no such closed-form contraction exists in finite
bar-degree, and one must pass to the coderived
category (Theorem~\ref{thm:bar-cobar-inversion-qi}, clause~(2)).
\end{remark}

\begin{lemma}[Filtered comparison; \ClaimStatusProvedHere]
\label{lem:filtered-comparison}
\index{spectral sequence!filtered comparison}
Let $(\cA, \cC, \tau, F_\bullet)$ be a chiral twisting datum
whose filtration is exhaustive, complete, and bounded below.
If $\varepsilon_\tau$ is a quasi-isomorphism, then:
\begin{enumerate}[label=\textup{(\alph*)}]
\item the associated graded datum
 $(\operatorname{gr} \cA, \operatorname{gr} \cC,
 \operatorname{gr} \tau)$ is a classical Koszul datum;
\item both twisted tensor products $K_\tau^L$ and $K_\tau^R$ are
 acyclic.
\end{enumerate}
\end{lemma}

\begin{proof}
Filter $\varepsilon_\tau$ by $F_\bullet$. Since the filtration
is bounded below and complete, the spectral sequence converges
strongly (Eilenberg--Moore). The associated graded
$\operatorname{gr} \varepsilon_\tau$ is a quasi-isomorphism
between graded objects, which by \cite[Theorem~3.2.1]{LV12}
identifies $\operatorname{gr} \tau$ as a classical Koszul
morphism. Classical Koszulity gives acyclicity of the
associated graded twisted tensor products. Strong convergence
then lifts this to acyclicity of $K_\tau^L$ and $K_\tau^R$
themselves.
\end{proof}

\begin{lemma}[Bar-side filtered comparison; \ClaimStatusProvedHere]
\label{lem:filtered-comparison-unit}
\index{spectral sequence!bar-side filtered comparison}
Let $(\cA, \cC, \tau, F_\bullet)$ be a chiral twisting datum
whose filtration is exhaustive, complete, and bounded below.
If $\eta_\tau\colon \cC \to \bar{B}_X(\cA)$ is a weak equivalence,
then:
\begin{enumerate}[label=\textup{(\alph*)}]
\item the associated graded datum
 $(\operatorname{gr} \cA, \operatorname{gr} \cC,
 \operatorname{gr} \tau)$ is a classical Koszul datum;
\item both twisted tensor products $K_\tau^L$ and $K_\tau^R$ are
 acyclic.
\end{enumerate}
\end{lemma}

\begin{proof}
This is the unit-side analogue of
Lemma~\ref{lem:filtered-comparison}: one replaces the cobar-side
filtration on $\Omega_X(\cC)$ by the induced filtration on the bar
coalgebra $\barB_X(\cA)$. Because $\barB_X(\cA)$ is the cofree tensor
coalgebra on $s^{-1}\bar\cA$ and $F_\bullet$ is compatible with the
bar structure, the associated graded bar coalgebra is canonically
the classical reduced bar construction:
\[
\operatorname{gr}_F \barB_X(\cA)
\;\cong\;
\barB(\operatorname{gr}_F \cA).
\]
Filter $\eta_\tau$ by $F_\bullet$. Since the filtration is bounded
below and complete, the induced spectral sequence converges
strongly, and the associated graded map
\[
\operatorname{gr}_F\eta_\tau \colon \operatorname{gr}_F\cC
\longrightarrow \operatorname{gr}_F\barB_X(\cA)
\cong \barB(\operatorname{gr}_F\cA)
\]
is a quasi-isomorphism of classical conilpotent coalgebras. By the
classical fundamental theorem of twisting morphisms
\cite[Theorem~2.3.1]{LV12}, the induced twisting morphism
$\operatorname{gr}_F\tau$ is classically Koszul, proving~(a).

For~(b), Lemma~\ref{lem:twisted-product-cone-unit} identifies
$K_\tau^R$ with $\operatorname{Cone}(\eta_\tau)[-1]$, so the weak
equivalence of $\eta_\tau$ makes $K_\tau^R$ acyclic. Filter
$K_\tau^L$ by $F_\bullet$; its associated graded complex is the
classical left twisted tensor product for $\operatorname{gr}_F\tau$,
which is acyclic by~(a). Strong convergence then lifts this
acyclicity to~$K_\tau^L$.
\end{proof}

\begin{theorem}[Fundamental theorem of chiral twisting morphisms; \ClaimStatusProvedHere]
\label{thm:fundamental-twisting-morphisms}
\index{twisting morphism!fundamental theorem|textbf}
\textup{[Regime: quadratic
\textup{(}Convention~\textup{\ref{conv:regime-tags})}.]}

\smallskip\noindent
This is the chiral analogue of \cite[Theorem~2.3.1]{LV12}.
For a chiral twisting datum
$(\cA, \cC, \tau, F_\bullet)$
\textup{(}Definition~\textup{\ref{def:chiral-twisting-datum})},
the following are equivalent:
\begin{enumerate}[label=\textup{(\roman*)}]
\item\label{ftm:koszul}
 $\tau$ is a chiral Koszul morphism
 \textup{(}Definition~\textup{\ref{def:chiral-koszul-morphism})}.
\item\label{ftm:counit}
 The canonical counit
 $\varepsilon_\tau\colon \Omega_X(\cC) \to \cA$
 is a quasi-isomorphism.
\item\label{ftm:unit}
 The canonical unit
 $\eta_\tau\colon \cC \to \bar{B}_X(\cA)$
 is a weak equivalence of conilpotent complete factorization
 coalgebras.
\item\label{ftm:acyclic}
 Both twisted tensor products $K_\tau^L(\cA,\cC)$ and
 $K_\tau^R(\cC,\cA)$ are acyclic.
\end{enumerate}
\end{theorem}

\begin{proof}
The equivalence follows the same logical structure as the
classical case \cite[Theorem~2.3.1]{LV12}, adapted to the
chiral/factorization setting.

\emph{\ref{ftm:koszul}$\Rightarrow$\ref{ftm:acyclic}:}
By definition (Definition~\ref{def:chiral-koszul-morphism},
condition~1).

\emph{\ref{ftm:acyclic}$\Rightarrow$\ref{ftm:counit}:}
By Lemma~\ref{lem:twisted-product-cone-counit},
$K_\tau^L \simeq \operatorname{Cone}(\varepsilon_\tau)[-1]$.
Acyclicity of $K_\tau^L$ is therefore equivalent to
$\varepsilon_\tau$ being a quasi-isomorphism.

\emph{\ref{ftm:acyclic}$\Rightarrow$\ref{ftm:unit}:}
By Lemma~\ref{lem:twisted-product-cone-unit},
$K_\tau^R \simeq \operatorname{Cone}(\eta_\tau)[-1]$.
Acyclicity of $K_\tau^R$ is therefore equivalent to
$\eta_\tau$ being a weak equivalence.

\emph{\ref{ftm:counit}$\Rightarrow$\ref{ftm:koszul}:}
Because the filtration in the twisting datum is exhaustive,
complete, and bounded below on both $\Omega_X(\cC)$ and
$\barB_X(\cA)$, the comparison spectral sequences converge
strongly. Lemma~\ref{lem:filtered-comparison} then gives classical
Koszulity of the associated graded datum and acyclicity of the
twisted tensor products. These are conditions~\textup{(1)} and
\textup{(2)} of Definition~\ref{def:chiral-koszul-morphism}; the
strong convergence just noted is condition~\textup{(3)}.

\emph{\ref{ftm:unit}$\Rightarrow$\ref{ftm:koszul}:}
The same completeness and bounded-below hypotheses give strong
convergence on the bar-side filtration. Lemma~\ref{lem:filtered-comparison-unit}
shows that the weak equivalence of $\eta_\tau$ forces the
associated graded datum
$(\operatorname{gr} \cA, \operatorname{gr} \cC,
\operatorname{gr} \tau)$ to be classically Koszul and lifts the
acyclicity of the classical twisted tensor products to
$K_\tau^L$ and $K_\tau^R$. Together with strong convergence, this
is exactly Definition~\ref{def:chiral-koszul-morphism}.

\end{proof}

\begin{corollary}[Three bijections for chiral twisting morphisms;
\ClaimStatusProvedHere]
\label{cor:three-bijections}
\index{twisting morphism!three bijections|textbf}
For a conilpotent complete factorization coalgebra $\cC$
and an augmented chiral algebra $\cA$ on a smooth curve~$X$,
there are natural bijections
\begin{align}\label{eq:three-bijections}
\mathrm{Tw}(\cC, \cA)
&\;\cong\;
\operatorname{Hom}_{\mathrm{alg}}(\Omega_X(\cC), \cA)
\notag\\
&\;\cong\;
\operatorname{Hom}_{\mathrm{coalg}}(\cC, \bar{B}_X(\cA))
\end{align}
where\/ $\mathrm{Tw}(\cC, \cA)$ is the set of degree~$+1$
morphisms $\tau\colon \cC \to \cA$ satisfying the
Maurer--Cartan equation $d\tau + \tau \star \tau = 0$.

The \emph{universal twisting morphism}
$\tau_{\mathrm{univ}}\colon \bar{B}_X(\cA) \to \cA$
(projection to cogenerators) represents the twisting
functor: every $\tau \in \mathrm{Tw}(\cC, \cA)$
factors uniquely through $\tau_{\mathrm{univ}}$.
Dually, the \emph{canonical twisting morphism}
$\iota\colon \cC \to \Omega_X(\cC)$
(inclusion of cogenerators) corepresents
twisting: every $\tau$ factors uniquely through~$\iota$.
\end{corollary}

\begin{proof}
The first bijection sends $\tau \in \mathrm{Tw}(\cC, \cA)$
to the algebra morphism $g_\tau\colon \Omega_X(\cC) \to \cA$
defined on cogenerators by $g_\tau(s^{-1} c) = \tau(c)$; the
MC equation for $\tau$ is equivalent to $g_\tau$ commuting
with differentials
(the same proof as \cite[Theorem~2.2.7]{LV12}).
The second bijection sends $\tau$ to the coalgebra morphism
$f_\tau\colon \cC \to \bar{B}_X(\cA)$ defined dually.
The universal property of the bar construction as a cofree
coalgebra ensures uniqueness of factorization.
\end{proof}

\begin{remark}[Why the MC equation is Stokes' theorem]
\label{rem:MC-is-Stokes}
\index{Maurer--Cartan equation!geometric origin}
The Maurer--Cartan equation $d\tau + \tau \star \tau = 0$ is
not an imposed axiom. It is forced by the geometry of
$\overline{C}_2(X)$. The convolution
$(\tau \star \tau)(x) = \mu \circ (\tau \otimes \tau)
\circ \Delta(x)$ evaluates the two-point OPE: it extracts
the residue of the propagator $\eta_{12}$ at the collision
divisor $D_{12} \subset \overline{C}_2(X)$, then multiplies
in~$\cA$. The term $d\tau$ computes the Stokes boundary
on $\overline{C}_2(X)$. The MC equation says:
\emph{on the compactified two-point configuration space,
the propagator residue is exact}.
At three points, this becomes the Arnold relation.
At $n$ points, $d_{\mathrm{bar}}^2 = 0$.

At genus~$g \geq 1$, the propagator acquires monodromy, the
Stokes argument picks up a curvature term
$\kappa(\cA) \cdot \omega_g$, and $\tau$ becomes curved
(Remark~\ref{rem:curved-twisting-higher-genus}).
\end{remark}

\begin{remark}[Theorem~A decomposition]\label{rem:theorem-A-decomposition}
\index{Theorem A!decomposition}
The fundamental theorem of chiral twisting morphisms
(Theorem~\ref{thm:fundamental-twisting-morphisms}) is
Theorem~$\mathrm{A}_0$ in the decomposition of the main
duality theorem:
\begin{itemize}
\item \emph{$\mathrm{A}_0$} (this theorem): for a single
 twisting datum, Koszulity $\Leftrightarrow$ counit qi
 $\Leftrightarrow$ unit we $\Leftrightarrow$ twisted tensor
 products acyclic.
\item \emph{$\mathrm{A}_1$}: bar concentration
 (Theorem~\ref{thm:bar-concentration}): for a Koszul morphism,
 $H^i(\bar{B}^{\mathrm{ch}}(\cA)) = 0$ for $i \neq 0$ and
 $H^0 \cong \cC$.
\item \emph{$\mathrm{A}_2$}: geometric bar-cobar duality
 (Theorem~\ref{thm:bar-cobar-isomorphism-main}): for a chiral
 dual pair, applying Verdier duality to the bar coalgebra of one side
 and then reading the result on the post-Verdier side produces the
 dual factorization algebra of the other.
\end{itemize}
$\mathrm{A}_0$ is the Koszul recognition theorem;
$\mathrm{A}_1$ extracts the dual coalgebra;
$\mathrm{A}_2$ links two Koszul data via Verdier geometry.
\end{remark}

\begin{proposition}[Universal property of bar construction {\cite{LV12}}; \ClaimStatusProvedElsewhere]
\label{prop:bar-universal-property}
\index{bar construction!universal property}
For any augmented chiral algebra~$\cA$, the bar construction
$\barB_X(\cA)$ together with its canonical twisting morphism
$\tau_{\mathrm{can}} \colon \barB_X(\cA) \to \cA$
is initial among pairs $(C, \tau)$ consisting of a conilpotent chiral
coalgebra~$C$ and a twisting morphism $\tau \colon C \to \cA$.
Concretely: for any twisting morphism $\alpha \colon C \to \cA$,
there exists a unique coalgebra map
$f_\alpha \colon C \to \barB_X(\cA)$
such that $\tau_{\mathrm{can}} \circ f_\alpha = \alpha$.
Dually, the cobar construction $\Omega_X(C)$ together with
$\iota_{\mathrm{can}} \colon C \to \Omega_X(C)$ is terminal among
pairs $(A, \iota)$ consisting of an augmented chiral algebra~$A$
and a twisting morphism $\iota \colon C \to A$.
\end{proposition}

\begin{proof}
This is the chiral analogue of
\cite[Theorem~2.2.8 and Proposition~2.2.7]{LV12}.
The universal property of the bar construction is the statement
that $\barB_X(\cA) = T^c_{\mathrm{chiral}}(s^{-1}\bar{\cA})$
is the cofree conilpotent coalgebra on the desuspended
augmentation ideal, equipped with the unique coderivation~$d$
whose projection to cogenerators recovers the chiral product.
Given $\alpha \colon C \to \cA$, the map $f_\alpha$ is the
unique coalgebra map $C \to T^c_{\mathrm{chiral}}(s^{-1}\bar{\cA})$
whose projection to cogenerators is $\alpha$; the MC equation
on~$\alpha$ ensures $f_\alpha$ is a chain map. The cobar
statement is dual.
\end{proof}

\begin{theorem}[Comparison lemma for chiral twisted products {\cite[Theorem~2.4.1]{LV12}}; \ClaimStatusProvedElsewhere]
\label{thm:chiral-comparison-lemma}
\index{comparison lemma!chiral|textbf}
\index{twisted tensor product!comparison}
Let $\alpha, \beta \colon C \to \cA$ be two chiral twisting morphisms
from the same conilpotent coalgebra~$C$ to the same algebra~$\cA$,
and let $f \colon C \to C$ be a coalgebra quasi-isomorphism such that
$\alpha = \beta \circ f$. Then the induced map on twisted tensor
products
\[
\mathrm{id}_\cA \otimes f \colon
\cA \otimes_\alpha C \;\xrightarrow{\;\sim\;}\; \cA \otimes_\beta C
\]
is a quasi-isomorphism. In particular, if
$\cA \otimes_\alpha C$ is acyclic, then so is
$\cA \otimes_\beta C$.

This is the chiral analogue of \textup{\cite[Theorem~2.4.1]{LV12}}.
The proof transfers directly: filter both sides by bar degree, observe
that $\mathrm{gr}(f)$ is an isomorphism (since $f$ is a quasi-isomorphism
of coalgebras), and conclude by spectral sequence comparison.
\end{theorem}

\subsection{Chiral Koszul pairs}

The Koszul-pair notion extends beyond the quadratic setting through a
set of \emph{recognition criteria} (checkable conditions on the pair)
from which the full bar-cobar identification follows as a theorem. No
quasi-isomorphism statement is built into the definition itself.

\begin{definition}[Chiral Koszul pair]\label{def:chiral-koszul-pair}
\index{Koszul pair!chiral|textbf}
A \emph{chiral Koszul pair} on a smooth projective curve~$X$
is a pair of chiral Koszul data
\textup{(}Definition~\textup{\ref{def:chiral-twisting-datum}},
Definition~\textup{\ref{def:chiral-koszul-morphism})}
\[
(\cA_1, \cC_1, \tau_1, F_\bullet), \qquad
(\cA_2, \cC_2, \tau_2, F_\bullet)
\]
in one fixed completed lane: the coalgebras~$\cC_i$ are complete
conilpotent factorization coalgebras, the filtrations are exhaustive
and separated, each conformal-weight/bar-length piece is finite over
the ground field or locally free of finite rank over~$\cO_X$, and the
transition maps satisfy the Mittag--Leffler condition used in
Lemma~\ref{lem:completion-convergence}. Verdier duals below are taken
in the holonomic/perfect $\cD$-module factorization category, or in the
equivalent continuous finite-window algebraic model. Without these
topological and finite-type hypotheses one has twisting data, not a
Verdier-Koszul pair. The pair is equipped with Verdier-compatible
factorization-algebra
identifications
\[
\mathbb{D}_{\operatorname{Ran}}(\cC_1) \simeq \Omega_X(\cC_2),
\qquad
\mathbb{D}_{\operatorname{Ran}}(\cC_2) \simeq \Omega_X(\cC_1),
\]
compatible with the twisting morphisms and filtrations.
In this situation we write $\cA_2 \simeq \cA_1^!$ and
$\cA_1 \simeq \cA_2^!$.
The notation names the post-Verdier algebra. The coalgebra~$\cC_i$
models the bar-dual coalgebra $\cA_i^{\mathrm i}$; it is not the
algebra $\cA_i^!$, and the cobar reconstruction
$\Omega_X(\cC_i)\simeq\cA_i$ is not a definition of Koszul duality.

\smallskip\noindent
These are \emph{antecedent hypotheses}: the twisting data
(Definition~\ref{def:chiral-twisting-datum}), the Koszulness
condition (Definition~\ref{def:chiral-koszul-morphism}),
and Verdier compatibility can all be verified without
invoking bar-cobar duality itself.
Theorem~\ref{thm:bar-cobar-isomorphism-main} then
\emph{proves} the full bar-cobar identification as a
consequence.

\smallskip\noindent
\emph{Standard construction.}
For the standard examples, the Koszul pair is constructed
as follows: $\cC_i = \bar{B}_X(\cA_i)$ with the canonical
twisting morphism, $F_\bullet$ is the PBW filtration
(Theorem~\ref{thm:pbw-koszulness-criterion}), acyclicity
follows from spectral sequence degeneration, and the
algebra-level Verdier compatibility
\[
\mathbb{D}_{\operatorname{Ran}}(\cC_1)
\simeq \Omega_X(\cC_2)
\]
is supplied by Theorem~\ref{thm:verdier-bar-cobar}, which is already
a statement about the factorization algebra obtained after applying
$\mathbb{D}_{\operatorname{Ran}}$ to the bar coalgebra, together with
the counit quasi-isomorphism
$\Omega_X(\cC_2) \xrightarrow{\sim} \cA_2$ from
Theorem~\ref{thm:fundamental-twisting-morphisms}.
The Heisenberg pair
$(\mathcal{H}_k, \mathrm{Sym}^{\mathrm{ch}}(V^*))$ is the
prototype (\S\ref{sec:frame-koszul-dual}).

\smallskip\noindent
In particular, each algebra~$\cA_i$ in a chiral Koszul pair is a
Koszul chiral algebra in the sense of
Definition~\textup{\ref{def:koszul-chiral-algebra}}, and hence lies on
the Koszul locus~$\operatorname{Kosz}(X)$
(Definition~\textup{\ref{def:koszul-locus}}).
\end{definition}

\begin{remark}[Five-object firewall]\label{rem:five-object-firewall-kp}
For a chiral algebra~$\cA$ there are five objects, and no two of them
are identified without an explicit functor:
\[
\cA,\qquad
\barB_X(\cA),\qquad
\cA^{\mathrm i},\qquad
\cA^!,\qquad
Z^{\mathrm{der}}_{\mathrm{ch}}(\cA).
\]
The bar entry has two ambient realizations. The symmetric bar
$\barB_X^\Sigma(\cA)$ is the Ran bar coalgebra used in
Theorem~A and in the comparison with Beilinson--Drinfeld factorization
homology. The ordered bar
$\barB_X^{\mathrm{ord}}(\cA)=T^c(s^{-1}\bar\cA)$ is the
$\Eone$ refinement which remembers the tensor order and carries the
Yangian, affine Yangian, and $r$-matrix constructions. It maps to the
symmetric bar only after imposing the appropriate locality, descent, or
averaging condition; it is not the same object before that passage.

The symbol $\cA^{\mathrm i}$ denotes the intrinsic bar-dual coalgebra,
computed by the relevant bar construction on the chosen lane. The symbol
$\cA^!$ denotes the post-Verdier Koszul-dual chiral algebra, obtained
from $\cA^{\mathrm i}$ by Verdier duality, or by continuous finite-type
linear duality in the affine algebraic model. Finally,
$Z^{\mathrm{der}}_{\mathrm{ch}}(\cA)$ is the cochain-level chiral
Hochschild closed-sector actor and is not a Koszul dual. It becomes a
physical bulk object only after a boundary local-operator comparison map
has been supplied and proved to be a quasi-isomorphism in the chosen
completion.

Thus the type-correct maps are
\[
\barB_X(\cA)\longrightarrow\cA^{\mathrm i}
\quad\text{a coalgebra quasi-isomorphism on the Koszul locus},
\]
\[
\mathbb D_{\operatorname{Ran}}\cA^{\mathrm i}\simeq\cA^!
\quad\text{as chiral algebras on the finite-type Verdier lane},\qquad
\Omega_X\barB_X(\cA)\simeq\cA\quad\text{on the Koszul locus}.
\]
The last equivalence is bar-cobar inversion. It is not the definition of
Koszul duality and it does not identify $\Omega_X\barB_X(\cA)$ with
$\cA^!$.
\end{remark}


\begin{example}[Heisenberg specialization]\label{ex:heisenberg-koszul-pair}
For the Heisenberg algebra $\mathcal{H}_k$ at level $k \neq 0$,
Definition~\ref{def:chiral-koszul-pair} yields the Koszul pair
$(\mathcal{H}_k, \mathrm{Sym}^{\mathrm{ch}}(V^*))$,
where $\mathcal{H}_k$ is $\chirLie$-type (generated by a single
bosonic field $a(z)$ with OPE $a(z)a(w) \sim k/(z-w)^2$)
and $\mathrm{Sym}^{\mathrm{ch}}(V^*)$ is $\chirCom$-type: the
chiral incarnation of classical $\mathrm{Lie}$--$\mathrm{Com}$ duality.
The explicit bar computation gives the bar coalgebra
$\bar{B}^{\mathrm{ch}}(\mathcal{H}_k) \simeq
\mathrm{coLie}^{\mathrm{ch}}(V^*)$.
Applying Theorem~\ref{thm:bar-cobar-isomorphism-main} to this bar
coalgebra identifies the Verdier dual of the resulting bar-dual
coalgebra with the Koszul dual algebra
$\mathcal{H}_k^! = \mathrm{Sym}^{\mathrm{ch}}(V^*)$:
the bar complex $\bar{B}(\mathcal{H}_k)$ was computed degree by
degree in \S\ref{sec:frame-bar-deg1}--\S\ref{sec:frame-bar-all},
producing the coalgebra whose finite dual is the commutative Koszul dual
(\S\ref{sec:frame-koszul-dual}).
Because $\mathcal{H}_k$ is abelian, the PBW spectral sequence
(Theorem~\ref{thm:pbw-koszulness-criterion}) is trivial and the
Koszul complex $K_\tau(\mathcal{H}_k,
\mathrm{Sym}^{\mathrm{ch}}(V^*))$ is manifestly acyclic.
\end{example}

\begin{example}[Explicit correspondence: \texorpdfstring{$bc$}{bc} ghost system and \texorpdfstring{$\beta\gamma$}{beta-gamma} system]
\label{ex:bc-betagamma-koszul}
On the split genus-$0$ finite-type quadratic lane, consider the
Verdier-compatible pair $(\mathcal{BC}, \mathcal{BG})$ where:
\begin{itemize}
\item $\mathcal{BC}$ is the $bc$ ghost system with fermionic fields
 $b(z), c(z)$ and OPE $b(z)c(w) \sim (z-w)^{-1}$;
\item $\mathcal{BG}$ is the $\beta\gamma$ system with bosonic fields
 $\beta(z), \gamma(z)$ and OPE
 $\beta(z)\gamma(w) \sim (z-w)^{-1}$.
\end{itemize}

The comparison has two type-distinct steps. First the bar computation
gives the intrinsic coalgebra
\[
\bar{B}^{\mathrm{ch}}(\mathcal{BC})
\xrightarrow{\;\sim\;}
\mathcal{BC}^{\mathrm i}
\]
whose corelations are residue-orthogonal to the exterior
relations. Second, after finite Verdier duality,
$(\mathcal{BC}^{\mathrm i})^\vee$ is identified with the
post-Verdier algebra~$\mathcal{BG}$. Thus the classical
$\Lambda^* \leftrightarrow \mathrm{Sym}^*$ exchange appears only after
the coalgebra step has been dualized.

\emph{Generator-relation exchange.}
\begin{align*}
\text{$bc$ OPE: } b(z)c(w) &\sim (z-w)^{-1}
 \quad \text{(fermionic exterior relations)}, \\
\text{$\beta\gamma$ OPE: } \beta(z)\gamma(w) &\sim (z-w)^{-1}
 \quad \text{(bosonic symmetric dual algebra)}.
\end{align*}

\emph{Geometric picture.}
The bar complex $\bar{B}^{\text{ch}}(\mathcal{BC})$ involves:
\[\bar{B}^{\text{ch}}(\mathcal{BC})_n = \Gamma\left(\overline{C}_n(X), (b \oplus c)^{\boxtimes n} \otimes \Omega^*_{\log}\right)\]
The residues at collision divisors define the coproduct and the
bar-dual corelations. Cobar reconstructs $\mathcal{BC}$ from this
coalgebra on the Koszul locus; finite Verdier duality, not cobar
reconstruction, identifies the post-Verdier algebra with
$\mathcal{BG}$.
\end{example}

\subsection{PBW deformation as a Koszulness criterion}
\label{sec:koszulness-verification}
\index{Koszul property!verification}
\index{PBW filtration!Koszulness verification}

A systematic method reduces
chiral Koszulness to classical Koszulness of the PBW-associated
graded via a flat deformation argument.

\begin{theorem}[PBW criterion for chiral Koszulness; \ClaimStatusProvedHere]
\label{thm:pbw-koszulness-criterion}
\index{Koszul property!PBW criterion|textbf}
\textup{[Regime: filtered-complete
\textup{(}Convention~\textup{\ref{conv:regime-tags})}.]}

Let $\cA$ be a chiral algebra with PBW filtration
$F_0 \subset F_1 \subset \cdots \subset \cA$
such that the associated graded
$\operatorname{gr}_F \cA$ is a commutative chiral algebra
(equivalently, a vertex Poisson algebra).
Suppose:
\begin{enumerate}
\item\label{item:pbw-flat} The filtration is \emph{flat}: each $F_p/F_{p-1}$ is a free
 $\cO_X$-module of finite rank in each conformal weight.
\item\label{item:pbw-classical-koszul} The associated graded $\operatorname{gr}_F \cA$
 is classically Koszul: the Koszul complex
 $\barBgeom(\operatorname{gr}_F \cA) \otimes_\tau \operatorname{gr}_F \cA$
 is acyclic in positive degrees.
\item\label{item:pbw-bounded} For each bar degree $n$ and conformal
 weight $h$, the chain group $\barBgeom^n_h(\cA)$ is finite-dimensional,
 the induced filtration on this bidegree is complete and separated, and
 the finite-window quotients form the Mittag--Leffler system of
 Lemma~\ref{lem:completion-convergence}.
\end{enumerate}
Then $\cA$ is chiral Koszul: the Koszul complex
$\barBgeom(\cA) \otimes_\tau \cA$ is acyclic in positive degrees.
\end{theorem}

\begin{proof}
Equip the Koszul complex $K = \barBgeom(\cA) \otimes_\tau \cA$ with
the filtration induced by $F$:
\[
F_p K^n = \sum_{i+j = n} F_i \barBgeom(\cA) \otimes_\tau F_j \cA.
\]
This is a bounded-below, exhaustive filtration compatible with the
differential (the bar differential and the twisting morphism $\tau$
both respect $F$ because the OPE is filtered by PBW degree; conformal
weight supplies the finite windows in
Hypothesis~\ref{item:pbw-bounded}). The resulting spectral sequence has:
\[
E_0^{p,q} = \operatorname{gr}^p_F K^{p+q},
\qquad
d_0 = \operatorname{gr}_F(d_K).
\]
By flatness~\ref{item:pbw-flat}, the associated graded of $K$ is
the Koszul complex of $\operatorname{gr}_F \cA$:
\[
\operatorname{gr}_F K \cong
\barBgeom(\operatorname{gr}_F \cA) \otimes_\tau \operatorname{gr}_F \cA.
\]
(This uses that the bar construction commutes with taking associated
graded for flat filtrations: the residue differential of the associated
graded sees only the leading-order OPE terms, i.e., the vertex Poisson
bracket.)

By hypothesis~\ref{item:pbw-classical-koszul}, the $E_1$ page
is concentrated in degree $0$:
\[
E_1^{p,q} = H^{p+q}(\operatorname{gr}^p K) =
\begin{cases}
\bC & p = q = 0, \\
0 & p + q > 0.
\end{cases}
\]
All higher differentials $d_r$ for $r \geq 1$ are therefore zero
(they map between zero groups), so the spectral sequence collapses
at $E_1$.

By Hypothesis~\ref{item:pbw-bounded}, the completion has no
$R^1\varprojlim$ contribution in any fixed bidegree, and
Lemma~\ref{lem:completion-convergence} gives strong convergence:
$E_\infty = \operatorname{gr} H^*(K)$.
Since $E_\infty$ is concentrated in degree~$0$,
we conclude $H^n(K) = 0$ for $n > 0$.
\end{proof}

\begin{remark}[Relation to classical Koszulness]
\label{rem:classical-to-chiral-koszul}
The PBW criterion reduces chiral Koszulness to the classical statement that $\operatorname{gr}_F \cA$ is Koszul.
For Kac--Moody and Virasoro algebras, $\operatorname{gr}_F \cA$ is a polynomial algebra (Koszul by Priddy's theorem); semisimplicity of $\mathfrak{g}$ is not required.
\end{remark}

\begin{theorem}[Affine Kac--Moody algebras are chiral Koszul; \ClaimStatusProvedHere]
\label{thm:km-chiral-koszul}
\index{Kac--Moody algebra!chiral Koszulness|textbf}
\textup{[Regime: curved-central
\textup{(}Convention~\textup{\ref{conv:regime-tags})}.]}

For any simple Lie algebra $\fg$ and generic level
$k \notin \Sigma(\fg)$ (the finite exceptional set of
Theorem~\textup{\ref{thm:bar-cohomology-level-independence}}),
the affine Kac--Moody vertex algebra
$\widehat{\fg}_k$ is chiral Koszul.
\end{theorem}

\begin{proof}
We verify the three hypotheses of
Theorem~\ref{thm:pbw-koszulness-criterion}.

\emph{Hypothesis~\ref{item:pbw-flat}.}
The PBW filtration on $\widehat{\fg}_k$ is defined by
$F_p = \operatorname{span}\{J^{a_1}_{-n_1} \cdots J^{a_r}_{-n_r}
|0\rangle : r \leq p\}$.
Each graded piece $F_p/F_{p-1}$ is a direct sum of
weight spaces of the symmetric algebra
$S^p(\fg \otimes t^{-1}\bC[t^{-1}])$, which is a free
$\bC$-module of finite rank in each conformal weight.

\emph{Hypothesis~\ref{item:pbw-classical-koszul}.}
The associated graded is:
\[
\operatorname{gr}_F \widehat{\fg}_k
\cong \operatorname{Sym}(\fg \otimes t^{-1}\bC[t^{-1}]).
\]
This is a polynomial algebra on the generators
$\{J^a_{-n} : a = 1, \ldots, \dim\fg,\; n \geq 1\}$.
By Priddy's theorem (\cite[Theorem~3.4.5]{LV12}),
polynomial algebras are Koszul, with Koszul dual the
exterior coalgebra $\Lambda^c(\fg^* \otimes t\bC[t])$.
Equivalently, the Koszul complex of $\operatorname{Sym}(V)$
is the de~Rham complex of the polynomial ring, which is
acyclic in positive degrees (Poincar\'{e} lemma for
formal power series).

\emph{Hypothesis~\ref{item:pbw-bounded}.}
The bar chain group $\barBgeom^n_h(\widehat{\fg}_k)$ is a
subspace of $(\fg \otimes t^{-1}\bC[t^{-1}])^{\otimes n}
\otimes \Omega^{n-1}(\overline{C}_n)$ with total conformal
weight~$h$. Since the number of monomials of weight~$h$
in $n$ tensor factors is bounded by $p(h)^n < \infty$,
each chain group is finite-dimensional.

All three hypotheses are verified.
By Theorem~\ref{thm:pbw-koszulness-criterion},
$\widehat{\fg}_k$ is chiral Koszul.
\end{proof}

\begin{theorem}[Virasoro chiral Koszulness; \ClaimStatusProvedHere]
\label{thm:virasoro-chiral-koszul}
\index{Virasoro algebra!chiral Koszulness|textbf}
\textup{[Regime: curved-central\textup{;}
Convention~\textup{\ref{conv:regime-tags}}.]}

For any central charge $c \in \bC$, the Virasoro vertex algebra
$\mathrm{Vir}_c$ is chiral Koszul.
\end{theorem}

\begin{proof}
The PBW filtration on $\mathrm{Vir}_c$ is defined by the
number of Virasoro mode applications:
\[
F_p = \operatorname{span}\{L_{-n_1} \cdots L_{-n_r}
|0\rangle : r \leq p,\; n_i \geq 2\}.
\]

\emph{Hypothesis~\ref{item:pbw-flat}.}
Each $F_p/F_{p-1} \cong S^p(V)$ where
$V = \bigoplus_{n \geq 2} \bC \cdot L_{-n}$
(the Virasoro generators modulo vacuum), which is
free of finite rank in each conformal weight.

\emph{Hypothesis~\ref{item:pbw-classical-koszul}.}
$\operatorname{gr}_F \mathrm{Vir}_c
\cong \operatorname{Sym}(V)$
is a polynomial algebra (the OPE reduces to a commutative
product on the associated graded: the non-linear terms
$T_{(0)}T = \partial T$, $T_{(1)}T = 2T$ are lower-order in the
PBW filtration, and the quartic pole $T_{(3)}T = c/2$ is a
scalar that contributes to the curvature, not the
associated graded bracket).
By Priddy's theorem, $\operatorname{Sym}(V)$ is Koszul.

\emph{Hypothesis~\ref{item:pbw-bounded}.}
Same argument as the Kac--Moody case: the number of
partitions of weight~$h$ into at most $n$ parts (each $\geq 2$)
is finite.

By Theorem~\ref{thm:pbw-koszulness-criterion},
$\mathrm{Vir}_c$ is chiral Koszul.
\end{proof}

\begin{corollary}[Bar cohomology computes the bar-dual coalgebra; \ClaimStatusProvedHere]
\label{cor:bar-cohomology-koszul-dual}
\index{bar complex!Koszul dual computation}
For a chiral Koszul algebra $\cA$ satisfying the hypotheses of
Theorem~\textup{\ref{thm:pbw-koszulness-criterion}},
the PBW spectral sequence on $\barBgeom(\cA)$ collapses
at $E_2$, and the bar cohomology is the intrinsic bar-dual
coalgebra~$\cA^{\mathrm i}$. On a finite-type Verdier lane its
continuous dual is the algebra~$\cA^!$; hence the finite graded
dimensions agree:
\[
\dim H^n(\barBgeom(\cA))
= \dim(\cA^{\mathrm i})_n
= \dim(\cA^!)_n.
\]
In particular:
\begin{enumerate}
\item For $\widehat{\fg}_k$ at generic level:
 $\dim H^n(\barBgeom(\widehat{\fg}_k))
 = \dim(\widehat{\fg}_k^{\mathrm i})_n$, and after finite
 continuous duality this equals $\dim(\widehat{\fg}_k^!)_n$
 \textup{(}Theorem~\textup{\ref{thm:universal-kac-moody-koszul})}.
 The Koszul dual $\widehat{\fg}_k^!$ has the same modular characteristic
 as $\widehat{\fg}_{-k-2h^\vee}$, but the two are distinct
algebras: $\widehat{\fg}_k^!$ is the linear dual of the bar
coalgebra, not the affine algebra at the dual level.
\item For $\mathrm{Vir}_c$:
 $\dim H^n(\barBgeom(\mathrm{Vir}_c))
 = \dim(\mathrm{Vir}_c^{\mathrm i})_n$; only after the finite-window
 Verdier passage is this also the dimension of
 $(\mathrm{Vir}_c^!)_n$.
\end{enumerate}
\end{corollary}

\begin{proof}
\emph{Step~1: The PBW filtration on the bar complex.}
Equip $\barBgeom(\cA)$ with the filtration inherited from
the PBW filtration $F$ on $\cA$:
\[
F_p \barBgeom^n(\cA) = \bigl\{
a_1 \otimes \cdots \otimes a_{n+1} \otimes \omega
\;\big|\; a_i \in F_{p_i}\cA,\;
\textstyle\sum p_i \leq p\bigr\}
\]
where $\omega \in \Omega^n_{\log}(\overline{C}_{n+1}(X))$.
This filtration is exhaustive (every bar element lies in
some $F_p$ by finiteness of conformal weight) and
compatible with the differential (the bar differential
respects conformal weight since the OPE does).

\emph{Step~2: The associated graded.}
The filtration $F_p$ is placed on the algebra factor $\cA$ in the bar construction $T^c(s^{-1}\bar{\cA})$
underlying $\barBgeom(\cA)$; the configuration-space form factor
$\Omega^\bullet_{\log}(\overline{C}_\bullet(X))$ carries PBW degree~$0$.
Since $T^c$ is graded cofree and Hypothesis~\ref{item:pbw-flat} makes each
$\operatorname{gr}_F^p\cA$ a free graded $\cO_X$-module in each conformal weight,
tensor-product flatness gives
$\operatorname{gr}_F T^c(s^{-1}\bar{\cA}) \cong T^c(s^{-1}\operatorname{gr}_F\bar{\cA})$.
The bar differential decomposes along the filtration because the OPE respects
conformal weight, giving
$\operatorname{gr}_F d_{\mathrm{bar}} = d_{\mathrm{bar},\,\operatorname{gr}_F\cA}$.
Therefore:
\begin{equation}\label{eq:pbw-associated-graded-bar}
\operatorname{gr}_F \barBgeom(\cA) \;\cong\;
\barBgeom(\operatorname{gr}_F \cA)
\;=\; \barBgeom(\operatorname{Sym}^{\mathrm{ch}}(V))
\end{equation}
where $V$ is the generating $\mathcal{D}_X$-module.
This identification holds because the bar differential on the
associated graded sees only the leading-order OPE terms.
For $\widehat{\fg}_k$: the singular OPE
$J^a(z) J^b(w) \sim f^{ab}_c J^c(w)/(z-w)
+ k\kappa^{ab}/(z-w)^2$
reduces on $\operatorname{gr}_F$ to the commutative product
(the bracket $f^{ab}_c J^c$ and level term $k\kappa^{ab}$ are
lower-order in the PBW filtration).
For $\mathrm{Vir}_c$: similarly, the Virasoro OPE reduces to a
commutative product on the associated graded.

\emph{Step~3: Computation of the $E_1$ page.}
The $E_0$ page is
$E_0^{p,q} = \operatorname{gr}^p_F \barBgeom^{p+q}(\cA)$,
with $d_0$ the associated-graded bar differential.
By~\eqref{eq:pbw-associated-graded-bar}, the $E_1$ page
is the bar cohomology of the commutative chiral algebra
$\operatorname{Sym}^{\mathrm{ch}}(V)$.

Since $\operatorname{Sym}^{\mathrm{ch}}(V)$ is a free
commutative algebra, it is Koszul with Koszul dual
the exterior coalgebra $\Lambda^{\mathrm{c}}(V^*)$
(Priddy's theorem applied to the chiral setting,
cf.\ Remark~\ref{rem:classical-to-chiral-koszul}).
Therefore:
\begin{equation}\label{eq:e1-page-exterior}
E_1^{p,q} = H^{p+q}(\barBgeom^p(\operatorname{Sym}^{\mathrm{ch}}(V)))
= \begin{cases}
\Lambda^p(V^*) & \text{if } q = 0, \\
0 & \text{if } q \neq 0.
\end{cases}
\end{equation}
The $E_1$ page is \emph{concentrated in the row $q = 0$}.
In particular, for $\widehat{\fg}_k$ with
$V = \fg \otimes t^{-1}\bC[t^{-1}]$:
\[
E_1^{n,0} = \Lambda^n(\fg^* \otimes t\bC[t])
\]
is the $n$-th exterior power of the dual loop algebra
generators (graded by conformal weight).

\emph{Step~4: $E_2$ collapse.}
Since $E_1^{p,q} = 0$ for $q \neq 0$, the higher differentials
$d_r\colon E_r^{p,q} \to E_r^{p+r, q-r+1}$ vanish for
$r \geq 2$: the target has second index
$q - r + 1 \leq 0 - 2 + 1 = -1 < 0$, placing it in the
zero region. The differential
$d_1\colon E_1^{p,0} \to E_1^{p+1,0}$ maps \emph{within}
the nonzero row $q = 0$ and is generally nonzero:
it is the Chevalley--Eilenberg differential of the loop
algebra $\fg \otimes t^{-1}\bC[t^{-1}]$ acting on
$\Lambda^*(\fg^* \otimes t\bC[t])$ by bracket contraction
(the leading PBW correction to the commutative bar
differential). Therefore $E_\infty = E_2$ and:
\[
H^n(\barBgeom(\cA)) \cong E_\infty^{n,0}
= E_2^{n,0}
= H^n_{\mathrm{CE}}\!\bigl(\fg \otimes t^{-1}\bC[t^{-1}],\, \bC\bigr)
= (\cA^{\mathrm i})_n.
\]
If the finite-type Verdier dual exists, the continuous dual of this
coalgebraic piece is the corresponding graded piece of~$\cA^!$.
For commutative chiral algebras (e.g.\ the Heisenberg algebra),
the Lie bracket of $\fg$ vanishes, $d_1 = 0$,
$E_\infty = E_2 = E_1 = \Lambda(V^*)$, and the bar cohomology
reduces to the exterior algebra. For non-abelian algebras
(Kac--Moody, Virasoro), $d_1 \neq 0$ and the bar
cohomology is computed by the CE cohomology, a proper
subquotient of $\Lambda(V^*)$.

\emph{Step~5: Verification for $\widehat{\mathfrak{sl}}_2$.}
$V = \bC^3 \otimes t^{-1}\bC[t^{-1}]$
(three generators $J^+, J^0, J^-$, each with modes
$J^a_{-n}$ for $n \geq 1$).
The $E_1$ page has
$\dim {E_1}^{n,0}_h = \dim \Lambda^n(\bC^3 \otimes t\bC[t])_h$.
For example: $\dim {E_1}^{1,0}_1 = 3$ (the three weight-$1$
generators), $\dim {E_1}^{2,0}_2 = \binom{3}{2} = 3$
(wedge products $J^{a*}_{+1} \wedge J^{b*}_{+1}$, $a < b$).
The $d_1$ differential is the CE differential of
$\mathfrak{sl}_2 \otimes t^{-1}\bC[t^{-1}]$;
the resulting $E_2$ page gives:
$\dim H^1 = 3$ (all at weight $h = 1$; $d_1 = 0$ at this
weight since there are no weight-$1$ chains in degree~$0$),
and $\dim H^2 = 5$ (concentrated at weight $h = 3$:
chain groups $\dim\Lambda^1_3 = 3$,
$\dim\Lambda^2_3 = 9$,
$\dim\Lambda^3_3 = 1$; differential ranks
$3$ and $1$; cohomology
$(9 - 1) - 3 = 5$).
The computation is verified by checking $d_1^2 = 0$
at each weight (Computation~\ref{comp:sl2-ce-verification}).
\end{proof}

\begin{remark}[CE cohomology vs exterior algebra dimensions]
\label{rem:ce-vs-exterior}
\index{bar complex!CE vs exterior algebra}
Bar cohomology of a non-abelian chiral Koszul algebra is the CE cohomology of $\fg \otimes t^{-1}\bC[t^{-1}]$, not the exterior algebra $\Lambda^n(V^*)$. For $\widehat{\mathfrak{sl}}_2$: $E_1$ totals $= \prod_{m \geq 1}(1+q^m)^3$, but $E_\infty = E_2 \neq E_1$ because the $d_1$ differential is nonzero for non-abelian $\fg$. In particular, $\dim H^2 = 5$ (not $R(5) = 6$; Appendix~\ref{app:combinatorial-frontier}).
\end{remark}

\begin{remark}[Chevalley--Eilenberg cohomology vs chiral bar cohomology]%
\label{rem:ce-vs-chiral-bar}%
\index{Chevalley--Eilenberg cohomology!vs chiral bar}%
\index{bar complex!vs CE cohomology}%
\index{Witt algebra!CE vs bar}%
\index{Orlik--Solomon form!contribution to discrepancy}%
The Chevalley--Eilenberg cohomology of the negative-mode
Lie algebra $\fg_{<0} = \fg \otimes t^{-1}\bC[t^{-1}]$
and the chiral bar cohomology $H^*(\barBgeom(\cA))$ give
\emph{different} Betti sequences, though both compute from
the same underlying OPE data.
For the Witt algebra $\mathrm{Witt} = \mathrm{Der}(\bC((t)))$,
the negative-mode Lie algebra is the Lie algebra of
polynomial vector fields vanishing to order~$\geq 2$
at the origin, and its CE cohomology has
$\dim H^1_{\mathrm{CE}}(\mathrm{Witt}_{<0}) = 2$
(the two independent cocycles $L_{-1}^*$ and~$L_{-2}^*$;
$L_{-3}^*$ is \emph{not} a cocycle since
$d(L_{-3}^*)(L_{-1}, L_{-2}) = L_{-3}^*([L_{-1}, L_{-2}])
= L_{-3}^*(L_{-3}) = 1 \neq 0$).
The chiral bar cohomology of the Virasoro algebra has
$\dim H^1(\barBgeom(\mathrm{Vir})) = 1$
(a single generator~$T$ of conformal weight~$2$).
The discrepancy $2 - 1 = 1$ quantifies the
Orlik--Solomon form contribution: the chiral bar differential
uses the full logarithmic residue $d\log(z_1 - z_2)$
(the Arnold relation on configuration space),
which collapses the three CE cocycles to the single
chiral primary.
\end{remark}

\begin{theorem}[Bar concentration for Koszul pairs; \ClaimStatusProvedHere]
\label{thm:bar-concentration}
\index{bar complex!concentration}
\textup{[Regime: quadratic on the Koszul locus
\textup{(}Convention~\textup{\ref{conv:regime-tags})}.]}

Let $(\cA_1, \cA_2)$ be a chiral Koszul pair
\textup{(}Definition~\textup{\ref{def:chiral-koszul-pair})}.
Then the augmented bar complex $\bar{B}^{\mathrm{ch}}(\cA_1)$,
equipped with its bigrading by bar degree~$p$ and bar-differential
cohomological degree~$q$, satisfies:
\[
H^{p,q}(\bar{B}^{\mathrm{ch}}(\cA_1)) = 0
\quad\text{for } q \neq 0,
\qquad
H^{p,0}(\bar{B}^{\mathrm{ch}}(\cA_1)) \cong (\cA_1^{\mathrm i})_p.
\]
Equivalently,
$\bar{B}^{\mathrm{ch}}(\cA_1) \simeq \cA_1^{\mathrm i}$ in the
derived category of graded chiral coalgebras, where
$\cA_1^{\mathrm i}=H^*(\barB^{\mathrm{ch}}(\cA_1))$ carries its
internal grading concentrated in bar-differential degree~$0$.
After continuous linear or Verdier duality on the finite-type lane,
$((\cA_1^{\mathrm i})^\vee)$ is the algebra denoted
$\cA_1^!\simeq\cA_2$.
In particular, the natural map
$\bar{B}^{\mathrm{ch}}(\cA_1) \to \cA_1^{\mathrm i}$ is a
quasi-isomorphism of graded chiral coalgebras.
\end{theorem}

\begin{proof}
Filter $\bar{B}^{\mathrm{ch}}(\cA_1)$ by bar degree.
By the PBW criterion (Theorem~\ref{thm:pbw-koszulness-criterion}),
the associated graded is the classical Koszul complex of
$\operatorname{gr}_F \cA_1$, which is acyclic in positive
degrees by hypothesis.
Corollary~\ref{cor:bar-cohomology-koszul-dual} then identifies the
graded pieces of $H^*(\barB^{\mathrm{ch}}(\cA_1))$ with those of the
bar-dual coalgebra $\cA_1^{\mathrm i}$.
It remains to show that the comparison is a coalgebra
quasi-isomorphism, not merely a graded vector-space identification.
The twisting morphism
$\tau\colon \bar{B}^{\mathrm{ch}}(\cA_1) \to \cA_2$
satisfies the Maurer--Cartan equation $\partial(\tau) + \tau \star \tau = 0$
in the chiral convolution algebra
$\mathrm{Hom}(\bar{B}^{\mathrm{ch}}(\cA_1), \cA_2)$
(Definition~\ref{def:chiral-koszul-morphism}), so the
twisted tensor products
$K_\tau^L(\cA_1, \cC_1)$ and $K_\tau^R(\cC_1, \cA_1)$
are well-defined chain complexes (twisted differentials square to
zero by the MC equation).
These are acyclic by hypothesis.
The canonical projection to cohomology induces a coalgebra map
\[
\rho\colon \barB^{\mathrm{ch}}(\cA_1)\longrightarrow \cA_1^{\mathrm i}
\]
because the bar comultiplication $\Delta_{\bar{B}}$ is a chain map.
Acyclicity of $K_\tau^L$ and $K_\tau^R$ upgrades this map from a
graded comparison to a coalgebra isomorphism
(the chiral analogue of \cite[Theorem~2.3.1]{LV12}; the
identical argument applies because acyclicity of the Koszul
complexes is the only input beyond the MC equation, and both
hold in the chiral setting by
Lemma~\ref{lem:bar-holonomicity} and
Theorem~\ref{thm:fundamental-twisting-morphisms}).
The algebra $\cA_2\simeq\cA_1^!$ is obtained only after applying
the continuous/Verdier dual to this coalgebra.
\end{proof}

%%% =======================================================
%%% INTRINSIC CHARACTERIZATIONS OF CHIRAL KOSZULNESS
%%% =======================================================

\subsection{Intrinsic characterizations of chiral Koszulness}
\label{sec:intrinsic-characterizations-koszulness}
\index{Koszul property!intrinsic characterizations}
\index{MC element!genus-0 formality}

The MC element $\Theta_\cA \in \MC(\mathfrak{g}_\cA^{\mathrm{mod}})$
lives in a modular convolution algebra bigraded by degree and genus.
Chiral Koszulness is a \emph{genus-$0$ formality property}:
it says the genus-$0$ component $\Theta_\cA^{(0)}$ is
determined by its binary part; equivalently, all genus-$0$
obstruction classes $o_{r+1}^{(0)}(\cA) = 0$ vanish for $r \geq 2$.
The shadow obstruction tower probes all genera; Koszulness constrains only
genus~$0$. Twelve characterizations and consequences of chiral
Koszulness split into seven intrinsic bidirectional equivalences,
the Barr--Beck--Lurie consequence of the bar-cobar counit, the
conditional factorization-homology comparison lane, the one-way
chiral Hochschild consequence on the Koszul locus, the
perfectness/non-degeneracy conditional Lagrangian criterion, and the
one-directional D-module purity implication. Each reads
$\Theta_\cA^{(0)}$ through a different invariant of the
convolution algebra, but only the first seven form the intrinsic
unconditional equivalence core.

\begin{remark}[One-loop exactness]
\label{rem:one-loop-exactness}
\index{PBW filtration!one-loop exactness}
\index{one-loop exactness}

The PBW spectral sequence
(Corollary~\ref{cor:bar-cohomology-koszul-dual}) admits a sharp
physics interpretation. The $E_0$ page is the tree-level bar complex
$\barBgeom(\operatorname{Sym}^{\mathrm{ch}}(V))$, with cohomology
$\Lambda^{\mathrm{c}}(V^*)$ by Priddy's theorem. The $d_1$
differential is the Chevalley--Eilenberg differential extracted from
the simple-pole OPE $a_{(0)}b$: a single operator collision.
Collapse $E_2 = E_\infty$ means that no multi-particle collisions
(the higher poles $a_{(n)}b$, $n \geq 2$, generating $d_r$ for
$r \geq 2$) contribute additional cohomology. Chiral Koszulness is
therefore \emph{one-loop exactness}: the Lie-algebraic data from the
single-pole OPE determines the bar cohomology completely.
\end{remark}

\begin{proposition}[Formality implies chiral Koszulness; \ClaimStatusProvedHere]
\label{prop:ainfty-formality-implies-koszul}
\index{Koszul property!$A_\infty$ characterization}
\index{$A_\infty$-algebra!minimal model obstruction}

Let $\cA$ be a chiral algebra on a smooth curve~$X$, and let
$(m_2, m_3, m_4, \ldots)$ denote the transferred
$A_\infty$-structure on the finite homotopy dual
$\cA^!_\infty = H^*(\barBgeom_X(\cA))^\vee$
obtained via a deformation retract
\[
\barBgeom_X(\cA) \;\underset{p}{\overset{i}{\rightleftarrows}}\;
H^*(\barBgeom_X(\cA))
\qquad \text{with contracting homotopy } h
\]
and the homological perturbation lemma
\textup{(}Kadeishvili \cite{Kad82}; cf.\ \cite{LPWZ09}\textup{)}.

If the minimal $A_\infty$-structure is \emph{formal}
($m_n = 0$ for all $n \geq 3$), then $\cA$ is chirally Koszul.
\end{proposition}

\begin{proof}
Formality of the $A_\infty$-structure means the bar complex
$\barBgeom_X(\cA)$ is formal as a dg coalgebra: the natural map to
its cohomology $H^*(\barBgeom_X(\cA))$ is a quasi-isomorphism of
$A_\infty$-coalgebras. The higher products $m_n$ ($n \geq 3$) are
identified with the PBW spectral sequence differentials $d_{n-1}$
via the standard comparison (cf.\ \cite[Theorem~3.4.1]{LPWZ09}):
$m_n = 0$ for $n \geq 3$ is equivalent to $E_2 = E_\infty$, which
is the PBW collapse criterion of
Theorem~\ref{thm:pbw-koszulness-criterion}.
\end{proof}

\begin{theorem}[Converse: chiral Koszulness implies formality;
\ClaimStatusConditional]
\label{thm:ainfty-koszul-characterization}%

The converse of Proposition~\textup{\ref{prop:ainfty-formality-implies-koszul}}
holds: $\cA$ is chirally Koszul if and only if $m_n = 0$ for all $n \geq 3$.
\end{theorem}

\begin{proof}
Forward: Proposition~\ref{prop:ainfty-formality-implies-koszul}.
Converse: if $m_n = 0$ for $n \geq 3$, the minimal $A_\infty$-model
of $\barBgeom_X(\cA)$ is a strict dg algebra with product~$m_2$.
The chiral $A_\infty$ structure on $H^*(\barBgeom_X(\cA))$ is
computed fiberwise on each FM stratum
(Proposition~\ref{prop:shadow-formality-low-degree}). On each fiber,
the HPL transfer is an ordinary $A_\infty$ transfer over a field, to
which Keller's classicality theorem applies. The PBW filtration is
defined fiberwise and compatible with the FM stratification, so
fiberwise $E_2$-collapse assembles to global $E_2$-collapse: all
transferred differentials $d_r = 0$ for $r \geq 2$ because
$m_n = 0$. Hence Koszulness.
\end{proof}

\begin{remark}[$A_\infty$ products as genus-$0$ shadows]
\label{rem:ainfty-genus0-shadow}

The sequence $(m_3, m_4, \ldots)$ is the genus-$0$ analogue of the
shadow obstruction tower $\Theta_\cA^{\leq 3}, \Theta_\cA^{\leq 4},
\ldots$ of Theorem~\ref{thm:mc2-bar-intrinsic}: the $m_n$ measure
nonlinearity of the bar differential at genus~$0$, while the shadows
$\Theta_\cA^{\leq r}$ measure it at all genera.
\end{remark}

\begin{example}[Truncated polynomials: the homotopy Koszul dual off the Koszul locus]
\label{ex:truncated-polynomial-ainfty}
\index{truncated polynomial algebra!$A_\infty$ dual}
\index{Koszul dual!homotopy, off the Koszul locus}

Consider $A = k[x]/(x^n)$ for $n \geq 3$ (non-Koszul; the Koszul
case $n = 2$ gives the exterior algebra). The bigraded bar
cohomology $H^{p,q}(\barB(A))$ has generators
\[
a \in H^{1,1}, \qquad b \in H^{2,n},
\]
and $H$ is \emph{not} concentrated on the diagonal $p = q$,
the signature of non-Koszulness
(Theorem~\ref{thm:ainfty-koszul-characterization}).

The transferred $A_\infty$-structure on $H^*(\barB(A))^\vee$
carries exactly one nontrivial higher operation:
\[
m_2 = \cdots = m_{n-2} = 0 \quad \text{on } a,
\qquad
m_{n-1}(\underbrace{a, \ldots, a}_{n-1}) = b.
\]
This is the Massey product detecting the off-diagonal class; it is
the first and only nontrivial higher operation
(Lu--Palmieri--Wu--Zhang \cite{LPWZ09}). The homotopy Koszul dual
$A^!_\infty = (H^*(\barB(A))^\vee, m_2, m_{n-1})$ is a genuine
$A_\infty$-algebra whose non-formality is measured by the single
operation~$m_{n-1}$. For $n = 2$ (exterior algebra), $m_{n-1} = m_1
= 0$ on $H$, recovering formality and hence Koszulness.

The chiral analogue: for a chiral algebra~$\cA$ whose bar
cohomology has off-diagonal bigraded classes, the homotopy Koszul
dual $\cA^!_\infty$ carries genuine higher chiral operations
measuring non-formality. The shadow obstruction tower
(Remark~\ref{rem:ainfty-genus0-shadow}) is the all-genera lift of
this phenomenon.
\end{example}

\begin{theorem}[Ext diagonal vanishing criterion; \ClaimStatusProvedHere]
\label{thm:ext-diagonal-vanishing}%
\index{Koszul property!Ext criterion}
\index{Ext algebra!diagonal vanishing}

Define the \emph{chiral Ext algebra}
$\operatorname{Ext}^{*,*}_\cA(\omega_X, \omega_X)
= \bigoplus_{p,q} \operatorname{Ext}^{p,q}_\cA(\omega_X, \omega_X)$,
bigraded by cohomological degree~$p$ and conformal weight~$q$.
Then $\cA$ is chirally Koszul if and only if
$\operatorname{Ext}^{p,q}_\cA(\omega_X, \omega_X) = 0$ for $p \neq q$.

\smallskip\noindent\textbf{Classical precedent.}
On the formal disk $\hat{D}$, this reduces to the
Beilinson--Ginzburg--Soergel criterion for $\operatorname{Zhu}(\cA)$
\cite{BGS96}.
On $\bP^1$, diagonal vanishing follows from PBW since
$H^1(\bP^1, \cO) = 0$.
At genus $g \geq 1$, the Hodge-theoretic contributions from
$H^0(X, \omega_X)$ produce the exceptional set $\Sigma(\fg)$ of
Theorem~\ref{thm:bar-cohomology-level-independence} where diagonal
vanishing may fail.

\end{theorem}

\begin{proof}
Forward: $E_2$-collapse concentrates $H^*(\barBgeom(\cA))$ on
the diagonal $p = q$; since the bar resolution computes Ext,
$\operatorname{Ext}^{p,q}(\omega_X, \omega_X) = 0$ for
$p \neq q$.
Converse: diagonal concentration forces all PBW differentials
$d_r$ ($r \geq 2$) to vanish, since a $d_r$-differential shifts
the bigrading by $(r, 1-r)$, producing off-diagonal classes
that contradict~$\operatorname{Ext}^{p,q} = 0$ for $p \neq q$.
Hence $E_2$-collapse, hence Koszulness.
\end{proof}

\begin{proposition}[PBW universality; \ClaimStatusProvedHere]
\label{prop:pbw-universality}
\index{Koszul property!PBW universality}
\index{vertex algebra!freely strongly generated}

A vertex algebra~$\cA$ is \emph{freely strongly generated} if it
admits strong generators $\{a^i\}_{i \in I}$ whose normally ordered
monomials
$:\!\partial^{n_1} a^{i_1} \cdots \partial^{n_r} a^{i_r}\!:$
form a PBW basis; equivalently,
$\operatorname{gr}_F \cA \cong \operatorname{Sym}^{\mathrm{ch}}(V)$
where $V$ is the $\cD_X$-module spanned by the generators.

Every freely strongly generated vertex algebra is chirally Koszul.
\end{proposition}

\begin{proof}
The three hypotheses of Theorem~\ref{thm:pbw-koszulness-criterion}
hold: (1)~flatness of $F_p/F_{p-1} \cong S^p(V)$ by the PBW basis;
(2)~classical Koszulness of $\operatorname{gr}_F \cA \cong
\operatorname{Sym}^{\mathrm{ch}}(V)$ by Priddy's theorem;
(3)~finite-dimensionality of bar chain groups in each bigrading by
the partition bound of Theorem~\ref{thm:km-chiral-koszul}.
\end{proof}

\begin{corollary}[Universal vertex algebras are chirally Koszul;
\ClaimStatusConditional]
\label{cor:universal-koszul}
\index{Koszul property!universal algebras}

The following are chirally Koszul:
\begin{enumerate}
\item $V_k(\fg)$ at all levels~$k$ \textup{(}including $k = -h^\vee$\textup{)},
 since its vacuum module is Verma with no null vectors.
 \emph{Caveat}: at $k = -h^\vee$ the PBW spectral sequence
 still collapses, but diagonal Ext concentration fails
 \textup{(}Remark~\textup{\ref{rem:pbw-vs-diagonal-critical}}\textup{)}:
 ``chirally Koszul'' here means PBW-Koszul
 \textup{(}condition~\textup{(ii)} of
 Theorem~\textup{\ref{thm:koszul-equivalences-meta}}\textup{)},
 not Ext-diagonal
 \textup{(}condition~\textup{(iv)}\textup{)}.
\item $\mathrm{Vir}_c$ as the universal vertex algebra at all~$c$,
 since the vacuum Verma module is freely strongly generated.
\item $\mathcal{W}^k(\fg)$ obtained by Drinfeld--Sokolov reduction,
 by the Feigin--Frenkel theorem \textup{(}the Miura images freely
 generate\textup{)}.
\end{enumerate}

\noindent\textbf{The universal vs.\ simple distinction.}
The universal algebra $V_k(\fg)$ is always chirally Koszul.
The simple quotient $L_k(\fg) = V_k(\fg)/I_k$ may fail Koszulness
at levels where the singular vector lies in the bar-relevant range
; the precise criterion is
Theorem~\textup{\ref{thm:kac-shapovalov-koszulness}}.
\end{corollary}

\begin{proof}
Proposition~\ref{prop:pbw-universality} applies to every freely
strongly generated vertex algebra. The universal affine algebra
$V_k(\fg)$ and the universal Virasoro algebra are freely strongly
generated by their current, respectively stress-tensor, fields, so
items~(1) and~(2) follow immediately. For
$\mathcal{W}^k(\fg)$, the Feigin--Frenkel Miura theorem supplies a
free strong generating set by the Miura images, so the same
proposition applies. The final paragraph records the
universal-versus-simple warning for these families.
\end{proof}

\begin{remark}[Critical level: PBW degeneration vs diagonal Ext concentration]
\label{rem:pbw-vs-diagonal-critical}
\index{Koszul property!critical level subtlety}
At the critical level $k = -h^\vee$, the PBW spectral sequence
degenerates at~$E_2$ (the $E_2$-collapse mechanism is
$k$-independent, since $\operatorname{gr}_F V_k(\fg) \cong
\operatorname{Sym}^{\mathrm{ch}}(\fg)$ for all~$k$).
However, the $E_2$ page is \emph{not} diagonally concentrated:
the bar cohomology
$H^*(\barB(V_{-h^\vee}(\fg))) = \Omega^*(\mathrm{Op}_{\fg^\vee}(D))$
\textup{(}Theorem~\textup{\ref{thm:langlands-bar-bridge}}\textup{)}
has $H^0 = \operatorname{Fun}(\mathrm{Op}) \neq \bC$,
contributing off-diagonal classes
$\operatorname{Ext}^{0,q} \neq 0$ for all even $q \geq 0$.
Hence the diagonal criterion of
Theorem~\textup{\ref{thm:pbw-koszulness-criterion}}
\textup{(}$\operatorname{Ext}^{p,q} = 0$ for
$p \neq q$\textup{)} \emph{fails} at $k = -h^\vee$.

At generic level ($\kappa \neq 0$), the curvature of the bar
differential forces $H^0(\barB) = \bC$ and concentrates bar
cohomology on the diagonal. The critical level is the unique
point where curvature vanishes, PBW degeneration persists, but
diagonal concentration is lost.

Accordingly, the parenthetical ``including $k = -h^\vee$'' in
item~(1) of
Corollary~\textup{\ref{cor:universal-koszul}}
should be understood as: PBW degeneration holds at all~$k$;
the full Koszulness
\textup{(}diagonal Ext concentration\textup{)}
holds at generic~$k$
but fails at $k = -h^\vee$.
Theorem~H
\textup{(Theorem~\ref{thm:hochschild-polynomial-growth})}
requires the diagonal criterion, so it does not apply at
critical level.
See Remark~\textup{\ref{rem:critical-level-lie-vs-chirhoch}}
for the Beilinson--Drinfeld comparison that replaces
Theorem~H at critical level.
\end{remark}

\begin{theorem}[Kac--Shapovalov criterion for simple quotients;
\ClaimStatusConditional]
\label{thm:kac-shapovalov-koszulness}
\index{Koszul property!Kac--Shapovalov criterion}
\index{Kac determinant!Koszulness criterion}
\index{Shapovalov form!Koszulness criterion}

Let $\cA$ be a conformal vertex algebra with Shapovalov form on the
vacuum module. The \emph{bar-relevant range} is the set of conformal
weights~$h$ with $\barBgeom^n_h(\cA) \neq 0$ for some $n > 0$.
Assume the PBW filtration is flat and split by conformal weight, the
bar spectral sequence is bidegreewise finite and strongly convergent,
and the PBW-to-bar comparison is detected by the Shapovalov pairing:
a nonzero class on the PBW page pairs nontrivially with some class in
the same bar-relevant weight. Under these hypotheses, $\cA$ is chirally
Koszul if and only if the Kac--Shapovalov determinant
$\det G_h \neq 0$ for all~$h$ in the bar-relevant range.

For universal algebras the PBW flatness and pairing hypotheses hold
by Corollary~\ref{cor:universal-koszul}. The genuine content is for
simple quotients: $L_k(\fg)$ has null vectors at
$h = h_{\mathrm{sing}}(k)$ by the Kac--Kazhdan theorem; at admissible
levels $k + h^\vee = p/q$, $(p,q)=1$, these lie in the bar-relevant
range and Koszulness may fail. For the Virasoro algebra, the Kac
determinant $\det G_h = \prod_{rs \leq h}(h - h_{r,s}(c))^{p(h-rs)}$
is nonvanishing in the bar-relevant range for the universal algebra,
but $L(c_{p,q}, 0)$ may fail at minimal-model values.
\end{theorem}

\begin{proof}
Forward: Koszulness gives $E_2$-collapse, hence the PBW
filtration on the vacuum module is strict:
$\operatorname{gr}_F M_0(\cA) \cong
\operatorname{Sym}^{\mathrm{ch}}(V)$. Strictness implies the
Shapovalov form~$G_h$ is non-degenerate in the bar-relevant
range (a degenerate $G_h$ would give a null vector in the
PBW-associated graded, contradicting strictness).

Converse: under the PBW/Shapovalov detection hypothesis,
$\det G_h \neq 0$ in the bar-relevant range means the
PBW-to-bar comparison map
$\operatorname{Sym}^{\mathrm{ch}}(V)_h \to \barBgeom_h(\cA)$ is
injective at every bar-relevant weight. Injectivity forces
$E_2$-collapse: a nonzero differential $d_r$ with $r \geq 2$ would
kill an element that pairs nontrivially with the Shapovalov form
(since the element is nonzero in $\operatorname{gr}_F$ and $G_h$ is
non-degenerate), contradicting the $d_r$-exactness. Strong convergence
then identifies the collapsed spectral sequence with bar cohomology.
\end{proof}

\begin{remark}[Rationality is not, by itself, a Koszul criterion]
\label{rem:rationality-not-koszul-criterion}
\index{Koszul property!rationality not sufficient}
\index{rational vertex algebra!limitations of rationality}
Rationality, $C_2$-cofiniteness, and semisimplicity of the Zhu
algebra are strong \emph{finiteness} inputs, but they do not by
themselves imply chiral Koszulness in the sense used here.
The obstruction is categorical: the ordinary positive-energy module
category $V$-$\mathrm{gmod}$ controls Zhu theory and conformal
blocks, whereas Ext groups in the
chiral/factorization module setting are computed via the bar complex
\textup{(}Corollary~\ref{cor:bar-computes-ext}\textup{)}.
No theorem here identifies these two Ext theories for general rational
vertex algebras.

Consequently, the implication
\[
\text{ordinary semisimplicity}
\;\Longrightarrow\;
H^{>0}\!\bigl(\barBgeom(V)\bigr)=0
\]
is \emph{not} established here. For simple quotients at
admissible level the Kac--Shapovalov route is visibly obstructed
by null vectors in the bar-relevant range
\textup{(}Theorem~\ref{thm:kac-shapovalov-koszulness}\textup{)},
while the rationality route stops at finiteness of the ordinary
representation theory.
\end{remark}

\begin{remark}[Admissible simple quotients]
\label{rem:admissible-koszul-status}
\index{admissible level!Koszulness status}
For the universal algebra $V_k(\fg)$, free strong generation
gives chiral Koszulness at every level
\textup{(}Corollary~\ref{cor:universal-koszul}\textup{)}.
For the simple quotient $L_k(\fg)$, null vectors enter the
bar-relevant range at admissible levels, so the PBW/Shapovalov
criterion no longer applies directly.

What \emph{is} proved for non-degenerate admissible levels is:
\begin{enumerate}[label=\textup{(\roman*)}]
\item rationality and semisimplicity of the ordinary module
 category \textup{(}Theorem~\ref{thm:admissible-rep-theory}(ii)\textup{)};
\item finite-page bar and Li--bar spectral sequences in the
 ordinary-character calculations: the worked
 $\widehat{\mathfrak{sl}}_2$, $k=-1/2$ bar derivation of the vacuum
 Kac--Wakimoto character \textup{(}Theorem~\ref{thm:kw-bar-spectral}\textup{)},
 together with weightwise finite-page control and the external
 general-rank Kac--Wakimoto formulas on the broader
 non-degenerate admissible lane
 \textup{(}Proposition~\ref{prop:bar-admissible} and
 Theorem~\ref{thm:kw-bar-general-rank}\textup{)}.
\end{enumerate}
These results are compatible with chiral Koszulness, but they do
not yet prove bar-cobar inversion or Ext-diagonal purity for the
simple quotient in general.

For $\fg = \mathfrak{sl}_2$, admissible Koszulness is settled:
the simple quotient $L_k(\mathfrak{sl}_2)$ is chirally Koszul at
\emph{every} admissible level $k = -2 + p/q$ ($p \geq 2$,
$q \geq 1$, $\gcd(p,q) = 1$). The argument is structural:
the null-vector ideal $I_k$ is generated in a single conformal
weight, the bar complex of $L_k$ is the quotient bar complex
$\barB(V_k)/\barB(I_k)$, and the quotient bar spectral sequence
degenerates at~$E_2$ by the Kac--Wakimoto character formula.

For $\fg = \mathfrak{sl}_3$ at denominator $q \geq 3$, the finite
admissible Li--bar model isolates a Cartan obstruction. When the
reduced associated graded in the bar-relevant window is the truncated
Cartan/root model displayed in
Theorem~\ref{thm:admissible-sl3-non-koszul-qge3}, the bar spectral
sequence has eight off-diagonal generator classes: six root classes
and two Cartan classes. The Poisson differential $d_1$ acts through
the Lie bracket; it kills the root-generator classes, while the two
Cartan classes survive the $d_1$ page because
$[\mathfrak h,\mathfrak h]=0$. If no non-reduced extension or hidden
higher differential hits those Cartan classes, they give
$H^2(\barB(L_k(\mathfrak{sl}_3)))\neq0$.
At denominator $q \leq 2$, the off-diagonal classes lie on the
$E_1$ diagonal where $d_1 = 0$, and no obstruction appears; the present
argument proves no non-Koszulness statement there.

The mechanism is uniform: for any simple $\fg$ of
$\mathrm{rk}(\fg) \geq 2$, the abelian Cartan subalgebra
$\mathfrak{h}$ produces $\mathrm{rk}(\fg)$ classes in
$E_1^{2,q}$ that are annihilated by the Lie-bracket derivation
$d_1$ (since $[\mathfrak{h}, \mathfrak{h}] = 0$). Persistence to
$E_\infty$ is the extra finite-model assertion isolated below.
\end{remark}

\begin{theorem}[Cartan obstruction in the finite admissible
Li--bar model for $L_k(\mathfrak{sl}_3)$]%
\label{thm:admissible-sl3-non-koszul-qge3}%
\ClaimStatusConditional
\index{admissible level!Koszulness obstruction}%
\index{Cartan subalgebra!bar obstruction}%
Let $k = -3 + p/q$ be a non-degenerate admissible level for
$\widehat{\mathfrak{sl}}_3$, with coprime $p \geq 3$, $q \geq 3$.
Assume the following finite admissible Li--bar package.
\begin{enumerate}[label=\textup{(A\arabic*)}]
\item The quotient bar complex is bidegreewise finite and the Li--bar
 spectral sequence converges strongly in each conformal weight.
\item In the bar-relevant window the reduced associated graded is
 \[
 R^{\mathrm{red}} \;\cong\;
 \mathbb{C}[H_1]/(H_1^{d_C}) \otimes
 \mathbb{C}[H_2]/(H_2^{d_C})
 \otimes
 \bigotimes_{\alpha > 0}
 \mathbb{C}[x_\alpha]/(x_\alpha^{d_R}),
 \qquad d_R=(p-2)q,\quad d_C=(p-1)q.
 \]
\item The $d_1$ differential on the resulting Li--bar page is the
 Kirillov--Kostant Poisson differential of
 Proposition~\ref{prop:li-bar-poisson-differential}.
\item The non-reduced nilradical and extension data of $R_{L_k}$ do
 not create a higher differential hitting the two Cartan classes
 $[H_1]$ and~$[H_2]$.
\item The periodic-CDG projection used in
 Theorem~\ref{thm:admissible-kl-bar-cobar-adjunction} is available
 and absorbs, rather than kills, the surviving Cartan classes.
\end{enumerate}
Then the simple quotient $L_k(\mathfrak{sl}_3)$ is \emph{not}
chirally Koszul on the original, non-periodic bar complex. More
precisely:
\begin{enumerate}[label=\textup{(\roman*)}]
\item $\dim H^2(\barB^{\mathrm{ch}}(L_k(\mathfrak{sl}_3)))
 \geq \mathrm{rk}(\mathfrak{sl}_3) = 2$.
\item The bar--cobar counit
 $\varepsilon \colon \Omega^{\mathrm{ch}}(\barB^{\mathrm{ch}}(L_k)) \to L_k$
 is \emph{not} a chain quasi-isomorphism on the original
 complex. It becomes a quasi-isomorphism only after the
 $2p$-periodic projection of
 Theorem~\ref{thm:periodic-cdg-is-koszul-compatible}.
\item In the twelve-equivalence meta-theorem
 \textup{(}Theorem~\ref{thm:koszul-equivalences-meta}\textup{)},
 $L_k(\mathfrak{sl}_3)$ fails conditions \textup{(i)}--\textup{(vi)}
 and \textup{(ix)}--\textup{(x)}, but satisfies the
 weakened periodic-CDG analogues of \textup{(i)} and
 \textup{(ii)} modulo $\sigma_{2p}$ via
 Theorem~\ref{thm:admissible-kl-bar-cobar-adjunction}.
\end{enumerate}
\end{theorem}

\begin{proof}
\emph{Step~1 (Finite admissible model).}
Hypothesis~\textup{(A1)} supplies the finite, strongly convergent
Li--bar spectral sequence. Arakawa finite-length results and
Finkelberg semisimplification are the standard sources of finite
admissible control, but the theorem uses them only through the explicit
finite package stated above.

\emph{Step~2 \textup{(}$E_1$ page via K\"unneth on the truncated
associated graded\textup{)}.}
Hypothesis~\textup{(A2)} gives the displayed reduced associated
graded, with $h_\theta=(p-2)q$ and $h_\alpha=(p-1)q$.
The K\"unneth decomposition of
$\mathrm{Tor}^{R^{\mathrm{red}}}(\mathbb{C}, \mathbb{C})$ produces
eight generator classes at $E_1^{2, *}$: six root-generator
classes at conformal weight $h_\theta$ and two Cartan classes
$[H_1], [H_2]$ at conformal weight $h_\alpha$.

\emph{Step~3 (Cartan obstruction: $d_1$ vanishes on Cartan
classes).}
By Hypothesis~\textup{(A3)} and
Proposition~\ref{prop:li-bar-poisson-differential}, the $d_1$
differential on the Li--bar spectral sequence is the Poisson
bracket on $R_V$, transported by the Killing form to the
Kirillov--Kostant bracket on $\mathfrak{sl}_3^*$. On
generator classes, $d_1([x_a])$ is the linearization of
$[a, -]$. Abelianness of the Cartan subalgebra,
$[\mathfrak{h}, \mathfrak{h}] = 0$, forces
\[
d_1([H_i]) = 0, \qquad i = 1, 2.
\]
The six root-generator classes at $E_1^{2, h_\theta}$ are
mapped by $d_1$ into $E_1^{1, h_\theta + 1}$ (a diagonal position),
where surjectivity follows from the non-degeneracy of the Killing
pairing on the root space; those six classes are killed. The two
Cartan classes persist to $E_2^{2, h_\alpha}$.

\emph{Step~4 (Higher differentials cannot resurrect the Cartan
classes).}
A higher differential $d_r$ ($r \geq 2$) from the position
$E_r^{2, h_\alpha}$ lands in $E_r^{2-r, h_\alpha + r - 1}$.
For $r = 2$ the target $E_2^{0, h_\alpha + 1}$ vanishes: the
origin $E_2^{0, *}$ is concentrated at conformal weight $0$, and
$h_\alpha + 1 = (p-1)q + 1 > 0$ for $q \geq 3$, forcing the
conformal-weight mismatch. For $r \geq 3$ the target filtration
degree $2 - r < 0$ is empty. Hence
\[
E_\infty^{2, h_\alpha} \;=\; E_2^{2, h_\alpha} \;\geq\; 2,
\]
with equality in finite windows where Hypothesis~\textup{(A4)} has
been checked directly.

The rank $\geq 2$ threshold is genuine: for
$\mathfrak{sl}_2$, the single Cartan class $[H]$ at
$E_1^{2, h_\alpha}$ co-resonates with a root target at
$E_1^{1, h_\alpha + 1}$ via the Poisson bracket $[H, F] = -2F$,
and $d_r$ cancellation through that channel restores Koszulness
(consistent with the sharp $L_k(\mathfrak{sl}_2)$ result at all
admissible levels recalled above). For $\mathfrak{sl}_3$ the
two-dimensional Cartan decouples because $[H_1, H_2] = 0$
removes the analogous cancellation.

\emph{Step~5 (Falsification of the bar--cobar counit).}
If $L_k(\mathfrak{sl}_3)$ were chirally Koszul, the counit
$\varepsilon$ would be a chain quasi-isomorphism and
$H^n(\barB^{\mathrm{ch}}(L_k))$ would vanish for $n \geq 2$ at
weights off the diagonal. Step~4 contradicts this at
$(n, \text{weight}) = (2, h_\alpha)$ with dimension $\geq 2$,
establishing~(i). Hence $\varepsilon$ fails on the original
complex at weight $h_\alpha$, establishing the first half
of~(ii). Under Hypothesis~\textup{(A5)}, the periodic-CDG adjunction
Theorem~\ref{thm:admissible-kl-bar-cobar-adjunction} recovers
a quasi-isomorphism modulo $\sigma_{2p}$: the surviving Cartan
classes are absorbed into the $2p$-periodic direct sum of Ext
bigradings, but they are not killed. This establishes the
second half of~(ii). The failure of
conditions~(i)--(vi) and~(ix)--(x) of
Theorem~\ref{thm:koszul-equivalences-meta} is a direct
consequence of $H^2 \neq 0$ and the non-isomorphism of
$\varepsilon$; the $2p$-periodic CDG analogues of~(i)
and~(ii) are the content of
Theorem~\ref{thm:admissible-kl-bar-cobar-adjunction}. This
establishes~(iii).
\end{proof}

\begin{conjecture}[Higher-rank admissible non-Koszulness, $\mathrm{rk}(\fg) \geq 3$]%
\label{conj:admissible-koszul-rank-obstruction}%
\ClaimStatusConjectured
\index{admissible level!higher-rank Koszulness obstruction}%
For any simple Lie algebra $\fg$ with $\mathrm{rk}(\fg) \geq 3$,
the simple quotient $L_k(\fg)$ at admissible level with
denominator $q \geq 3$ is \emph{not} chirally Koszul, with
$\dim H^2(\barB(L_k(\fg))) \geq \mathrm{rk}(\fg)$. At
$\mathrm{rk}(\fg) = 2$ this is
Theorem~\ref{thm:admissible-sl3-non-koszul-qge3}; for
$\mathrm{rk}(\fg) \geq 3$ subtle $d_2$ contributions from
triple-root brackets remain to be controlled.
\end{conjecture}

\begin{remark}[$W(2)$ at $c = -2$: Koszulness status]
\label{rem:w2-koszulness}
\index{W(2) algebra@$W(2)$ algebra!Koszulness}%
\index{Kazhdan filtration!$W(2)$ Koszulness}%
\index{symplectic fermion!orbifold}%
The $W(2)$ algebra at $c = -2$ (the triplet algebra at $p = 2$,
$\mathbb{Z}/2$-orbifold of the symplectic fermion
algebra~$\mathrm{SF}$) has Koszulness status \emph{open}.
PBW universality
\textup{(}Proposition~\ref{prop:pbw-universality}\textup{)} does
not apply directly: $W(2)$ is \emph{not} freely strongly
generated, since the vacuum module carries a singular vector at
conformal weight~$5$
\textup{(}Adamovi\'c--Milas;
Feigin--Gainutdinov--Semikhatov--Tipunin\textup{)},
which lies in the bar-relevant range ($h = 5 \geq 2 h_{\min} = 4$).
The Kac-determinant zero at weight~$15$ is a
Virasoro phenomenon, not a $W(2)$ null vector; but the
weight-$5$ singular vector is genuine and obstructs the PBW route.

A potential alternative route uses the Kazhdan filtration:
$W(2)$ embeds into the lattice VOA~$V_{2\mathbb{Z}}$ as the
momentum-zero sector. If the Kazhdan-filtered associated graded
of $W(2)$ (not of $V_{2\mathbb{Z}}$) is polynomial, then
Koszulness would follow. This requires verifying that
the weight-$5$ singular vector is absorbed by the Kazhdan
filtration, which is not established.
\end{remark}

\begin{remark}[Arithmetic shadows: Koszulness as harmonic metric]
\label{rem:kac-shapovalov-arithmetic-shadow}
The Kac--Shapovalov non-degeneracy criterion is the analytic shadow of a deeper structure: by
Theorem~\ref{thm:general-nahc} (Chapter~\ref{chap:arithmetic-shadows}), Koszulness of $\cA$
is equivalent to the nondegeneracy of the Shapovalov form on the
shadow triple $(\Theta_\cA, M_\cA, h_\cA)$: the propagator
$P = h^{-1}$ mediates between the flat side (MC element) and the
Higgs side (spectral curve).
\end{remark}

\begin{remark}[Deformation quantization compatibility]
\label{rem:dq-koszul-compatibility}
\index{Koszul property!deformation quantization}
\index{deformation quantization!Koszul compatibility}

Let $\cA_\hbar$ be a flat deformation of a classically Koszul
commutative chiral algebra~$\cA_0$, with
$\cA_\hbar/\hbar\cA_\hbar \cong \cA_0$.
Proposition~\ref{prop:pbw-universality} implies: $\cA_\hbar$ is
chirally Koszul whenever the PBW filtration deforms flatly
($\operatorname{gr}_F \cA_\hbar \cong \cA_0$). All standard
families satisfy this: for $V_k(\fg)$ with $\hbar = k^{-1}$, the
bracket $[J^a, J^b]_\star = f^{ab}_c J^c + k\kappa^{ab}$ lies in
$F_{\geq 1}$; Virasoro and $\mathcal{W}$-algebras are analogous
with $\hbar = c^{-1}$. Flat PBW deformation is therefore the
mechanism producing chirally Koszul algebras from commutative seeds.
\end{remark}

\subsection{The Li--bar spectral sequence and geometric Koszulness}
\label{subsec:li-bar-geometric-koszulness}

The Kac--Shapovalov criterion
(Theorem~\ref{thm:kac-shapovalov-koszulness}) detects Koszulness
by non-degeneracy of a bilinear form weight by weight. A geometric
refinement is as follows: the Li filtration on~$\cA$ induces a
spectral sequence on the bar complex whose $E_1$~page is the bar
complex of a \emph{commutative Poisson algebra}, and Koszulness
is controlled by the geometry of the associated variety~$X_\cA$.

\begin{construction}[Li--bar spectral sequence]
\label{constr:li-bar-spectral-sequence}
\index{Li filtration!bar spectral sequence}
\index{bar complex!Li filtration spectral sequence}
Let $V$ be a vertex algebra with Li filtration~$F_p V$
(Definition~\ref{def:associated-variety}). The bar complex
$\barBgeom(V) = \bigoplus_n \barBgeom^n(V)$ inherits a
multiplicative filtration
\[
F_p \barBgeom^n(V)
= \sum_{p_1 + \cdots + p_n = p}
 F_{p_1} V \otimes \cdots \otimes F_{p_n} V
 \cap \barBgeom^n(V).
\]
The Li filtration is multiplicative with respect to chiral
operations: $a_{(m)} b \in F_{p+q}$ when $a \in F_p$,
$b \in F_q$, and $m \leq -2$ (the defining condition
of~$F_p$), while $a_{(-1)} b \in F_{p+q}$ (product) and
$a_{(0)} b \in F_{p+q-1}$ (Poisson bracket lowers filtration
by~$1$). The bar differential $d_{\barB}$, assembled from
the chiral operations $\mu_n$, respects the filtration:
$d_{\barB}(F_p \barBgeom^n) \subset F_{p-1} \barBgeom^{n-1}$.
This gives rise to a convergent spectral sequence
\begin{equation}\label{eq:li-bar-ss}
E_1^{p,q}
= H^{p+q}\!\bigl(\operatorname{gr}^F_p \barBgeom(V)\bigr)
\;\Longrightarrow\;
H^{p+q}\!\bigl(\barBgeom(V)\bigr).
\end{equation}
The $E_0$~page is the bar complex of the associated
graded~$R_V := \operatorname{gr}^F V$, which is a commutative
vertex algebra with Poisson bracket $\{-,-\}$ induced by the
$0$-product (Definition~\ref{def:associated-variety}).
The $E_0$~differential is the bar differential of~$R_V$ as a
commutative algebra; the $d_1$~differential arises from the
Poisson bracket.
\end{construction}

\begin{proposition}[Poisson differential on the Li--bar $E_1$ page;
\ClaimStatusProvedHere]
\label{prop:li-bar-poisson-differential}
\index{Li filtration!Poisson differential on bar}
\index{Poisson bracket!bar complex differential}

The $d_1$ differential on the Li--bar spectral
sequence~\eqref{eq:li-bar-ss} is determined by the Poisson
bracket on $R_V = \operatorname{gr}^F V$:
\[
d_1 \colon E_1^{p,q} \to E_1^{p-1,q},
\qquad
d_1 = \sum_{i < j}
\{-,-\}_{ij} \circ \operatorname{contract}_{ij}
\]
where $\{-,-\}_{ij}$ applies the Poisson bracket to the
$i$-th and $j$-th tensor factors. In particular:
\begin{enumerate}[label=\textup{(\roman*)}]
\item When $R_V \cong \operatorname{Sym}(V_0)$ is a polynomial
 algebra with trivial Poisson bracket, $d_1 = 0$.
 Any further degeneration beyond~$E_1$ requires separate
 control of the higher filtration-lowering pieces of the
 vertex-algebra differential.
\item When $R_V = \mathcal{O}(Y)$ for an affine Poisson
 variety~$(Y, \pi)$, the $E_1$ page is the ordinary
 commutative bar cohomology
 $H^*_{\mathrm{bar}}(\mathcal{O}(Y))$, and $d_1$ is the
 differential induced there by~$\pi$. Equivalently, the
 $E_2$ page is the homology of this commutative
 bar-cohomology page with respect to that bracket-induced
 differential.
\end{enumerate}
\end{proposition}

\begin{proof}
The $d_1$ differential is the component of $d_{\barB}$ that
lowers the Li filtration by exactly~$1$. In the OPE
$a_{(n)} b$ on a vertex algebra, the filtration behavior is:
$a_{(-1)} b$ preserves the filtration (giving the commutative
product in $\operatorname{gr}^F$), $a_{(0)} b$ lowers it
by~$1$ (giving the Poisson bracket), and $a_{(n)} b$ with
$n \geq 1$ lowers it by $n \geq 1$, contributing to
$d_r$ with $r \geq 2$.

Thus $d_1$ on $\operatorname{gr}^F \barBgeom(V)$ is exactly the
component from the $0$-th product, which on the associated
graded is the Poisson bracket
$\{\bar{a}, \bar{b}\} = \overline{a_{(0)} b}$.

For~(i): if $\{-,-\} = 0$, then $d_1 = 0$.
This removes only the first Poisson correction on the
$E_1$ page; higher filtration-lowering OPE terms can still
contribute to later differentials $d_r$ with $r \geq 2$.
For the universal VOA $V_k(\fg)$,
$\operatorname{gr}^F V_k(\fg) \cong
\operatorname{Sym}(\fg[t^{-1}]t^{-1})$ with
Poisson bracket from the Lie bracket of~$\fg$;
here $d_1$ is the Chevalley--Eilenberg differential,
which is nonzero but has concentrated cohomology
by Whitehead vanishing \textup{(}for semisimple~$\fg$\textup{)}.

For~(ii): when $R_V = \mathcal{O}(Y)$ with $Y$ Poisson, the
$E_0$ page is the ordinary commutative bar complex of
$\mathcal{O}(Y)$, and the filtration-$1$ component of the bar
differential induces on
$E_1 = H^*_{\mathrm{bar}}(\mathcal{O}(Y))$ the differential
coming from the Poisson bracket~$\pi$. This is bar-side
Poisson data on the Li--bar spectral sequence, not ordinary
Lichnerowicz--Poisson cohomology.
\end{proof}

\begin{theorem}[Associated-variety criterion for Koszulness;
\ClaimStatusProvedHere]
\label{thm:associated-variety-koszulness}
\index{Koszul property!associated variety criterion}
\index{associated variety!Koszulness criterion}
\index{Li filtration!Koszulness criterion}

Let $V$ be a vertex algebra with Li filtration~$F_p V$,
associated graded $R_V = \operatorname{gr}^F V$, and
associated variety $X_V = \operatorname{Specm}(R_V/\!\sqrt{0})$.
Consider the Li--bar spectral sequence~\eqref{eq:li-bar-ss}.
\begin{enumerate}[label=\textup{(\roman*)}]
\item \emph{Sufficiency from diagonal concentration of the
 Li--bar $E_2$ page.}
 If the Li--bar $E_2$ page of~$R_V$
 is concentrated on the diagonal
 \textup{(}$H^{p,q}(E_1, d_1) = 0$ for
 $p \neq q$\textup{)}, then the Li--bar spectral sequence
 degenerates at~$E_2$ and $V$~is chirally Koszul.
\item \emph{Orbit criterion.}
 Suppose $R_V \cong \mathcal{O}(\overline{\mathbb{O}})$
 for a nilpotent orbit closure
 $\overline{\mathbb{O}} \subset \fg^*$, equipped with
 the Kirillov--Kostant Poisson structure. Then the Li--bar
 $E_1$ page is the commutative bar cohomology of
 $\mathcal{O}(\overline{\mathbb{O}})$, and its $d_1$
 differential is induced by that Poisson bracket. In
 particular, diagonal concentration of the resulting reduced
 Li--bar $E_2 = H(E_1,d_1)$ page is a sufficient condition
 for $V$ to be chirally Koszul.
\end{enumerate}
\end{theorem}

\begin{proof}
For~(i): diagonal concentration of the $E_2$ page
$H^{p,q}(E_1, d_1)$ forces all higher differentials
$d_r \colon E_r^{p,q} \to E_r^{p-r, q+r-1}$ to vanish
for degree reasons: sources and targets lie on different
diagonals. Hence $E_2 = E_\infty$, and the bar cohomology
$H^*(\barBgeom(V))$ inherits diagonal concentration. By
Theorem~\ref{thm:pbw-koszulness-criterion}, this is
Koszulness.

For~(ii): if $R_V \cong \mathcal{O}(\overline{\mathbb{O}})$ as
Poisson algebras, then
Construction~\ref{constr:li-bar-spectral-sequence}
identifies the $E_0$ page with the commutative bar complex of
$\mathcal{O}(\overline{\mathbb{O}})$, so the $E_1$ page is its
ordinary commutative bar cohomology. Then
Proposition~\ref{prop:li-bar-poisson-differential}
identifies the induced $d_1$ differential on that $E_1$ page
with the operator coming from the Kirillov--Kostant Poisson
bracket on the reduced Poisson algebra. Thus the Li--bar
$E_2$ page is exactly the homology of that reduced
bar-cohomology page under the induced $d_1$. Diagonal
concentration of this reduced Li--bar $E_2$ page implies
chiral Koszulness by~(i).
\end{proof}

\begin{remark}[What this theorem does not identify]
\label{rem:li-bar-no-poisson-identification}
\index{Li filtration!not ordinary Poisson cohomology}

Theorem~\ref{thm:associated-variety-koszulness} is a theorem
about the Li--bar spectral sequence and its $d_1$ complex on the
commutative bar construction of~$R_V$. It does \emph{not}
identify that $E_2$ page with ordinary Lichnerowicz--Poisson
cohomology, nor with
$H^*(\widetilde{\mathbb{O}}, \mathcal{O}_{\widetilde{\mathbb{O}}})$
for a symplectic resolution
$\widetilde{\mathbb{O}} \to \overline{\mathbb{O}}$.
Any such comparison would require extra geometric input and is
not part of the proved surface here.
\end{remark}

\begin{corollary}[Minimal-orbit levels under reducedness and
reduced Li--bar concentration; \ClaimStatusConditional]
\label{cor:minimal-orbit-koszul}
\index{minimal orbit!Koszulness}
\index{admissible level!minimal orbit Koszulness}

Let $\fg$ be a simple Lie algebra and $k$ an admissible level
with $X_{L_k(\fg)} = \overline{\mathbb{O}}_{\min}$
\textup{(}the closure of the minimal nilpotent orbit\textup{)}.
Assume in addition that the Li associated graded is reduced:
\[
R_{L_k(\fg)}=\operatorname{gr}^F L_k(\fg)
\;\cong\;
\mathcal{O}(\overline{\mathbb{O}}_{\min})
\]
as Poisson algebras, and assume that the resulting reduced
Li--bar $E_2$ page is diagonally concentrated.
Then $L_k(\fg)$ is chirally Koszul.
\end{corollary}

\begin{proof}
This is exactly the orbit-reduction criterion of
Theorem~\ref{thm:associated-variety-koszulness}(ii)
applied to the reduced Poisson algebra
$\mathcal{O}(\overline{\mathbb{O}}_{\min})$.
\end{proof}

\begin{proposition}[Nilradical obstruction at degenerate
admissible levels; \ClaimStatusProvedHere]
\label{prop:large-orbit-obstruction}
\index{associated variety!nilradical obstruction}
\index{nilpotent cone!Koszulness obstruction}

Let $k$ be a degenerate admissible level for~$\widehat{\fg}$
with $X_{L_k(\fg)} = \overline{\mathbb{O}} \neq \{0\}$.
The associated graded $R_{L_k} = \operatorname{gr}^F L_k$
is a Poisson algebra whose reduced quotient gives
$\mathcal{O}(\overline{\mathbb{O}})$. The Koszulness
obstruction at degenerate levels comes from two sources:
\begin{enumerate}[label=\textup{(\roman*)}]
\item \emph{Reduced Li--bar surface.}
 The reduced quotient
 $\mathcal{O}(\overline{\mathbb{O}})=R_{L_k}^{\mathrm{red}}$
 controls only the reduced Li--bar $E_2$ page furnished by
 Theorem~\ref{thm:associated-variety-koszulness}(ii). Any
 diagonal-concentration statement established on that reduced
 surface is therefore a statement only for the reduced quotient,
 not for the full non-reduced algebra~$R_{L_k}$.
\item \emph{Nilradical obstruction.}
 The nilradical $\mathcal{I} = \sqrt{0} \subset R_{L_k}$
 carries additional bar-cohomological data not visible
 on the associated variety~$X_{L_k}$. When
 $\mathcal{I} \neq 0$, the quotient
 $R_{L_k} \twoheadrightarrow
 \mathcal{O}(\overline{\mathbb{O}})$ forgets tensors in the
 bar complex with at least one nilpotent factor. These extra
 summands can support differentials and off-diagonal classes
 invisible on the reduced quotient, so reduced associated-variety
 data alone do not determine chiral Koszulness of~$L_k(\fg)$.
\end{enumerate}
\end{proposition}

\begin{proof}
For~(i): this is exactly the distinction in
Theorem~\ref{thm:associated-variety-koszulness}(ii) between the
reduced Poisson algebra
$\mathcal{O}(\overline{\mathbb{O}})$ and the full
non-reduced Li graded~$R_{L_k}$.

For~(ii): the quotient map
$R_{L_k} \twoheadrightarrow
\mathcal{O}(\overline{\mathbb{O}})$ induces a surjection of
commutative bar complexes, and its kernel contains bar tensors
with at least one factor in~$\mathcal{I}$. These nilpotent
summands are absent from the reduced quotient, so the reduced
Li--bar $E_2$ page cannot by itself determine the full
bar-cohomological behavior of~$R_{L_k}$.
\end{proof}

\begin{remark}[The orbit hierarchy of Koszulness]
\label{rem:orbit-hierarchy-koszulness}
\index{nilpotent orbit!Koszulness hierarchy}

Theorem~\ref{thm:associated-variety-koszulness} organizes the
Koszulness question for simple quotients into a hierarchy indexed
by nilpotent orbits, but only after the full Li associated
graded~$R_{L_k}$ has been identified. The associated variety
$X_{L_k}=(\operatorname{Specm}R_{L_k})_{\mathrm{red}}$ records only
the reduced geometry, so nilradical contributions remain a
separate obstruction.
For $\fg = \mathfrak{sl}_N$, the orbits are indexed by
partitions $\lambda \vdash N$:

\smallskip
\begin{center}
\begin{tabular}{llll}
\toprule
\textbf{Orbit} & \textbf{Partition} &
 \textbf{Reduced Li--bar surface} & \textbf{Koszulness} \\
\midrule
$\{0\}$ & $(1^N)$ &
 finite-dim.\ reduced quotient & nilradical-dependent \\
$\mathbb{O}_{\min}$ & $(2, 1^{N-2})$ &
 reduced Li--bar concentration assumed &
 conditional on reducedness and reduced
 Li--bar concentration
 \textup{(}Cor~\ref{cor:minimal-orbit-koszul}\textup{)} \\
$\mathbb{O}_{\mathrm{subreg}}$ & $(N{-}1, 1)$ &
 reduced Li--bar concentration open & nilradical-dependent \\
$\mathcal{N}$ & $(N)$ &
 reduced Li--bar concentration open & nilradical-dependent \\
\bottomrule
\end{tabular}
\end{center}

\smallskip\noindent
The reduced orbit geometry can still provide useful input, but
the present theorem does not identify the reduced Li--bar
$E_2$ page with ordinary Poisson cohomology or with
$H^*(\widetilde{\mathbb{O}}, \mathcal{O}_{\widetilde{\mathbb{O}}})$
for a symplectic resolution. What it does show is that any
reduced-level diagonal-concentration statement constrains only the
reduced Li--bar page and leaves nilradical contributions as a
separate obstruction.

A conjectural type-$A$ orbit-duality principle
\textup{(}Conjecture~\textup{\ref{conj:orbit-duality}}\textup{)}
would exchange the associated varieties of $L_k$ and~$L_{k'}$:
$\overline{\mathbb{O}} \leftrightarrow
\overline{d_{\mathrm{BV}}(\mathbb{O})}$. If that principle holds,
it imposes a \emph{compatibility constraint}: the Li--bar
obstructions for~$L_k$ and~$L_{k'}$ should be
Barbasch--Vogan dual.
\end{remark}

\begin{remark}[Comparison with the Kac--Shapovalov criterion]
\label{rem:li-bar-vs-kac-shapovalov}
\index{Li filtration!comparison with Kac--Shapovalov}

Theorem~\ref{thm:associated-variety-koszulness} refines
Theorem~\ref{thm:kac-shapovalov-koszulness} as follows.
The Shapovalov form~$G_h$ detects null vectors weight by weight.
The Li filtration organizes these null vectors \emph{geometrically}:
a null vector at weight~$h$ in the vacuum module corresponds
to a relation in $R_V = \operatorname{gr}^F V$ at Li
filtration level~$p$, and the Poisson bracket determines
whether this relation propagates to the bar complex via~$d_1$.
The Shapovalov criterion asks: are there null vectors in the
bar-relevant range? The Li--bar criterion asks: does the
\emph{geometry} of the associated variety~$X_V$ force those
null vectors to produce off-diagonal bar cohomology?

The gain is maximal when $X_V \neq \{0\}$ and one can
separately establish diagonal concentration on the reduced
Li--bar page, possibly guided by geometric input from a
symplectic resolution. Minimal-orbit levels are the first
test case where the reduced geometry suggests concentration, but
an unconditional Koszul theorem still requires the extra
reducedness and reduced Li--bar concentration hypotheses recorded in
Corollary~\ref{cor:minimal-orbit-koszul}.
\end{remark}

\begin{remark}[Koszul Reflection orientation: the twelve equivalences as properties of $K$]
\label{rem:koszul-reflection-orientation}
\index{Koszul reflection!the twelve equivalences as properties of K}
\index{Theorem A!unified form pulls back the twelve tests}
The twelve characterisations packaged in
Theorem~\ref{thm:koszul-equivalences-meta} are not twelve parallel
statements about $\cA$; they are twelve diagnostics for a single
involutive structure carried by the chiral bar functor. The unified
\emph{Koszul Reflection} of
Theorem~\ref{thm:koszul-reflection}, $K = \Bbarch_X$ with inverse
$K^{-1} = \Omegach_X$, is a symmetric-monoidal adjoint equivalence
of $(\infty,1)$-categories
\[
 K\;\colon\;\Alg^{\mathrm{fact},\mathrm{aug},\mathrm{comp}}_X
 \;\rightleftarrows\;
 \CoAlg^{\mathrm{fact},\mathrm{conil},\mathrm{co}}_X
 \;\colon\;K^{-1},
\]
satisfying $K^{2}\simeq\mathrm{id}$ on the Koszul locus
$\Kosz(X)\subset\Alg^{\mathrm{fact},\mathrm{aug},\mathrm{comp}}_X$ (chain-level
quasi-isomorphism for both unit and counit) and
$K^{2}\simeq\mathrm{id}$ in
$D^{\co}(\Bbarch\text{-}\mathrm{CoFact})$ off
$\Kosz(X)$ (coderived equivalence under the covariantly-weighted
Positselski package). Each entry of
Theorem~\ref{thm:koszul-equivalences-meta} reads as a coordinate
test for the involution:
\begin{itemize}
\item items \textup{(i)}, \textup{(v)} register the unit
 $\eta\colon\mathrm{id}\Rightarrow K^{-1}K$ and counit
 $\varepsilon\colon KK^{-1}\Rightarrow\mathrm{id}$ as
 quasi-isomorphisms;
\item items \textup{(ii)}, \textup{(iii)}, \textup{(x)} register the
 $E_2$-collapse of the bar spectral sequence that promotes the
 coderived involutivity of $K^{2}\simeq\mathrm{id}$ to chain level on
 $\Kosz(X)$;
\item item \textup{(iv)} is the diagonal-vanishing fingerprint of
 $K$ at the Hochschild level
 $\Ext^{p,q}_\cA(\omega_X,\omega_X)$;
\item item \textup{(vi)} is the Barr--Beck--Lurie monadicity test
 on the comodule fibre of $K$;
\item item \textup{(vii)} is the factorization-homology shadow of
 $K^{2}\simeq\mathrm{id}$ over $\mathbb{P}^1$ only after the
 comparison hypotheses of Proposition~\ref{prop:bar-fh}; on the
 uniform-weight lane the same conditional comparison may be imposed
 over every $\Sigma_g$;
\item item \textup{(viii)} is the structural output
 $K^{2}\!\!\restriction\!\Kosz$ on the Hochschild centre tower
 $\ChirHoch^*(\cA)$, with cohomological amplitude $[0,2]$
 reflecting the curve dimension $\dim_\mathbb{C} X = 1$ and the
 $[2]$-Verdier amplitude of the bigraded model;
\item item \textup{(ix)} is the Kac--Shapovalov non-degeneracy
 witness that places $\cA$ on $\Kosz(X)$;
\item item \textup{(xi)} is the conditional Lagrangian-transversality
 upgrade $K^{2}\simeq\mathrm{id}$ at the level of the
 $(-1)$-shifted symplectic deformation space
 $\cM_{\mathrm{comp}}$ (perfectness hypothesis);
\item item \textup{(xii)} is the one-directional purity refinement of
 $K$ as a mixed-Hodge-module morphism along FM strata.
\end{itemize}
The coordinate organisation matters: the symmetry $K \leftrightarrow
K^{-1}$ exchanges the unit and counit clauses, exchanges
$\Kosz(X)$ with the conilpotent locus $\Conil(X)$, and exchanges
$\cA$ with the Verdier-Koszul dual
$\cA^! := \Omegach_X(D_{\Ran}\Bbarch_X(\cA))$ whenever the completed
cobar construction on the Verdier-dual bar coalgebra converges. The
Hochschild clause~\textup{(viii)} is the trace identity
$K^{2}\!\!\restriction\!\Kosz$ reads on the chiral derived centre
$Z^{\mathrm{der}}_{\mathrm{ch}}(\cA) = \ChirHoch^*(\cA,\cA)$; the
factorization-homology clause~\textup{(vii)} is the conditional trace
identity that the same involution reads on $\int_{\Sigma_g}\cA$
after the Ran comparison has been installed. The
Climax of Vol~I, $\dbar = \KZ^*(\nabla_{\Arn})$ together with
$\kappa(\cA) = -c_{\mathrm{ghost}}(\BRST(\cA))$, packages the
coordinate dependence into the universal Arnold connection and the
universal ghost functor (Chapter~\ref{chap:climax-platonic};
Chapter~\ref{chap:universal-conductor}); the twelve coordinate tests
of Theorem~\ref{thm:koszul-equivalences-meta} are the
$\KZ$-pullback shadows of these climax structures along the unique
adjunction whose involutivity is $K^{2}\simeq\mathrm{id}$.
\end{remark}

\begin{remark}[Honest count of the meta-theorem]
\label{rem:koszul-equivalences-meta-scope}
The content of
Theorem~\ref{thm:koszul-equivalences-meta} separates exactly as
$7 + 1 + 1 + 1 + 1 + 1 = 12$ numbered items. Items
\textup{(i)--(v)} and \textup{(ix)--(x)} are seven
\emph{genuinely independent} bidirectional equivalences on the
standard landscape at
\emph{generic} parameters: non-critical level $k \neq -h^\vee$,
non-admissible level for lattice and $\mathcal{W}_N$ algebras,
and non-minimal-model central charge for Virasoro and
$\mathcal{W}$ algebras. Item~\textup{(vi)} (Barr--Beck--Lurie
monadicity) is a proved consequence of item~\textup{(v)}
(bar-cobar quasi-isomorphism) via the standard comonadic
reduction, not an independent equivalence; it is listed because
it is the $\infty$-categorical formulation most useful for
applications, not because it supplies a logically new route.
Item~\textup{(vii)} (factorization homology concentration) is a
conditional comparison consequence: Proposition~\ref{prop:bar-fh}
supplies only a filtered model under locally constant,
constructible, compact Ran hypotheses, and a converse from
factorization homology concentration to Koszulness requires the
extra detection hypothesis stated in the theorem. It is therefore
not counted as an unconditional equivalence.
Item~\textup{(viii)} (chiral Hochschild concentration and
duality) is a proved \emph{one-way} consequence on the Koszul
locus; the converse from~\textup{(viii)} to the bar-cobar counit
is not proved here. Item~\textup{(xi)} (the Lagrangian criterion)
is conditional on perfectness of the bar-cobar normal complex and
non-degeneracy of the ambient tangent complex. On the standard
landscape these hypotheses are discharged by
Proposition~\ref{prop:lagrangian-perfectness}; outside it they remain
part of the statement. Item~\textup{(xii)}
($\mathcal{D}$-module purity) is a one-directional
implication~\textup{(xii)}$\Rightarrow$\textup{(x)}; the converse
is proved for affine Kac--Moody at non-critical level only after the
localization-weight package of
Proposition~\ref{prop:d-module-purity-km}; it is open for Virasoro
and $\mathcal{W}$ algebras.
At special parameter values some equivalences require individual
verification or fail; in particular the critical level
$k = -h^\vee$ for affine Kac--Moody is excluded from the generic
scope, where $\kappa$ vanishes and the bar cohomology reshapes to
the opers complex. A separately stated characterization,
tropical Koszulness
(Theorem~\ref{thm:tropical-koszulness} in
Chapter~\ref{chap:higher-genus}), provides a further
bidirectional criterion via a weight spectral sequence on the
tropical real blowup; it is proved there as a standalone theorem
but is \emph{not} integrated into the unified list of the present
meta-theorem.
\end{remark}

\begin{theorem}[Equivalences and consequences of chiral Koszulness; \ClaimStatusConditional]
\label{thm:koszul-equivalences-meta}
\index{Koszul property!equivalences of characterizations|textbf}

Let $\cA$ be a chiral algebra on a smooth projective curve~$X$
with PBW filtration $F_\bullet$. The intrinsic equivalence core is
read on the split finite-window PBW lane: conformal-weight/bar-length
pieces are finite, the induced filtrations are Hausdorff, complete, and
strongly convergent, and a conformal-weight splitting identifies
$E_\infty$ with its associated graded. Outside this lane the theorem
records only the implications whose stated hypotheses supply the
missing convergence or comparison input.
For a smooth projective curve~$\Sigma$, write
$\mathsf{FH}_\Sigma(\cA)$ for the comparison package of
Proposition~\ref{prop:bar-fh}: the Ran object attached to~$\cA$ is
locally constant on configuration strata, constructible on
Fulton--MacPherson compactifications, compact enough for the Ran
colimit to commute with the bar-length filtration, and the resulting
comparison
\[
\barB^{\mathrm{geom}}_\Sigma(\cA)
\;\simeq\;
\textstyle\int_\Sigma \overline{\cA}^{\mathrm{Lie}}
\]
is filtered for bar length versus the Goodwillie filtration and
strongly convergent; in condition~\textup{(vii)} this right-hand
side is abbreviated by $\int_\Sigma\cA$. Write
$\mathsf{FH}^{\mathrm{det}}_\Sigma(\cA)$ for the detecting
strengthening in which off-diagonal PBW classes in the bar model
cannot be killed by cancellation after passage through the
Goodwillie realization; equivalently, degree-zero concentration of
$\int_\Sigma\cA$ reflects degree-zero concentration of
$\barB^{\mathrm{geom}}_\Sigma(\cA)$.
These are extra hypotheses, not consequences of chiral
Koszulness.
Conditions \textup{(i)--(v)} and \textup{(ix)--(x)} below are seven
genuinely independent bidirectional equivalences; condition~\textup{(vi)}
(Barr--Beck--Lurie comparison) is a proved consequence of
condition~\textup{(v)} rather than an independent route, and is listed
for completeness.
Condition~\textup{(vii)} is a conditional factorization-homology
consequence under $\mathsf{FH}_{\mathbb{P}^1}(\cA)$; it becomes a
genus-$0$ criterion only under the stronger
$\mathsf{FH}^{\mathrm{det}}_{\mathbb{P}^1}(\cA)$.
On the uniform-weight lane, condition~\textup{(vii)} strengthens to
the all-genera concentration statement only under
$\mathsf{FH}_{\Sigma_g}(\cA)$ for the relevant curves.
Condition~\textup{(viii)} is a proved one-way consequence of
the intrinsic core on the Koszul locus: it yields chiral Hochschild duality,
cohomological amplitude~$[0,2]$, and polynomial
Hilbert series. Any free-polynomial description of the cup-product
algebra is conditional on a Massey-vanishing consequence of
Fulton--MacPherson formality.
Under the additional perfectness and non-degeneracy hypotheses on the
ambient tangent complex, condition~\textup{(xi)} is also equivalent to
the intrinsic core. Condition~\textup{(xii)} implies condition~\textup{(x)}
(Remark~\ref{rem:d-module-purity-content}); the converse is open.

\smallskip
\noindent\textbf{Intrinsic equivalence core and monadic consequence:}
\begin{enumerate}
\item[\textup{(i)}] $\cA$ is chirally Koszul
 \textup{(}Definition~\textup{\ref{def:chiral-koszul-morphism})}.
\item[\textup{(ii)}] The PBW spectral sequence on
 $\barBgeom(\cA)$ collapses at $E_2$
 \textup{(}Theorem~\textup{\ref{thm:pbw-koszulness-criterion})}.
\item[\textup{(iii)}] The minimal $A_\infty$-model of
 $\barBgeom(\cA)$ is formal: $m_n = 0$ for $n \geq 3$
 \textup{(}Proposition~\textup{\ref{prop:ainfty-formality-implies-koszul}}
 and Theorem~\textup{\ref{thm:ainfty-koszul-characterization})}.
\item[\textup{(iv)}] Ext diagonal vanishing:
 $\operatorname{Ext}^{p,q}_\cA(\omega_X, \omega_X) = 0$ for
 $p \neq q$.
\item[\textup{(v)}] The bar-cobar counit
 $\Omega(\barBgeom(\cA)) \to \cA$ is a
 quasi-isomorphism
 \textup{(}Theorem~\textup{\ref{thm:koszul-reflection}},
 reconstruction clause\textup{)}. This is the intrinsic counit of
 the bar-cobar adjunction, not the factorization-cohomology model of
 Proposition~\textup{\ref{prop:cobar-fh}}.
\item[\textup{(vi)}] The Barr--Beck--Lurie comparison functor attached
 to the bar-cobar adjunction is an equivalence after restricting to the
 Koszul bar-cobar component generated by~$\cA$.
\item[\textup{(ix)}] The Kac--Shapovalov determinant
 $\det G_h \neq 0$ in the bar-relevant range
 \textup{(}Theorem~\textup{\ref{thm:kac-shapovalov-koszulness})}.
\item[\textup{(x)}] The restriction
 $i_S^!\,\barB_n(\cA)$ to every FM boundary stratum~$S$ is
 acyclic outside degree~$0$.
\end{enumerate}

\smallskip
\noindent\textbf{Conditional factorization-homology comparison:}
\begin{enumerate}
\item[\textup{(vii)}] If
 $\mathsf{FH}_{\mathbb{P}^1}(\cA)$ holds, then chiral Koszulness
 implies genus-$0$ factorization homology concentration:
 \[
 H^k\!\left(\textstyle\int_{\mathbb{P}^1}\cA\right)=0
 \qquad (k \neq 0).
 \]
 If the detecting strengthening
 $\mathsf{FH}^{\mathrm{det}}_{\mathbb{P}^1}(\cA)$ also holds, this
 genus-$0$ concentration statement is equivalent to
 conditions~\textup{(i)--(v)} and \textup{(ix)--(x)}. If $\cA$
 lies on the uniform-weight lane and
 $\mathsf{FH}_{\Sigma_g}(\cA)$ holds for every smooth projective
 curve $\Sigma_g$ of genus $g \geq 0$, then chiral Koszulness
 implies
 \[
 H^k\!\left(\textstyle\int_{\Sigma_g}\cA\right)=0
 \qquad (k \neq 0)
 \]
 for every $g \geq 0$, and for $g \geq 1$ the surviving scalar
 obstruction class is
 \[
 \mathrm{obs}_g(\cA)=\kappa(\cA)\cdot\lambda_g
 \qquad \textup{(}UNIFORM-WEIGHT\textup{)}
 \]
 \textup{(UNIFORM-WEIGHT; Theorem~\textup{\ref{thm:genus-universality})}.}
 No converse from factorization homology concentration to
 Koszulness is asserted without
 $\mathsf{FH}^{\mathrm{det}}_{\mathbb{P}^1}(\cA)$.
\end{enumerate}

\smallskip
\noindent\textbf{One-way chiral Hochschild consequence on the Koszul locus:}
\begin{enumerate}
\item[\textup{(viii)}] $\mathrm{ChirHoch}^*(\cA)$ has cohomological
 amplitude $[0,2]$ and satisfies the duality
 \[
 \ChirHoch^n(\cA)
 \cong \ChirHoch^{2-n}(\cA^!)^\vee \otimes \omega_X
 \]
 with Hochschild--Hilbert series
 \[
 P_\cA(t)=\dim Z(\cA)+\dim \ChirHoch^1(\cA)\cdot t
 + \dim Z(\cA^!)\cdot t^2,
 \]
 and any free-polynomial description of the graded-commutative
 cup-product algebra is conditional on a Massey-vanishing
 consequence of Fulton--MacPherson formality
 \textup{(}Theorem~H,
 Theorem~\textup{\ref{thm:main-koszul-hoch}},
 Theorem~\textup{\ref{thm:hochschild-polynomial-growth}},
 cf.\ Proposition~\textup{\ref{prop:e2-formality-hochschild}}\textup{)}.
\end{enumerate}

\smallskip
\noindent\textbf{Conditional and partial:}
\begin{enumerate}
\item[\textup{(xi)}] The Lagrangian criterion: under perfectness,
 non-degeneracy, and the tangent-complex comparison identifying the
 completed twisted tensor product with the derived intersection,
 $\cM_\cA$ and $\cM_{\cA^!}$ are transverse Lagrangians in the
 $(-1)$-shifted symplectic deformation space~$\cM_{\mathrm{comp}}$.
\item[\textup{(xii)}] $\mathcal{D}$-module purity: each
 $\barBgeom_n(\cA)$ is pure of weight~$n$ as a mixed Hodge module,
 with characteristic variety aligned to FM boundary strata.
\end{enumerate}
\end{theorem}

\begin{proof}
\textsc{The core circuit}
\textup{(i)}$\Leftrightarrow$\textup{(ii)}$\Leftrightarrow$%
\textup{(iii)}$\Leftrightarrow$\textup{(v)}:

\smallskip\noindent
\textup{(i)}$\Leftrightarrow$\textup{(ii)}:
Koszulness (Definition~\ref{def:chiral-koszul-morphism}) is
acyclicity of $\barBgeom(\cA) \otimes_\tau \cA$ in positive degrees.
Theorem~\ref{thm:pbw-koszulness-criterion} identifies this with
$E_2$-collapse; the converse is immediate.

\smallskip\noindent
\textup{(ii)}$\Rightarrow$\textup{(v)}:
$E_2$-collapse gives bar concentration
(Theorem~\ref{thm:bar-concentration}), which is the input for
bar-cobar inversion (Theorem~\ref{thm:bar-cobar-inversion-qi}).

\smallskip\noindent
\textup{(v)}$\Rightarrow$\textup{(i)}:
Condition~\textup{(v)} is the statement that the canonical counit
\[
\varepsilon_{\tau_{\mathrm{univ}}}\colon
\Omega_X(\barB_X(\cA)) \longrightarrow \cA
\]
of the universal twisting morphism
$\tau_{\mathrm{univ}}\colon \barB_X(\cA)\to \cA$
\textup{(}Corollary~\ref{cor:three-bijections}\textup{)}
is a quasi-isomorphism.
The fundamental theorem of chiral twisting morphisms
\textup{(}Theorem~\ref{thm:fundamental-twisting-morphisms}\textup{)}
identifies this directly with
$\tau_{\mathrm{univ}}$ being a chiral Koszul morphism.
With the PBW filtration fixed in the statement of the theorem,
that is exactly condition~\textup{(i)}.
No factorization-cohomology comparison is used in this implication;
Proposition~\ref{prop:cobar-fh} gives an external completed model
only under its own constructibility and Mittag--Leffler hypotheses.

\medskip
\textsc{chiral Hochschild consequence}
\textup{(v)}$\Rightarrow$\textup{(viii)}:

\smallskip\noindent
The quasi-isomorphism
$\Omega(\barBgeom(\cA)) \xrightarrow{\sim} \cA$ identifies the
bar-cobar resolution with the chiral Hochschild complex. On the
Koszul locus, Theorem~H,
Theorem~\ref{thm:main-koszul-hoch},
Theorem~\ref{thm:hochschild-polynomial-growth}, and
Proposition~\ref{prop:e2-formality-hochschild} give the stated
cohomological concentration, duality, and polynomial Hilbert
series. Proposition~\ref{prop:e2-formality-hochschild} isolates the
Fulton--MacPherson formality input governing higher braces and
Massey products. Any further free-polynomial refinement of the
graded-commutative cup-product algebra requires that extra
Massey-vanishing package and is not part of the unconditional
equivalence core.
No converse from
\textup{(viii)} to the bar-cobar counit is proved here.

\medskip
\textsc{$A_\infty$ formality}
\textup{(i)}$\Leftrightarrow$\textup{(iii)}:

\smallskip\noindent
\textup{(ii)}$\Rightarrow$\textup{(iii)}:
$E_2$-collapse means all PBW differentials $d_r = 0$ for $r \geq 2$.
The $A_\infty$ products $m_n$ are identified with $d_{n-1}$ under
the HPL transfer, so $m_n = 0$ for $n \geq 3$
(Proposition~\ref{prop:ainfty-formality-implies-koszul}).

\smallskip\noindent
\textup{(iii)}$\Rightarrow$\textup{(ii)}:
If $m_n = 0$ for $n \geq 3$, the minimal $A_\infty$-model is
strictly associative: $H^*(\barBgeom(\cA))$ is a dg algebra with
$m_2$ as the product and all higher operations trivial. The
bar spectral sequence of this formal $A_\infty$-algebra collapses
at~$E_2$ by the classicality theorem for formal $A_\infty$-algebras
(Keller). The chiral $A_\infty$ structure is computed fiberwise on
each FM stratum (Proposition~\ref{prop:shadow-formality-low-degree});
on each fiber the HPL transfer is an ordinary $A_\infty$ transfer
over a field, to which Keller's theorem applies directly. The PBW
filtration is defined fiberwise and compatible with the FM
stratification, so fiberwise $E_2$-collapse assembles to global
$E_2$-collapse: all transferred differentials $d_r = 0$ for
$r \geq 2$. Hence the spectral sequence
on~$\barBgeom(\cA)$ also collapses at~$E_2$.

\medskip
\textsc{Conditional factorization-homology comparison}
\textup{(i)}$\Rightarrow$\textup{(vii)} under
$\mathsf{FH}_{\mathbb{P}^1}$, with a converse only under
$\mathsf{FH}^{\mathrm{det}}_{\mathbb{P}^1}$
\textup{(}the uniform-weight all-genera statement is a conditional
refinement. Cf.\ Remark~\ref{rem:fh-vii-uniform-weight-scope}\textup{)}:

\smallskip\noindent
\textup{(i)}$\Rightarrow$\textup{(vii)}:
The intrinsic statement is bar-model concentration. By
\textup{(i)}$\Leftrightarrow$\textup{(ii)}, the PBW spectral
sequence of $\barBgeom(\cA)$ collapses at~$E_2$, hence
Theorem~\ref{thm:bar-concentration} places the bar model in
cohomological degree~$0$. If
$\mathsf{FH}_{\mathbb{P}^1}(\cA)$ holds, then
Proposition~\ref{prop:bar-fh} identifies this filtered bar model
with the Goodwillie-filtered factorization homology object over
$\mathbb{P}^1$. Strong convergence of the comparison transfers
degree-zero concentration to $\int_{\mathbb{P}^1}\cA$, giving
$H^k(\int_{\mathbb{P}^1}\cA)=0$ for $k \neq 0$.

If $\cA$ lies on the uniform-weight lane and
$\mathsf{FH}_{\Sigma_g}(\cA)$ holds for every $\Sigma_g$, the same
filtered comparison is available at each genus. The loop-order
spectral sequence of the higher-genus bar complex collapses on the
modular Koszul locus
(Theorem~\ref{thm:loop-order-collapse}), so after the vertexwise
genus-$0$ PBW collapse every loop-order stratum remains
cohomologically concentrated in degree~$0$. Theorem~\ref{thm:genus-universality}
then identifies the surviving scalar obstruction class by
\[
\mathrm{obs}_g(\cA)=\kappa(\cA)\cdot\lambda_g
\qquad (g \geq 1)
\]
\textup{(UNIFORM-WEIGHT)}.
Because no mixed-weight channels exist on that lane, the
cross-channel correction $\delta F_g^{\mathrm{cross}}$ is absent,
so no positive-degree cohomology reappears. Hence
$H^k(\int_{\Sigma_g}\cA)=0$ for $k \neq 0$ for all $g \geq 0$ on
the uniform-weight lane.

\smallskip\noindent
\textup{(vii)}$\Rightarrow$\textup{(i)}:
Suppose $H^k(\int_{\mathbb{P}^1}\cA) = 0$ for $k \neq 0$ and assume
$\mathsf{FH}^{\mathrm{det}}_{\mathbb{P}^1}(\cA)$. The detected
comparison reflects degree-zero concentration back to the bar model:
an off-diagonal class on the PBW $E_2$ page would survive as an
off-degree Goodwillie class, contradicting concentration of
$\int_{\mathbb{P}^1}\cA$. Hence the PBW spectral sequence collapses
at~$E_2$, and Theorem~\ref{thm:pbw-koszulness-criterion} gives
chiral Koszulness.

Without $\mathsf{FH}^{\mathrm{det}}_{\mathbb{P}^1}(\cA)$, the
external factorization homology object may have cancellations in its
Goodwillie realization that are invisible to the PBW page. Therefore
factorization homology concentration alone is not used here to prove
Koszulness.

\begin{remark}[All-genera form of condition~\textup{(vii)}]
\label{rem:fh-vii-uniform-weight-scope}
Condition~\textup{(vii)} has an intrinsic bar-model theorem and a
conditional factorization-homology avatar. The unconditional
bar-model statement is
\[
H^k\!\bigl(|\barBgeom_{\mathbb{P}^1}(\cA)|\bigr)=0
\qquad (k \neq 0)
\quad\Longleftrightarrow\quad
\cA\in\Kosz(X),
\]
by Theorem~\ref{thm:bar-concentration} and
Theorem~\ref{thm:pbw-koszulness-criterion}. The replacement of this
bar model by $\int_{\mathbb{P}^1}\cA$ requires
$\mathsf{FH}_{\mathbb{P}^1}(\cA)$, and the converse requires
$\mathsf{FH}^{\mathrm{det}}_{\mathbb{P}^1}(\cA)$.

On the uniform-weight lane the higher-genus refinement is proved.
If $\mathsf{FH}_{\Sigma_g}(\cA)$ holds, Proposition~\ref{prop:bar-fh}
supplies the bar comparison model for factorization homology on the
curve~$\Sigma_g$. On the Koszul locus, the genus-$0$ PBW collapse
forces each vertex contribution to be concentrated in degree~$0$, and
Theorem~\ref{thm:loop-order-collapse} prevents the loop-order
spectral sequence from reintroducing positive-degree cohomology.
Definition~\ref{def:scalar-lane} leaves only one channel, so
mixed-channel boundary graphs do not occur. The surviving scalar
class is then identified by Theorem~\ref{thm:genus-universality}:
\[
\mathrm{obs}_g(\cA)=\kappa(\cA)\cdot\lambda_g
\qquad \textup{(UNIFORM-WEIGHT)}.
\]
Hence
\[
H^k\!\bigl(\textstyle\int_{\Sigma_g}\cA\bigr)=0
\qquad (k \neq 0)
\]
for every $g \geq 0$ on the uniform-weight lane. This is the
all-genera form of condition~\textup{(vii)}.

The scalar identity in that argument has an independent proof path.
Remark~\ref{rem:theorem-d-alt-grr} applies
Grothendieck--Riemann--Roch to the family bar complex over the
universal curve and recovers the same class
$\kappa(\cA)\lambda_g$ without using the shadow tower, the
Maurer--Cartan element, or
Theorem~\ref{thm:genus-universality}. This does not replace the
conditional bar/factorization-homology concentration argument above,
but it does give a second proof of the scalar part of the
all-genera uniform-weight statement.

For multi-weight algebras, the correct higher-genus replacement of
the second sentence of condition~\textup{(vii)} is not another
concentration equivalence. The correct all-genera reformulation
passes to the scalar projection of the genus-$g$ obstruction
package: one replaces higher-genus concentration by the
diagonal-plus-cross-channel identity
\[
F_g(\cA)
\;=\;
\kappa(\cA)\cdot\lambda_g^{\mathrm{FP}}
\;+\;
\delta F_g^{\mathrm{cross}}(\cA)
\qquad
\textup{(ALL-WEIGHT, with cross-channel correction
$\delta F_g^{\mathrm{cross}}$)}
\]
for $g \geq 1$, with
$\delta F_1^{\mathrm{cross}}=0$ universally and
$\delta F_g^{\mathrm{cross}}=0$ on the uniform-weight lane
\textup{(}Theorem~\ref{thm:multi-weight-genus-expansion}\textup{)}.
In this form, condition~\textup{(vii)} survives at all genera as a
statement about the scalar trace rather than about full
cohomological concentration.
The uniform-weight hypothesis therefore cannot be removed in
general: the interacting algebra~$\cW_3$ leaves the scalar locus
already at $g=2$, while
Proposition~\ref{prop:free-field-scalar-exact} identifies the
free-field exception where the correction vanishes genus by genus.

Geometrically, $\delta F_g^{\mathrm{cross}}$ is the scalar trace of
mixed-channel boundary residues. The edges still carry the universal
weight-$1$ propagator $d\log E(z,w)$; the correction comes from the
vertices, where OPE residues couple distinct conformal-weight
sectors on stable-graph boundary strata. Construction~\ref{constr:cross-channel-graph-sum}
makes this precise: $\delta F_g^{\mathrm{cross}}$ is the sum over
stable graphs whose edge-labelling
$\sigma\colon E(\Gamma)\to\{1,\ldots,r\}$ uses at least two
weights. At genus~$1$ the unique boundary graph has one edge,
hence no mixed assignment and no correction. At genus~$g \geq 2$,
multi-edge strata admit genuine mixed assignments, and their
integrated OPE residues produce
$\delta F_g^{\mathrm{cross}}$. In this precise sense, the
cross-channel term is the higher-genus shadow of mixed-weight OPE
residues.
\end{remark}

\begin{remark}[The content of \textup{(viii)}]
Theorem~\ref{thm:main-koszul-hoch} and
Theorem~\ref{thm:hochschild-polynomial-growth} prove duality,
cohomological amplitude~$[0,2]$, and
polynomial Hilbert growth on the Koszul locus.
The interval $[0,2]$ is the curve-level chiral Hochschild amplitude:
it uses the one-dimensional Fulton--MacPherson boundary calculus and
the Verdier shift on a curve. It is not a statement about product
ambients, categorical Hochschild cochains of Calabi--Yau categories, or
topological Hochschild theory. Product or external tensor ambients
require a K\"unneth comparison and may have larger total amplitude.
A further identification of $\ChirHoch^*(\cA)$ as a free
polynomial algebra is conditional on a Massey-vanishing
consequence of Fulton--MacPherson formality and is not used in
the equivalence core.
\end{remark}

\medskip
\textsc{FM boundary acyclicity}
\textup{(i)}$\Leftrightarrow$\textup{(x)}:

\smallskip\noindent
\textup{(i)}$\Rightarrow$\textup{(x)}:
Each FM boundary stratum $S_T \cong \prod_{v \in V(T)}
\overline{\operatorname{Conf}}_{|v|}(X_v)$ is indexed by a tree~$T$.
The residue component $d_{\mathrm{res}}|_{S_T}$ of the bar
differential restricts to the tensor product of lower-degree bar
complexes. $E_2$-collapse at each vertex degree
(Theorem~\ref{thm:pbw-koszulness-criterion}) gives acyclicity:
$H^k(i_{S_T}^!\,\barB_n(\cA)) = 0$ for $k \neq 0$.

\smallskip\noindent
\textup{(x)}$\Rightarrow$\textup{(i)}:
At the deepest stratum $S = D_{\{1,2\}} \cong X \times
\overline{\operatorname{Conf}}_{n-1}(X)$ (binary collision), the
restriction $i_S^!\,\barB_n(\cA)$ computes the bar-complex
contribution from the OPE $a_{(k)}b$ at a single collision point.
Acyclicity of $i_S^!$ at all strata forces the residues at every
collision to be exact, which is the PBW condition:
$d_1(e_i) = 0$ on the associated graded, hence $E_2$-collapse.

\medskip
\textsc{Barr--Beck--Lurie monadicity}
\textup{(v)}$\Rightarrow$\textup{(vi)} and
\textup{(vi)}$\Rightarrow$\textup{(i)}:

\smallskip\noindent
The Barr--Beck--Lurie theorem (\cite[Theorem~4.7.3.5]{HA}) states
that the comparison functor $\Phi$ is an equivalence if and only if
$\barB_\kappa$ is conservative and preserves totalizations of
$\barB_\kappa$-split cosimplicial objects.
Conservativity: $\barB_\kappa$ detects quasi-isomorphisms on the
Koszul locus by bar-cobar inversion~(v).
Totalization preservation: on the Koszul bar-cobar component generated
by~$\cA$, the bar-cobar
adjunction is a Quillen equivalence
(Theorem~\ref{thm:quillen-equivalence-chiral}). The derived
adjunction on the underlying $\infty$-categories satisfies the
BBL conditions because the Quillen-equivalence comparison functor
is an equivalence restricted to the Koszul locus; this is a
formal consequence of Quillen equivalences satisfying monadic
descent (\cite[Proposition~4.7.3.16]{HA}).
Therefore $\Phi$ is an equivalence if and only if~(v) holds,
i.e., if and only if $\cA$ is chirally Koszul.

\smallskip\noindent
\textup{(vi)}$\Rightarrow$\textup{(i)}:
If the BBL comparison functor $\Phi$ is an equivalence, then
$\barB_\kappa$ is conservative (detecting quasi-isomorphisms).
Conservativity of the bar functor forces bar-cobar inversion,
which is~(v), which is~(i).

\medskip
\textsc{Kac--Shapovalov determinant}
\textup{(i)}$\Leftrightarrow$\textup{(ix)}:

\smallskip\noindent
\textup{(i)}$\Rightarrow$\textup{(ix)}:
If $\cA$ is chirally Koszul, the PBW filtration on the vacuum
module is strict: $\operatorname{gr}_F M_0(\cA) \cong
\operatorname{Sym}^{\mathrm{ch}}(V)$. Strictness implies the
Shapovalov form $G_h$ is non-degenerate in the bar-relevant range.

\smallskip\noindent
\textup{(ix)}$\Rightarrow$\textup{(i)}:
Non-degeneracy $\det G_h \neq 0$ in the bar-relevant range
means the PBW-to-bar comparison map
$\iota \colon \operatorname{Sym}^{\mathrm{ch}}(V)_h \to
\barB_h(\cA)$ is injective at every bar-relevant weight~$h$.
Suppose $d_r \neq 0$ for some $r \geq 2$. Then there exists
$x \in E_r^{p,q}$ with $d_r(x) \neq 0$, represented by
$\tilde{x} \in F^p\barB$ with $d_{\mathrm{bar}}(\tilde{x})
\in F^{p+r}\barB$. If $\tilde{x} \in \operatorname{im}(\iota)$,
it represents a nonzero class in
$\operatorname{gr}^p(\barB) \cong \operatorname{Sym}^p(V)$,
and $d_r(\tilde{x})$ represents a nonzero class in
$\operatorname{gr}^{p+r}$. By PBW strictness (injectivity
of~$\iota$ at all bar-relevant weights), this contradicts the
acyclicity of $\operatorname{Sym}(V)$ as a Koszul complex.
Hence $d_r = 0$, giving~(ii) and hence~(i).

\medskip
\textsc{Ext diagonal vanishing}
\textup{(i)}$\Leftrightarrow$\textup{(iv)}:

\smallskip\noindent
\textup{(i)}$\Rightarrow$\textup{(iv)}:
$E_2$-collapse places $H^*(\barBgeom(\cA))$ on the diagonal
$p = q$; the bar resolution computes Ext, so
$\operatorname{Ext}^{p,q}(\omega_X, \omega_X) = 0$ for $p \neq q$.

\smallskip\noindent
\textup{(iv)}$\Rightarrow$\textup{(i)}:
Diagonal Ext-vanishing means $\barBgeom(\cA)$ has cohomology
concentrated on the line $p = q$ in the bigrading. The PBW
spectral sequence has $E_1^{p,q} = H^{p+q}(\operatorname{gr}_F^p
\barBgeom)$ converging to $H^{p+q}(\barBgeom)$; diagonal
concentration of the abutment forces all differentials
$d_r$ ($r \geq 2$) to vanish (a $d_r$-differential moves
off-diagonal by $(r, 1-r)$, which would produce off-diagonal
classes contradicting~(iv)). The PBW spectral sequence converges
conditionally in the sense of Boardman: the filtration is Hausdorff
and complete (the PBW filtration on a positively graded algebra is
bounded below). In characteristic zero with a split filtration (the
PBW filtration splits by conformal weight), the extension problem is
trivial: $E_\infty = E_2$ as bigraded objects, not merely as
filtered objects. Hence $E_2$-collapse, hence~(i).

\medskip
\textsc{Lagrangian criterion}
\textup{(i)}$\Leftrightarrow$\textup{(xi)}
\emph{(only under perfectness, non-degeneracy, and tangent-complex
comparison; the standard landscape uses
Proposition~\textup{\ref{prop:lagrangian-perfectness}} where its
hypotheses have been verified)}:

\smallskip\noindent
\textup{(i)}$\Rightarrow$\textup{(xi)}:
On the Koszul locus, the bar-cobar adjunction provides a free
resolution. Theorem~C supplies the Verdier-side homotopy eigenspace
decomposition
$\mathbf{Q}_g(\cA) \oplus \mathbf{Q}_g(\cA^!)
\simeq H^*(\overline{\cM}_g, \cZ_\cA)$
(not a negation of MC elements and not a scalar identity beyond its
trace projection). Under perfectness and non-degeneracy this
decomposition identifies $\cM_\cA$ and $\cM_{\cA^!}$ as complementary
subspaces of $\cM_{\mathrm{comp}}$. The shifted-symplectic structure
(Theorem~\ref{thm:ambient-complementarity-fmp})
makes these subspaces isotropic; the complementarity splitting
gives maximal dimension, hence Lagrangian.

\smallskip\noindent
\textup{(xi)}$\Rightarrow$\textup{(i)}:
If $\cM_\cA$ and $\cM_{\cA^!}$ are transverse Lagrangians and the
tangent-complex comparison identifies their derived intersection with
the completed twisted tensor product,
their derived intersection
$\cM_\cA \times^h_{\cM_{\mathrm{comp}}} \cM_{\cA^!}$
is discrete (transverse Lagrangian intersection in a $(-1)$-shifted
symplectic space has expected dimension~$0$). Under the comparison
hypothesis this derived intersection computes the twisted tensor product
$K_\tau(\cA, \cA^!)$; its acyclicity (i.e., discreteness of the
intersection) is Koszulness
(Definition~\ref{def:chiral-koszul-morphism}).
\end{proof}

\begin{remark}[Proof web redundancy]
\label{rem:koszul-proof-web-redundancy}
The proof of Theorem~\ref{thm:koszul-equivalences-meta} is not a
single chain.
Its unconditional core already has several proof lanes that do not
route through the same intermediate condition:
\[
\textup{(i)} \Longleftrightarrow \textup{(v)}
\quad\text{(universal twisting morphism and bar filtration),}\qquad
\textup{(i)} \Longleftrightarrow \textup{(iii)}
\quad\text{($A_\infty$ transfer and Keller classicality),}
\]
\[
\textup{(i)} \Longleftrightarrow \textup{(iv)}
\quad\text{(Ext diagonal vanishing),}
\]
\[
\textup{(i)} \Longleftrightarrow \textup{(ix)}
\quad\text{(Kac--Shapovalov non-degeneracy),}\qquad
\textup{(i)} \Longleftrightarrow \textup{(x)}
\quad\text{(FM boundary acyclicity),}
\]
and on the monadic side
\[
\textup{(v)} \Rightarrow \textup{(vi)} \Rightarrow \textup{(i)}
\quad\text{(Barr--Beck--Lurie monadicity).}
\]

Three direct cross-links are the load-bearing redundancy.

\smallskip\noindent
\textup{(v)}$\Rightarrow$\textup{(i)} is the shortest loop.
For the universal twisting morphism
$\tau_{\mathrm{univ}}\colon \barB_X(\cA)\to \cA$,
condition~\textup{(v)} is the quasi-isomorphism of the counit
$\varepsilon_{\tau_{\mathrm{univ}}}$, and
Theorem~\ref{thm:fundamental-twisting-morphisms}
identifies that directly with
$\tau_{\mathrm{univ}}$ being a chiral Koszul morphism.
This uses the universal bar twisting datum and the bar filtration,
not the route through \textup{(ii)} or \textup{(iii)}.

\smallskip\noindent
\textup{(i)}$\Rightarrow$\textup{(viii)} is also direct.
Theorem~\ref{thm:main-koszul-hoch},
Theorem~\ref{thm:hochschild-polynomial-growth}, and
Proposition~\ref{prop:e2-formality-hochschild}
start from a chiral Koszul datum and compute
$\ChirHoch^*(\cA)$ by the bar-cobar Hochschild resolution.
This yields duality, cohomological amplitude~$[0,2]$, and
polynomial Hilbert growth. Any free-polynomial upgrade of the
cup-product algebra is conditional on the Massey-vanishing
Fulton--MacPherson formality package, so it is not routed through
\textup{(ii)}--\textup{(vii)} inside the meta-theorem.

\smallskip\noindent
\textup{(iii)}$\Rightarrow$\textup{(x)} is fiberwise and geometric.
For a boundary stratum $S_T$, the restriction
$i_{S_T}^!\,\barB_n(\cA)$ decomposes as the tensor product of the
vertexwise collision complexes, and
Remark~\ref{rem:iterated-residues-ainfty}
identifies those iterated residues with the transferred
$A_\infty$ operations.
If $m_k = 0$ for $k \geq 3$, only the strict binary residue survives
at each vertex.
Fiberwise Keller classicality, exactly as in
Theorem~\ref{thm:ainfty-koszul-characterization},
then gives degree-zero concentration for every vertex factor, so
the K\"unneth spectral sequence yields
\[
H^k\bigl(i_{S_T}^!\,\barB_n(\cA)\bigr)=0
\qquad (k \neq 0).
\]
This reaches \textup{(x)} without passing through the shadow tower.

\smallskip\noindent
The same argument gives further direct arrows already implicit in
the proof:
\[
\textup{(iv)} \Rightarrow \textup{(ii)},\qquad
\textup{(ix)} \Rightarrow \textup{(ii)},\qquad
\textup{(v)} \Longleftrightarrow \textup{(vi)}.
\]
Under $\mathsf{FH}_{\mathbb{P}^1}$ the arrow
\textup{(ii)}$\Rightarrow$\textup{(vii)} may be added to this web;
under $\mathsf{FH}^{\mathrm{det}}_{\mathbb{P}^1}$ the converse
arrow may also be added. These conditional arrows are deliberately
not displayed in the unconditional core.
Hence the unconditional core
\[
\{\textup{(i)},\textup{(ii)},\textup{(iii)},\textup{(iv)},
\textup{(v)},\textup{(vi)},\textup{(ix)},
\textup{(x)}\}
\]
forms a web rather than a chain.

\begin{center}
\begin{tikzcd}[column sep=3em,row sep=2.8em]
\textup{(iv)} \arrow[r] \arrow[dr] &
\textup{(ii)} \arrow[d, shift left=.5ex] \arrow[r, shift left=.5ex] &
\textup{(v)} \arrow[l, shift left=.5ex] \arrow[dl] \arrow[dr] &
\\
&
\textup{(i)} \arrow[r, shift left=.5ex] \arrow[d, shift left=.5ex] \arrow[dr] \arrow[rr, bend left=12] &
\textup{(iii)} \arrow[l, shift left=.5ex] \arrow[dl] &
\textup{(vi)} \arrow[ll, bend left=12]
\\
\textup{(ix)} \arrow[uur, bend left=16] \arrow[r] &
\textup{(x)} \arrow[u, shift left=.5ex] \arrow[ur] &
\textup{(viii)} &
\end{tikzcd}
\end{center}

Here \textup{(viii)} is a satellite consequence, not part of the
equivalence core.
The node \textup{(vii)} has been omitted because it is a conditional
factorization-homology comparison lane. The uniform-weight all-genera
statement recorded in
Theorem~\ref{thm:koszul-equivalences-meta}\textup{(vii)} is a
conditional strengthening under the relevant
$\mathsf{FH}_{\Sigma_g}$ hypotheses, not an extra unconditional
equivalence.
Let $G_{\mathrm{core}}$ be the undirected graph obtained from the
displayed core subweb by forgetting arrow directions and deleting the
satellite node~\textup{(viii)}.
Every displayed core edge lies on one of the cycles
\[
(\textup{(i)},\textup{(ii)},\textup{(v)}),\qquad
(\textup{(i)},\textup{(ii)},\textup{(iv)}),
\]
\[
(\textup{(i)},\textup{(iii)},\textup{(x)}),\qquad
(\textup{(i)},\textup{(v)},\textup{(vi)}),\qquad
(\textup{(i)},\textup{(ii)},\textup{(ix)},\textup{(x)}).
\]
Therefore $G_{\mathrm{core}}$ has no bridges.
Equivalently, deleting any one displayed implication still leaves the
unconditional core connected: the implication graph has no single
point of failure.
\end{remark}

\begin{remark}[Lagrangian and purity refinements]
Condition~\textup{(xi)} remains conditional on perfectness and
non-degeneracy of the ambient tangent complex. Condition~\textup{(xii)}
implies condition~\textup{(x)} (via Saito's theory of mixed Hodge
modules; see Remark~\textup{\ref{rem:d-module-purity-content}}).
The converse direction (Koszulness implies D-module purity) is
reduced to a single gap
(Remark~\ref{rem:d-module-purity-content}).
The converse of~\textup{(xii)} is proved for the affine Kac--Moody
lineage under the localization-weight package
(Proposition~\ref{prop:d-module-purity-km}); the non-affine case
(Virasoro, $\cW$-algebras) remains open.
Beyond these twelve conditions, Sklyanin--Poisson rigidity
(Theorem~\ref{thm:koszulness-from-sklyanin}) provides a conditional
Poisson-geometric route to Koszulness restricted to the affine
Kac--Moody subclass and to the Sklyanin--bar comparison package.
\end{remark}

\begin{proposition}[Closure of chiral Koszulness under
tensor, dualization, and base change;
\ClaimStatusConditional]
\label{prop:koszul-closure-properties}
\index{Koszul property!closure under operations|textbf}

The class of chirally Koszul algebras is closed under the following
operations when the stated finiteness and comparison hypotheses hold:
\begin{enumerate}[label=\textup{(\alph*)}]
\item \emph{Chiral tensor product.}
 If $\cA, \cB$ are chirally Koszul on the curve $X$ and satisfy the
 split mixed-residue hypotheses of
 Proposition~\textup{\ref{prop:koszul-dual-tensor-product}}, then
 $\cA \boxtimes \cB$ is chirally Koszul, and its bar-dual coalgebra is
 $\cA^{\mathrm i}\boxtimes\cB^{\mathrm i}$.
\item \emph{Koszul dualization} \textup{(}involutivity\textup{)}.
 If $\cA$ is chirally Koszul and the finite-type Verdier dual algebra
 $\cA^!$ exists with convergent dual bar construction, then $\cA^!$ is
 chirally Koszul and $(\cA^!)^! \simeq \cA$ canonically.
\item \emph{Smooth base change.}
 If $f \colon Y \to X$ is a smooth morphism of curves and
 $\cA$ is chirally Koszul on $X$, with holonomicity, finite-type
 duality, and PBW convergence preserved by $f^*$, then $f^*\cA$ is
 chirally Koszul on $Y$, and
 $(f^*\cA)^! \simeq f^*(\cA^!)$.
\end{enumerate}
\end{proposition}

\begin{proof}
\textup{(a)} is Proposition~\ref{prop:koszul-dual-tensor-product}.
The K\"unneth formula applies only after the mixed-color hypotheses in
that proposition are imposed; it then identifies the colored bar
complex with $\barB(\cA)\boxtimes\barB(\cB)$ and gives
$(\cA \boxtimes \cB)^{\mathrm i}\simeq
\cA^{\mathrm i}\boxtimes\cB^{\mathrm i}$. Finite duality gives the
algebra statement when the duals exist.

\textup{(b)}: Involutivity is a consequence of
Theorem~\ref{thm:fundamental-twisting-morphisms} (the bar-cobar
adjunction is a Quillen equivalence on the Koszul locus) and the
biduality of holonomic $\mathcal D$-modules under Verdier duality.
The bar-cobar counit reconstructs $\cA$ from its bar-dual coalgebra,
while Verdier duality identifies the finite dual algebra
$\cA^!$ with $(\cA^{\mathrm i})^\vee$. Applying the same construction
to $\cA^!$ and using Verdier biduality gives
$(\cA^!)^!\simeq\cA$. This is the chiral finite-type analogue of the
standard quadratic fact $(A^!)^!\simeq A$
(Loday--Vallette~\cite{LV12}, Theorem~3.4.6).

\textup{(c)}: Smooth base change uses
Lemma~\ref{lem:pushforward-preserves-qi} to verify that
$f^*$ preserves quasi-isomorphisms of chiral algebras, combined
with Theorem~\ref{thm:bar-cobar-inversion-functorial} to verify
that $f^*$ commutes with bar-cobar. The witnessing datum for
chiral Koszulness pulls back along $f$, and the
$E_2$-collapse of the PBW spectral sequence is preserved under
flat pullback (since smooth morphisms are flat). The duality
identification $(f^*\cA)^! \simeq f^*(\cA^!)$ follows from the
compatibility of $f^*$ with Verdier duality on the finite-type
holonomic lane.
\end{proof}

\begin{remark}[Closure does not extend to bare quotients]
\label{rem:koszul-closure-not-quotients}
The closure properties of
Proposition~\ref{prop:koszul-closure-properties} do \emph{not} extend
to general quotients $\cA / \cI$ by chiral ideals: the quotient of a
chirally Koszul algebra by a non-Koszul ideal is generically non-Koszul.
The W-algebra construction $\cW^k(\fg, f)$ via Drinfeld--Sokolov
reduction is the principal example where the quotient \emph{does}
preserve Koszulness, and this is a substantive theorem of the chapter
(Theorem~\ref{thm:w-algebra-koszul}), not a corollary of closure.
\end{remark}

\begin{remark}[$A_\infty$ formality and Gaiotto--Kulp--Wu]
\label{rem:ainfty-formality-gkw}
\index{Gaiotto--Kulp--Wu!formality characterization}
The $A_\infty$ formality characterization
(item~(iii) of Theorem~\ref{thm:koszul-equivalences-meta})
is the bar-complex statement of the formality theorem of
Gaiotto--Kulp--Wu~\cite{GKW2025}: the transferred operations
$m_k = 0$ for $k \geq 3$ on bar cohomology if and only if
$\cA$ is chirally Koszul. The shadow depth classification
$\mathbf{G}/\mathbf{L}/\mathbf{C}/\mathbf{M}$
(Remark~\ref{rem:shadow-depth-gkw-refinement}) refines their
binary formal/non-formal dichotomy into a four-class stratification
governed by the discriminant $\Delta = 8\kappa S_4$.
\end{remark}

\begin{remark}[Loop-exactness ordering $G < L < C < M$]
\label{rem:loop-exactness-ordering}
\index{shadow depth!loop-exactness ordering}
The four shadow-depth classes are ordered by the number of loop corrections needed to determine the full $A_\infty$ structure on~$\cA^!_{\mathrm{line}}$:
\begin{itemize}
\item Class~$G$ (Gaussian, $r_{\max}=2$): tree-level exact.
\item Class~$L$ (Lie/tree, $r_{\max}=3$): one-loop exact (DNP non-renormalization, Theorem~\ref{thm:non-renormalization-tree}(iii)).
\item Class~$C$ (Contact, $r_{\max}=4$): two-loop exact.
\item Class~$M$ (Mixed, $r_{\max}=\infty$): requires all-loop resummation.
\end{itemize}
This ordering refines the generic Koszul locus: the standard
generic/universal representatives in all four classes are chirally
Koszul (bar $E_2$-collapse), but they differ in Swiss-cheese formality
depth.
\end{remark}

\begin{proposition}[Swiss-cheese non-formality by shadow class; \ClaimStatusConditional]
\label{prop:swiss-cheese-nonformality-by-class}
\index{Swiss-cheese operad!non-formality by shadow class|textbf}
\index{shadow depth!Swiss-cheese formality}
\index{Ainfty@$A_\infty$!non-formality for class M}
The Swiss-cheese operations $m_k^{\mathrm{SC}}$ on~$\cA$ itself
\textup{(}not on the bar cohomology $H^*(\barB(\cA))$, which is
always $A_\infty$-formal for Koszul algebras by
Theorem~\textup{\ref{thm:koszul-equivalences-meta}(iii)}\textup{)}
are classified by shadow depth as follows.
\begin{center}
\renewcommand{\arraystretch}{1.3}
\begin{tabular}{@{}lcccc@{}}
\toprule
\textbf{Class}
& \textbf{Families}
& $m_3^{\mathrm{SC}}$
& $m_4^{\mathrm{SC}}$
& $m_k^{\mathrm{SC}},\; k \geq 5$ \\
\midrule
$G$ \textup{(}$r_{\max}=2$\textup{)}
& Heisenberg, lattice VOA, free fermion
& $= 0$
& $= 0$
& $= 0$ \\
$L$ \textup{(}$r_{\max}=3$\textup{)}
& affine KM
& $\neq 0$
& $= 0$
& $= 0$ \\
$C$ \textup{(}$r_{\max}=4$\textup{)}
& $\beta\gamma$
& $= 0$
& $\neq 0$
& $= 0$ \\
$M$ \textup{(}$r_{\max}=\infty$\textup{)}
& Virasoro, $\cW_N$
& $\neq 0$
& $\neq 0$
& nonzero for infinitely many $k \geq 5$
 \textup{(}Theorem~\ref{thm:shadow-archetype-classification}\textup{)} \\
\bottomrule
\end{tabular}
\end{center}

\medskip\noindent
\emph{Mechanism.}
For class~$G$\textup{:}
Heisenberg, lattice VOAs, and free fermions are Gaussian/free-field:
their ordered bar data is generated by the binary two-point kernel,
so every connected genus-$0$ tree with at least three external legs
factors through pairings and carries no primitive higher vertex.
Equivalently, the shadow tower truncates at degree~$2$:
$S_r(\cA)=0$ for $r \geq 3$
\textup{(}Theorem~\ref{thm:shadow-archetype-classification}; for
lattices see Corollary~\ref{cor:lattice-postnikov-termination}\textup{)}.
Under the averaging map
$\operatorname{av}(\Theta_\cA^{\Eone}) = \Theta_\cA$
\textup{(}Theorem~\ref{thm:e1-primacy}\textup{)},
the mixed Swiss-cheese tree of degree~$r$ projects to the degree-$r$
shadow. Since there is no primitive tree for $r \geq 3$,
all higher mixed-sector operations vanish.
For class~$L$\textup{:} the affine KM OPE has at most a double pole; the
three-channel tree sum over $\overline{\cM}_{0,4}$ is nonzero
\textup{(}$S_3 = 2h^\vee/(k+h^\vee) \neq 0$\textup{)} but the
tower terminates at depth~$3$ because $S_r = 0$ for $r \geq 4$.
For class~$C$\textup{:}
the cubic shadow vanishes by weight parity of the $\gamma$
generator \textup{(}conformal weight~$0$\textup{)}, while the
quartic contact invariant
$Q^{\mathrm{contact}} \neq 0$ from the
$\beta\gamma\beta\gamma$ self-contraction; rank-one rigidity of
the contact stratum kills all $m_k^{\mathrm{SC}}$ for $k \geq 5$.
For class~$M$\textup{:}
the quartic pole in the Virasoro OPE
\textup{(}$T_{(3)}T = c/2$\textup{)} produces $S_3 = 2$
\textup{(}$c$-independent\textup{)} and the critical discriminant
$\Delta = 8\kappa S_4 \neq 0$ forces all higher operations to be
nonzero.
\end{proposition}

\begin{proof}
Class~$G$: for Heisenberg, centrality of the bracket kills all nested
tree compositions, so $m_k^{\mathrm{SC}} = 0$ for $k \geq 3$.
For lattice VOAs, bosonization reduces connected genus-$0$ amplitudes
to Wick pairings of the underlying Heisenberg field, and
Corollary~\ref{cor:lattice-postnikov-termination} identifies the
shadow tower with its weight-$2$ truncation. For the free fermion,
Wick's theorem gives Pfaffian factorization by the basic contraction
$\psi(z)\psi^*(w) \sim (z-w)^{-1}$, so again no connected tree with
$r \geq 3$ survives after extracting the binary pairings. In all three
cases the mixed Swiss-cheese degree-$r$ operation is the open-colored
lift of the same genus-$0$ tree sum whose closed projection is
$S_r(\cA)$; equivalently,
$\operatorname{av}(\Theta_\cA^{\Eone}) = \Theta_\cA$
\textup{(}Theorem~\ref{thm:e1-primacy}\textup{)}.
Because class~$G$ means $S_r(\cA) = 0$ for all $r \geq 3$,
the mixed-sector operations vanish:
$m_k^{\mathrm{SC}} = 0$ for $k \geq 3$.
The weaker condition $\Delta = 0$ is not sufficient:
affine Kac--Moody algebras are class~$L$ with $\Delta = 0$ but
$m_3^{\mathrm{SC}} \neq 0$.
Class~$L$: the cubic shadow
$S_3 = 2h^\vee/(k + h^\vee)$ is computed from the OPE coefficient
ratio (the structure constant divided by the level); the quartic
and higher shadows vanish because the Jacobi identity and rank
constraints eliminate all $\overline{\cM}_{0,r+1}$ tree sums for
$r \geq 4$ when the OPE has at most double poles.
Class~$C$: the cubic shadow vanishes by the parity obstruction on
the weight-$0$ generator~$\gamma$; the quartic contact invariant
$Q^{\mathrm{contact}}$ is nonzero by the explicit
$\beta\gamma\beta\gamma$ channel computation; termination at
depth~$4$ follows from rank-one rigidity
(Theorem~\ref{thm:riccati-algebraicity}).
Class~$M$:
$S_3 = (\text{scalar projection of } T_{(1)}T)/T_{(3)}T
= 2\kappa/\kappa = 2$, independent of~$c$;
$S_4 = 10/[c(5c{+}22)]$ and $\Delta = 8\kappa S_4 = 40/(5c{+}22) \neq 0$
for generic~$c$, so the Riccati algebraicity theorem
forces the tower to be infinite.
Computational verification:
\texttt{theorem\_ainfty\_nonformality\_class\_m\_engine.py}
(three independent methods for $S_3$,
quartic shadow $S_4$ and $Q^{\mathrm{contact}}$, quintic
shadow $S_5 = -48/(c^2(5c+22))$, all four classes verified).
\end{proof}

The four classes exhaust the standard landscape, but the
relationship between SC-formality and the classification is
sharper than the table suggests. Proposition~\ref{prop:swiss-cheese-nonformality-by-class}
already isolates the first nonzero Swiss-cheese operation in each
non-Gaussian class and, on the Gaussian locus, identifies the
mixed-sector tree formulas with the degree-by-degree shadow tower
after averaging. The criterion below packages this as an equivalence
between SC-formality and truncation of the shadow tower at degree~$2$.

\begin{proposition}[SC-formality characterises class~$G$; \ClaimStatusConditional]
\label{prop:sc-formal-iff-class-g}
\index{Swiss-cheese operad!formality characterisation}
\index{class G@class~$G$!characterised by SC-formality}
Let~$\cA$ be a chiral algebra in the standard landscape.
Then~$\cA$ is Swiss-cheese formal
\textup{(}$m_k^{\mathrm{SC}} = 0$ for all $k \geq 3$\textup{)}
if and only if~$\cA$ belongs to class~$G$.
\end{proposition}

\begin{proof}
The proof is operadic: both directions use the genus-$0$
tree-shadow correspondence
\textup{(}Theorem~\ref{thm:shadow-formality-identification}\textup{)}
and the averaging identity
$\operatorname{av}(\Theta_\cA^{\Eone}) = \Theta_\cA$
\textup{(}Theorem~\ref{thm:e1-primacy}\textup{)}.
No invariant bilinear form or metric is required;
the argument applies uniformly to all families in the standard
landscape, including $\beta\gamma$ systems.

\medskip
\textsc{Forward} \textup{(}class~$G$ $\Longrightarrow$ SC-formal\textup{)}:

\smallskip\noindent
Class~$G$ is the Gaussian locus
\textup{(}Theorem~\ref{thm:shadow-archetype-classification}\textup{)}:
Heisenberg, lattice VOA, and free fermion. For these families,
the shadow tower truncates at degree~$2$:
$S_r(\cA) = 0$ for every $r \geq 3$.

The genus-$0$ operadic transfer
\textup{(}Theorem~\ref{thm:shadow-formality-identification};
Proposition~\ref{prop:shadow-tower-three-lenses}(ii)\textup{)}
identifies $S_r(\cA)$ with the transferred $L_\infty$ bracket
$\ell_r^{(0),\mathrm{tr}}$ evaluated on the truncated MC element.
The Swiss-cheese mixed-sector operation $m_r^{\mathrm{SC}}$ of
degree~$r$ is the open-colour lift of the \emph{same} genus-$0$
tree sum: the mixed tree and the shadow tree share their internal
edge propagators and vertex operations; they differ only in the
output colour (open vs closed), with passage to the closed sector
given by the averaging morphism~$\operatorname{av}$.

Since class~$G$ means the genus-$0$ transfer is generated by the
binary two-point kernel alone, every connected tree of degree
$r \geq 3$ factors through binary pairings and carries no
primitive higher vertex. The tree sum vanishes at the open level
(no connected tree survives), so
$m_k^{\mathrm{SC}} = 0$ for all $k \geq 3$.

For completeness, we record the family-by-family mechanism.
Heisenberg: centrality of the bracket kills all nested tree
compositions.
Lattice VOA: bosonization reduces connected genus-$0$ amplitudes
to Wick pairings of the underlying Heisenberg field
\textup{(}Corollary~\ref{cor:lattice-postnikov-termination}\textup{)}.
Free fermion: Pfaffian factorization by the basic contraction
$\psi(z)\psi^*(w) \sim (z{-}w)^{-1}$ ensures no connected tree
with $r \geq 3$ external legs survives after extracting binary
pairings. In all three cases the mechanism is Gaussian: only the
two-point kernel generates.

\medskip
\textsc{Converse} \textup{(}SC-formal $\Longrightarrow$ class~$G$\textup{)}:

\smallskip\noindent
Suppose~$\cA$ is SC-formal, so
$m_k^{\mathrm{SC}} = 0$ for all $k \geq 3$.

\emph{Step~1: SC-formality forces $S_r(\cA) = 0$ for $r \geq 3$.}
By the tree-shadow correspondence, the degree-$r$ shadow
$S_r(\cA) = \operatorname{Sh}_r(\cA)$ is the closed-colour
projection of the same genus-$0$ transferred tree that defines
the mixed-sector operation $m_r^{\mathrm{SC}}$.
Precisely: the averaging identity
$\operatorname{av}(\Theta_\cA^{\Eone}) = \Theta_\cA$
\textup{(}Theorem~\ref{thm:e1-primacy}(ii)\textup{)}
sends the degree-$r$ component of the $\Eone$ MC element
to the degree-$r$ shadow. On the $\Eone$ side, this component
is the same transferred tree operation that, before averaging,
determines the mixed-sector Swiss-cheese operation.

The logical chain is:
\begin{equation}
\label{eq:sc-formal-shadow-vanishing}
m_r^{\mathrm{SC}} = 0
\;\;\Longrightarrow\;\;
\text{degree-}r\text{ tree} = 0
\;\;\Longrightarrow\;\;
S_r(\cA) = \operatorname{av}(\text{degree-}r\text{ tree}) = 0.
\end{equation}
The first implication is the two-colour transfer hypothesis: the
bulk-to-boundary structure map is injective on connected genus-$0$ trees
of positive degree, so vanishing of the open-colour output forces
vanishing of the underlying ordered tree. The second implication is the
averaging identity: a zero tree has zero $\Sigma_n$-coinvariant
projection.

\begin{remark}[Two-colour transfer hypothesis]
\label{rem:sc-formal-open-transfer-frontier}
The implication ``open-colour output vanishes $\Longrightarrow$
underlying $\Eone$ tree vanishes'' is used here as an operational
lemma supported by the injectivity of the bulk-to-boundary
structure map on connected genus-$0$ trees of positive degree.
It is the analogue of Theorem~\ref{thm:e1-primacy} in the
ordered-to-mixed direction rather than the ordered-to-symmetric
direction. Theorem~\ref{thm:e1-primacy}(iv) establishes that
$\ker(\operatorname{av})$ is non-trivial (contains the Drinfeld
associator at degree~$3$), so the transfer from mixed-sector vanishing
to ordered-tree vanishing requires an injectivity statement genuinely
different from the averaging map. The converse direction of
Proposition~\ref{prop:sc-formal-iff-class-g} is conditional on this
two-colour transfer lemma; the forward direction is unconditional.
\end{remark}

\emph{Step~2: $S_r = 0$ for all $r \geq 3$ is precisely class~$G$.}
The discriminant $\Delta = 0$ is necessary but not sufficient:
class~$L$ has $\Delta = 0$ but $S_3 \neq 0$.
The full vanishing criterion $S_r = 0$ for all $r \geq 3$
identifies exactly class~$G$ among the four standard classes:
\begin{center}
\small
\renewcommand{\arraystretch}{1.15}
\begin{tabular}{@{}lcccl@{}}
\toprule
\textbf{Class} & $S_3$ & $S_4$ & $S_r\;(r \geq 5)$
& \textbf{Ruled out by} \\
\midrule
$G$ & $=0$ & $=0$ & $=0$ &; \\
$L$ & $\neq 0$ & $=0$ & $=0$ & $S_3 \neq 0$ \\
$C$ & $=0$ & $\neq 0$ & $=0$ & $S_4 \neq 0$ \\
$M$ & $\neq 0$ & $\neq 0$ & nonzero for infinitely many $r$ & $S_3 \neq 0$ \\
\bottomrule
\end{tabular}
\end{center}
By Theorem~\ref{thm:shadow-archetype-classification}, the
standard-landscape locus with $S_r = 0$ for all $r \geq 3$ is
precisely class~$G$. By Step~1, an SC-formal algebra has
this vanishing pattern; hence it belongs to class~$G$.
\end{proof}

\begin{remark}[Independence from invariant bilinear forms]
\label{rem:sc-formal-no-metric}
\index{Swiss-cheese operad!metric-independence of formality proof}
The proof of
Proposition~\ref{prop:sc-formal-iff-class-g} is purely
operadic: it uses the tree-shadow correspondence
\textup{(}Theorem~\ref{thm:shadow-formality-identification}\textup{)},
the averaging identity
\textup{(}Theorem~\ref{thm:e1-primacy}\textup{)}, and the shadow
archetype classification
\textup{(}Theorem~\ref{thm:shadow-archetype-classification}\textup{)}.
No invariant bilinear form, Killing metric, or Casimir tensor
enters at any step. This is essential because the class-$G$ locus
includes the free fermion and lattice VOAs, and the class-$C$
family $\beta\gamma$ (which must be ruled out by the converse)
admits no invariant metric and no Sugawara construction.
The operadic structure of the genus-$0$ tree transfer, not the
representation theory of any Lie algebra, is the correct
mechanism.
\end{remark}

\begin{remark}[Livernet--LV comparison]
\label{rem:livernet-lv-comparison}
\index{Swiss-cheese operad!comparison with Livernet and Loday--Vallette}
The Swiss-cheese statements in
Propositions~\ref{prop:swiss-cheese-nonformality-by-class}
and~\ref{prop:sc-formal-iff-class-g}
concern transferred mixed operations on a fixed chiral algebra~$\cA$.
They should be distinguished from operadic Koszul duality of the
two-colored operad itself.
The comparison with the classical literature is color by color.

The convention check is decisive.
In the dual calculation of
Proposition~\ref{prop:sc-koszul-dual-three-sectors},
the closed color of $\SCchtop$ is first replaced by its algebraic
commutative shadow.
The resulting quadratic model has closed color $\operatorname{Com}$,
open color $\operatorname{Ass}$, and mixed sector governed by
shuffles.
This is the same distributive-law package as in the classical
Swiss-cheese literature around Livernet, made explicit in
Loday--Vallette
\cite[\S7.1, Chapter~8, \S13.3]{LV12}.
Accordingly,
\[
(\SCchtop)^! \cong (\mathrm{Lie},\, \mathrm{Ass},\, \text{shuffle-mixed}),
\]
with closed dimensions $(n{-}1)!$, open dimensions $m!$, and mixed
dimensions $(k{-}1)!\binom{k+m}{m}$
\textup{(}Proposition~\ref{prop:sc-koszul-dual-three-sectors}\textup{)}.
Operadic self-duality already fails in the classical algebraic model:
the closed $\operatorname{Com}$ sector is one-dimensional in each
input degree, whereas the dual closed $\operatorname{Lie}$ sector has
dimension $(n{-}1)!$.
Hence the duality functor is involutive, but the operad is not
self-dual.

What changes in the chiral--topological setting is the proof method.
Before passage to the commutative shadow, the closed color is
represented by the dg operad $C_\bullet(\FM_k(\CC))$.
One therefore does not prove strict quadratic Koszulity directly on
chains.
Instead one transports the classical Koszul statement along
Kontsevich formality for the closed color together with the bar-cobar
transfer theorems of Loday--Vallette
\cite[Theorems~11.3.3 and~11.4.1]{LV12}.
In this sense the genus-zero comparison is compatible with the
Livernet--LV lane:
classical Swiss-cheese provides the strict Koszul input, while
$\SCchtop$ enters only through the weaker homotopy-Koszul
consequence.
\end{remark}

\begin{remark}[D-module purity: reduction to a single gap]
\label{rem:d-module-purity-content}
\index{D-module purity!reduction to Saito weight}
\index{Saito weight filtration!and PBW filtration}
The converse direction of condition~\textup{(xii)} (Koszulness
$\Longrightarrow$ D-module purity of~$\barBgeom_n(\cA)$) reduces to a
weight-comparison problem. The first three steps are formal or
PBW-local; the fourth is the family-specific proof obligation.

\begin{enumerate}[label=\textup{Step \arabic*.}]
\item
The PBW filtration $F_\bullet^{\mathrm{PBW}}$ on
$\barBgeom_n(\cA)$ is a filtration by $\cD$-submodules on
$\FM_n(X)$, with associated graded given by the FG bar complex
(the bar of the chiral Lie algebra).
\textup{Proved:} the PBW filtration is defined by order of OPE
pole, compatible with $\cD$-module structure.

\item
On the Koszul locus, $\mathrm{gr}^F_\bullet \barBgeom_n$ is
a direct sum of rank-one $\cD$-modules supported on the
strata of the FM boundary, each pure of weight determined
by the stratum codimension.
\textup{Proved:} follows from PBW concentration
\textup{(ii)} and the explicit identification of strata.

\item
Purity of the associated graded implies purity of the
filtered module if and only if the PBW filtration
coincides with the Saito weight filtration from mixed
Hodge module theory.
\textup{Proved:} this is a formal consequence of Saito's
strictness theorem for filtered $\cD$-modules.

\item
\textbf{The PBW filtration on $\barBgeom_n(\cA)$ equals
the Saito weight filtration from mixed Hodge modules on
$\FM_n(X)$.}
For affine Kac--Moody algebras this is the localization-weight
comparison: chiral localization must identify $\barBgeom_n$ with the
corresponding mixed Hodge module on the Beilinson--Drinfeld
Grassmannian and identify the PBW filtration with the
Saito--Kashiwara weight filtration. For Virasoro and $\cW$-algebras it
is open: the BPZ equations do not presently supply a polarizable
variation of Hodge structure controlling the bar filtration.

\item
Given steps~1--4, purity of $\barBgeom_n$ in weight~$n$
follows by the strictness of the weight filtration.
\textup{Proved:} formal.
\end{enumerate}

\noindent
Thus the exact proof obligation is the equality
$F_\bullet^{\mathrm{PBW}}=W_\bullet^{\mathrm{Saito}}$ on the bar
mixed Hodge module. Once this equality is known, Koszulness forces
purity by strictness. Without it, PBW collapse proves bar
concentration but not purity. The affine Kac--Moody lane is treated
below under the localization-weight package; the non-affine lineage
(Virasoro, $\cW$-algebras) remains open.
\end{remark}

\begin{proposition}[$\cD$-module purity for affine Kac--Moody
under localization weights]
\label{prop:d-module-purity-km}
\ClaimStatusProvedHere
\index{D-module purity!affine Kac--Moody|textbf}
For $\cA = V_k(\fg)$ at generic non-critical level~$k$, assume the
localization-weight package: Frenkel--Gaitsgory chiral localization
identifies each $\barBgeom_n(\cA)$ with the corresponding
Beilinson--Drinfeld Grassmannian mixed Hodge module, the Hecke
operations realizing the $n$ insertions are proper on the relevant
finite windows, and the PBW filtration agrees with the
Saito--Kashiwara weight filtration there. Then
$\barBgeom_n(\cA)$ is pure of weight~$n$ as a mixed Hodge
module on $\FM_n(X)$.
\end{proposition}

\begin{proof}
Under the localization-weight package, $\barBgeom_n(\cA)$ is represented
by a mixed Hodge module on the Beilinson--Drinfeld Grassmannian whose
restriction to each finite Hecke window is pure before the $n$ bar
insertions. Proper Hecke pushforward preserves purity, and the
Saito--Kashiwara weight comparison identifies the PBW degree of each
insertion with one unit of Saito weight. Therefore the $n$-fold bar
piece is pure of weight~$n$ on the corresponding FM stratum, and the
proper descent to $\FM_n(X)$ preserves that weight.
\end{proof}

\subsection*{D-module purity: reduction and family dependence}
\label{subsec:d-module-purity-reduction}

Under the localization-weight package the Kac--Moody case strengthens
to a full equivalence. For other families, the two-way implication
remains conditional or open.

\begin{proposition}[Kac--Moody equivalence under
Saito--Kashiwara weight comparison;
\ClaimStatusConditional]
\label{prop:d-module-purity-km-equivalence}
\index{D-module purity!Kac--Moody equivalence|textbf}
\index{Saito weight filtration!PBW comparison}
\index{Kashiwara filtration!Kac--Moody}
Let $\cA = V_k(\fg)$ be the universal affine Kac--Moody
vertex algebra at generic level~$k$ on a smooth projective
curve~$X$. Assume the localization-weight package of
Proposition~\ref{prop:d-module-purity-km}. Then the following are
equivalent:
\begin{enumerate}[label=\textup{(\alph*)}]
\item $\cA$ is chirally Koszul
 \textup{(}condition~\textup{(i)} of
 Theorem~\textup{\ref{thm:koszul-equivalences-meta}}\textup{)}.
\item The PBW filtration $F_\bullet^{\mathrm{PBW}}$ on each
 $\barBgeom_n(\cA)$ is strict and concentrated in
 PBW-degree~$n$
 \textup{(}condition~\textup{(ii)}\textup{)}.
\item Each $\barBgeom_n(\cA)$ is pure of weight~$n$ as a
 mixed Hodge module on~$\FM_n(X)$
 \textup{(}condition~\textup{(xii)}\textup{)}.
\end{enumerate}
\end{proposition}

\begin{proof}
The implications \textup{(a)}$\Leftrightarrow$\textup{(b)}
are part of the core circuit
\textup{(i)}$\Leftrightarrow$\textup{(ii)} of
Theorem~\ref{thm:koszul-equivalences-meta}.
The implication \textup{(c)}$\Rightarrow$\textup{(b)}
is the one-directional part of the meta-theorem
(\textup{(xii)}$\Rightarrow$\textup{(x)}): purity
of~$\barBgeom_n$ forces Saito strictness on the weight
filtration, and Saito strictness gives FM stratum
acyclicity, which in turn forces PBW strictness by the
Kac--Shapovalov implication
\textup{(ix)}$\Rightarrow$\textup{(ii)}.

The new content is the converse \textup{(b)}$\Rightarrow$\textup{(c)}
under the stated localization-weight package. Proposition~\ref{prop:d-module-purity-km}
supplies \textup{(a)}$\Rightarrow$\textup{(c)} after that package is
installed. The converse uses the same Saito--Kashiwara weight
comparison: the PBW filtration on~$V_k(\fg)$ is identified with the
Saito weight filtration on the associated mixed Hodge module after
chiral localization to the affine Grassmannian.

First, Frenkel--Gaitsgory chiral localization identifies
$\barBgeom_n(V_k(\fg))$ with a $\cD$-module on
$\mathrm{Gr}_G^{(n)}$ (the $n$-th Beilinson--Drinfeld
Grassmannian), carrying a canonical variation of mixed
Hodge structure on the chosen finite windows.
Second, the PBW filtration on the vacuum module lifts to
a Hodge-theoretic filtration on this variation; the required
identification is the Saito--Kashiwara comparison supplied by the
hypothesis, in the same circle as Mochizuki's proof of Saito's
decomposition theorem for pure twistor $\cD$-modules
\textup{(}\cite{Mochizuki2015,Sabbah2014}\textup{)},
where the Kashiwara filtration along a boundary divisor
(here, the FM stratification) gives a weight filtration
that, in the presence of a polarization, agrees with the
Saito weight filtration by strictness.
Third, for Kac--Moody the polarization is supplied by the
Kac--Shapovalov form
\textup{(}Theorem~\textup{\ref{thm:kac-shapovalov-koszulness}}\textup{):}
$G_h$ is non-degenerate in the bar-relevant range at
generic level, providing the positivity data that forces
Saito strictness. PBW strictness~\textup{(b)} is then
transported, by the identification above, to strictness
of the Saito weight filtration, which is equivalent
to~\textup{(c)} by Saito's theorem.

The three ingredients (Frenkel--Gaitsgory localization,
Mochizuki--Sabbah twistor strictness, Kac--Shapovalov
positivity) are each family-specific: each relies on
properties of~$V_k(\fg)$ that are not shared by Virasoro
or~$\beta\gamma$. The non-affine case is precisely the
open direction recorded in
Remark~\ref{rem:d-module-purity-content}, Step~4.
\end{proof}

\begin{remark}[Family dependence of the
$\cD$-module purity equivalence]
\label{rem:d-module-purity-family-dependence}
\index{D-module purity!family dependence|textbf}
\index{Virasoro algebra!D-module purity}
\index{beta-gamma system@$\beta\gamma$ system!D-module purity}
\index{W-algebra!principal Drinfeld--Sokolov reduction}
The Kac--Moody equivalence of
Proposition~\ref{prop:d-module-purity-km-equivalence} is not available
for the full standard landscape without additional Hodge-theoretic
input. The missing input is family-specific and is visible in the
shadow class of the algebra.

\emph{Virasoro} \textup{(}class~$\mathbf{M}$,
$r_{\max} = \infty$\textup{):}
$\mathrm{Vir}_c$ is chirally Koszul at generic~$c$
(Table~\ref{tab:koszulness-landscape}), so condition~(ii)
holds. The shadow depth is infinite: the shadow tower
does not terminate, and the higher Massey products
$S_r \neq 0$ for infinitely many~$r$ obstruct the naive assignment of
one Saito weight per PBW bar insertion. The BPZ differential equations
do not carry a known polarization of twistor type, so Mochizuki's
strictness theorem does not apply. Thus PBW concentration is proved,
but $\cD$-module purity is an extra open statement rather than a
consequence of condition~(ii).
This is the \emph{shadow depth~$\infty$ vs.\ PBW
concentration} gap.

\emph{$\beta\gamma$} \textup{(}class~$\mathbf{C}$,
$r_{\max} = 4$\textup{):}
The contact shadow $Q^{\mathrm{contact}} \neq 0$ at degree~$4$
(Proposition~\ref{prop:swiss-cheese-nonformality-by-class})
marks a parity obstruction candidate for the Saito weight: the
weight-$0$ generator~$\gamma$ gives a nontrivial quartic Massey product
in the bar transfer. A proof that this Massey product is incompatible
with purity would require a mixed-Hodge model for the $\beta\gamma$
bar complex. As for Virasoro, condition~(ii) holds, but
condition~(xii) is not proved from it.

\emph{$\cW_N$ at principal nilpotent, generic central
charge} \textup{(}class~$\mathbf{M}$\textup{):}
\begin{conjecture}[Principal $\cW_N$ equivalence on the
generic locus; \ClaimStatusConjectured]
\label{conj:d-module-purity-wn-principal}
Let $\cW_N^k$ be the universal principal $\cW$-algebra
obtained from $V_k(\mathfrak{sl}_N)$ by principal
Drinfeld--Sokolov reduction, at generic central charge.
On the generic locus of the FM configuration space,
condition~(ii) of
Theorem~\ref{thm:koszul-equivalences-meta} is equivalent
to condition~(xii), conditional on the preservation of
Saito weights under principal Drinfeld--Sokolov reduction.
\end{conjecture}
The conditional hypothesis is the conjectured compatibility
of the principal nilpotent BRST differential with the
Kashiwara filtration of the ambient Kac--Moody mixed Hodge
module. Only on the generic locus (away from the FM deep
strata where higher shadow products interact) is the
equivalence plausible; at the boundary, class~$\mathbf{M}$
behavior reintroduces the Virasoro obstruction.

The upshot is that condition~(xii) is family-dependent:
it depends on more than $\dim(\fg)$ or Euler characteristic data, and
the equivalence of (xii) with the Koszulness hierarchy is proved here
only on the affine Kac--Moody lineage under the localization-weight
package. For the non-affine families, the one-way
implication (xii)$\Rightarrow$(x) holds unconditionally
(Remark~\ref{rem:d-module-purity-content}, Steps 1--3, 5),
but the converse requires family-specific input. The working
hypothesis is that the shadow obstruction tower and the mixed-Hodge
weight filtration measure the same obstruction through different
functors; the direct bridge is presently available only for
Kac--Moody.
\end{remark}

\begin{remark}[Computational verification landscape]
\label{rem:koszulness-computational-landscape}
\index{Koszul property!computational verification}
Items \textup{(i)--(v)}, \textup{(ix)}, and \textup{(x)}
\textup{(}seven genuinely independent bidirectional equivalences
\textup{)}, together with item~\textup{(vi)} as a listed
consequence of~\textup{(v)} and the Hochschild
consequence~\textup{(viii)}, have been checked for a landscape of
17 algebras
(Table~\textup{\ref{tab:koszulness-landscape}}):
13 proved Koszul via equivalences~\textup{(ii)} and~\textup{(ix)},
1 proved not Koszul (Ising minimal model,
$L(c_{3,4}, 0)$: Kac--Shapovalov null vector at $h = 6$
in the bar-relevant range, so~\textup{(ix)} fails), 1 conditional
finite-model obstruction
\textup{(}$L_k(\mathfrak{sl}_3)$ at admissible $q \geq 3$:
$\mathrm{rk}(\fg) = 2$ Cartan classes survive in $H^2$ under
Theorem~\textup{\ref{thm:admissible-sl3-non-koszul-qge3}}\textup{)},
2 open ($L_1(\mathfrak{sl}_2)$ and triplet $\cW(2)$).
The principal proof mechanism is PBW universality
(Proposition~\textup{\ref{prop:pbw-universality}}), which
applies to all freely strongly generated algebras: Heisenberg,
$V_k(\fg)$ at all levels including critical, $\mathrm{Vir}_c$
at all central charges, universal $\cW$-algebras, $\beta\gamma$,
free fermion, and the symplectic fermion
(Remark~\textup{\ref{rem:symplectic-logarithmic}}).
Three additional mechanisms are needed: lattice filtration
(Theorem~\textup{\ref{thm:lattice:koszul-morphism}}) for
lattice VOAs, MC4 completion tower
(Theorem~\textup{\ref{thm:completed-bar-cobar-strong}}) for
$\cW_{1+\infty}$, and Kac--Shapovalov null vector detection
for disproving Koszulness of minimal models.
Cross-family consistency checks: $\kappa$ additivity,
complementarity consistency, and shadow depth independence
from Koszulness (shadow classes G, L, C, M each contain
Koszul algebras). The factorization-homology clause~\textup{(vii)}
is used only after the relevant $\mathsf{FH}_\Sigma$ comparison
model has been verified; it is not a computational test for
Koszulness in this table.
\end{remark}

\begin{remark}[Lattice examples and the same firewall]
\label{rem:lattice-firewall-kp}
For an even lattice~$\Lambda$ of rank~$d$, the Heisenberg currents in
$V_\Lambda$ satisfy
\[
J^a(z)J^b(w)\sim \frac{(a,b)}{(z-w)^2},
\qquad
\kappa(V_\Lambda)=d.
\]
The lattice proof of Koszulness uses bosonization and the Gaussian
symmetric Ran bar: connected genus-$0$ amplitudes reduce to Wick
pairings, so the shadow tower stops at degree~$2$. This does not identify
the bar coalgebra with the post-Verdier algebra. The first object is the
Gaussian conilpotent coalgebra $V_\Lambda^{\mathrm i}$, carrying the
lattice momentum grading and the discriminant-sector grading; only
after Verdier or continuous finite-sector duality does one get
$V_\Lambda^!$. In the unimodular case this finite-sector duality returns
the same lattice VOA.

The quantum lattice family $\Vlat_\Lambda^{N,q}$ lies on the strict
ordered $\Eone$ lane. Its $q\mapsto -q$ pairing is therefore a statement
about completed ordered finite windows before averaging; it is not a
symmetric Ran-bar identity until the ordered data have descended.
\end{remark}

\begin{table}[ht]
\centering
\caption{Koszulness landscape: 17 algebras}
\label{tab:koszulness-landscape}
\index{Koszul property!landscape table}
\index{Koszul property!universal vs.\ simple quotient}
\smallskip
{\small
\begin{tabular}{@{}llccll@{}}
\toprule
\textbf{Algebra} & \textbf{Type}
 & \textbf{Gens} & \textbf{Status}
 & \textbf{Mechanism} & \textbf{Shadow} \\
\midrule
\multicolumn{6}{@{}l}{\textit{Free fields}} \\[2pt]
Heisenberg $\cH_k$
 & free field & $1$ & Proved & PBW univ.\ & $G$, $r_{\max}=2$ \\
$\beta\gamma$ ($\lambda{=}1$)
 & free field & $2$ & Proved & PBW univ.\ & $C$, $r_{\max}=4$ \\
Free fermion $\psi$
 & free field & $1$ & Proved & PBW univ.\ & $G$, $r_{\max}=2$ \\
Symplectic fermion${}^{\dagger}$
 & free field & $2$ & Proved & PBW univ.\ & $C$, $r_{\max}=4$ \\
\midrule
\multicolumn{6}{@{}l}{\textit{Affine Kac--Moody}} \\[2pt]
$V_k(\mathfrak{sl}_2)$, generic $k$
 & universal & $3$ & Proved & PBW univ.\ & $L$, $r_{\max}=3$ \\
$V_{-2}(\mathfrak{sl}_2)$, critical
 & logarithmic & $4$ & Open & not known & unknown \\
$V_k(\mathfrak{sl}_2)$, admissible${}^{\ddagger}$
 & universal & $3$ & Proved & PBW univ.\ & $L$, $r_{\max}=3$ \\
$L_1(\mathfrak{sl}_2)$, integrable${}^{\ddagger}$
 & simple quot.\ & $3$ & Open & null vect.\ & $L$, $r_{\max}=3$ \\
$L_k(\mathfrak{sl}_3)$, adm.\ $q \geq 3$${}^{\ddagger}$
 & simple quot.\ & $8$ & Cond.\ not K.\ & Cartan $H^2$ & $L$, $r_{\max}=3$ \\
\midrule
\multicolumn{6}{@{}l}{\textit{Virasoro and $\cW$-algebras}} \\[2pt]
$\mathrm{Vir}_c$, generic $c$
 & universal & $1$ & Proved & PBW univ.\ & $M$, $r_{\max}=\infty$ \\
$\mathrm{Vir}$ at $c = 0$
 & universal & $1$ & Proved & PBW univ.\ & $M$, $r_{\max}=\infty$ \\
$L(c_{3,4}, 0)$ Ising${}^{\ddagger}$
 & simple quot.\ & $1$ & \textbf{Not K.}\ & null vect.\ & $M$, $r_{\max}=\infty$ \\
$\cW_3$ (universal)
 & universal & $2$ & Proved & PBW univ.\ & $M$, $r_{\max}=\infty$ \\
$\cW_{1+\infty}$
 & limit & $\infty$ & Proved & MC4 compl.\ & $M$, $r_{\max}=\infty$ \\
$N{=}2$ SCA${}^{\S}$
 & extended & $4$ & Proved & PBW univ.\ & $M$, $r_{\max}=\infty$ \\
\midrule
\multicolumn{6}{@{}l}{\textit{Lattice and logarithmic}} \\[2pt]
$V_{D_4}$ lattice VOA
 & lattice & $28$ & Proved & lattice filt.\ & $G$, $r_{\max}=2$ \\
Triplet $\cW(2)$
 & logarithmic & $4$ & Open & not known & unknown \\
\bottomrule
\end{tabular}
}

\smallskip\noindent
\textbf{Legend.}
\textit{Status}: Proved = chirally Koszul
(Definition~\ref{def:chiral-koszul-morphism});
Open = neither proved nor disproved;
Not K.\ = proved not chirally Koszul; Cond.\ not K.\ = non-Koszul
under the finite admissible Li--bar package of
Theorem~\ref{thm:admissible-sl3-non-koszul-qge3}.
\textit{Mechanism}: PBW univ.\ =
Proposition~\ref{prop:pbw-universality};
MC4 compl.\ =
Theorem~\ref{thm:completed-bar-cobar-strong};
lattice filt.\ =
Theorem~\ref{thm:lattice:koszul-morphism};
null vect.\ = Kac--Shapovalov obstruction
(Theorem~\ref{thm:kac-shapovalov-koszulness});
Cartan~$H^2$ = abelian Cartan subalgebra classes
surviving in $H^2(\barB)$ under the finite admissible Li--bar package
(Remark~\ref{rem:admissible-koszul-status}).
\textit{Shadow}: archetype class $G$/$L$/$C$/$M$
(Theorem~\ref{thm:nms-archetype-trichotomy})
and shadow depth $r_{\max}$;
all four classes contain Koszul algebras
(shadow depth $\neq$ Koszulness).

\smallskip\noindent
${}^{\dagger}$\,The symplectic fermion
($\beta\gamma$ at $\lambda = \tfrac{1}{2}$,
$c = -1$) \emph{is} chirally Koszul:
the algebra is freely strongly generated, and
logarithmic phenomena appear in the module category,
not in the bar complex
(Remark~\ref{rem:symplectic-logarithmic}).

\smallskip\noindent
${}^{\ddagger}$\,\textbf{Universal vs.\ simple quotient.}
The universal algebra $V_k(\fg)$ is chirally Koszul at
\emph{all} levels~$k$, including critical $k = -h^\vee$
and admissible (Corollary~\ref{cor:universal-koszul}).
The simple quotient $L_k(\fg)$ may fail: null vectors
in the bar-relevant range (conformal weight
$h \geq 2 \cdot h_{\min}$, where $h_{\min}$ is the
minimal generator weight) block the PBW filtration.
At integrable level $k = 1$, the first null vector of
$L_1(\mathfrak{sl}_2)$ lies at $h = 2$, exactly at the
bar-relevant threshold: Koszulness is open.
For minimal models $L(c_{p,q}, 0)$ with coprime $p < q$,
the first null at $h = pq - p - q + 1$ lies in the
bar-relevant range whenever $pq - p - q + 1 \geq 4$
(i.e., for $(p,q) \neq (2,3)$): these are proved
not chirally Koszul by
Theorem~\ref{thm:kac-shapovalov-koszulness}.

\smallskip\noindent
${}^{\S}$\,\textbf{CE complex vs.\ chiral bar.}
The $N{=}2$ SCA has $H^2_{\mathrm{CE}} \neq 0$
at the mode-algebra level, but $H^2(\barB) = 0$:
the CE complex is a coarse upper bound on bar cohomology,
and PBW universality resolves the apparent paradox
(Remark~\ref{rem:n2-sca-koszulness}).
\end{table}

\begin{proposition}[Koszulness as formality of the convolution algebra;
\ClaimStatusConditional]
\label{prop:koszulness-formality-equivalence}
\index{Koszul property!convolution algebra formality}
\index{formality!convolution algebra characterization}

Assume the genus-$0$ obstruction-completeness package: the transferred
$L_\infty$ minimal model of
$\mathfrak{g}^{(0)}_\cA$ exists in the completed filtered category, the
filtration is complete and separated with no $\varprojlim^1$ contribution
in each conformal bidegree, and every finite-degree transferred bracket is
detected by the corresponding shadow obstruction class. Then a chiral
algebra $\cA$ is chirally Koszul if and only if the modular
convolution algebra $\gAmod$ satisfies:
\begin{enumerate}[label=\textup{(\roman*)}]
\item The shadow algebra
 $\cA^{\mathrm{sh}} = H_\bullet(\Defcyc^{\mathrm{mod}}(\cA))$ is
 concentrated in degree~$2$:
 $\cA^{\mathrm{sh}}_{r,g} = 0$ for $r \geq 3$, $g = 0$.
\item The bar spectral sequence
 $E_1^{p,q} = H^q(F^p\gAmod/F^{p+1}\gAmod)$ degenerates
 at~$E_2$.
\item The genus-$0$ convolution algebra
 $\mathfrak{g}^{(0)}_\cA$ is formal as a dg Lie algebra.
\end{enumerate}
These three conditions are equivalent to each other and to the
unconditional core consisting of
\textup{(i)--(vi)} and \textup{(ix)--(x)} in
Theorem~\textup{\ref{thm:koszul-equivalences-meta}}.
The factorization-homology clause~\textup{(vii)} joins this list only
after the comparison and detection hypotheses
$\mathsf{FH}^{\mathrm{det}}_{\mathbb{P}^1}(\cA)$ are imposed.
\end{proposition}

\begin{proof}
Condition~(i) states $\cA^{\mathrm{sh}}_{r,0} = 0$ for $r \geq 3$,
i.e.\ all genus-$0$ obstruction classes
$o_{r+1}^{(0)}(\cA) = 0$ for $r \geq 2$. This is the PBW
concentration condition~(ii) of
Theorem~\ref{thm:koszul-equivalences-meta}, establishing
(i)$\Leftrightarrow$(ii).
Condition~(ii) (bar spectral sequence $E_2$-collapse) is the same
statement expressed in spectral sequence language.
For~(i)$\Leftrightarrow$(iii): in the filtered obstruction-complete
category just assumed, formality of a dg Lie algebra is equivalent to the
vanishing of the transferred higher $L_\infty$ brackets on its minimal
model.
Proposition~\ref{prop:shadow-formality-low-degree} identifies
the genus-$0$ obstruction classes with the transferred $L_\infty$
brackets $\ell_{r+1}^{(0)}$ at degrees $r = 2,3,4$; at each
higher degree the same inductive mechanism applies (the
obstruction at degree $r{+}1$ is the next transferred bracket
evaluated on the truncated MC element). The detection hypothesis is
exactly what promotes these finite-degree identifications to formality of
the completed convolution algebra.
\end{proof}

\begin{remark}[Two transverse slices of the MC element]
\label{rem:koszulness-vs-shadow-depth}
\index{Koszul property!shadow depth independence}
\index{shadow obstruction tower!depth vs.\ Koszulness}
\index{MC element!genus-0 vs all-genera slices}

The MC element $\Theta_\cA \in \MC(\gAmod)$ carries two
transverse families of invariants.
\emph{Koszulness} reads the genus-$0$ slice:
$o_{r+1}^{(0)}(\cA) = 0$ for all $r \geq 2$
(the transferred $A_\infty$~products $m_r = 0$ for $r \geq 3$;
Proposition~\ref{prop:shadow-formality-low-degree}).
\emph{Shadow depth} reads the all-genera projection:
$r_{\max} := \min\{r : o_{s+1}(\cA) = 0 \text{ for all } s > r\}$.
These are independent:
\begin{center}
\begin{tabular}{lll}
\hline
\textbf{Family} & \textbf{Shadow depth $r_{\max}$}
 & \textbf{Archetype} \\
\hline
Heisenberg & $2$ (finite) & Gaussian \\
Affine $\widehat{\fg}_k$ & $3$ (finite) & Lie \\
$\beta\gamma$ system & $4$ (finite) & Contact \\
Virasoro $\mathrm{Vir}_c$ & $\infty$ & Mixed \\
$\mathcal{W}_N$ & $\infty$ & Mixed \\
\hline
\end{tabular}
\end{center}
The generic or universal representatives in these five rows are chirally
Koszul (genus-$0$ slice is formal), while special simple quotients can
leave the Koszul locus. In the generic/universal rows the
Virasoro algebra has infinite shadow depth
($\mathfrak{o}^{(5)}_{\mathrm{Vir}} \neq 0$
at genus~$\geq 1$). The archetype classification
(Theorem~\ref{thm:nms-archetype-trichotomy}) stratifies the
Koszul locus by the all-genera termination behavior of
$\Theta_\cA$, a transverse invariant invisible to genus~$0$.
\end{remark}

\begin{remark}[Shadow classes under direct sum]
\label{rem:shadow-class-direct-sum}
\index{shadow depth!direct sum}
\index{shadow obstruction tower!join-semilattice}

The shadow classes $\{G, L, C, M\}$ form a join-semilattice under
the ordering $G < L < C < M$, with the join operation realised by
direct sum of chiral algebras. For $\cA_1$, $\cA_2$ in the
standard landscape, the OPE of $\cA_1 \oplus \cA_2$ has no mixed
singular terms between the two sectors: the $\cA_1$-fields and
$\cA_2$-fields are mutually local with regular OPE. The shadow
tower of the direct sum therefore decomposes sector by sector, and
the shadow depth satisfies
\begin{equation}
\label{eq:shadow-depth-direct-sum}
 r_{\max}(\cA_1 \oplus \cA_2)
 \;=\;
 \max\bigl(r_{\max}(\cA_1),\, r_{\max}(\cA_2)\bigr).
\end{equation}
In particular:
\begin{center}
\renewcommand{\arraystretch}{1.3}
\begin{tabular}{@{}llcl@{}}
\toprule
$\cA_1$ & $\cA_2$ & $r_{\max}(\cA_1 \oplus \cA_2)$ & Class \\
\midrule
$G$ & $G$ & $\max(2,2) = 2$ & $G$ \\
$G$ & $L$ & $\max(2,3) = 3$ & $L$ \\
$L$ & $L$ & $\max(3,3) = 3$ & $L$ \\
$G$ & $M$ & $\max(2,\infty) = \infty$ & $M$ \\
$C$ & $M$ & $\max(4,\infty) = \infty$ & $M$ \\
\bottomrule
\end{tabular}
\end{center}
The entry $L \oplus L = L$ (not~$C$) is the decisive case: two
affine Kac--Moody factors at levels $k_1$, $k_2$ produce no
cross-sector shadow terms at depth~$4$, because the quartic
obstruction $S_4$ is intrinsic to each factor's OPE and the
mixed channels contribute only regular terms.

This join-semilattice structure is specific to direct sums.
Constructions that create new OPE singularities between
sectors can raise the shadow class strictly: a quantum
Drinfeld--Sokolov reduction of an affine algebra (class~$L$)
produces a $\mathcal{W}$-algebra (class~$M$), and a coset
$\cA / \cA'$ can land in a higher class than either factor.
The shadow depth is monotone under direct sum but not under
quotient or reduction.
\end{remark}


\begin{remark}[Class~$\mathbf{W}$: root multiplicities and the Koszul boundary]
\label{rem:class-w-koszul-frontier}
\index{wild quiver!root multiplicity mechanism}
\index{Koszul property!boundary beyond class M}
\index{Kronecker quiver!bar spectral sequence non-collapse}
\index{root multiplicity!bar cohomology obstruction}

The $\mathbf{G}/\mathbf{L}/\mathbf{C}/\mathbf{M}$ classification
(Remark~\ref{rem:koszulness-vs-shadow-depth}) presupposes chiral
Koszulness: $E_2$-collapse of the bar spectral sequence
(Theorem~\ref{thm:koszul-equivalences-meta}). Beyond this boundary
lies class~$\mathbf{W}$ (wild), where $E_2$-collapse fails and the
unconditional equivalence core loses force simultaneously; seven
genuinely independent bidirectional equivalences (items
\textup{(i)--(v)} and \textup{(ix)--(x)}) and the listed
consequence~\textup{(vi)} -- and the chiral Hochschild
consequence~\textup{(viii)} -- no longer follow. The
factorization-homology clause~\textup{(vii)} has no independent
force in this regime: even when an external integration model exists,
the converse to Koszulness requires the detecting comparison
$\mathsf{FH}^{\mathrm{det}}_{\mathbb{P}^1}$, and the PBW diagonal
already fails.
The mechanism is controlled by the Kronecker quiver $K_m$ (two
vertices, $m$ parallel arrows) and its root multiplicities.

\emph{Root multiplicities and higher differentials.}
The symmetric Euler form of $K_m$ has determinant $4 - m^2$.
For $m \leq 2$, this form is positive semi-definite: every
positive root has multiplicity~$1$, the indecomposable
representations are classified (finite type for $m = 1$, tame
for $m = 2$), and the bar spectral sequence collapses at~$E_2$
by the Kac--Moody concentration argument. For $m \geq 3$, the
form is indefinite. The dimension vectors
$(d_0, d_1)$ with $d_0^2 + d_1^2 - m d_0 d_1 < 0$ support
\emph{families} of indecomposable representations: the moduli
space $\mathcal{M}(d_0, d_1)$ has positive dimension
$1 - d_0^2 - d_1^2 + m d_0 d_1$, and the root multiplicity
exceeds~$1$. These multiple roots produce non-split extensions
in the bar complex $\barB(\cA)$: the $E_2^{p,q}$ page acquires
off-diagonal classes supported on the excess moduli, and the
higher differentials $d_r$ ($r \geq 2$) act nontrivially between
them. The signed Euler series
$\prod_{n \geq 1}(1 - t^n)^{m+2}$ acquires negative
coefficients at weight~$2$, confirming the failure of bar
concentration.

\emph{Why the unconditional equivalences break.}
Each of the seven genuinely independent bidirectional
equivalences \textup{(i)--(v)} and \textup{(ix)--(x)} in
Theorem~\ref{thm:koszul-equivalences-meta}, together with the
listed consequence \textup{(vi)}, requires, at its core,
that bar cohomology $H^\bullet(\barB(\cA))$ be concentrated on the
PBW diagonal. In the wild regime, the excess moduli inject
off-diagonal classes into $H^2$ (and higher) that persist to
$E_\infty$: condition~(ii) (PBW concentration) fails, and with
it the equivalence chain. The conditional Lagrangian criterion
\textup{(xi)} and the $\mathcal{D}$-module purity refinement
\textup{(xii)} are no longer consequences of the Koszul core. This is
only a failure of implication; no mixed-Hodge purity failure is asserted
without a separate Saito-weight computation.

\emph{Concrete sources.}
The following proven and conditional sources produce class~$\mathbf{W}$
behavior:
\begin{enumerate}[label=\textup{(\roman*)}]
\item \emph{Virasoro minimal models.}
 The simple quotient $L(c_{p,q}, 0)$ is non-Koszul by
 Proposition~\ref{prop:minimal-model-non-koszul}: the null
 quotient at weight $w_0 = (p{-}1)(q{-}1)$ enlarges
 $\ker(d\colon \barB^2 \to \barB^1)$, producing $H^2 \neq 0$.
\item \emph{Admissible-level simple quotients at higher rank.}
 For $\fg$ with $\mathrm{rk}(\fg) \geq 2$ and denominator
 $q \geq 3$, the abelian Cartan subalgebra $\mathfrak{h}$
 contributes $\mathrm{rk}(\fg)$ candidate classes to $E_1^{2,q}$.
 At $\mathrm{rk}(\fg)=2$ their survival is conditional on the finite
 admissible Li--bar package of
 Theorem~\ref{thm:admissible-sl3-non-koszul-qge3}; in higher rank it is
 conjectural
 (Conjecture~\ref{conj:admissible-koszul-rank-obstruction}).
\item \emph{Non-principal $\mathcal{W}$-algebras at wild nilpotent type.}
 Under the expected finite DS--Li comparison, quantum
 Drinfeld--Sokolov reduction at a nilpotent $e$ whose associated
 Slodowy slice has wild representation type should produce a
 $\mathcal{W}$-algebra whose Ext quiver has $m \geq 3$ arrows, so the
 resulting bar complex would inherit the wild root structure.
\end{enumerate}
The replacement for the shadow obstruction tower in
class~$\mathbf{W}$ is the motivic Donaldson--Thomas theory of
Kontsevich--Soibelman: the DT invariants $\Omega(d_0, d_1)$
grow exponentially in $\max(d_0, d_1)$ for $m \geq 3$, and
their wall-crossing structure encodes the information that the
shadow tower would carry in the Koszul world
(Remark~\ref{rem:wild-quiver-boundary}).
\end{remark}

\subsubsection{Koszulness from Sklyanin Poisson cohomology}
\label{subsubsec:sklyanin-poisson-koszulness}
\index{Koszul property!Sklyanin Poisson rigidity}
\index{Poisson cohomology!Sklyanin bracket}
\index{Sklyanin bracket!Poisson rigidity}

The Koszulness meta-theorem
(Theorem~\ref{thm:koszul-equivalences-meta}) characterizes chiral
Koszulness intrinsically, through the bar complex and its spectral
sequences. The next proposition gives a Poisson-geometric
characterization that applies on the affine Kac--Moody locus after an
explicit Sklyanin--bar comparison, and is then logically independent of
the PBW spectral sequence.

The classical $r$-matrix
$r^{\mathrm{cl}}(z) = k\Omega/z$ of an affine algebra
$\cA = \widehat{\fg}_k$ defines a Sklyanin--Poisson bracket on
$\fg^*$ via the Semenov--Tian--Shansky
construction~\cite{STS83}. The infinitesimal rigidity of this
bracket (vanishing of its second Poisson cohomology) becomes a
Koszulness criterion only after the finite-window comparison identifies
the PBW associated graded of the ordered bar deformation complex with
the Lichnerowicz--Poisson complex.

\begin{theorem}[Koszulness from Sklyanin--Poisson rigidity;
\ClaimStatusConditional{} for affine KM]
\label{thm:koszulness-from-sklyanin}
\index{Koszul property!Sklyanin Poisson rigidity|textbf}
\index{Poisson cohomology!vanishing|textbf}
\index{Whitehead's second lemma!application to Poisson rigidity}

For $\cA = \widehat{\fg}_k$ with $\fg$ semisimple at generic
non-critical level, the
Sklyanin--Poisson bracket
$\{-,-\}_{\mathrm{STS}}$ on $\fg^*$ determined by the classical
$r$-matrix $r^{\mathrm{cl}}(z) = k\Omega/z$ has vanishing second
Poisson cohomology:
\[
H^2_\pi\bigl(\fg^*, \{-,-\}_{\mathrm{STS}}\bigr) = 0.
\]
Consequently the Sklyanin bracket is infinitesimally rigid: every
first-order Poisson deformation is trivial. Assume, in addition, the
Sklyanin--bar comparison package: in each finite PBW window the
associated graded of the ordered bar deformation complex is the
polynomial Lichnerowicz complex
$(\mathfrak{X}^\bullet(\fg^*),[\pi_{\mathrm{STS}},-]_{\mathrm{SN}})$,
and degree-two Poisson rigidity detects all off-diagonal bar classes.
Under this package, rigidity gives a Poisson-geometric proof of
$E_2$-collapse independent of the PBW spectral sequence.
\end{theorem}

\begin{proof}
The Sklyanin bracket
$\pi_{\mathrm{STS}} = \pi_{\mathrm{LP}} + \pi_r$ decomposes as the
Lie--Poisson bracket $\pi_{\mathrm{LP}}$ (the Kirillov--Kostant
structure on~$\fg^*$) plus a constant-coefficient bivector~$\pi_r$
from the symmetric part of the $r$-matrix. The Lichnerowicz--Poisson
complex $(\mathfrak{X}^\bullet(\fg^*), \delta_\pi =
[\pi_{\mathrm{STS}}, -]_{\mathrm{SN}})$ computes the Poisson
cohomology $H^\bullet_\pi$.

Filter the Poisson complex by polynomial degree on~$\fg^*$. The
$E_0$ page is the Chevalley--Eilenberg complex
$C^\bullet_{\mathrm{CE}}(\fg, \mathrm{Sym}^p(\fg^*))$ in each
polynomial stratum~$p$; the constant correction $\pi_r$ contributes
only at the $E_0$ level. Whitehead's second lemma gives
$H^2_{\mathrm{CE}}(\fg, M) = 0$ for every finite-dimensional
$\fg$-module~$M$ when $\fg$ is semisimple
(\ClaimStatusProvedElsewhere, Weibel~\cite{Weibel94}, Theorem~7.8.10).
In particular $H^2_{\mathrm{CE}}(\fg, \mathrm{Sym}^p(\fg^*)) = 0$
for all~$p \geq 0$. The spectral sequence degenerates at~$E_1$
and yields $H^2_\pi(\fg^*, \pi_{\mathrm{STS}}) = 0$.

The explicit verification for $\fg = \mathfrak{sl}_2$ (the
three-dimensional Lichnerowicz complex on~$\mathbb{C}^3$, carried out
in
\texttt{compute/lib/theorem\_sklyanin\_poisson\_cohomology\_engine.py};
57~tests, three independent verification paths: direct
kernel/image, CE comparison, and polynomial-degree spectral
sequence) confirms $H^0_\pi = \mathbb{C}[C_2]$ (Casimir ring),
$H^1_\pi = H^2_\pi = H^3_\pi = 0$.

Comparison with Koszulness:
the vanishing $H^2_\pi = 0$ means the Sklyanin bracket on
$\fg^*$ has no nontrivial infinitesimal deformations.
The bar-dual coalgebra $\barBgeom(\cA)^{\mathrm i}$ carries
$\fg^*$ as its semiclassical shadow; rigidity of the
Sklyanin bracket is the Poisson-geometric avatar of bar $E_2$-collapse
(condition~(ii) of
Theorem~\ref{thm:koszul-equivalences-meta}). The Sklyanin--bar
comparison hypothesis is precisely the assertion that this
semiclassical shadow is faithful in finite PBW windows: the Poisson
complex $(\mathfrak{X}^\bullet, \delta_\pi)$ is the associated graded of
the ordered bar deformation complex, and degree-two Poisson cohomology
controls the first possible off-diagonal bar obstruction. Under that
comparison, $H^2_\pi=0$ rules out the obstruction and gives the claimed
independent route to $E_2$-collapse.
\end{proof}

\begin{remark}[Scope and the class~$M$ obstruction]
\label{rem:sklyanin-koszulness-scope}
\index{Koszul property!Sklyanin rigidity scope}
\index{Poisson cohomology!differential Poisson brackets}
The Sklyanin rigidity proof applies to affine Kac--Moody algebras
(shadow classes~$G$ and~$L$), where the classical $r$-matrix is
rational of order~$1$ and the resulting Poisson bracket is
algebraic (polynomial coefficients on~$\fg^*$). For class~$M$
algebras (Virasoro, $\cW_N$), the collision residue has higher-order
poles producing \emph{differential} Poisson brackets of order
$\geq 2$ in the sense of~\cite{LPV13}. The Lichnerowicz complex
for such brackets lies outside the scope of Whitehead's lemma:
$H^2_{\mathrm{CE}}$ controls only the algebraic (order-$0$) sector.
Extending the Sklyanin rigidity proof to class~$M$ requires a theory
of differential Poisson cohomology that does not yet exist.
This is a conditional 14th characterization of chiral Koszulness,
restricted to the affine KM subclass and to the Sklyanin--bar comparison
package. On that subclass it is logically independent of the PBW
spectral sequence and provides a second, Poisson-geometric route to
$E_2$-collapse.
\end{remark}


\subsection{The Koszulness moduli scheme}
\label{subsec:koszulness-moduli}
\index{Koszulness moduli scheme|textbf}
\index{Grothendieck--Teichmuller group!torsor structure|textbf}
\index{Tamarkin Phi-transfer|textbf}

The classical meta-theorem presents Koszulness through a list of
equivalent conditions. The list reorganizes as a moduli
scheme: the conditions are \emph{charts} on a
$\mathrm{GRT}_1$-torsor $\cM_{\mathrm{Kosz}}(\cA)$, and the
$\mathrm{GRT}_1$-action on Drinfeld associators is the torsor
structure. This is the Kontsevich--Tamarkin principle made
operational for a fixed chiral algebra.

\begin{definition}[Koszulness moduli scheme]
\label{def:koszulness-moduli-scheme-kp}
\index{Koszulness moduli scheme!definition}
Let $\cA$ be a chiral algebra on~$X$ with PBW filtration. The
\emph{Koszulness moduli scheme} $\cM_{\mathrm{Kosz}}(\cA)$ is the
affine $\bQ$-scheme representing the functor
\[
R \longmapsto
\bigl\{(\Phi, c_\Phi) :
\Phi \in \mathrm{Assoc}(R),\;
c_\Phi \in \mathrm{Chart}_\Phi(\cA;R)\bigr\}
\]
where $\mathrm{Assoc}(R)$ is the Drinfeld associator
torsor~\cite{Drinfeld90} and
$\mathrm{Chart}_\Phi(\cA;R) = \{(\Pi_\Phi, \mathrm{wit}_\Phi,
\gamma_\Phi)\}$ is the chart set: a Koszulness predicate
$\Pi_\Phi$ depending on $(\cA, \Phi)$, a witness
$\mathrm{wit}_\Phi(\cA)$, and a Tamarkin transition cocycle
$\gamma_\Phi\colon \Phi \leadsto \Phi'$ to every other
associator, satisfying the coherence
$\Pi_\Phi(\cA) \iff \Pi_{\Phi \cdot g}(\cA)$
for $g \in \mathrm{GRT}_1(R)$.
\end{definition}

\begin{theorem}[Koszulness moduli; \ClaimStatusConditional]
\label{thm:koszulness-moduli-kp}
\index{Koszulness moduli scheme!torsor theorem}
Let $\cA$ be a chirally Koszul chiral algebra. Then
$\cM_{\mathrm{Kosz}}(\cA)$ is an affine $\bQ$-scheme with:
\begin{enumerate}[label=\textup{(A\arabic*)},leftmargin=3em]
\item
 $\cM_{\mathrm{Kosz}}(\cA) \neq \emptyset$ if and only if
 $\cA$ is chirally Koszul;
\item
 the projection
 $\cM_{\mathrm{Kosz}}(\cA) \to \mathrm{Spec}(\mathrm{Assoc})$
 is a $\mathrm{GRT}_1$-torsor: transitions between charts are
 Tamarkin $\Phi$-transfer maps, and the cocycle computation
 closes in $\mathrm{GRT}_1$;
\item
 the fourteen classical characterizations of
 Theorem~\ref{thm:koszul-equivalences-meta} form a finite affine
 atlas: charts $U_1, \ldots, U_{10}$ at $\Phi_{\mathrm{KZ}}$;
 $U_{11}$ at $\Phi_{\mathrm{AT}}$;
 $U_{12}$ at $\Phi_{\mathrm{dR,B}}$;
 $U_{13}$ at $\Phi_{\mathrm{ell}}$;
 $U_{14}$ at $\Phi_{\mathrm{Kon}}$;
\item
 the shadow-class stratification $G \sqcup L \sqcup C \sqcup M
 \sqcup FF$ selects a distinguished chart on each stratum,
 respectively $\Phi_{\mathrm{KZ}}$, $\Phi_{\mathrm{ell}}$,
 $\Phi_{\mathrm{AT}}$, $\Phi_\infty = \Phi_{\mathrm{Kon}}$,
 and $\Phi_{\mathrm{dR,B}}$.
\end{enumerate}
\end{theorem}

\begin{proof}
The argument has four parts, one for each axiom.

\emph{(A1).} Definition~\ref{def:chiral-koszul-morphism} is
$\mathrm{GRT}_1$-invariant: acyclicity of the twisted tensor
product $K_\tau^L(\cA, \barBgeom(\cA))$ is preserved under the
Tamarkin transfer $\tau \mapsto \Phi \cdot \tau$ (the transfer
is by construction a quasi-isomorphism between
$\Ainfty$-structures). Hence the chart set is non-empty exactly
when Koszulness holds.

\emph{(A2).} Fix the KZ associator as base point. For any
$\Phi \in \mathrm{Assoc}$, there is a unique
$g_\Phi \in \mathrm{GRT}_1$ with
$\Phi = \Phi_{\mathrm{KZ}} \cdot g_\Phi$ (the associator torsor
is principal). The Tamarkin transfer
$g_\Phi^*\colon
\mathrm{Chart}_{\Phi_{\mathrm{KZ}}}(\cA) \xrightarrow{\;\sim\;}
\mathrm{Chart}_\Phi(\cA)$
is a bijection of chart sets, and strict associativity of
associator composition in $\mathrm{GRT}_1$ gives the cocycle
condition
$g_{\Phi''} = g_{\Phi'} \cdot h_{\Phi' \to \Phi''}$.
This is exactly the torsor property.

\emph{(A3).} Conditions (i)--(x) of
Theorem~\ref{thm:koszul-equivalences-meta} are stated in the KZ
language (PBW, bar, $\Ainfty$, BBL, FM) and their coherence data
is the KZ associator; they cover a Zariski-open neighbourhood
of $\Phi_{\mathrm{KZ}}$. The four supplementary charts
$U_{11}, \ldots, U_{14}$ are established by direct argument at
their own associators in the standalone paper
\cite[Propositions~2.5--2.8]{LorgatMKD1_stand}, which imports
the Alekseev--Torossian associator~\cite{AT12}, the de~Rham--Betti
associator~\cite{Deligne89}, the Enriquez elliptic
associator~\cite{Enriquez14}, and the Kontsevich integral
associator~\cite{Kontsevich93}.

\emph{(A4).} The shadow class is an intrinsic invariant of the
MC element $\Theta_\cA$. The selection of a distinguished chart
on each stratum is determined by the computational complexity
of the shadow algebra bracket $[-,-]^{\mathrm{sh}}$: class~G
admits a trivial KZ coordinate (abelian shadow); class~L
activates the genus-one elliptic coupling
(Enriquez~$\Phi_{\mathrm{ell}}$); class~C requires the AT
Kashiwara--Vergne diagonalization of the quartic bracket; class~M
needs the full Kontsevich expansion; and the Feigin--Frenkel
locus needs Hodge--Betti purity.
\end{proof}

\begin{remark}[Kontsevich--Tamarkin reformulated]
\label{rem:kt-reformulation-kp}
\index{Kontsevich--Tamarkin!reformulation}
The Kontsevich--Tamarkin formality theorem~\cite{Kontsevich99,
Tamarkin03} asserts that formality of little disks holds
universally but depends on an associator choice. The moduli
scheme $\cM_{\mathrm{Kosz}}(\cA)$ gives this content at the level
of a single algebra: chiral Koszulness is a
$\mathrm{GRT}_1$-\emph{invariant property}; the classical
characterizations are $\mathrm{GRT}_1$-\emph{coordinate charts}.
Switching between two classical criteria is an application of
the Tamarkin $\Phi$-transfer.
\end{remark}

\begin{remark}[Comparison packages]
\label{rem:fm-closures-kp}


The moduli scheme isolates two comparison packages.
\textup{} (Virasoro Koszulness, circular bar--cobar):
the Hodge--Betti chart $U_{12}$ would provide a non-circular proof
via Feigin--Fuchs BGG resolution~\cite{FeiginFuchs90,RCW82} and
Kac determinant non-degeneracy~\cite{KacBook} after the polarizable
Hodge--Betti realization of the Virasoro bar complex is constructed
(Theorem~\ref{thm:virasoro-koszulness-non-circular-kp}). The unconditional
proof remains the PBW proof of Theorem~\ref{thm:virasoro-chiral-koszul}.
\textup{} (Yangian chiral lift): the associative
Koszulness of $Y_\hbar(\fg)$
(Proposition~\ref{prop:yangian-koszul-general}) is refined to
an open embedding
$\cM_{\mathrm{Kosz}}^{\mathrm{assoc}}(Y_\hbar) \hookrightarrow
 \cM_{\mathrm{Kosz}}^{\mathrm{chiral}}(Y_\hbar)$;
the open question is chart-surjectivity
(Theorem~\ref{thm:yangian-chart-inclusion-kp}).
\end{remark}

\begin{theorem}[Hodge--Betti lane for Virasoro Koszulness;
\ClaimStatusConditional]
\label{thm:virasoro-koszulness-non-circular-kp}
\index{Virasoro algebra!non-circular Koszulness proof}
\index{Feigin--Fuchs BGG resolution|textbf}
For $\Vir_c$ at $c$ outside the minimal-model lattice
$\{c_{p,q}\}$, assume a polarizable Hodge--Betti realization of the
Feigin--Fuchs resolution whose induced weight filtration on
$\barBgeom_n(\Vir_c)$ agrees with the PBW bar filtration in every finite
window. Then the bar cohomology weight components satisfy
\[
H^\bullet(\Vir_c, \Vir_c)_{\mathrm{weight}\,r+1} = 0
\qquad \text{for all } r \geq 2,
\]
without appeal to bar--cobar inversion.
\end{theorem}

\begin{proof}
Work at the Hodge--Betti chart $U_{12}$, centered at
$\Phi_{\mathrm{dR,B}}$~\cite{Deligne89}.

\emph{Step 1: Feigin--Fuchs BGG resolution.}
For each highest weight $h$ in the vacuum Verma module
$M_0(c)$ of $\Vir_c$, the Feigin--Fuchs
theorem~\cite{FeiginFuchs90}, elaborated by
Rocha-Caridi and Wallach~\cite{RCW82}, constructs a resolution
of the vacuum irreducible by Verma modules indexed by an affine
Weyl group orbit:
\[
\cdots \to
\bigoplus_{w \in W^{(2)}} M_{w \cdot 0}(c) \to
\bigoplus_{w \in W^{(1)}} M_{w \cdot 0}(c) \to
M_0(c) \to L(c,0) \to 0,
\]
with differentials given by singular-vector embeddings. At
$c \notin \{c_{p,q}\}$, every singular-vector embedding is
nonzero and the resolution is exact.

\emph{Step 2: Kac determinant non-degeneracy.}
The Kac determinant~\cite{KacBook,FeiginFuchs90} at weight~$h$ is
\[
\det G_h
\;=\; \prod_{\substack{p,q \geq 1 \\ pq \leq h}}
(h - h_{p,q}(c))^{P(h - pq)},
\]
where $h_{p,q}(c)$ are the Kac values and $P$ is the partition
function. At $c \notin \{c_{p,q}\}$, $\det G_h \neq 0$ for every
bar-relevant weight~$h$; the BGG resolution embeds faithfully in
the bar complex.

\emph{Step 3: Conditional purity by Saito stability.}
Under the Hodge--Betti hypothesis, each Verma module $M_{w\cdot 0}(c)$
has a polarizable mixed-Hodge realization whose PBW associated graded is
pure of weight equal to its conformal weight. Saito's stability
theorem~\cite{Saito} then propagates purity through the exact sequence:
the iterated bar complex $\barBgeom_n(\Vir_c)$ is pure of weight~$n$ in
each finite window.

\emph{Step 4: Off-diagonal weight vanishing.}
Purity of weight~$n$ on $\barBgeom_n$ forces
$H^\bullet(\Vir_c, \Vir_c)_{\mathrm{weight}\,r+1} = 0$ for every
$r \geq 2$ off the Koszul diagonal: any nonzero class in weight
$r + 1 \neq n$ would violate purity.

The four steps use Feigin--Fuchs (1), Kac (2), the assumed
Hodge--Betti/Saito comparison (3), and purity (4). Bar--cobar inversion
is not invoked. The open input is exactly the polarizable
Hodge--Betti realization and its equality with the PBW bar filtration;
without it this theorem does not prove $\mathcal D$-module purity for
Virasoro.
\end{proof}

\begin{theorem}[Yangian chart inclusion; \ClaimStatusProvedHere]
\label{thm:yangian-chart-inclusion-kp}
\index{Yangian!associative-to-chiral chart embedding}
For the Yangian $Y_\hbar(\fg)$ with $\fg$ simple, the associative
Koszulness moduli embeds as an open subscheme of the chiral
Koszulness moduli:
\[
\cM_{\mathrm{Kosz}}^{\mathrm{assoc}}(Y_\hbar)
\;\hookrightarrow\;
\cM_{\mathrm{Kosz}}^{\mathrm{chiral}}(Y_\hbar).
\]
The embedding is $\mathrm{GRT}_1$-equivariant; the open
question is chart-surjectivity: whether every
chiral chart lifts an associative one.
\end{theorem}

\begin{proof}
At $\Phi = \Phi_{\mathrm{KZ}}$, the associative chart
$U_1^{\mathrm{assoc}}$ detects Yangian Koszulness in the
classical sense of
Proposition~\ref{prop:yangian-koszul-general}:
PBW flatness $Y_\hbar \rightsquigarrow \mathrm{Sym}(V)$,
RTT quadraticity, the QYBE for the quantum line kernel
$R(z)=1+\hbar r(z)+O(\hbar^2)$, and local finiteness.
The Yangian carries a natural $\Eone$-chiral structure on the
formal disk via its evaluation modules $V_u$; the PBW filtration
is compatible with the chiral bracket. The associative chart
coordinate
$(\Phi_{\mathrm{KZ}}, c^{\mathrm{assoc}}_{\mathrm{KZ}})$ maps to
$(\Phi_{\mathrm{KZ}}, c^{\mathrm{chiral}}_{\mathrm{KZ}})$: both
detect $\Etwo$-collapse of the same PBW spectral sequence (the
associative one on $Y_\hbar$, the chiral one on its
$\Eone$-chiral avatar). The classical residue $r(z)$ satisfies the
CYBE by the Arnold relation, but this checks only the semiclassical
part of the chart. The chart inclusion uses the RTT-complete quantum
lift and the QYBE input; a bare CYBE solution does not define a
Yangian chart. The map is defined on the whole
associative moduli by $\mathrm{GRT}_1$-equivariance.
Openness: the chart $U_1^{\mathrm{assoc}}$ is the complement of
the Shapovalov zero locus; its preimage in the chiral moduli is
the complement of the same zero locus, hence open.
\end{proof}

\begin{remark}[Chart-surjectivity]
\label{rem:fm161-honest}

For $\fg = \mathfrak{sl}_2$ the embedding is an isomorphism
(the Yangian has no genuinely chiral Koszul data beyond the
associative PBW). For higher rank, surjectivity is open and is
the precise content of~: do there exist chiral charts
$(\Phi, c^{\mathrm{chiral}}_\Phi)$ detecting Koszulness that
have no associative preimage? The moduli reformulation strips
the question of all phrasing ambiguity.
\end{remark}

\subsection{Koszulness and the conformal bootstrap}
\label{subsec:koszulness-bootstrap}

\begin{theorem}[Koszulness implies bootstrap closure;
\ClaimStatusConditional]
\label{thm:koszulness-bootstrap}
\index{conformal bootstrap!from Koszulness|textbf}
\index{Koszul property!bootstrap closure|textbf}
Let $\cA$ be a chirally Koszul chiral algebra. Then:
\begin{enumerate}[label=\textup{(\roman*)}]
\item The Maurer--Cartan moduli $\MC(\gAmod)$ is discrete
 \textup{(}zero-dimensional\textup{)}: the MC element
 $\Theta_\cA$ is the unique solution, determined by the
 genus-$0$ OPE data.
\item The shadow depth $r_{\max}(\cA)$ equals the bootstrap
 closure order: the number of OPE channels at which the
 crossing equation
 $D\Theta + \tfrac{1}{2}[\Theta,\Theta] = 0$ at degree
 $n = 4$ is satisfied.
\item The shadow depth classification corresponds to bootstrap
 complexity:
 class~$G$ has no OPE corrections beyond leading order,
 class~$L$ has one cubic correction
 \textup{(}the pentagon identity\textup{)},
 class~$C$ terminates at the quartic,
 class~$M$ requires the infinite tower.
\end{enumerate}
The MC equation at $(g{=}0, n{=}4)$ is the crossing equation
for the four-point function.
\end{theorem}

\begin{proof}
(i) By $A_\infty$~formality
(item~(iii) of Theorem~\ref{thm:koszul-equivalences-meta}),
the transferred $A_\infty$ operations on $H^*(\barB(\cA))$ satisfy
$m_k = 0$ for $k \geq 3$ on the Koszul locus. The MC moduli of a
formal $A_\infty$ algebra is discrete
(Kontsevich's formality implies the Deligne groupoid is
contractible). Uniqueness of $\Theta_\cA$ then follows
from the bar-intrinsic construction
(Theorem~\ref{thm:mc2-bar-intrinsic}).

(ii) The MC equation projected to degree~$r{+}1$ determines
$o_{r+1}(\cA)$ from $\Theta_\cA^{\leq r}$
(Lemma~\ref{lem:nms-euler-inversion}). The shadow depth
$r_{\max}$ is the smallest~$r$ such that
$o_{r+1} = 0$ identically. At degree~$4$, the MC
equation is the crossing constraint for the four-point
function $\langle TTTT \rangle$: both express the
vanishing of the obstruction class in
$H^2(\Defcyc^{\mathrm{mod}})$ at weight~$4$.

(iii) Class~$G$: $S_3 = S_4 = 0$, so the crossing
equation is automatically satisfied (no corrections).
Class~$L$: $S_3 \neq 0$, $S_4 = 0$; one cubic crossing
constraint (the pentagon identity on five FM boundary
strata of $\overline{\mathcal{M}}_{0,5}$).
Class~$C$: $S_4 \neq 0$, $S_5 = 0$; two crossing levels.
Class~$M$: all $S_r \neq 0$; infinite tower.
\end{proof}

\begin{proposition}[Minimal model non-Koszulness;
\ClaimStatusProvedHere]
\label{prop:minimal-model-non-koszul}
\index{minimal model!non-Koszulness mechanism|textbf}
\index{null vector!bar complex effect|textbf}
\index{Koszul property!failure mechanism}
Let $L(c_{p,q}, 0)$ be a non-trivial unitary minimal model
with first vacuum singular vector at
weight $w_0 = (p{-}1)(q{-}1)$. The simple quotient
$L(c_{p,q}, 0)$ is \emph{not} chirally Koszul: the bar
cohomology
$H^2(\barB(L(c_{p,q},0)))_{w_0} \neq 0$.

The mechanism is structural:
the null quotient removes a generator from
$\barB^1$ at weight~$w_0$ while leaving the domain
$\barB^2$ at weight~$w_0$ unchanged
\textup{(}all tensor factors have weight $< w_0$\textup{)}.
The differential $d\colon \barB^2 \to \barB^1$
therefore acquires a larger kernel, producing new
cocycles in~$H^2$.
\end{proposition}

\begin{proof}
At weight $w_0$, the universal Virasoro Verma module has
$p(w_0)$ states (unrestricted partitions of~$w_0$ into parts
$\geq 2$), and the simple quotient removes one null state.
Thus $\dim \barB^1_{w_0}(V_{c}) = p(w_0)$ while
$\dim \barB^1_{w_0}(L) = p(w_0) - 1$. The domain
$\barB^2_{w_0}$ consists of tensor products
$s^{-1}a \otimes s^{-1}b$ with $|a| + |b| = w_0$ and
$|a|, |b| \geq 2$, which is determined by weights
$< w_0$ and is therefore unaffected by the null quotient.
For the Ising model $M(4,3)$: $w_0 = 6$, $p(6) = 4$,
$\dim \barB^1_6(L) = 3$, $\dim \barB^2_6 = 5$, and the
Euler characteristic shifts from~$0$ (universal) to~$1$
(simple quotient), confirming $H^2_6 \neq 0$.
\end{proof}

\subsection{Completion kinematics: the primitive cumulant spectrum}
\label{subsec:completion-kinematics-cumulant}
\index{primitive cumulant spectrum|textbf}
\index{completion kinematics}
\index{Koszul radius}

The shadow obstruction tower probes the \emph{dynamic} axis of the
completed bar theory: the OPE structure, curvature, and modular
anomaly. The complementary
\emph{kinematic} axis depends only on the vacuum character
of~$\cA$. The central objects (the primitive cumulant quotients,
their generating series, and the completion entropy) are
character-level invariants that control the combinatorial growth
of the completed bar coalgebra, independently of any curvature
parameter.

\begin{definition}[Primitive cumulant quotient; \ClaimStatusProvedHere]
\label{def:primitive-cumulant-quotient}
\index{primitive cumulant quotient|textbf}
\index{quasi-primary space}

Let $\cA = \mathbf{1} \oplus \bar{\cA}$ be a conformal chiral
algebra with finite-dimensional weight spaces $\cA[h]$,
$h \geq 0$. Assume the translation operator
$\partial \colon \cA[h-1] \to \cA[h]$ is injective on positive
weights (away from null loci, i.e.\ when
$\det G_h \neq 0$ in the Kac--Shapovalov sense of
Theorem~\textup{\ref{thm:kac-shapovalov-koszulness}}).
The \emph{primitive cumulant quotient} at weight~$h$ is
\begin{equation}
\label{eq:primitive-cumulant-quotient}
Q_h(\cA) \;:=\; \cA[h]\,\big/\,\partial\cA[h-1].
\end{equation}
These are the quasi-primary (or primitive) cumulant spaces: the
new degrees of freedom at each conformal weight that are not
descendants of lower-weight fields.
\end{definition}

\begin{definition}[Primitive generating series; \ClaimStatusProvedHere]
\label{def:primitive-generating-series}
\index{primitive generating series|textbf}
\index{vacuum character!primitive extraction}

Let $a_h := \dim \cA[h]$ and write the vacuum character as
$\chi_\cA(q) = 1 + \sum_{h \geq 1} a_h\, q^h$.
Set $g_r := \dim Q_{r+1}(\cA)$. The \emph{primitive generating
series} is
\begin{equation}
\label{eq:primitive-generating-series}
G_\cA(t) \;:=\; \sum_{r \geq 1} g_r\, t^r
\;=\; \frac{1-t}{t}\,\bigl(\chi_\cA(t) - 1\bigr).
\end{equation}
Equivalently, $g_r = a_{r+1} - a_r$ for all $r \geq 1$.

\smallskip\noindent\textbf{Canonical extraction.}
The primitive spectrum is read off directly from the vacuum
character: $G_\cA$ depends only on the sequence $(a_1, a_2,
a_3, \ldots)$ and no further algebraic data.
\end{definition}

\begin{definition}[Completion Hilbert series; \ClaimStatusProvedHere]
\label{def:completion-hilbert-series}
\index{completion Hilbert series|textbf}
\index{bar coalgebra!reduced-weight window}

The \emph{completion Hilbert series} of $\cA$ is the
associated-graded window series of the completed cumulant bar
coalgebra:
\begin{equation}
\label{eq:completion-hilbert-series}
H_\cA(t) \;:=\; \frac{G_\cA(t)}{1 - G_\cA(t)}
\;=\; \sum_{q \geq 1} h_q\, t^q,
\end{equation}
where $h_q$ counts exact reduced-weight-$q$ bar words in the ideal
cumulant model. This is the nonempty ordered-word series on the
primitive alphabet $\{g_1, g_2, g_3, \ldots\}$: the coefficient
$h_q$ enumerates ordered concatenations of primitive generators
whose reduced weights sum to~$q$.
\end{definition}

\begin{definition}[Primitive defect series; \ClaimStatusProvedHere]
\label{def:primitive-defect-series}
\index{primitive defect series|textbf}
\index{strong generators!primitive defect}

For a chiral algebra with declared strong generators of conformal
weights $d_1, \ldots, d_s$, the \emph{primitive defect series} is
\begin{equation}
\label{eq:primitive-defect-series}
\Delta^{\mathrm{prim}}_\cA(t) \;:=\;
G_\cA(t) - \sum_{i=1}^s t^{d_i - 1}.
\end{equation}
The defect vanishes identically for quadratic (or free-like)
behaviour: $\Delta^{\mathrm{prim}}_\cA = 0$ when the
declared generators account for all quasi-primaries at every
weight. Nonzero defect detects precisely those weights at
which nonquadratic OPEs force new primitive cumulants beyond the
declared generators.
\end{definition}

\begin{definition}[Koszul radius and completion entropy;
\ClaimStatusProvedHere]
\label{def:completion-entropy}
\index{Koszul radius|textbf}
\index{completion entropy|textbf}
\index{bar coalgebra!exponential growth rate}

Define the \emph{Koszul radius} $\rho_\cA$ as the smallest positive
real root of
\begin{equation}
\label{eq:koszul-radius}
G_\cA(\rho_\cA) \;=\; 1.
\end{equation}
The \emph{completion entropy} is
\begin{equation}
\label{eq:completion-entropy}
h_K(\cA) \;:=\; \log\bigl(\rho_\cA^{-1}\bigr).
\end{equation}
This is the exponential growth rate of the ideal reduced-weight bar
windows: $h_q \sim e^{q \cdot h_K(\cA)}$ as $q \to \infty$. The
completion entropy is \emph{kinematic}: it depends only on the
primitive spectrum $(g_1, g_2, \ldots)$, hence only on the vacuum
character, and not on curvature parameters such as $c$ or~$k$.
\end{definition}

\begin{remark}[Kinematics versus dynamics]
\label{rem:kinematics-dynamics-split}
\index{completion kinematics!versus dynamics}
\index{modular dynamics}

The completed bar theory separates into two orthogonal axes:
\begin{enumerate}[label=\textup{(\roman*)}]
\item \emph{Completion kinematics}:
 $Q(\cA)$, $G_\cA(t)$, $H_\cA(t)$, $h_K(\cA)$; depends
 only on the vacuum character.
\item \emph{Modular dynamics}:
 $D_\cA$, $\kappa(\cA)$, $\Theta_\cA$, $r_\cA(z)$; depends
 on the full OPE structure, curvature, anomaly.
\end{enumerate}
The Virasoro involution $c \mapsto 26 - c$ changes dynamic
data (it sends $\kappa_c$ to $\kappa_{26-c}$) but does not
alter the completion entropy:
$h_K(\mathrm{Vir}_c) = h_K(\mathrm{Vir}_{26-c})$
for all~$c$, since both algebras share the same vacuum
character. The same-family shadow $\mathrm{Vir}_{26-c}$
lives on the dynamic axis
(Remark~\ref{rem:virasoro-resonance-model});
kinematically, even generic Virasoro has infinitely many
primitive cumulants, with $g_r = p(r+1) - p(r)$ growing
without bound.

\smallskip
This separation is structurally clean: the primitive spectrum
determines \emph{how large} the completed bar coalgebra is
(how many words at each reduced weight), while the OPE data
determines \emph{what the differential does} on those words.
Both axes are needed for the full MC4 theory
(Theorem~\textup{\ref{thm:completed-bar-cobar-strong}}).
\end{remark}

\begin{proposition}[Cumulant-to-window inversion; \ClaimStatusProvedHere]
\label{prop:cumulant-window-inversion}
\index{cumulant-to-window inversion|textbf}
\index{primitive cumulant spectrum!recovery from windows}

The cumulant law is invertible: from the window counts
$(h_1, h_2, \ldots)$ one recovers the hidden primitive spectrum
$(g_1, g_2, \ldots)$ by the recursion
\begin{equation}
\label{eq:cumulant-window-inversion}
g_1 = h_1, \qquad
g_q \;=\; h_q - \sum_{r=1}^{q-1} g_r\, h_{q-r}
\qquad (q \geq 2).
\end{equation}
\end{proposition}

\begin{proof}
The identity $H_\cA = G_\cA/(1 - G_\cA)$ is equivalent to
$G_\cA = H_\cA/(1 + H_\cA)$, or at the level of formal power
series,
\[
G_\cA(t) \;=\; H_\cA(t) - G_\cA(t)\,H_\cA(t).
\]
Extracting the coefficient of $t^q$ gives
$g_q = h_q - \sum_{r=1}^{q-1} g_r\, h_{q-r} - g_q\, h_0$;
since $h_0 = 0$ (no nonempty words of weight~$0$), the formula
follows by induction on~$q$.
\end{proof}

\begin{remark}[The primitive defect as a Koszulness probe]
\label{rem:primitive-defect-probe}
\index{primitive defect series!Koszulness probe}

For a quadratic chiral algebra $\cA$ with generators of
weights $d_1, \ldots, d_s$, the PBW basis ensures that all
quasi-primaries are generated by the declared strong generators,
so $\Delta^{\mathrm{prim}}_\cA = 0$ identically. The
Virasoro algebra $\mathrm{Vir}_c$ has a single strong generator
of weight~$2$ and $g_r = p(r+1) - p(r)$ for all $r \geq 1$,
giving primitive defect $\Delta^{\mathrm{prim}}_{\mathrm{Vir}}(t)
= G_{\mathrm{Vir}}(t) - t$, which is nonzero starting at
$t^3$. This nonvanishing defect is the kinematic signature of
nonquadraticity and feeds directly into the MC4 completion
package
(\S\textup{\ref{subsec:shadow-postnikov-tower}}).
\end{remark}

\subsection{Geometric bar-cobar duality (Theorem~A)}
\label{sec:theorem-a-statement}

Under the holonomic/perfect hypotheses of
Definition~\ref{def:chiral-koszul-pair}, the homotopy Koszul dual
$\cA^!_\infty := \mathbb{D}_{\Ran}\barB_X(\cA)$ is a factorization
algebra (Chapter~\ref{chap:NAP-koszul-derivation}); outside that lane it
is only a continuous/pro-object until dualizability is proved.
Theorem~A identifies the conditions under which $\cA^!_\infty$ is
\emph{formal}, quasi-isomorphic to the classical dual~$\cA^!$.

\begin{theorem}[Geometric bar--cobar duality; \ClaimStatusProvedHere]
\label{thm:bar-cobar-isomorphism-main}
\index{bar-cobar duality!main theorem|textbf}
\index{Theorem A|textbf}
\textup{[}Regime: quadratic on the Koszul locus
\textup{(}Convention~\textup{\ref{conv:regime-tags})].}
The unit maps in part~\textup{(1)} are equivalences in the derived
category of conilpotent complete factorization coalgebras on
$\operatorname{Ran}(X)$. The counit maps in part~\textup{(1)} and all
post-Verdier identifications in parts~\textup{(2)}--\textup{(3)} are
equivalences in the derived category
$D^b(\mathrm{Fact}^{\mathrm{aug}}(X))$ of augmented
factorization algebras on~$X$.

\smallskip\noindent
The Heisenberg instance was verified in~\S\ref{sec:frame-inversion}.
In general:

Let $(\cA_1, \cC_1, \tau_1, F_\bullet)$ and
$(\cA_2, \cC_2, \tau_2, F_\bullet)$ be a chiral Koszul pair in
the sense of Definition~\textup{\ref{def:chiral-koszul-pair}}.
Then:
\begin{enumerate}
\item the canonical unit maps are weak equivalences of conilpotent
 complete factorization coalgebras, and the canonical counit maps are
 quasi-isomorphisms of factorization algebras:
\[
\cC_i \xrightarrow{\;\sim\;} \bar{B}_X(\cA_i), \qquad
\Omega_X(\cC_i) \xrightarrow{\;\sim\;} \cA_i
\qquad (i = 1, 2);
\]

\item the homotopy Koszul dual factorization algebras are
 identified with the opposite member of the pair:
\[
\mathbb{D}_{\operatorname{Ran}} \bar{B}_X(\cA_1)
\simeq \Omega_X(\cC_2) \simeq \cA_2, \qquad
\mathbb{D}_{\operatorname{Ran}} \bar{B}_X(\cA_2)
\simeq \Omega_X(\cC_1) \simeq \cA_1;
\]

\item if $\cA_2$ is denoted by $\cA_1^!$, then
\[
(\cA_1)^!_\infty := \mathbb{D}_{\operatorname{Ran}} \bar{B}_X(\cA_1)
\simeq \cA_1^!, \qquad
(\cA_2)^!_\infty := \mathbb{D}_{\operatorname{Ran}} \bar{B}_X(\cA_2)
\simeq \cA_2^!.
\]
\end{enumerate}
\end{theorem}

\begin{proof}
Part~(1) is the fundamental theorem of chiral twisting morphisms
(Theorem~\ref{thm:fundamental-twisting-morphisms}): the Koszulness
conditions
(Definition~\ref{def:chiral-koszul-morphism}) are equivalent to
the unit maps being weak equivalences of factorization coalgebras and
the counit maps being quasi-isomorphisms of factorization algebras.
For part~(2), the Verdier compatibility in
Definition~\ref{def:chiral-koszul-pair} identifies the algebra-level
Verdier dual $\mathbb{D}_{\operatorname{Ran}}(\cC_1)$ with the
factorization algebra $\Omega_X(\cC_2)$, and likewise with the
indices reversed. In the standard examples, this hypothesis is
supplied by Theorem~\ref{thm:verdier-bar-cobar}, whose output already
lives on the post-Verdier factorization-algebra side rather than on a
coalgebra side. Let
\[
\eta_i \colon \cC_i \xrightarrow{\sim} \barB_X(\cA_i),
\qquad
\varepsilon_i \colon \Omega_X(\cC_i) \xrightarrow{\sim} \cA_i
\]
be the unit and counit equivalences from part~(1). Applying
$\mathbb{D}_{\operatorname{Ran}}$ to $\eta_1$ gives an
equivalence of factorization algebras
\[
\mathbb{D}_{\operatorname{Ran}}\barB_X(\cA_1)
\xrightarrow{\sim}
\mathbb{D}_{\operatorname{Ran}}(\cC_1).
\]
Composing this with the Verdier-compatibility identification
$\mathbb{D}_{\operatorname{Ran}}(\cC_1)\simeq \Omega_X(\cC_2)$ and
then with $\varepsilon_2$ yields
\[
\mathbb{D}_{\operatorname{Ran}} \bar{B}_X(\cA_1)
\simeq \Omega_X(\cC_2) \simeq \cA_2,
\]
and similarly for the other index. Every object after
$\mathbb{D}_{\operatorname{Ran}}$ is therefore a factorization
algebra, not a coalgebra.
Part~(3) is the same statement after naming $\cA_2$ as the Koszul
dual~$\cA_1^!$.
\end{proof}

\begin{remark}[Alternative proof via derived algebraic geometry]
There is a second proof of
Theorem~\ref{thm:bar-cobar-isomorphism-main} in the
$\infty$-categorical formulation of
Remark~\ref{rem:theorem-a-model}. In the stable presentable
$\infty$-category of augmented factorization algebras on
$\Ran(X)$, the role of the terminal algebra~$k$ is played by the unit
object~$\omega_X$. For an augmentation
$\varepsilon\colon\cA\to\omega_X$, the associated simplicial bar object
$\operatorname{Bar}_\bullet(\omega_X,\cA,\omega_X)$ is the
\v{C}ech nerve of~$\varepsilon$, and its geometric realization is the
bar coalgebra~$\barB_X(\cA)$~\cite[Chapter~5, especially
\S\S5.2.2--5.2.4]{HA}. Dually, a coaugmented conilpotent
coalgebra~$\cC$ has a cosimplicial cobar object obtained from the dual
\v{C}ech nerve, and its totalization is~$\Omega_X(\cC)$ by the same
bar\slash cobar formalism~\cite[Chapter~5, especially
\S\S5.2.3--5.2.4]{HA}. The general realization\slash totalization
adjunction therefore produces the functor pair
\[
\barB_X \colon
\operatorname{Alg}_{\mathrm{aug}}(\operatorname{Fact}(X))
\rightleftarrows
\operatorname{CoAlg}_{\mathrm{conil}}(\operatorname{Fact}(X))
\,:\!\Omega_X.
\]
In this $\infty$-categorical formulation, the Koszul locus of the
chapter is the locus where these \v{C}ech and co-\v{C}ech resolutions
are effective, so the unit and counit are equivalences there. This
proves part~\textup{(1)} without the filtered comparison from the
simplicial bar object to the explicit dg model.

For part~\textup{(2)}, Serre--Grothendieck duality on the relevant
holonomic factorization $\cD$-modules over the derived Ran prestack
sends geometric realizations to the totalizations of the dual
cosimplicial objects and exchanges the left and right module structures
carried by the augmentation object~\cite[Chapter~IV.5]{GR17}. Applied
to~$\cA_1$, this identifies
$\mathbb{D}_{\Ran}\barB_X(\cA_1)$ with the cobar
factorization algebra attached to the Verdier-dual datum. For a chiral
Koszul pair, Definition~\ref{def:chiral-koszul-pair} identifies that
dual datum with~$\Omega_X(\cC_2)$; composing with the counit
equivalence $\Omega_X(\cC_2) \xrightarrow{\sim} \cA_2$ from
part~\textup{(1)} gives
\[
\mathbb{D}_{\Ran}\barB_X(\cA_1)
\simeq \Omega_X(\cC_2) \simeq \cA_2,
\]
and similarly with the indices reversed.

This route is independent of the filtered-comparison lemmas. It proves
the adjunction in part~\textup{(1)} and the Verdier intertwining in
part~\textup{(2)} before choosing an explicit dg presentation. The
remaining input is family-specific: the Verdier-dual datum must be
represented by the opposite member of the chosen Koszul pair. That
verification is exactly the role isolated in
Theorem~\ref{thm:verdier-bar-cobar}.
\end{remark}

\begin{proposition}[Relative extension from relative Verdier base change;
\ClaimStatusProvedHere]
\label{prop:bar-cobar-relative-extension}
\index{bar-cobar duality!relative extension}
Let $\pi \colon \mathcal{X} \to S$ be a family of smooth projective
curves. Suppose given relative factorization data
\[
(\cA_{1,S}, \cC_{1,S}, \tau_{1,S}, F_\bullet),
\qquad
(\cA_{2,S}, \cC_{2,S}, \tau_{2,S}, F_\bullet)
\]
together with relative algebra-level Verdier identifications,
where $\Omega^{\mathrm{rel}}$ denotes the cobar functor formed over
the family $\pi$:
\[
\mathbb{D}_{\operatorname{Ran}(\mathcal{X}/S)}(\cC_{1,S})
\simeq \Omega^{\mathrm{rel}}(\cC_{2,S}),
\qquad
\mathbb{D}_{\operatorname{Ran}(\mathcal{X}/S)}(\cC_{2,S})
\simeq \Omega^{\mathrm{rel}}(\cC_{1,S}),
\]
whose geometric fibers are chiral Koszul pairs in the sense of
Definition~\textup{\ref{def:chiral-koszul-pair}}. Assume:
\begin{enumerate}[label=\textup{(\alph*)}]
\item each relative bar component is represented on the relative
 Fulton--MacPherson spaces by a bounded complex of relative
 holonomic $\cD$-modules;
\item for every geometric point $s \to S$, the canonical base-change
 morphism
 \[
 s^* \mathbb{D}_{\operatorname{Ran}(\mathcal{X}/S)}
 \barB_{\mathcal{X}/S}(\cA_{i,S})
 \longrightarrow
 \mathbb{D}_{\operatorname{Ran}(X_s)}
 \barB_{X_s}(\cA_{i,s})
 \]
 is an isomorphism;
\item the unit and counit quasi-isomorphisms of
 Theorem~\ref{thm:bar-cobar-isomorphism-main}\textup{(1)} commute
 with pullback to geometric fibers.
\end{enumerate}
Then the equivalences of
Theorem~\ref{thm:bar-cobar-isomorphism-main} assemble into
equivalences of relative factorization algebras over~$S$, functorial
under base change in~$S$.

For the universal smooth curve over $\mathcal{M}_{g,n}$, the additional
inputs beyond Theorem~\ref{thm:bar-cobar-isomorphism-main} are the lift
of the displayed relative Verdier identifications and the verification
of \textup{(a)} and \textup{(b)} on the relative Ran space. Over
$\overline{\mathcal{M}}_{g,n}$ the curve is nodal at the boundary, so an
annular sewing comparison is an additional hypothesis rather than a
formal consequence of this proposition.
\end{proposition}

\begin{proof}
Hypothesis~\textup{(c)} promotes part~\textup{(1)} of
Theorem~\ref{thm:bar-cobar-isomorphism-main} from a fiberwise
statement to a compatible relative unit/counit package. Applying
$\mathbb{D}_{\operatorname{Ran}(\mathcal{X}/S)}$ to the relative unit
map and then pulling back to a geometric fiber $X_s$ gives, by
hypothesis~\textup{(b)}, the algebra-level Verdier equivalence of
Theorem~\ref{thm:bar-cobar-isomorphism-main}\textup{(2)} on that
fiber. The displayed relative Verdier identifications provide the
algebra target on the total space, and hypothesis~\textup{(a)} is the
exact condition needed to take Verdier duals degreewise in the
relative holonomic category. These ingredients glue the fiberwise
identifications over~$S$. The universal smooth curve over
$\mathcal M_{g,n}$ is covered by the same argument once
\textup{(a)}--\textup{(c)} are verified there; the compactification
$\overline{\mathcal M}_{g,n}$ has nodal boundary and requires the
additional annular sewing comparison isolated below.
\end{proof}

\begin{remark}[Proof obligations for the modular-family extension]
\label{rem:bar-cobar-relative-programme}
\index{bar-cobar duality!modular-family proof obligations}
To prove the full modular-family statement over
$\overline{\mathcal{M}}_{g,n}$, one must establish the two relative
inputs isolated in
Proposition~\ref{prop:bar-cobar-relative-extension}:
\begin{enumerate}[label=\textup{(\arabic*)}]
\item construct the relative bar components as bounded holonomic
 $\cD$-complexes on the relative Fulton--MacPherson spaces;
\item prove the relative Verdier/base-change isomorphism on the
 relative Ran space, lift the algebra-level Verdier-compatible pair
 identifications, and check compatibility with the factorization
 gluing maps.
\end{enumerate}
The fixed-curve theorem above is the strongest unconditional result
proved in this chapter.
\end{remark}

\begin{remark}[Analysis of the three relative inputs]
\label{rem:relative-ran-analysis}
\index{bar-cobar duality!relative inputs analysis}
The three inputs required by
Proposition~\ref{prop:bar-cobar-relative-extension} admit the
following analysis.

\begin{enumerate}[label=\textup{(\alph*)}]
\item \emph{Relative holonomic bar.}
For each geometric fiber~$X_s$, the bar complex
$\barB_{X_s}(\cA_s)$ is a bounded holonomic $\cD$-complex by
the fixed-curve holonomicity established above. Kashiwara's
base-change theorem for holonomic $\cD$-modules
\textup{(}Kashiwara~\cite{Kas84}, Theorem~4.3.1\textup{)}
preserves holonomicity under proper base change. To obtain
hypothesis~\textup{(a)} one must still assemble these fiberwise
complexes into a bounded relative holonomic complex on the relative
Fulton--MacPherson stages and check the bar-length transition maps.
The ind-properness of
$\operatorname{Ran}(\mathcal{X}/S)\to S$ for a smooth proper family is
the geometric input; it is not by itself the completed relative bar
object.

\item \emph{Relative Verdier and base change on Ran.}
Verdier duality commutes with proper base change by Verdier's original
theorem. In the ind-scheme setting of the Ran space,
Gaitsgory--Rozenblyum~\textup{(}\cite{GR17},
Vol.~II, Chapter~7, Theorem~7.3.4\textup{)} supply the ambient
base-change theorem for ind-proper morphisms. To apply it to the
completed bar complex one must still verify finite-stage compatibility:
the bar truncations must be represented by holonomic complexes on the
relative Fulton--MacPherson stages, the transition maps must satisfy the
Mittag--Leffler condition, and Verdier duality must commute with the
completed inverse limit. These verifications are precisely
hypothesis~\textup{(b)} of
Proposition~\ref{prop:bar-cobar-relative-extension}; they are not
automatic from smoothness alone.

\item \emph{Lifted relative Verdier pair.}
The Koszul pair $(\cA, \cA^!)$ on a fixed curve lifts to a relative
Koszul pair over~$S$ provided the twisting morphism
$\tau \colon \barB(\cA) \to \cA$ is a natural transformation of
$\cD$-module-valued functors on
$\operatorname{Ran}(\mathcal{X}/S)$. This is a separate pair-lift
obligation: one must construct the OPE residue morphisms on the relative
Fulton--MacPherson boundary strata, prove their holomorphic variation in
the fiber parameter of~$\pi$, and check compatibility with factorization
gluing. Fiberwise existence of~$\tau$ does not by itself give the
relative natural transformation.
\end{enumerate}

\noindent
\emph{Conclusion.}
For smooth families $\pi \colon \mathcal{X} \to S$ satisfying the
three inputs above, the modular-family extension of
Theorem~\textup{A} holds over~$S$. This is a relative Verdier/base-change
statement, not an automatic consequence of smoothness alone. Over
$\overline{\mathcal{M}}_{g,n}$, the extension requires additional
care at the boundary: the bar complex over a nodal fiber involves
the annular bar $\barB^{\mathrm{ann}}$, which requires a separate
sewing argument. The sewing comparison is partially established
in Chapter~\ref{chap:higher-genus} \textup{(}the analytic
Hilbert-space sewing theorem, proved for all genera\textup{)}; the
remaining algebraic chain-level identification at nodal fibers is the
proof obligation recorded in
Remark~\ref{rem:bar-cobar-relative-programme}. The strongest
unconditional result is: fixed smooth curve; smooth relative families
after the displayed relative holonomicity, base-change, and pair-lift
inputs; $\overline{\mathcal{M}}_{g,n}$ conditional on the nodal sewing
comparison.
\end{remark}

\begin{remark}[Geometric substrate (Volume~II)]
\label{rem:theorem-a-lagrangian}
\index{Lagrangian self-intersection!Theorem A}
Volume~II reveals that this adjunction is the groupoid
comodule-module adjunction for the Lagrangian self-intersection
$\mathfrak{S} = \mathcal{L} \times_{\mathcal{M}} \mathcal{L}$:
the bar complex $\barB_X(\cA)$ is the structure sheaf of the
derived self-intersection of the boundary Lagrangian inside a
$(-2)$-shifted symplectic stack, and the cobar functor reconstructs
the Lagrangian from its groupoid algebra.
\end{remark}

\begin{remark}[Twisting morphism content of Theorem~A]
\label{rem:theorem-a-tau}
\index{twisting morphism!Theorem A content}
The unit $\eta_{\tau_i} \colon \cC_i \to \barB_X(\cA_i)$ and
counit $\varepsilon_{\tau_i} \colon \Omega_X(\cC_i) \to \cA_i$
in part~(1) are induced by the Koszul morphism~$\tau_i$
of the underlying twisting datum
(Definition~\ref{def:chiral-twisting-datum}).
Part~(1) says each~$\tau_i$ is acyclic;
part~(2) says that after applying Verdier duality to the bar
coalgebra of one side one obtains the post-Verdier dual factorization algebra
recovered from the other;
and Proposition~\ref{prop:bar-cobar-relative-extension} isolates
the extra relative Verdier/base-change input needed to transport
these identifications over a base.
\end{remark}

\begin{remark}[Elementary model presentation;
Convention~\ref{conv:proof-architecture}]
\label{rem:theorem-a-model}
\index{bar-cobar duality!model presentation}
\emph{$\infty$-categorical adjunction.}
The bar functor $\mathbf{B}_X$ and cobar functor
$\boldsymbol{\Omega}_X$ form an adjoint pair
\[
\mathbf{B}_X \colon
\operatorname{Alg}_{\mathrm{aug}}(\operatorname{Fact}(X))
\rightleftarrows
\operatorname{CoAlg}_{\mathrm{conil}}(\operatorname{Fact}(X))
\,:\! \boldsymbol{\Omega}_X
\]
on the stable $\infty$-category of factorization (co)algebras on
$\operatorname{Ran}(X)$. On the Koszul locus, the unit and counit
are equivalences: $\boldsymbol{\Omega}_X \mathbf{B}_X(\cA)
\simeq \cA$. Verdier duality identifies the bar coalgebra with the
homotopy dual factorization algebra via the post-Verdier equivalence
$\mathbb{D}_{\operatorname{Ran}} \mathbf{B}_X(\cA) \simeq
(\cA)^!_\infty \simeq \cA^!$ on the Koszul locus. This is the
theorem; what follows are presentations of the same objects.

\emph{Explicit dg model.}
The explicit dg model uses the reduced bar complex
$\bar{B}_X(\cA) = (T^c(s^{-1}\bar{\cA}),\, d_{\mathrm{bar}})$,
with $d_{\mathrm{bar}}$ encoded by configuration-space integrals
on $\overline{C}_n(X)$. The universal twisting morphism
$\tau \colon \bar{B}_X(\cA) \to \cA$ is the degree-$+1$
OPE extraction against the log form on $\overline{C}_2(X)$. The
post-Verdier comparison with~$\cA^!$ is a separate duality statement.
The proof reduces to showing the unit and counit are
quasi-isomorphisms, which follows from the Koszul conditions
(Definition~\ref{def:chiral-koszul-morphism}).

\emph{Model comparison.}
Any two admissible dg presentations of the same factorization
object are connected by a contractible space of
quasi-isomorphisms
(Proposition~\ref{prop:model-independence}). Hence the dg
quasi-isomorphisms represent the $\infty$-categorical equivalence
up to a contractible ambiguity.

\emph{Cohomology.}
The bar coalgebra $\bar{B}_X(\cA)$ has cohomology
$\cA^{\mathrm i}$ concentrated on the Koszul diagonal. When the
finite-type Verdier dual exists, the post-Verdier algebra~$\cA^!$ is
characterized by Verdier intertwining
$\mathbb{D}_{\operatorname{Ran}}\bar{B}_X(\cA) \simeq
\cA^!_\infty$ (factorization \emph{algebra}, not coalgebra;
Convention~\ref{conv:bar-coalgebra-identity}), with
dimensions controlled by the spectral
discriminant~$\Delta_{\cA}(x)$.

\emph{Cohomology extraction.}
Taking cohomology of the dg model extracts the bar-dual coalgebra.
Concretely, the PBW spectral sequence
$E_1^{p,q} \Rightarrow H^{p+q}(\bar{B}_X(\cA))$ collapses
at~$E_2$ on the Koszul locus
(Corollary~\ref{cor:bar-cohomology-koszul-dual}),
and the resulting Hilbert series
$\sum_n \dim H^n(\bar{B}_X(\cA))\, x^n$ is the Hilbert series
of~$\cA^{\mathrm i}$. After finite Verdier duality it is also the
Hilbert series of~$\cA^!$, with radius of convergence determined
by~$\Delta_{\cA}(x)$.
\end{remark}


% ================================================================
% SECTION 8.3: YANGIANS AND AFFINE YANGIANS - COMPLETE TREATMENT
% ================================================================

\section{Yangians and affine Yangians: self-duality and Koszul theory}
\label{sec:yangians-complete}

The Yangian $Y(\mathfrak{g})$ and affine Yangian
$Y_{\hbar}(\widehat{\mathfrak{g}})$ provide examples where the
finite-window ordered quadratic shadow is self-dual. The full
post-Verdier algebra-level statement requires the completed ordered bar
coalgebra, continuous duality, and the RTT/QYBE quantum line-kernel
datum; the CYBE collision residue is only its semiclassical
codimension-$2$ shadow.

For the type-A RTT definition, the Drinfeld presentation in all simple
types, the Hopf structure, affine Yangians, and the physical
interpretation, see Chapter~\ref{chap:yangians}. We record here only
the Koszul-duality-specific results.

\subsection{Self-duality}

\begin{theorem}[Type-A Yangian quadratic shadow; \ClaimStatusProvedHere]\label{thm:yangian-self-dual}
Let $Y^{\mathrm{RTT}}_\hbar(\mathfrak{gl}_N)$ be the completed RTT
Yangian with Yang matrix
\[
R(u;\hbar)=1-\frac{\hbar P}{u},
\]
where $P$ is the permutation of $\mathbb C^N\otimes\mathbb C^N$. Fix a
finite level window and let $Q_R$ be the homogeneous quadratic RTT
algebra generated by the corresponding coefficients of $T(u)$, with
relation
\begin{equation}\label{eq:rtt-relation}
R_{12}(u-v)\,T_1(u)\,T_2(v)
=T_2(v)\,T_1(u)\,R_{12}(u-v).
\end{equation}
Then the quadratic dual algebra of $Q_R$ is the RTT algebra with
projectively inverse Yang matrix:
\[
Q_R^{!_{\mathrm{quad}}}\cong Q_{R^{-1}}.
\]
Equivalently, since
\[
R(u;\hbar)^{-1}
=\left(1-\frac{\hbar^2}{u^2}\right)^{-1}
\left(1+\frac{\hbar P}{u}\right)
\sim R(u;-\hbar),
\]
and scalar rational factors cancel from the RTT relation, the type-A
quadratic shadow is self-dual after the parameter change
$\hbar\mapsto-\hbar$.

This is a statement about the ordered $\Eone$ quadratic shadow. It does
not identify the full symmetric Ran bar coalgebra with an RTT algebra,
and it does not say that the Verdier-dual algebra $Y_\hbar(\mathfrak g)^!$
is obtained by replacing $R$ with $R^{-1}$ outside the type-A RTT lane.
\end{theorem}

\begin{proof}[Proof of Theorem~\ref{thm:yangian-self-dual}]
Work in a finite level window, so all vector spaces are finite over
$\mathbb C[[\hbar]]$ and continuous duals are ordinary duals in each
bidegree. Write the RTT relation as the image of the linear map
\[
\Phi_R(X)=R_{12}X-\sigma(X)R_{12}
\]
on the quadratic tensor space of two matrix coefficients. The trace
pairing on $\End(\mathbb C^N)^{\otimes 2}$ identifies the orthogonal
complement of $\operatorname{Im}\Phi_R$ with the kernel of the adjoint
map. Moving $R_{12}$ across the trace changes it to the inverse
braiding, hence the adjoint relation is the RTT relation for
$R^{-1}_{21}$, which is projectively the same Yang matrix in this
normalization. Therefore the quadratic dual relation space is exactly
the RTT relation space attached to $R^{-1}$.

For the Yang matrix, $P^2=1$, so
\[
\left(1-\frac{\hbar P}{u}\right)^{-1}
=\left(1-\frac{\hbar^2}{u^2}\right)^{-1}
\left(1+\frac{\hbar P}{u}\right).
\]
The scalar factor $\left(1-\hbar^2/u^2\right)^{-1}$ multiplies both
sides of~\eqref{eq:rtt-relation} by the same invertible rational
function and is invisible to the RTT ideal. Thus the inverse matrix
defines the same quadratic algebra as $R(u;-\hbar)$. The bar-dual object
computed here is the quadratic coalgebra $Q_R^{\mathrm i}$; the notation
$Q_R^{!_{\mathrm{quad}}}$ appears only after finite continuous duality.
\end{proof}

\begin{proposition}[Yangian ordered-bar Koszulness in finite windows;
\ClaimStatusProvedHere]\label{prop:yangian-koszul-general}
\index{Yangian!Koszulness!general type}
Let $\mathfrak g$ be a finite-dimensional simple Lie algebra and let
$Y_\hbar(\mathfrak g)$ be Drinfeld's Yangian with augmentation and level
filtration. For every finite level window $Y_{\hbar,\le m}(\mathfrak g)$,
completed $\hbar$-adically and filtered by bar length, the ordered bar
complex is bidegreewise finite and its PBW spectral sequence is diagonal:
\[
H^*(\barB_X^{\mathrm{ord}}(Y_{\hbar,\le m}(\mathfrak g)))
\cong Y_{\hbar,\le m}(\mathfrak g)^{\mathrm i}
\]
as a conilpotent ordered chiral coalgebra. Passing to the inverse limit
over~$m$ gives the completed intrinsic coalgebra
$Y_\hbar(\mathfrak g)^{\mathrm i}$ whenever the finite windows form the
Mittag--Leffler system of Lemma~\ref{lem:completion-convergence}.
Only after continuous duality of these finite windows does one obtain a
post-Verdier algebra $Y_\hbar(\mathfrak g)^!$.
\end{proposition}

\begin{proof}
Drinfeld's PBW theorem \cite{Drinfeld85} gives the filtered
identification
$\operatorname{gr}Y_\hbar(\mathfrak g)\cong U(\mathfrak g[t])$.
For the classical matrix types this is compatible with the RTT
presentations treated by Molev \cite{Molev07,molev-yangians}; for the
exceptional and uniform Drinfeld presentations the required PBW basis is
the theorem of Guay--Regelskis--Wendlandt \cite[Thm.~3.2]{GRW18}. In a
fixed level window all homogeneous pieces are finite-dimensional.

Apply the Polishchuk--Positselski PBW theorem for locally finite
quadratic-linear filtered algebras \cite[Ch.~4]{PP05}. The associated
graded algebra $U(\mathfrak g[t]_{\le m})$ is PBW, and its further
monomial-length associated graded is
$\operatorname{Sym}(\mathfrak g[t]_{\le m})$, whose Koszul complex is
the classical symmetric-exterior complex. The PBW theorem transfers this
acyclicity to $Y_{\hbar,\le m}(\mathfrak g)$ and gives diagonal collapse
of the ordered bar spectral sequence.

The transition maps between finite windows are surjective in each fixed
bar length and level. Hence the inverse system is bidegreewise
Mittag--Leffler, so Lemma~\ref{lem:completion-convergence} identifies
the completed inverse limit with the completed ordered bar coalgebra.
The construction produces $Y_\hbar(\mathfrak g)^{\mathrm i}$ first. The
algebra $Y_\hbar(\mathfrak g)^!$ is obtained only by the continuous
Verdier-dual passage.
\end{proof}

% Hopf structure, physical interpretation, explicit computations, and quantum group
% connections for Yangians are developed in Chapter~\ref{chap:yangians}.

\section{Intrinsic bar-dual coalgebras and comparison}
\label{sec:three-stage-construction}

\subsection{The type-correct comparison problem}

\begin{problem}[Non-circular duality data]\label{prob:circularity}
The expression
\[
\bar{B}^{\mathrm{ch}}(\mathcal{A}_1)\simeq\mathcal{A}_2^!
\]
is not type-correct as stated. The left-hand side is a bar
coalgebra; $\mathcal{A}_2^!$ denotes a post-Verdier algebra on the
finite-type Koszul lane. The correct comparison has two separate
assertions:
\[
\widehat{\bar{B}^{\mathrm{ch}}(\mathcal{A}_1)}
\xrightarrow{\;\sim\;}
\mathcal{A}_1^{\mathrm i}
\quad\text{in conilpotent chiral coalgebras,}
\]
and, only after continuous linear or Verdier duality,
\[
(\mathcal{A}_1^{\mathrm i})^\vee\simeq\mathcal{A}_1^!\simeq
\mathcal{A}_2
\quad\text{in chiral algebras.}
\]
Thus a proof of chiral Koszul duality must supply:
\begin{enumerate}
\item an intrinsic candidate for the bar-dual coalgebra
 $\mathcal{A}_1^{\mathrm i}$ from generators, relations,
 OPE residues, and the residue pairing;
\item a proof that this candidate carries the required strict or
 completed $A_\infty$ chiral coalgebra structure;
\item a filtered comparison
 $\widehat{\bar{B}^{\mathrm{ch}}(\mathcal{A}_1)}
 \to\mathcal{A}_1^{\mathrm i}$ whose associated graded is the
 classical Koszul comparison.
\end{enumerate}
For quadratic algebras this is the usual orthogonality condition
$R_1^\perp=R_2$. For PBW but non-quadratic chiral algebras such as the
universal Virasoro and universal principal $\mathcal W$ families, the
comparison is filtered and completed; the uncompleted bar complex is
not identified with the dual algebra.
\end{problem}

\begin{remark}
A chiral coalgebra is the object controlled by collision residues in
the opposite direction. Its coproduct records how a configuration
splits along boundary strata; its differential is dual to the OPE
residue calculus after the desuspension convention of
Appendix~\ref{app:signs-shifts}.
\end{remark}

%================================================================
% INTRINSIC BAR-DUAL COALGEBRA
%================================================================

\subsection{Intrinsic bar-dual coalgebra}

\begin{definition}[Intrinsic bar-dual chiral coalgebra]\label{def:koszul-dual-coalgebra}
\label{def:intrinsic-koszul-dual}
\label{def:koszul-dual-chiral}
\index{Koszul dual!intrinsic|textbf}
Let $\mathcal{A}=T_{\mathrm{ch}}(\mathcal{V})/(R_{\geq 2})$ be a
finite-type PBW chiral algebra on~$X$, with $\mathcal{V}$ a locally free
$\mathcal{D}_X$-module of generators, completed relation ideal
$R_{\geq 2}$, and OPE residues
\[
\mu_{r}\colon j_*j^*(\mathcal{V}^{\boxtimes r})
\longrightarrow \Delta_*\mathcal{V}\otimes
\Omega_{\log}^{r-1}
\qquad (r\geq 2)
\]
obtained from the transferred chiral product. Put
\[
\mathcal{V}^{\vee}:=
\mathcal{H}om_{\mathcal{O}_X}(\mathcal{V},\omega_X),
\qquad
\widehat{T}^{c}_{\mathrm{ch}}(s^{-1}\mathcal{V}^{\vee})
:=\prod_{n\geq 0}
\pi_{n*}\bigl(j_*j^*(s^{-1}\mathcal{V}^{\vee})^{\boxtimes n}\bigr)_{\Sigma_n},
\]
with coinvariants replaced by the ordered tensor coalgebra in the
$\Eone$ lane. The \emph{candidate bar-dual coalgebra}
$\mathcal{A}^{\mathrm i}_{\mathrm{cand}}$ is the closed conilpotent
subcoalgebra of $\widehat{T}^{c}_{\mathrm{ch}}
(s^{-1}\mathcal{V}^{\vee})$ cut out by the residue-orthogonal
corelations $R_{\geq2}^{\perp}$.

Its coproduct is the completed deconcatenation coproduct, with counit
given by projection to the length-zero summand. Its structure
coderivation
\[
q=\sum_{r\geq 1}q_r
\]
is the unique continuous coderivation whose projection to cogenerators
is dual, under Serre--residue pairing and the desuspension convention of
Appendix~\ref{app:signs-shifts}, to the transferred OPE residues
$\{\mu_r\}_{r\geq 1}$. The binary part is the dual of the chiral product
residue along the collision divisor; higher parts record the transferred
higher residues. If $q_r=0$ for $r\geq3$, this is a strict dg chiral
coalgebra. In general it is a completed conilpotent
$A_\infty$ chiral coalgebra, and the completed cobar functor is taken in
that $A_\infty$ sense.
\end{definition}

\begin{remark}[Independence]\label{rem:independence-construction}
The construction uses only the presentation of~$\mathcal{A}$, the
transferred OPE residues, and the residue pairing. It does not define
the Verdier-Koszul dual algebra. The algebra is obtained, when the
finite-type dual exists, from
$(\mathcal{A}^{\mathrm i}_{\mathrm{cand}})^\vee$ or equivalently from
the Verdier-dual bar package of Theorem~A.
\end{remark}

%================================================================
% COALGEBRA AXIOMS
%================================================================

\subsection{Coalgebra identities}

\begin{theorem}[Coalgebra structure on \texorpdfstring{$\mathcal{A}^{\mathrm i}_{\mathrm{cand}}$}{A-i-cand}; \ClaimStatusProvedHere]
\label{thm:coalgebra-axioms-verified}
Assume the transferred residues $\{\mu_r\}$ satisfy the Borcherds
identities and the Arnold--Orlik--Solomon relations with the signs of
Appendix~\ref{app:signs-shifts}, and assume the corelations
$R_{\geq2}^{\perp}$ are preserved by the coderivation~$q$ of
Definition~\ref{def:intrinsic-koszul-dual}. Then
\[
(\mathcal{A}^{\mathrm i}_{\mathrm{cand}},\Delta,\epsilon,q)
\]
is a completed conilpotent chiral coalgebra. In the strict case
$q_r=0$ for $r\geq3$, it satisfies
\begin{enumerate}
\item $(\Delta\boxtimes \mathrm{id})\Delta
 =(\mathrm{id}\boxtimes\Delta)\Delta$;
\item $(\epsilon\boxtimes\mathrm{id})\Delta
 =\mathrm{id}=(\mathrm{id}\boxtimes\epsilon)\Delta$;
\item $\Delta q=(q\boxtimes\mathrm{id}+\mathrm{id}\boxtimes q)\Delta$;
\item $q^2=0$.
\end{enumerate}
\end{theorem}

\begin{proof}
The completed tensor coalgebra is cofree conilpotent, hence
deconcatenation is coassociative and the projection to length~$0$ is a
counit. Since $q$ is defined as the continuous coderivation extending
its projection to cogenerators, the coderivation identity follows from
the cofree universal property. No separate four-term calculation on
primitive generators is needed; such a calculation misses the middle
deconcatenation terms unless one has already extended by the cofree
coderivation rule.

It remains to prove the square-zero identity. On a cofree conilpotent
coalgebra a coderivation is zero if and only if its projection to
cogenerators is zero. Hence $q^2=0$ is checked after projection to
$s^{-1}\mathcal{V}^{\vee}$. The projected square is the dual of the
sum of all iterated residue operations obtained by first colliding one
cluster of points and then a second cluster. In arity~$3$ these are the
three codimension-two boundary contributions in
$\overline{C}_3(X)$, and their logarithmic forms satisfy the Arnold
relation
\[
\eta_{12}\wedge\eta_{23}
\;+\;\eta_{23}\wedge\eta_{31}
\;+\;\eta_{31}\wedge\eta_{12}=0.
\]
With the desuspension signs fixed in Appendix~\ref{app:signs-shifts},
this is exactly the Borcherds identity for the OPE. Higher arities are
the same boundary-of-boundary identity on the Fulton--MacPherson
strata, written in Stasheff form. The preservation of
$R_{\geq2}^{\perp}$ makes the coderivation descend to the closed
subcoalgebra. Therefore $q^2=0$ in the strict case, and the general
case satisfies the completed $A_\infty$ coalgebra identities.
\end{proof}

\begin{remark}
The theorem proves only the coalgebra side. Passing from
$\mathcal{A}^{\mathrm i}_{\mathrm{cand}}$ to the algebra
$\mathcal{A}^!$ is a separate finite-type Verdier-duality operation.
\end{remark}

%================================================================
% BAR COMPARISON
%================================================================

\subsection{Bar comparison}

\begin{theorem}[Bar computes the intrinsic bar-dual coalgebra; \ClaimStatusProvedHere]
\label{thm:bar-computes-koszul-dual-complete}
\label{thm:chiral-koszul-duality}
\textup{[Regime: quadratic on the Koszul locus;
filtered-complete with completion
\textup{(}Convention~\textup{\ref{conv:regime-tags})}.]}

Let $(\mathcal{A}_1,\mathcal{A}_2)$ be a finite-type PBW chiral Koszul
pair, and let
$\mathcal{A}_{1,\mathrm{cand}}^{\mathrm i}$ be the intrinsic
bar-dual coalgebra of Definition~\ref{def:intrinsic-koszul-dual}.
Assume the bar-length/PBW filtration is complete, separated,
exhaustive, and Mittag--Leffler in every fixed conformal bidegree. Then
there is a natural filtered quasi-isomorphism of completed conilpotent
chiral coalgebras
\[
\Phi\colon
\widehat{\bar{B}^{\mathrm{ch}}(\mathcal{A}_1)}
\xrightarrow{\;\sim\;}
\mathcal{A}_{1,\mathrm{cand}}^{\mathrm i}.
\]
After continuous linear or Verdier duality on the finite-type lane this
gives the dual algebra
\[
(\mathcal{A}_{1,\mathrm{cand}}^{\mathrm i})^\vee
\simeq \mathcal{A}_1^!\simeq \mathcal{A}_2.
\]
The map $\Phi$ satisfies:
\begin{enumerate}
\item it respects coproduct, counit, and differential;
\item it is functorial for morphisms preserving the PBW filtration and
 residue pairing;
\item in the quadratic finite case it is the classical chiral
 quadratic Koszul comparison with no completion;
\item in PBW non-quadratic cases it is a completed comparison, not an
 identification of the uncompleted bar complex with the dual algebra.
\end{enumerate}
\end{theorem}

\begin{proof}
The bar complex has bar-length filtration, and
$\mathcal{A}_{1,\mathrm{cand}}^{\mathrm i}$ has the cogenerator-count
filtration. The residue--Serre pairing on each collision stratum gives
filtered maps
\[
\Phi_n\colon
\bar{B}^{\mathrm{ch}}_n(\mathcal{A}_1)
\longrightarrow
(\mathcal{A}_{1,\mathrm{cand}}^{\mathrm i})^{(n)}.
\]
These maps assemble because restriction to a boundary stratum in
$\overline{C}_{n}(X)$ is dual to deconcatenation, while residue along a
collision divisor is dual to the corresponding OPE component of~$q$.
Stokes' theorem on Fulton--MacPherson strata, together with the Arnold
relation and the signs of Appendix~\ref{app:signs-shifts}, gives
\[
\Phi\,d_{\mathrm{bar}}=q\,\Phi.
\]
Thus $\Phi$ is a filtered morphism of completed chiral coalgebras.

On associated graded objects the PBW filtration replaces
$\mathcal{A}_1$ by its classical quadratic associated graded algebra.
The target becomes the residue-orthogonal quadratic coalgebra
$(\operatorname{gr}\mathcal{A}_1)^{\mathrm i}$, and the associated
graded of~$\Phi$ is the usual quadratic Koszul comparison. Since
$(\mathcal{A}_1,\mathcal{A}_2)$ is a chiral Koszul pair, the PBW
spectral sequence collapses on the diagonal
(Theorem~\ref{thm:pbw-koszulness-criterion}), so this associated
graded comparison is a quasi-isomorphism. The comparison theorem for
complete filtered complexes then gives a quasi-isomorphism on the
completed objects. The Mittag--Leffler hypothesis is exactly the
condition that passage to the inverse limit does not create
$\varprojlim^1$ cohomology.
\end{proof}

\begin{lemma}[Completion convergence; \ClaimStatusProvedHere]\label{lem:completion-convergence}
Let $C$ be a bar-length filtered chiral complex with completion
\[
\widehat C=\varprojlim_n C/F^nC.
\]
Assume that for every fixed bar length and conformal weight the
filtered pieces are finite-dimensional, the transition maps in
$\{C/F^nC\}_n$ are eventually surjective, the filtration is complete
and separated, and the differential is continuous. Then
$R^1\varprojlim H^*(C/F^nC)=0$, the completed spectral sequence
converges strongly, and a filtered quasi-isomorphism on associated
graded objects induces a quasi-isomorphism after completion.
\end{lemma}

\begin{proof}
For fixed bidegree the inverse system is Mittag--Leffler, so
$\varprojlim^1=0$ by the standard derived-limit criterion
\cite[Proposition~3.5.7]{Weibel94}. Completeness and separatedness
identify $C$ with the limit of its quotients, and continuity of the
differential makes the completed object a complex. Strong convergence
of the completed spectral sequence follows from the same
Mittag--Leffler condition in each bidegree. A filtered map that is a
quasi-isomorphism on associated graded pages is therefore a
quasi-isomorphism on the completed abutments.
\end{proof}

\begin{remark}[Completion package]
Finite generation and polynomial growth are useful sufficient
conditions only when they imply the bidegreewise finite-dimensionality
and Mittag--Leffler stabilization stated in
Lemma~\ref{lem:completion-convergence}. They are not, by themselves, a
substitute for completion convergence. The full derived completion
package is developed in Appendix~\ref{app:nilpotent-completion}.
\end{remark}

\begin{corollary}[Circularity-free Koszul duality; \ClaimStatusProvedHere]
\label{cor:circularity-free-koszul}
Let $\mathcal{A}$ be a finite-type PBW chiral algebra on a smooth
curve~$X$ satisfying:
\begin{enumerate}[label=\textup{(\roman*)}]
\item\label{hyp:fg} $\mathcal{A} = T_{\mathrm{chiral}}(\mathcal{V})/(R)$ is finitely generated over $\mathcal{D}_X$,
\item\label{hyp:poly} the transferred OPE residues preserve the
 residue-orthogonal corelations and satisfy the Borcherds and
 Arnold identities with the signs of Appendix~\ref{app:signs-shifts},
\item\label{hyp:pbw-koszul} the associated graded quadratic chiral
 algebra for the PBW filtration is Koszul, so the PBW bar spectral
 sequence is diagonal at $E_2$,
\item\label{hyp:smooth} the bar-length completion is
 bidegreewise Mittag--Leffler in the sense of
 Lemma~\ref{lem:completion-convergence}.
\end{enumerate}
Then the intrinsic coalgebra
$\mathcal{A}^{\mathrm i}_{\mathrm{cand}}$ of
Definition~\ref{def:intrinsic-koszul-dual} is defined without
reference to the bar construction, and there is a natural
quasi-isomorphism
\[
\Phi\colon
\widehat{\bar{B}^{\mathrm{ch}}(\mathcal{A})}
\xrightarrow{\;\sim\;}
\mathcal{A}^{\mathrm i}_{\mathrm{cand}}
\]
of completed conilpotent chiral coalgebras. If the finite-type dual
exists, $(\mathcal{A}^{\mathrm i}_{\mathrm{cand}})^\vee$ is the
Verdier-Koszul dual algebra. For PBW non-quadratic families the
completed statement is the theorem; the uncompleted comparison is a
separate finite-length condition and is not automatic.
\end{corollary}

\begin{proof}
Definition~\ref{def:intrinsic-koszul-dual} constructs
$\mathcal{A}^{\mathrm i}_{\mathrm{cand}}$ from the presentation and
residue pairing. Theorem~\ref{thm:coalgebra-axioms-verified} gives its
coalgebra identities. Hypothesis~\ref{hyp:pbw-koszul} supplies the
associated-graded Koszul comparison, and
Hypothesis~\ref{hyp:smooth} removes the possible
$\varprojlim^1$ obstruction. Theorem~\ref{thm:bar-computes-koszul-dual-complete}
then supplies the completed bar comparison.
\end{proof}

\begin{remark}[Comparison with classical Koszul duality]\label{rem:comparison-classical}
When $\mathcal{A}$ is quadratic with relation space
$R \subset j_*j^*(\mathcal{V}^{\boxtimes 2})$, the bar-dual coalgebra is
the subcoalgebra of the cofree coalgebra cogenerated by
$R^\perp$, and its finite dual algebra is the quotient
$T_{\mathrm{ch}}(s\mathcal{V}^{\vee})/(R^\perp)$. Thus quotients on the
algebra side become subcoalgebras on the coalgebra side. In this
finite quadratic case the completion is discrete and the statement
reduces to the Ginzburg--Kapranov/Loday--Vallette quadratic
comparison. PBW non-quadratic examples require the filtered completed
comparison above.
\end{remark}

\begin{remark}[BD comparison]\label{rem:bd-comparison-koszul}
Theorem~\ref{thm:chiral-koszul-duality} is compatible with the
chiral-homology results of Beilinson--Drinfeld, but it is not stated in
the same category. BD work on the symmetric Ran-space side. The
symmetric bar $\barB_X(\cA)$ entering Theorem~A is the object
comparable to BD's factorization coalgebra; the ordered bar of this
monograph is an additional $\Eone$ refinement.

On the pole-free BD-commutative subclass,
\cite[Theorem~4.6.1]{BD04} identifies chiral homology with de~Rham
cohomology, and \cite[Theorem~4.8.1]{BD04} identifies the chiral
homology of a chiral envelope with the Chevalley complex of
$R\Gamma_{\mathrm{DR}}(X,L)$. Our Koszul statement uses the same
symmetric bar input in those lanes and then adds a duality statement not
present in BD: on the Koszul locus, the bar coalgebra resolves back to
the original chiral algebra and is paired with its Verdier-dual Koszul
partner.

The terminology bridge is strict. BD Chapter~4 ``commutative'' means a
pole-free commutative $\cD_X$-algebra. It does not mean the full
$\Einf$-chiral class used here. Thus the comparison with BD is exact on
the pole-free subclass and on the chiral-envelope/Chevalley lane; for
local vertex algebras with OPE poles, the present theorem extends past
BD's Chapter~4 scope by using the residue differential and the
Fulton--MacPherson boundary geometry.
\end{remark}

%================================================================
% EXPLICIT CALCULATIONS: W-ALGEBRAS
%================================================================

\section{Explicit calculations: $\mathcal{W}$-algebras and beyond}
\label{sec:w-algebras-explicit-completion}

%================================================================
% SUBSECTION: VIRASORO (WARM-UP)
%================================================================

\subsection{Virasoro Koszul-duality calculation}

\begin{example}[Virasoro: first non-quadratic example]
\label{ex:virasoro-completion-chiral}
The Virasoro algebra is generated by the stress tensor $T(z)$ with
OPE
\[
T(z)T(w)\sim
\frac{c/2}{(z-w)^4}
+\frac{2T(w)}{(z-w)^2}
+\frac{\partial T(w)}{z-w}.
\]
The quartic pole is not discarded by multiplying the OPE by a logarithmic
one-form and taking an ordinary one-variable residue. In the chiral bar
complex the Fulton--MacPherson boundary map first expands the fields
along the diagonal and then applies the residue operation. Equivalently,
the binary part of the bar coderivation records the Virasoro products
\[
T_{(3)}T=\frac{c}{2}\mathbf 1,\qquad
T_{(1)}T=2T,\qquad
T_{(0)}T=\partial T
\]
with the suspension signs fixed in Appendix~\ref{app:signs-shifts}.

Thus the completed bar-dual coalgebra is not the algebra
$\mathrm{Vir}_c^!$. It is the conilpotent coalgebra
\[
\mathrm{Vir}_c^{\mathrm i}
=H^*\!\left(
\widehat{\barB_X^\Sigma(\mathrm{Vir}_c)},
d_{\mathrm{bar}}
\right),
\]
or the corresponding ordered coalgebra on the $\Eone$ lane. It is
presented inside the completed cofree coalgebra on the descendants
\[
T,\partial T,\partial^2T,\ldots
\]
with a continuous coderivation whose arity-two components are the three
displayed OPE modes and their $\partial$-translates. In a fixed conformal
weight and bar length there are only finitely many descendants, and the
bar-length quotients satisfy Lemma~\ref{lem:completion-convergence}.

The post-Verdier algebra is
\[
\mathrm{Vir}_c^!=
\mathbb D_{\operatorname{Ran}}(\mathrm{Vir}_c^{\mathrm i}),
\]
or the continuous finite-window dual in the affine model. The scalar
central channel is normalized by
\[
\kappa(\mathrm{Vir}_c)=\frac{c}{2};
\]
it is a curvature component only after the curved genus package is
turned on. The slogan identifying $\mathrm{Vir}_c^!$ with the completed
tensor coalgebra $\widehat{T^c(T^*)}$ is therefore type-wrong: that
coalgebra is a presentation of the bar-dual object before Verdier
duality, not the Koszul-dual algebra itself.
\end{example}

%================================================================
% SUBSECTION: W_3 ALGEBRA - COMPLETE
%================================================================

\subsection{\texorpdfstring{$W_3$ algebra: explicit calculation}{W-3 algebra: explicit calculation}}

\begin{example}[\texorpdfstring{$W_3$}{W3} algebra: full completion]
\label{ex:w3-completion-full}
The $W_3$ algebra is strongly generated by $T(z)$ of weight~$2$ and
$W(z)$ of weight~$3$. Its complete OPEs are those of
Definition~\ref{def:w3-algebra} and Theorem~\ref{thm:w-w-ope-complete}.
The composite field
\[
\Lambda={:}TT{:}-\frac{3}{10}\partial^2T
\]
is the quasi-primary of weight~$4$, with
\[
\langle\Lambda|\Lambda\rangle=\frac{c(5c+22)}{10}
\]
as fixed in the landscape census. The $W$-$W$ OPE contains the
load-bearing channel
\[
\frac{32}{22+5c}\frac{\Lambda(w)}{(z-w)^2}
+\frac{16}{22+5c}\frac{\partial\Lambda(w)}{z-w},
\]
together with the central term $c/3\,(z-w)^{-6}$ and the Virasoro
descendant terms.

The completed bar object is
\[
\widehat{\barB_X^\Sigma(W_3)}
=\varprojlim_n\barB_X^\Sigma(W_3)/F^n_{\mathrm{bar}},
\]
and its differential is the continuous coderivation whose binary
components are the OPE modes of $T,T$, $T,W$, and $W,W$. No pole is
discarded by an ordinary residue test on
$(dz-dw)/(z-w)$. The sextic central pole is one of the scalar
components of the coderivation, while the $\Lambda$ channel supplies a
genuine quadratic composite:
\[
\Delta(\Lambda^\vee)
=\Lambda^\vee\boxtimes 1+1\boxtimes\Lambda^\vee
+T^\vee\boxtimes T^\vee
\]
up to the standard suspension and normal-ordering signs.

The intrinsic bar-dual coalgebra is
\[
W_3^{\mathrm i}
=H^*\!\left(\widehat{\barB_X^\Sigma(W_3)},d_{\mathrm{bar}}\right),
\]
or the ordered version in the $\Eone$ refinement. It is represented by a
closed subcoalgebra of the completed cofree conilpotent coalgebra on the
duals of
\[
T,\partial T,\partial^2T,\ldots,\qquad
W,\partial W,\partial^2W,\ldots
\]
together with the composite coalgebra coordinates forced by the OPE,
beginning with $\Lambda^\vee$. The Koszul-dual algebra is not this
cofree coalgebra. It is
\[
W_3^!=\mathbb D_{\operatorname{Ran}}(W_3^{\mathrm i})
\]
when the Verdier finite-type hypotheses are satisfied. Completion
converges only under the bidegreewise finite-dimensional and
Mittag--Leffler hypotheses of Lemma~\ref{lem:completion-convergence};
finite strong generation alone is not a substitute.
\end{example}

%================================================================
% SUBSECTION: GENERAL W_N
%================================================================

\subsection{\texorpdfstring{General $W_N$ algebras}{General W-N algebras}}

\begin{construction}[$W_N$ bar-dual coalgebra]
\label{const:wn-general}

For $W_N$ with strong generators
$\{W^{(2)},W^{(3)},\ldots,W^{(N)}\}$ of weights
$2,3,\ldots,N$, the completed bar complex is
\[
\widehat{\barB_X^\Sigma(W_N)}
=\varprojlim_n \barB_X^\Sigma(W_N)/F^n_{\mathrm{bar}},
\]
with
\[
\barB^k(W_N)=
\bigoplus_{\substack{i_1,\ldots,i_{k+1}\in\{2,\ldots,N\}\\
m_1,\ldots,m_{k+1}\ge0}}
\Gamma\!\left(
\overline C_{k+1}(X),
\partial^{m_1}W^{(i_1)}\boxtimes\cdots\boxtimes
\partial^{m_{k+1}}W^{(i_{k+1})}\otimes\Omega^k_{\log}
\right).
\]
The coderivation is the continuous extension of all OPE modes among the
strong generators and their descendants.

The intrinsic coalgebra is
\[
W_N^{\mathrm i}
=H^*\!\left(\widehat{\barB_X^\Sigma(W_N)},d_{\mathrm{bar}}\right).
\]
It is not the algebra $W_N^!$. In a presentation it sits inside the
completed conilpotent coalgebra cogenerated by the duals of the strong
generators, their $\partial$-descendants, and the normal-ordered
composite coordinates forced by the OPE closure relations. Its structure
has the following features.
\begin{enumerate}
\item The strong-generator coordinates $(W^{(s)})^\vee$ have dual
 weight $-s$ before the suspension shift.
\item The coproduct is primitive only on genuine cogenerators; composite
 coordinates have reduced coproducts determined by normal ordering.
\item The differential records all OPE structure constants, including
 central high-pole channels and composite fields.
\item In the curved genus package the scalar normalization is governed
 by
 \[
 \kappa(\Walg_N)=c(H_N-1),\qquad
 H_N=\sum_{j=1}^{N}\frac1j,
 \]
 and this vanishes at $c=0$, not at the minimal-model values in
 general.
\end{enumerate}
The post-Verdier algebra is
\[
W_N^!=\mathbb D_{\operatorname{Ran}}(W_N^{\mathrm i})
\]
under the same finite-type and completion hypotheses as in
Corollary~\ref{cor:circularity-free-koszul}.
\end{construction}

\begin{table}[ht]
\centering
\caption{$W_N$ completion complexity}
\begin{tabular}{|c|c|c|c|}
\hline
$N$ & Generators & Max pole order & $\dim(\bar{B}^1/I^2)$ \\
\hline
2 (Virasoro) & 1 & 4 & 1 \\
3 & 2 & 6 & 3 \\
4 & 3 & 8 & 6 \\
5 & 4 & 10 & 10 \\
$N$ & $N-1$ & $2N$ & $\binom{N}{2}$ \\
\hline
\end{tabular}
\end{table}

%================================================================
% SUBSECTION: BEYOND W-ALGEBRAS
%================================================================

\subsection{Beyond $\mathcal{W}$-algebras: other non-quadratic examples}

\begin{example}[Affine Yangian]\label{ex:affine-yangian-kp}
\label{ex:yangian-completion}

The affine Yangian $Y_\hbar(\widehat{\mathfrak{gl}}_n)$ belongs to the
ordered $\Eone$ lane. Its spectral-parameter relations are read by
$\barB_X^{\mathrm{ord}}$, not by the symmetric Ran bar before averaging.
In a finite Drinfeld-current or RTT window the generators may be written
schematically as $T^a_{(r)}(u)$, with $a=1,\ldots,n^2$ and $r\ge0$; the
relations have rational coefficients in $u-v$ and $\hbar$ and are
completed at the chosen spectral pole.

The correct completed bar object is the finite-window inverse limit
\[
\widehat{\barB_X^{\mathrm{ord}}(Y_\hbar)}
=\varprojlim_{p,q,m}
\barB_X^{\mathrm{ord}}(Y_\hbar)/
\bigl(F_{\mathrm{bar}}^p+\hbar^q+\mathfrak m_u^m\bigr),
\]
where $\mathfrak m_u$ is the spectral-parameter maximal ideal at the
expansion point. The quotient is by the sum of the three closed
filtration ideals, not by the product $I^p\hbar^q$. The intrinsic object
obtained from this complex is
$Y_\hbar(\widehat{\mathfrak{gl}}_n)^{\mathrm i}$. A statement about
$Y_\hbar(\widehat{\mathfrak{gl}}_n)^!$ requires continuous Verdier
duality on these finite windows and the bidegreewise
Mittag--Leffler convergence of Lemma~\ref{lem:completion-convergence}.
\end{example}

\begin{example}[Bershadsky--Polyakov algebra]
\label{ex:bp-algebra}

The $W_3^{(2)}$ algebra (also called Bershadsky--Polyakov) is the quantum Drinfeld--Sokolov reduction of $\widehat{\mathfrak{sl}}_3$ at the subregular nilpotent. It has generators:
\begin{itemize}
\item A weight-1 field $J(z)$
\item Two weight-$3/2$ fermionic fields $G^{\pm}(z)$
\item The Virasoro field $T(z)$ (weight 2)
\end{itemize}

The structure constants involve \emph{irrational functions} of the central charge (e.g., $\sqrt{c}$ or $\sqrt{5c+22}$ appear in OPE coefficients), although the OPE poles themselves are always integer-order as required by vertex algebra locality. This requires:
\begin{itemize}
\item Working over $\mathbb{C}[c^{1/2}]$ (not $\mathbb{C}$) when treating $c$ as a formal parameter
\item Completion in the $\mathfrak{m}$-adic topology where $\mathfrak{m} = (c^{1/2} - c_0^{1/2})$
\item Careful treatment of the fermionic generators and their $\mathbb{Z}/2$-grading
\end{itemize}
\end{example}

\begin{example}[Superconformal algebras]
\label{ex:n2-superconformal}

The $\mathcal{N}=2$ superconformal algebra has:
\begin{itemize}
\item Bosonic: $T(z)$ weight 2, $J(z)$ weight 1
\item Fermionic: $G^+(z)$, $G^-(z)$ weight $3/2$
\end{itemize}

The $\mathbb{Z}/2\mathbb{Z}$-grading introduces fermionic signs:
\[G^+(z)G^+(w) \sim 0, \quad G^+(z)G^-(w) \sim \frac{c/3}{(z-w)^3} + \frac{J(w)}{(z-w)^2} + \cdots\]

\emph{Completion.} Requires \emph{super-coalgebra} structure:
\[\Delta(G^+) = G^+ \boxtimes 1 + 1 \boxtimes G^+\]
(generators are primitive, as in Step~2 above; the fermionic sign enters in the \emph{cocommutativity} constraint $\tau \circ \Delta = \Delta$ where $\tau(x \boxtimes y) = (-1)^{|x||y|} y \boxtimes x$).
\end{example}

%================================================================
% SUMMARY TABLE
%================================================================

\begin{table}[ht]
\centering
\caption{Non-quadratic examples}
\begin{tabular}{|l|c|c|l|}
\hline
\textbf{Chiral Algebra} & \textbf{Max Pole} & \textbf{Completion Type} & \textbf{Feature} \\
\hline
Virasoro & 4 & I-adic & Quartic central term \\
$W_3$ & 6 & I-adic & Sextic pole, composite fields \\
$W_N$ & $2N$ & I-adic & Complexity $\sim N^2$ \\
Affine Yangian & $\infty$ & Double (I + $\hbar$) & Spectral parameter \\
Bershadsky--Polyakov & 4 & I + $\sqrt{c}$-adic & Fractional exponents \\
$\mathcal{N}=2$ Super & 3 & I-adic (super) & Fermion signs \\
\hline
\end{tabular}
\end{table}

\begin{remark}[Feynman diagrams]\label{sec:feynman_genus_g}
The Feynman diagram interpretation of the bar-cobar complex (the
identification of bar chains with off-shell amplitudes and cobar chains
with on-shell propagator kernels) is developed in
Chapter~\ref{ch:v1-feynman}.
\end{remark}

% ================================================================
% SECTION: E1-CHIRAL KOSZUL DUALITY
% ================================================================
\section{\texorpdfstring{The $\Eone$-chiral Koszul duality theorem}{The E1-chiral Koszul duality theorem}}

The preceding sections developed the bar-cobar adjunction and Koszul
pairs for $\Einf$-chiral algebras (vertex algebras), where the
underlying operad is the commutative chiral operad $\chirCom$.
The \emph{associative} case treats algebras over the
associative chiral operad $\chirAss$
\index{chiral associative operad|textbf}
(Definition~\ref{def:chiral-ass-operad}). This is the
setting of nonlocal vertex algebras
(Definition~\ref{def:e1-chiral}), and the resulting Koszul duality
is self-dual by Proposition~\ref{prop:chirAss-self-dual}.
The $\infty$-categorical bar-cobar formalism for algebras over operads
that underlies the construction is developed in~\cite[§3.1]{HA}.

\subsection{\texorpdfstring{Augmented and pro-nilpotent $\Eone$-chiral algebras}{Augmented and pro-nilpotent E1-chiral algebras}}

\begin{definition}[Augmented \texorpdfstring{$\Eone$}{E1}-chiral algebra]
\label{def:augmented-e1}
An $\Eone$-chiral algebra $\mathcal{A}$ is \emph{augmented} if it
is equipped with a morphism of $\Eone$-chiral algebras
$\varepsilon: \mathcal{A} \to \omega_X$ (the unit chiral algebra).
The \emph{augmentation ideal} is
$\bar{\mathcal{A}} := \ker(\varepsilon)$.
\end{definition}

\begin{definition}[Pro-nilpotent and conilpotent]
\label{def:pronilpotent-conilpotent}
An augmented $\Eone$-chiral algebra $\mathcal{A}$ is
\emph{pro-nilpotent} if the inverse system
\[
\mathcal{A} \;\twoheadrightarrow\;
\mathcal{A}/\bar{\mathcal{A}}^{\chirtensor 2} \;\twoheadrightarrow\;
\mathcal{A}/\bar{\mathcal{A}}^{\chirtensor 3} \;\twoheadrightarrow\;
\cdots
\]
is pro-isomorphic to $\mathcal{A}$, where
$\bar{\mathcal{A}}^{\chirtensor n}$ denotes the $n$-fold chiral
tensor power of the augmentation ideal.

Dually, a coaugmented $\Eone$-chiral coalgebra $\mathcal{C}$ is
\emph{conilpotent} if
$\mathcal{C} = \bigcup_{n \geq 0} F_n\mathcal{C}$
where $F_n\mathcal{C}$ is the coradical filtration
(the largest sub-coalgebra on which the iterated reduced coproduct
$\bar{\Delta}^{(n)}$ vanishes).
\end{definition}

\begin{remark}[Convergence role of pro-nilpotence]\label{rem:convergence-pro-nil}
\label{rem:convergence-pronilpotent}
Pro-nilpotence ensures the cobar differential
$d_\Omega: \Omega(\mathcal{C})_n \to \Omega(\mathcal{C})_{n-1}$
is well-defined: the formula involves a sum over all ways to
decompose a tensor factor via the coproduct, which in general is an
infinite series. Conilpotence ensures this sum is finite in each
degree. On the algebra side, pro-nilpotence ensures the bar
construction is degree-wise finite-dimensional.
See Appendix~\ref{app:nilpotent-completion} for the completion
theory.
\end{remark}

\subsection{\texorpdfstring{The bar-cobar equivalence for $\chirAss$}{The bar-cobar equivalence for Ass ch}}

\begin{theorem}[\texorpdfstring{$\Eone$}{E1}-chiral Koszul duality; \ClaimStatusConditional]
\label{thm:e1-chiral-koszul-duality}
\index{E1-Koszul duality@$\mathbb{E}_1$-Koszul duality|textbf}
\textup{[Regime: quadratic on the Koszul locus;
filtered-complete via PBW degeneration
\textup{(}Convention~\textup{\ref{conv:regime-tags})}.]}

On the finite-type complete lane, the bar-cobar adjunction restricts to
an equivalence of $\infty$-categories. For the underlying $E_1$ bar
construction in the $\infty$-categorical setting, see
\cite[§5.2.2]{HA}.

The equivalence
is
\[
\bar{B}^{\mathrm{ch}}:
\chirAss\text{-}\mathrm{Alg}^{\mathrm{aug,\, pronil,\, ft}}
\;\rightleftarrows\;
\chirAss\text{-}\mathrm{CoAlg}^{\mathrm{coaug,\, conil,\, ft}}
:\Omega^{\mathrm{ch}},
\]
where:
\begin{itemize}
\item $\chirAss\text{-}\mathrm{Alg}^{\mathrm{aug,\, pronil,\, ft}}$
 is the $\infty$-category of augmented, pro-nilpotent
 $\Eone$-chiral algebras on $X$ whose bar filtration is complete,
 separated, bidegreewise finite, and compatible with the chiral tensor
 product;
\item $\chirAss\text{-}\mathrm{CoAlg}^{\mathrm{coaug,\, conil,\, ft}}$
 is the $\infty$-category of coaugmented, conilpotent
 $\Eone$-chiral coalgebras on $X$ satisfying the dual cobar
 convergence conditions.
\end{itemize}

The unit and counit of the adjunction are quasi-isomorphisms:
\begin{align}
\Omega^{\mathrm{ch}}(\bar{B}^{\mathrm{ch}}(\mathcal{A}))
&\xrightarrow{\;\sim\;} \mathcal{A}
\quad \text{for all } \mathcal{A} \in
\chirAss\text{-}\mathrm{Alg}^{\mathrm{aug,\, pronil,\, ft}},
\label{eq:e1-counit} \\
\mathcal{C}
&\xrightarrow{\;\sim\;}
\bar{B}^{\mathrm{ch}}(\Omega^{\mathrm{ch}}(\mathcal{C}))
\quad \text{for all } \mathcal{C} \in
\chirAss\text{-}\mathrm{CoAlg}^{\mathrm{coaug,\, conil,\, ft}}.
\label{eq:e1-unit}
\end{align}
\end{theorem}

\begin{proof}
The argument is the operadic bar-cobar theorem transported to the
finite-type chiral tensor category.

\emph{Step 1: Reduction to the classical operadic statement.}
By Proposition~\ref{prop:chirAss-self-dual}, the chiral associative
operad $\chirAss$ is the chiral lift of the classical associative
operad $\operatorname{Ass}$. The bar-cobar adjunction for a Koszul
operad $\mathcal{P}$ in a stable presentably symmetric monoidal
$\infty$-category $\mathcal{V}$ is an equivalence
$\mathcal{P}\text{-Alg}^{\mathrm{aug,pronil}}(\mathcal{V})
\simeq
\mathcal{P}^!\text{-CoAlg}^{\mathrm{coaug,conil}}(\mathcal{V})$
whenever $\mathcal{P}$ is Koszul \cite[Theorem~11.4.1]{LV12}.
We apply this with $\mathcal{V} = \mathcal{D}\text{-mod}(X)$
(the stable $\infty$-category of $\mathcal{D}$-modules on~$X$,
equipped with the chiral tensor product) and
$\mathcal{P} = \operatorname{Ass}$.

\emph{Step 2: Koszulity of $\operatorname{Ass}$ in the chiral
tensor category.}
By Lemma~\ref{lem:operadic-koszul-transfer} below
(applied with $\mathcal{P} = \operatorname{Ass}$ and
$\mathcal{V} = \mathcal{D}\text{-mod}(X)$),
the classical Koszulity of
$\operatorname{Ass}$ \cite[Theorem~7.4.1]{LV12} transfers to the
chiral tensor category: the operadic bar complex
$B_{\operatorname{Ass}}(k)$ is quasi-isomorphic to
$\operatorname{Ass}^!$ after base change to~$\mathcal{V}$.

\emph{Step 3: Geometric realization via configuration spaces.}
The bar complex $\bar{B}^{\mathrm{ch}}(\mathcal{A})$ has an
explicit geometric model (Definition~\ref{def:geometric-bar}):
\[
\bar{B}^{\mathrm{ch}}_n(\mathcal{A})
= \Gamma\!\left(\overline{C}_n(X),\;
j_* j^*\, \mathcal{A}^{\boxtimes n}
\otimes \Omega^{n-1}_{\overline{C}_n(X)}(\log D)\right).
\]
The key difference from the $\Einf$-chiral case is that the
$\Sigma_n$-action is \emph{not} imposed: the $n$-th bar component
is indexed by \emph{ordered} configurations $(z_1, \ldots, z_n)$,
not unordered ones. The differential is the same alternating sum of
face maps (collisions of adjacent points), reflecting the
non-$\Sigma$ nature of the associative operad.

For the cobar construction, by
Corollary~\ref{cor:cobar-nilpotence-verdier}, we have
$d_{\mathrm{cobar}}^2 = 0$ via Verdier duality. The counit map
$\Omega^{\mathrm{ch}}(\bar{B}^{\mathrm{ch}}(\mathcal{A}))
\to \mathcal{A}$ is a quasi-isomorphism by the acyclicity of the
augmented bar complex (a standard contraction argument using the
augmentation $\varepsilon$, as in
Theorem~\ref{thm:bar-cobar-inversion-qi}), and the unit map is a
quasi-isomorphism by the dual argument for conilpotent coalgebras.
\end{proof}

\begin{lemma}[Operadic Koszulness transfer \cite{LV12}; \ClaimStatusProvedElsewhere]
\label{lem:operadic-koszul-transfer}
\index{Koszul operad!transfer to symmetric monoidal categories}
Let\/ $\mathcal{P}$ be a quadratic operad over a field\/~$k$, and let\/
$\mathcal{V}$ be a stable presentably symmetric monoidal\/
$\infty$-category over\/~$k$. If\/ $\mathcal{P}$ is Koszul in the
classical sense \textup{(}i.e., over the category of\/ $k$-modules\textup{)},
then\/ $\mathcal{P}$ is Koszul as an operad internal to\/~$\mathcal{V}$:
the operadic bar construction\/
$B_{\mathcal{P}}(k) \xrightarrow{\;\sim\;} \mathcal{P}^!$
is a quasi-isomorphism of cooperads in\/ $\Sigma$-objects of\/~$\mathcal{V}$.
\end{lemma}

\begin{proof}
Koszulness of a quadratic operad $\mathcal{P} = \mathcal{F}(V)/(R)$
is equivalent to the acyclicity of the Koszul complex
$\mathcal{P} \circ_\kappa \mathcal{P}^!$, where $\kappa \colon
\mathcal{P}^! \to \mathcal{P}$ is the universal twisting morphism
\cite[Theorem~7.4.2]{LV12}. We explain why this property, once
established over~$k$\nobreakdash-$\mathrm{Mod}$, holds automatically
in~$\mathcal{V}$.

\emph{Step~1: The Koszul complex is built from operadic composition.}
The Koszul complex $\mathcal{P} \circ_\kappa \mathcal{P}^!$ is
defined as the composition product
$\bigoplus_{n \geq 0} \mathcal{P}(n) \otimes_{\Sigma_n}
(\mathcal{P}^!)^{\otimes n}$ equipped with a differential twisted
by~$\kappa$. Both the composition product and the
twisting differential are determined by the \emph{quadratic data}
$(V, R)$: the generators~$V$, the space of
relations~$R \subset \mathcal{F}(V)(3)$, and the combinatorics
of planar trees indexing operadic compositions
\cite[§7.4]{LV12}. No structure of the ambient symmetric monoidal
category enters beyond the tensor product and direct sums indexed
by these combinatorial data.

\emph{Step~2: Base change preserves acyclicity.}
Given a symmetric monoidal functor $k\text{-}\mathrm{Mod} \to
\mathcal{V}$ (the unit functor $M \mapsto M \otimes \mathbf{1}_{\mathcal{V}}$),
the Koszul complex over~$\mathcal{V}$ is the image
$(\mathcal{P} \circ_\kappa \mathcal{P}^!) \otimes_k \mathbf{1}_{\mathcal{V}}$.
Since~$\mathcal{V}$ is $k$-linear (presentably enriched over~$k$),
this base-change functor is exact: it preserves all colimits and
finite limits (being a left adjoint between stable categories),
hence preserves quasi-isomorphisms. The acyclicity of the
Koszul complex over $k$\nobreakdash-$\mathrm{Mod}$
\cite[Theorem~7.4.1]{LV12} therefore implies acyclicity
over~$\mathcal{V}$.

\emph{Step~3: From Koszul complex to bar-cobar resolution.}
By \cite[Theorem~11.3.3]{LV12}, acyclicity of the Koszul complex
is equivalent to the bar-cobar resolution being a quasi-isomorphism:
$\Omega(\mathcal{P}^{\text{\textexclamdown}}) \xrightarrow{\sim}
\mathcal{P}$ as operads. This equivalence is formal
(it holds in any abelian category where the Koszul complex is
defined), so it transfers together with Step~2.
The $\infty$-categorical upgrade follows from
\cite[Theorem~11.4.1]{LV12} (see also
\cite[§16]{Fresse-operads} for the model-categorical treatment
of the bar-cobar adjunction for algebras over Koszul operads in
general symmetric monoidal model categories).
\end{proof}

\begin{remark}[Verification of hypotheses]\label{rem:e1-koszul-hypotheses}
Steps~1--2 of the proof invoke \cite[Theorem~11.4.1]{LV12} and
Lemma~\ref{lem:operadic-koszul-transfer} with
$\mathcal{V} = \mathcal{D}\text{-mod}(X)$ equipped with the chiral
tensor product~$\otimes^{\mathrm{ch}}$. We verify the required
hypotheses on~$\mathcal{V}$:
\begin{enumerate}
\item \emph{Stable presentable $\infty$-category.}
The $\infty$-category $\mathcal{D}\text{-mod}(X)$ is stable and
presentable by \cite[Chapter~I.3]{GR17}; see also
\cite[Proposition~5.5.3.6]{HA} and \cite[§1.1]{HA}.
\item \emph{Symmetric monoidal structure.}
The chiral tensor product endows $\mathcal{D}\text{-mod}(X)$ with a
symmetric monoidal structure by \cite[3.4.10]{BD04}. While
$\otimes^{\mathrm{ch}}$ does not in general coincide with the
$!$-tensor product, it does define a symmetric monoidal structure on
the $\infty$-categorical enhancement, as established in
\cite[Chapter~IV.5]{GR17}.
\item \emph{Exactness of base change
\textup{(}Lemma~\textup{\ref{lem:operadic-koszul-transfer}},
Step~\textup{2)}.}
The unit functor $k\text{-}\mathrm{Mod} \to \mathcal{D}\text{-mod}(X)$
is exact because $\mathcal{D}\text{-mod}(X)$ is $k$-linear and
stable. More concretely, $\mathcal{D}\text{-mod}(X)$ carries a
$t$-structure for which $\otimes^{\mathrm{ch}}$ is right-exact on each
factor \cite[3.4.12]{BD04}, and the bar filtration produces complexes
bounded below (in the $t$-structure sense), so the derived tensor
product agrees with the underived one on each filtered piece.
\end{enumerate}
\end{remark}

\begin{example}[Heisenberg specialization]\label{ex:heisenberg-e1-duality}
For the Heisenberg algebra $\mathcal{H}_k$,
Theorem~\ref{thm:e1-chiral-koszul-duality} simplifies dramatically:
$\mathcal{H}_k$ is $\Einf$-chiral as a local vertex algebra, but it is
not BD-commutative: the OPE $J(z)J(w) \sim k/(z-w)^2$ still has a
double pole. Hence the $\Eone$ ordered bar descends to the symmetric
$\Einf$ bar through the OPE-derived $R$-matrix datum; only for the
pole-free BD-commutative subclass is the $\Sigma_n$-quotient on ordered
configurations a literal isomorphism. The bar components
$\bar{B}_n^{\mathrm{ch}}(\mathcal{H}_k)$ are therefore indexed by
\emph{unordered} configurations after imposing this descent datum.
The counit quasi-isomorphism
$\Omega^{\mathrm{ch}}(\bar{B}^{\mathrm{ch}}(\mathcal{H}_k))
\xrightarrow{\sim} \mathcal{H}_k$
was verified explicitly in
Chapter~\ref{ch:heisenberg-frame}:
the bar complex
(\S\ref{sec:frame-bar-deg1}--\S\ref{sec:frame-bar-all})
recovers the intrinsic bar-dual coalgebra
$\mathcal H_k^{\mathrm i}$; finite Verdier duality identifies the
post-Verdier algebra with $\mathrm{Sym}^{\mathrm{ch}}(V^*)$, and the
cobar of the coalgebra reconstructs $\mathcal{H}_k$.
The spectral discriminant is trivial ($\Delta_{\mathcal{H}_k}(x)=1$,
rank~$1$), reflecting that no non-trivial higher operations arise.
\end{example}

\begin{corollary}[\texorpdfstring{$\Eone$}{E1}--\texorpdfstring{$\Eone$}{E1} Self-Duality; \ClaimStatusConditional]
\label{cor:e1-self-duality}
\tolerance=1500
Since
$(\chirAss)^! \cong \chirAss \otimes \operatorname{sgn}$
(Proposition~\ref{prop:chirAss-self-dual}), the ordered bar
of an augmented pro-nilpotent finite-type
$\Eone$\nobreakdash-chiral algebra~$\mathcal{A}$ is its
bar-dual coalgebra
\[
\mathcal{A}^{\mathrm i}_{\mathrm{ord}}
\coloneqq \bar{B}^{\mathrm{ch}}(\mathcal{A}).
\]
This coalgebra is not the dual algebra. When the continuous graded dual
exists, set
\[
\mathcal{A}^{!}_{\mathrm{ord}}
\coloneqq (\mathcal{A}^{\mathrm i}_{\mathrm{ord}})^\vee .
\]
The double-dual map
\[
\mathcal{A} \;\xrightarrow{\;\sim\;}\;
(\mathcal{A}^{!}_{\mathrm{ord}})^!
\]
is a quasi-isomorphism. This is the $\Eone$-chiral manifestation
of the classical self-duality $\operatorname{Ass}^! \cong
\operatorname{Ass} \otimes \operatorname{sgn}$.
\end{corollary}

\begin{proof}
The double-dual map factors as
$\mathcal{A} \xrightarrow{\sim}
\Omega^{\mathrm{ch}}(\bar{B}^{\mathrm{ch}}(\mathcal{A}))
= \Omega^{\mathrm{ch}}(\mathcal{A}^{\mathrm i}_{\mathrm{ord}})
\xrightarrow{\sim}(\mathcal{A}^{!}_{\mathrm{ord}})^!$,
where the first map is the counit quasi-isomorphism
\eqref{eq:e1-counit}. The second is the finite-type identification
obtained by applying continuous duality to the bar-dual coalgebra and
then using the unit-counit equivalence of
Theorem~\ref{thm:e1-chiral-koszul-duality}.
\end{proof}

\begin{remark}[Comparison with Francis--Gaitsgory]\label{rem:FG-comparison}
Francis--Gaitsgory~\cite{FG12} establish chiral Koszul duality for
the $\chirLie$--$\chirCom$ pair (mechanism~(i) of
Remark~\ref{rem:three-koszul-mechanisms}).
Theorem~\ref{thm:e1-chiral-koszul-duality} is the \emph{associative}
counterpart, operating within the single operadic type $\chirAss$.
The two results are complementary:
\begin{center}
\begin{tabular}{c|c|c}
& \textbf{Operadic pair} & \textbf{Algebra $\leftrightarrow$ Coalgebra} \\
\hline
Francis--Gaitsgory & $\chirLie \dashv \chirCom$ &
 $\chirLie$-alg $\simeq$ $\chirCom$-coalg \\
This work & $\chirAss \dashv \chirAss$ &
 $\chirAss$-alg $\simeq$ $\chirAss$-coalg
\end{tabular}
\end{center}
The self-dual case is simpler (the bar and cobar
land in the same operadic type) but has not been treated in the
chiral setting in the existing literature.
\end{remark}
\begin{remark}[Disk bar model and wrapped Fukaya categories; \ClaimStatusConditional]
\label{rem:fukaya-bar-koszul}
Abouzaid~\cite{AbouzaidGeneration} proves that a split-generating
Lagrangian~$L$ in the wrapped Fukaya category $\mathcal{W}(M)$ has its
wrapped endomorphism algebra computed from a disk model.
Under a chosen wrapped-generation and orientation package, the ordered
bar \(B^{\mathrm{ord}}(\mathcal A)\) on the disk is an algebraic model
for the Floer polygon complex: the bar differential records the same
$A_\infty$ polygon operations that Fukaya--Oh--Ohta--Ono~\cite{FOOO09}
define by counts of pseudoholomorphic $(n+1)$-gons with boundary
on~$L$.
Ganatra--Pardon--Shende~\cite{GPSSectorialDescent} show that sectorial
descent for wrapped Fukaya categories identifies this bar model with
$\mathcal{W}(M)$ itself under their hypotheses. In this identification
the bar coalgebra controls generation:
$\Omega(B^{\mathrm{ord}}(\mathcal{A})) \xrightarrow{\sim} \mathcal{A}$
becomes Floer-theoretic generation. A further homological-mirror
identification between compact and wrapped sectors requires the
corresponding HMS equivalence and split-generation assumptions
\cite{KontsevichHMS,SeidelBook}; it is not a formal consequence of the
chiral bar-cobar adjunction alone.
The objects remain distinct:
\(\mathcal A\) is the open-string $A_\infty$ algebra,
\(B^{\mathrm{ord}}(\mathcal A)\) is the bar coalgebra,
\(\mathcal A^{\mathrm i}\) is the bar-dual coalgebra,
\(\mathcal A^!\) is the finite-type Verdier/Koszul dual algebra when
that dual exists, and
$Z^{\mathrm{der}}_{\mathrm{ch}}(\mathcal{A})=\ChirHoch^*(\mathcal A)$
is the cochain-level Hochschild closed-sector actor. It is not the
physical closed-string bulk unless a separate open-closed comparison
identifies the physical local operators with this cochain object.
\end{remark}

% ================================================================
% SECTION 8.5: CATEGORIES OF MODULES AND DERIVED EQUIVALENCES
% ================================================================

\section{Categories of modules and derived equivalences}

\subsection{The fundamental theorem for chiral Koszul pairs}

\begin{remark}[General chiral module-category duality remains frontier]
\label{rem:module-category-frontier}
The broad ordinary-derived package for a general
chiral Koszul pair $(\cA_1,\cA_2)$
\textup{(}derived equivalence, Ext exchange with derived Hom on the
bar-comodule side and finite-type Tor reformulations, simple/projective
transport, and Hochschild transport\textup{)}
is \emph{not} proved on this general surface. The bar-cobar machinery
in Part~\ref{part:bar-complex} gives the intrinsic bar-coalgebra
comparison, and the proved ordinary module-level statement is the
$\Eone$ theorem
Theorem~\ref{thm:e1-module-koszul-duality}, restricted to the
quadratic genus-$0$ complete/conilpotent lane.

Any extension of that package to arbitrary chiral Koszul pairs would
require extra hypotheses and a module-level comparison theorem beyond
what is currently established here.
Nor does Hochschild transport reconstruct a physical bulk
factorization algebra: it gives the cochain-level closed-sector actor
$Z^{\mathrm{der}}_{\mathrm{ch}}(\cA)$ only after the chiral Hochschild
comparison is installed. A physical bulk requires a separate
bulk-to-boundary local-operator map and a quasi-isomorphism statement.
\end{remark}

\subsection{\texorpdfstring{$\Eone$-chiral module category Koszul duality}{E1-chiral module category Koszul duality}}

\begin{theorem}[\texorpdfstring{$\Eone$}{E1}-module category Koszul duality;
\ClaimStatusConditional]
\label{thm:e1-module-koszul-duality}
\index{module Koszul duality|textbf}
Let $\mathcal{A}$ be a Koszul $\Eone$-chiral algebra with Koszul
dual bar coalgebra
$C_{\mathcal{A}} := \bar{B}^{\mathrm{ch}}(\mathcal{A})$
\textup{(}Theorem~\ref{thm:e1-chiral-koszul-duality},
Corollary~\ref{cor:e1-self-duality}\textup{)}. When the finite-type
graded dual exists, we write $\mathcal{A}^! := C_{\mathcal{A}}^\vee$
for the corresponding dual algebra. Then the bar-cobar adjunction
restricts to an equivalence of derived categories between
\emph{complete} (pro-nilpotent) $\mathcal{A}$-modules and
\emph{conilpotent} $C_{\mathcal{A}}$-comodules:
\[
D(\mathrm{Mod}^{\Eone,\,\mathrm{compl}}_{\mathcal{A}})
\;\simeq\;
D(\mathrm{CoMod}^{\Eone,\,\mathrm{conil}}_{C_{\mathcal{A}}})
\]
given by the bar-cobar induction-restriction adjunction:
\begin{align*}
\mathrm{Ind}\colon \mathrm{Mod}^{\Eone,\,\mathrm{compl}}_{\mathcal{A}}
&\;\longrightarrow\;
\mathrm{CoMod}^{\Eone,\,\mathrm{conil}}_{C_{\mathcal{A}}}, &
\mathcal{M} &\;\longmapsto\;
\bar{B}^{\mathrm{ch}}(\mathcal{A}) \otimes_{\mathcal{A}} \mathcal{M},
\\
\mathrm{Res}\colon \mathrm{CoMod}^{\Eone,\,\mathrm{conil}}_{C_{\mathcal{A}}}
&\;\longrightarrow\;
\mathrm{Mod}^{\Eone,\,\mathrm{compl}}_{\mathcal{A}}, &
\mathcal{N} &\;\longmapsto\;
\Omega^{\mathrm{ch}}(C_{\mathcal{A}}) \otimes_{C_{\mathcal{A}}}
\mathcal{N}.
\end{align*}
Under this equivalence:
\begin{enumerate}[label=\textup{(\roman*)}]
\item For a simple $\mathcal{A}$-module~$L$, the bar image
 $\mathrm{Ind}(L)$ is a cofree, hence injective,
 $C_{\mathcal{A}}$-comodule. On finite-type dualized lanes, its
 graded dual is projective as an $\mathcal{A}^!$-module.
\item Derived extension groups are exchanged with derived morphisms on
 the bar-comodule side:
 \[
 \mathrm{Ext}^i_{\mathcal{A}}(L, M) \cong
 \operatorname{Hom}_{D(\mathrm{CoMod}^{\Eone,\,\mathrm{conil}}_{C_{\mathcal{A}}})}
 (\mathrm{Ind}(L), \mathrm{Ind}(M)[i]).
 \]
\item On any lane where the finite-type graded-dual algebra
 interpretation $C_{\mathcal{A}}^\vee \cong \mathcal{A}^!$ is
 available, character formulas for $\mathcal{A}$-modules
 (Chapter~\ref{chap:chiral-modules}) are computed via the Koszul
 resolution by $\mathcal{A}^!$.
\end{enumerate}
\end{theorem}

\begin{proof}
The argument follows the same geometric bar-resolution strategy
discussed above, but now on the proved $\Eone$
complete/conilpotent lane:

\emph{Step 1.}
The bar complex $\bar{B}^{\mathrm{ch}}(\mathcal{A})$ is computed
using \emph{ordered} configurations (no $\Sigma_n$-quotient), so
$\mathrm{Ind}(\mathcal{M})$ tensors $\mathcal{M}$ with
$\bar{B}^{\mathrm{ch}}_n(\mathcal{A})$ over ordered
configuration spaces.

\emph{Step 2.}
The Koszulity of $\chirAss$
(Proposition~\ref{prop:chirAss-self-dual}) ensures the bar
resolution is a cofibrant replacement in $\Eone$-chiral modules.
The counit $\Omega^{\mathrm{ch}} \circ \bar{B}^{\mathrm{ch}}
\xrightarrow{\sim} \mathrm{id}$
(Theorem~\ref{thm:e1-chiral-koszul-duality}, equation
\eqref{eq:e1-counit}) extends to modules:
$\mathrm{Res} \circ \mathrm{Ind}(\mathcal{M}) \simeq \mathcal{M}$.

\emph{Step 3.}
The unit
$\mathrm{id} \xrightarrow{\sim}
\mathrm{Ind} \circ \mathrm{Res}$
follows from the dual argument on conilpotent comodules.

We verify each property in turn.

\emph{Property (i): simple-to-cofree transport.}
Let $L$ be a simple $\mathcal{A}$-module, concentrated in conformal
weight~$h$. The functor $\mathrm{Ind}$ sends $L$ to
$\bar{B}^{\mathrm{ch}}(\mathcal{A}) \otimes_{\mathcal{A}} L$.
By the Koszul property, the Ext groups
$\mathrm{Ext}^n_{\mathcal{A}}(\omega_X, L)$ are concentrated in
internal degree~$n$ (this is the defining condition of Koszulity:
the diagonal purity of Ext; see~\cite{BGS96}, Theorem~2.10.1).
Therefore
$\mathrm{Ind}(L) = \bigoplus_n \mathrm{Ext}^n_{\mathcal{A}}(\omega_X, L)$
is a cofree $C_{\mathcal{A}}$-comodule: it equals
$C_{\mathcal{A}} \otimes L_h$ where $L_h$ is the lowest weight space.
Cofree comodules are injective, establishing the first sentence
of~(i). On any finite-type dualized lane, graded duality sends this
cofree comodule to a projective $\mathcal{A}^!$-module, giving the
second sentence.

\emph{Property (ii): Ext exchange.}
The displayed derived equivalence is covariant, so it preserves
morphism spaces with the same order of arguments:
\[
\operatorname{Hom}_{D(\mathrm{Mod}^{\Eone,\,\mathrm{compl}}_{\mathcal{A}})}
(L, M[i])
\cong
\operatorname{Hom}_{D(\mathrm{CoMod}^{\Eone,\,\mathrm{conil}}_{C_{\mathcal{A}}})}
(\mathrm{Ind}(L), \mathrm{Ind}(M)[i]).
\]
By definition, the left-hand side is
$\mathrm{Ext}^i_{\mathcal{A}}(L,M)$, giving~(ii).

\emph{Property (iii): character computation.}
On any lane where $C_{\mathcal{A}}^\vee \cong \mathcal{A}^!$ is
available, the usual dual-algebra Koszul resolution may be used.
The Koszul resolution of any
$\mathcal{A}$-module $\mathcal{M}$ by free
$\mathcal{A}$-modules has terms
$\mathcal{P}_n = \mathcal{A} \otimes (\mathcal{A}^!)_n
\otimes \mathcal{M}$
(Theorem~\ref{thm:koszul-resolution-module} in
Chapter~\ref{chap:chiral-modules}); the character formula then
follows from the acyclicity of the Koszul resolution combined
with the Hilbert series identity $h_{\mathcal{A}}(t) \cdot
h_{\mathcal{A}^!}(-t) = 1$.
\end{proof}

\begin{proposition}[Finite quadratic Hilbert identity; \ClaimStatusConditional]
\label{prop:koszul-character-identity}
\index{Koszul property!character identity|textbf}
\index{character identity!Koszul|textbf}

Let $\cA$ lie on the split finite-type quadratic lane, and assume its
bar-dual coalgebra~$\cA^{\mathrm i}$ admits the finite Verdier dual
algebra~$\cA^!=(\cA^{\mathrm i})^\vee$. Then the Hilbert series of the
algebra and of the post-Verdier dual satisfy the classical Koszul
identity
\begin{equation}\label{eq:koszul-character-identity}
 h_\cA(t)\,h_{\cA^!}(-t)=1.
\end{equation}
Here $h_{\cA^!}$ means the Hilbert series of
$(\cA^{\mathrm i})^\vee$. It is not the vacuum quasi-primary character
of~$\cA$, and equation~\eqref{eq:koszul-character-identity} is not a
formula for quasi-primary counts outside the split finite-type
quadratic lane.
\end{proposition}

\begin{proof}
On the finite quadratic lane the Koszul resolution of the unit has
terms $\cA\otimes(\cA^{\mathrm i})_n$. Taking the Euler characteristic
of this exact resolution gives
\[
h_\cA(t)\sum_{n\ge0}(-1)^n\dim(\cA^{\mathrm i})_n\,t^n=1.
\]
Finite Verdier duality preserves the graded dimensions, so the signed
series is $h_{\cA^!}(-t)$. This proves
\eqref{eq:koszul-character-identity}. For PBW non-quadratic or
completed families, the same expression is only a finite-window Euler
identity after the Mittag--Leffler hypotheses have been imposed.
\end{proof}

\begin{proposition}[Bar cohomology $\neq$ quasi-primary count;
\ClaimStatusProvedHere]
\label{prop:bar-neq-quasiprimary}
\index{bar cohomology!quasi-primary count distinction|textbf}
\index{quasi-primary space!bar cohomology distinction|textbf}
The bar cohomology dimension $\dim H^1(\barBgeom(\cA))_h$
counts cogenerators of the bar-dual coalgebra~$\cA^{\mathrm i}$ at
weight~$h$; after finite Verdier duality it counts generators of the
post-Verdier algebra~$\cA^!$. The quasi-primary count
$d_{\mathrm{qp}}(h) = p_{\geq 2}(h) - p_{\geq 2}(h-1)$
counts quasi-primary fields of the \emph{original}
algebra~$\cA$. These are different sequences for any
chirally Koszul algebra whose generators have weight~$\geq 2$.
\end{proposition}

\begin{proof}
Two independent obstructions.

\emph{(a) Different sequences.}
For the Virasoro algebra, $\dim H^1(\barBgeom(\mathrm{Vir}))_h
= M(h+1) - M(h)$ where $M(n)$ is the $n$-th Motzkin number
(Theorem~\ref{thm:virasoro-chiral-koszul}); this sequence is
$0, 1, 2, 5, 12, 30, 76, \ldots$ and is independent of~$c$.
The quasi-primary count
$d_{\mathrm{qp}}(h) = p_{\geq 2}(h) - p_{\geq 2}(h{-}1)$
grows as a partition function, with
$d_{\mathrm{qp}}(4) = 1$, $d_{\mathrm{qp}}(6) = 2$,
$d_{\mathrm{qp}}(8) = 3$. At weight~$4$:
$\dim H^1(\barBgeom(\mathrm{Vir}))_4 = 12$ (from
$M(5) - M(4) = 21 - 9$) while
$d_{\mathrm{qp}}(4) = 1$. The sequences already diverge at
the first nontrivial weight above the generator.

\emph{(b) Bar cohomology sees the Koszul dual, not
the original.}
By the Koszul property,
$H^1(\barBgeom(\cA))_h \cong (\cA^{\mathrm i})_h^{\mathrm{cogen}}$:
bar cohomology in degree~$1$ counts cogenerators of the
\emph{bar-dual} coalgebra, and only after finite duality generators of
the dual algebra. For Kac--Moody:
$\dim H^1(\barBgeom(\widehat{\mathfrak{g}}))_1
= \dim\mathfrak{g}$
(both the currents $J^a$ and the dual generators live at
weight~$1$). The coincidence at weight~$1$ is a consequence
of the KM algebra being generated in weight~$1$, so generators
of~$\cA$ and generators of~$\cA^!$ share the same weight.
For Virasoro, with a weight-$2$ generator, the sequences
diverge starting at weight~$4$.
\end{proof}

\begin{remark}[Why the coincidence for free fields is tempting]
\label{rem:bar-quasiprimary-free-field}
\index{bar cohomology!free-field coincidence}
For free-field algebras (Heisenberg, free fermion), the
quasi-primary count at weight~$h$ equals the bar cohomology
dimension at weight~$h$, because the algebra is generated by
a weight-$1$ field and the Koszul dual is the
symmetric/exterior algebra on the same space. This is a
special-case coincidence: it holds when the algebra is
both free and generated in weight~$1$, which forces the Koszul
dual to have the same Hilbert series as the quasi-primary tower.
For algebras with generators of weight~$\geq 2$, the Koszul
dual's weight distribution diverges from the quasi-primary
distribution.
\end{remark}

\begin{remark}[Categorified Koszul duality and geometric representation theory]
\label{rem:categorified-koszul-kl-bfm}
\ClaimStatusHeuristic
\index{Koszul duality!categorified}
\index{Kazhdan--Lusztig conjecture!bar complex interpretation}
The numerical complementarity $\kappa(\cA) + \kappa(\cA^!) = K$
(Theorem~\ref{thm:quantum-complementarity-main}) has a categorical
shadow only after the finite-type Verdier and completion hypotheses
have been fixed. The theorem proved here is the bar-comodule
equivalence on the complete $\Eone$ lane:
$D(\mathrm{Mod}_{\cA}^{\mathrm{compl}}) \simeq
D(\mathrm{CoMod}_{B(\cA)}^{\mathrm{conil}})$
(Theorem~\ref{thm:e1-module-koszul-duality}). Its Euler
decategorification can be compared with the scalar conductor identity
only after identifying the bar-dual coalgebra with a finite
post-Verdier algebra and inserting the family-specific conductor
normalization
$\kappa + \kappa^! = \varrho_{\cA}\,K$
(with $\varrho_{\cA} = 0$ for Kac--Moody and free-field families,
$\varrho_{\mathrm{Vir}} = 1/2$, $\varrho_{\cW_N} = H_N - 1$ for principal
$\cW_N$, and $\varrho_{\BP} = 1/6$). The module equivalence alone does
not prove those scalar constants.

This structure has three classical shadows in geometric
representation theory.

\begin{enumerate}[label=\textup{(\roman*)}]
\item \emph{BGS Koszul duality.}
 For category~$\mathcal{O}$ of a semisimple Lie algebra~$\fg$,
 the Ext algebra of simples is Koszul, and the
 Koszul dual category is
 $\mathcal{O}^!(\fg^\vee)$~\cite{BGS96}. The
 Kazhdan--Lusztig polynomials $P_{y,w}(q)$ are read from the
 $E_2$ page of the bar spectral sequence
 (Remark~\ref{rem:kl-bar}); the Kazhdan--Lusztig
 conjecture~\cite{KL79}, proved by
 Beilinson--Bernstein~\cite{BB81} and
 Brylinski--Kashiwara~\cite{BK81}, is the statement that
 these polynomials compute composition multiplicities. In our
 framework the bar complex $\barBgeom(\cA)$ provides the
 universal resolution underlying this spectral sequence.

\item \emph{Kazhdan--Lusztig tensor equivalence.}
 At positive integer level~$k$, the
 Kazhdan--Lusztig equivalence~\cite{KL93}
 $\mathcal{O}_k(\widehat{\fg}) \simeq
 \mathcal{C}(U_q(\fg))$ identifies affine category~$\mathcal{O}$
 with a quantum group category. Passing to the Koszul dual
 level $k' = -(k + 2h^\vee)$ exchanges the two sides of
	 this equivalence; the numerical shadow is the affine Kac--Moody
	 complementarity: the sum
	 $\kappa(V_k(\fg)) + \kappa(V_{k'}(\fg))$ vanishes.
 The chiral bar-cobar adjunction is the operadic skeleton
 that governs this level-exchange: the bar differential at
 level~$k$ and the cobar differential at level~$k'$ are
 intertwined by the Verdier functor of Theorem~A.

\item \emph{Bezrukavnikov--Finkelberg--Mirkovi\'{c}.}
 The BFM programme~\cite{BFM04} constructs a derived
 equivalence between (mixed) perverse sheaves on the affine
 Grassmannian $\mathrm{Gr}_G$ and modules over the Langlands
 dual Lie algebra~$\fg^\vee$, extending the geometric Satake
 correspondence to a Koszul duality statement. Their
 equivalence is a derived incarnation of the same pattern:
 the bar resolution of the equivariant intersection cohomology
 complexes on $\mathrm{Gr}_G$ computes $\mathrm{Ext}$ groups
 that are formal and Koszul, and the resulting dual category
 is governed by~$\fg^\vee$.
 Arkhipov--Bezrukavnikov--Ginzburg~\cite{ABG04} establish the
 parallel statement for the quantum group, completing the
 triangle \emph{affine~$\mathcal{O}$}
 $\leftrightarrow$ \emph{quantum group}
 $\leftrightarrow$ \emph{perverse sheaves on~$\mathrm{Gr}$}.
\end{enumerate}

In each case the shared mechanism is Koszul-type formality: the
relevant Ext algebra is generated in degree~$1$ or, in the chiral
setting, the bar cohomology is concentrated along the diagonal. The
chiral Koszul pair
$(\cA,\cA^!)$ of Definition~\ref{def:chiral-koszul-pair} axiomatises
this bar-formality mechanism after the bar-dual coalgebra has been
separated from its post-Verdier dual algebra.

BGS formality, the Kazhdan--Lusztig tensor equivalence, and the BFM
derived Satake correspondence require their own geometric hypotheses:
flag-variety localization, integral-level affine category~$\mathcal O$,
or affine-Grassmannian constructibility. The chiral bar-cobar
formalism supplies the common Koszul skeleton, not a replacement for
those inputs. When the external hypotheses are present, the bar
spectral sequence records the same graded multiplicities; without
them the paragraph is only a heuristic comparison.
\end{remark}


% ================================================================
% SECTION 8.6: INTERCHANGE OF STRUCTURES
% ================================================================

\section{Interchange of structures under Koszul duality}

\subsection{Generators and relations}

\begin{theorem}[Structure exchange on the finite quadratic lane; \ClaimStatusProvedHere]\label{thm:structure-exchange}
Let $(\mathcal{A}_1,\mathcal{A}_2)$ be a finite-type quadratic chiral
Koszul pair, with $\mathcal{A}_2\simeq\mathcal{A}_1^!$ obtained by
finite Verdier duality from the bar-dual coalgebra
$\mathcal{A}_1^{\mathrm i}$. Then:
\begin{enumerate}
\item \emph{Generators $\leftrightarrow$ Relations:}
\[\mathrm{Gen}(\mathcal{A}_1) \leftrightarrow \mathrm{Rel}(\mathcal{A}_2)^{\perp}\]
\[\mathrm{Rel}(\mathcal{A}_1) \leftrightarrow \mathrm{Gen}(\mathcal{A}_2)^{\perp}\]

\item \emph{Products $\leftrightarrow$ Coproducts:}
The multiplication of $\mathcal{A}_1$ is dual to the differential on
$\mathcal{A}_1^{\mathrm i}$, while the coproduct of
$\mathcal{A}_1^{\mathrm i}$ is the deconcatenation coproduct of the bar
coalgebra.

\item \emph{Syzygy ladder:}
\[
\mathrm{Syz}^n_{\mathcal{A}_1}(\mathbf 1)
\leftrightarrow
\mathrm{CoSyz}^{n+1}(\mathcal{A}_1^{\mathrm i})
\]
under the minimal Koszul resolution.
\end{enumerate}
\end{theorem}

\begin{proof}
Suppose $\mathcal{A}_1$ is presented as
$\mathcal{A}(V,R)$ with
$R \subset j_*j^*(V\boxtimes V)$. The quadratic bar-dual coalgebra is
the subcoalgebra of $T^c_{\mathrm{ch}}(s^{-1}V^\vee)$ cogenerated by
the residue-orthogonal corelations $R^\perp$. Its finite dual algebra
is
\[
\mathcal{A}_2\simeq
T_{\mathrm{ch}}(sV^\vee\omega_X^{-1})/(R^\perp),
\]
where $R^\perp$ is taken under the residue pairing
\[
\langle\, \cdot\,, \cdot\, \rangle_{\mathrm{res}} : j_*j^*(V \boxtimes V) \otimes j_*j^*(V^\vee \boxtimes V^\vee) \to \mathbb{C}, \quad
\langle \alpha, \beta \rangle_{\mathrm{res}} = \mathrm{Res}_{z_1 = z_2}\, \alpha(z_1, z_2) \cdot \beta(z_1, z_2).
\]
This proves the generator-relation exchange; the reverse direction is
the same statement applied to the finite quadratic dual.

The bar construction $\bar{B}(\mathcal{A}_1)$ is a coaugmented dg
coalgebra with deconcatenation coproduct. Its differential encodes the
product of $\mathcal{A}_1$:
\[
d_{\bar{B}}(a_1 | \cdots | a_n) = \sum_{i=1}^{n-1} \pm\, (a_1 | \cdots | \mu(a_i, a_{i+1}) | \cdots | a_n)
\]
with the sign fixed by the desuspension convention of
Appendix~\ref{app:signs-shifts}. Therefore multiplication in
$\mathcal{A}_1$ is dual to the coderivation on
$\mathcal{A}_1^{\mathrm i}$, not to the assertion
$\mathcal{A}_2=\Omega\bar{B}(\mathcal{A}_1)$. The latter cobar object
reconstructs $\mathcal{A}_1$ by Theorem~A.

Finally, the minimal Koszul resolution has
$P_n\simeq\mathcal{A}_1\otimes(\mathcal{A}_1^{\mathrm i})_n$. Hence
the $n$-th syzygy of the unit module is controlled by the
$(n+1)$-st cogenerator layer of the bar-dual coalgebra, which is the
displayed cosyzygy statement.
\end{proof}

\subsection{\texorpdfstring{$A_\infty$ operations exchange}{A-infinity operations exchange}}

\begin{theorem}[\texorpdfstring{$A_\infty$}{A-infinity} operations under Verdier pairing; \ClaimStatusProvedHere]
\label{thm:ainfty-duality-exchange}
Let $(\mathcal{A}_1,\mathcal{A}_2)$ be a finite-type quadratic chiral
Koszul pair with non-degenerate Verdier pairing. Then:
\begin{itemize}
\item transferred higher operations on $\mathcal{A}_1$ pair with
 higher cooperations on $\mathcal{A}_1^{\mathrm i}$;
\item after finite duality, these cooperations become higher
 operations on $\mathcal{A}_2\simeq\mathcal{A}_1^!$;
\item Massey products on $\mathcal{A}_1$ pair with the corresponding
 comassey classes on $\mathcal{A}_1^{\mathrm i}$.
\end{itemize}
\end{theorem}

\begin{proof}
The Verdier duality pairing on configuration spaces
(Theorem~\ref{thm:verdier-config}) gives
\[
\langle m_k^{(1)}, n_k^{(2)} \rangle = \int_{\overline{C}_k(X)} \omega_{m_k} \wedge \delta_{n_k}
\]
where $\omega_{m_k}$ represents the $k$-ary operation on
$\mathcal{A}_1$ and $\delta_{n_k}$ represents the dual cooperation on
the bar-dual coalgebra. Non-degeneracy detects non-vanishing of the
paired operation. Applying finite continuous duality turns the
cooperation on $\mathcal{A}_1^{\mathrm i}$ into an operation on the dual
algebra~$\mathcal{A}_2$. Massey and comassey classes are the cohomology
classes represented by these higher operations and cooperations, so the
same pairing gives the final assertion.
\end{proof}

% ================================================================
% SECTION 8.6: FILTERED AND CURVED EXTENSIONS
% ================================================================

\section{Filtered and curved extensions}

\subsection{Filtered and curved structures}

The central extensions of Virasoro, Kac--Moody, and Yangian algebras
produce curvature $m_0 \neq 0$ in their bar complexes. To handle
this systematically:

\begin{definition}[Filtered chiral algebra]
A filtered chiral algebra has an exhaustive filtration:
\[0 = F_{-1}\mathcal{A} \subset F_0\mathcal{A} \subset F_1\mathcal{A} \subset \cdots\]
with $\mu(F_i \otimes F_j) \subset F_{i+j}$ and $\mathcal{A} = \varprojlim \mathcal{A}/F_n\mathcal{A}$.
\end{definition}

\begin{definition}[Curved \texorpdfstring{$A_\infty$}{A-infinity}]\label{def:curved-ainfty-kp}
A curved $A_\infty$ structure has operations $m_k$ for $k \geq 0$ with curvature $m_0 \in F_{\geq 1}\mathcal{A}[2]$ satisfying the Maurer--Cartan equation.
\end{definition}

\begin{proposition}[Uniqueness of the Feigin--Frenkel involution; \ClaimStatusProvedHere]
\label{prop:ff-involution-uniqueness}
\index{Feigin--Frenkel involution!uniqueness}
Let $\mathfrak{g}$ be a finite-dimensional simple Lie algebra and
consider the family of affine Kac--Moody algebras
$\{\widehat{\mathfrak{g}}_k\}_{k \in \bC}$. The Feigin--Frenkel
involution $\iota_{\mathrm{FF}}: k \mapsto -k - 2h^\vee$ is the
\emph{unique} involution $\iota$ of the level parameter satisfying:
\begin{enumerate}[label=\textup{(\roman*)}]
\item $\iota$ is an affine-linear involution: $\iota(k) = ak + b$
 with $\iota^2 = \mathrm{id}$;
\item $\iota$ is compatible with Koszul duality:
 the Koszul dual $\widehat{\mathfrak{g}}_k^!$ has the same
 modular characteristic as $\widehat{\mathfrak{g}}_{\iota(k)}$,
 i.e.\ $\kappa(\widehat{\mathfrak{g}}_k^!) =
 \kappa(\widehat{\mathfrak{g}}_{\iota(k)})$;
\item $\iota$ fixes the critical level: $\iota(-h^\vee) = -h^\vee$.
\end{enumerate}
\end{proposition}

\begin{proof}
An affine-linear involution $\iota(k) = ak + b$ satisfies
$\iota^2(k) = a^2k + ab + b = k$ for all~$k$, giving $a^2 = 1$
and $b(a+1) = 0$. If $a = 1$, then $b = 0$ and $\iota = \mathrm{id}$,
which contradicts~(ii) (a non-trivial duality). Therefore $a = -1$
and $\iota(k) = -k + b$.

Condition~(iii) requires $\iota(-h^\vee) = h^\vee + b = -h^\vee$,
so $b = -2h^\vee$. This gives $\iota(k) = -k - 2h^\vee$, the
Feigin--Frenkel involution.

That this involution satisfies~(ii) is the content of
Theorem~\ref{thm:universal-kac-moody-koszul}: the bar construction
at level~$k$ produces a Koszul dual whose modular characteristic
equals $\kappa(\widehat{\mathfrak{g}}_{-k-2h^\vee})$.
\end{proof}

\subsection{Curved Koszul duality}

\begin{theorem}[Curved conductor compatibility; \ClaimStatusConditional]\label{thm:curved-koszul-pairs}
\index{curved Ainfinity@curved $A_\infty$!Koszul pairs}
\textup{[Regime: curved-central
\textup{(}Convention~\textup{\ref{conv:regime-tags})}.]}

Let $(\cA_1,\cA_2)$ be filtered-complete chiral algebras with modular
curvatures $\kappa_1,\kappa_2$. Put
\[
\cC_i=\widehat{\barB}^{\mathrm{curv}}_X(\cA_i),\qquad
\cA_i^{\mathrm i}=H^*(\cC_i),
\]
so $\cC_i$ is the completed curved bar coalgebra and
$\cA_i^{\mathrm i}$ is the intrinsic bar-dual coalgebra. When the latter
is finite-type Verdier dualizable, write
$\cA_i^!=\mathbb D_{\operatorname{Ran}}\cA_i^{\mathrm i}$. The Hochschild
object $Z^{\mathrm{der}}_{\mathrm{ch}}(\cA_i)
=C^\bullet_{\mathrm{ch}}(\cA_i,\cA_i)$ is not part of this pair datum.
Assume:
\begin{enumerate}
\item the associated graded pair
$(\mathrm{gr}\,\cA_1,\mathrm{gr}\,\cA_2)$ is a classical Koszul pair;
\item the filtrations are complete, separated, and bidegreewise
Mittag--Leffler in the sense of Lemma~\ref{lem:completion-convergence};
\item the coalgebras $\cA_i^{\mathrm i}$ are finite-type Verdier
dualizable on the chosen lane, and the candidate Verdier comparison maps
\[
\mathbb D_{\operatorname{Ran}}\cA_1^{\mathrm i}\longrightarrow \cA_2,
\qquad
\mathbb D_{\operatorname{Ran}}\cA_2^{\mathrm i}\longrightarrow \cA_1
\]
are maps of completed factorization algebras whose associated graded
maps are the classical Koszul-duality maps from \textup{(1)};
\item the curved twisting morphism exists in the Positselski
completed category and its curvature scalars satisfy the family
conductor identity
\[
\kappa_1+\kappa_2=K^\kappa_{\mathrm{fam}}
\]
under these candidate pairings;
\item the spectral sequence associated to the filtration degenerates at
$E_2$.
\end{enumerate}
Then the completed curved bar-cobar comparison identifies
$\Omega_X^{\mathrm{curv}}\cC_i$ with~$\cA_i$, and the two maps in
\textup{(3)} are the post-Verdier Koszul-duality comparisons
$\cA_1^!\simeq\cA_2$ and $\cA_2^!\simeq\cA_1$ on this lane. The scalar
identity $\kappa_1+\kappa_2=K^\kappa_{\mathrm{fam}}$ is only the
curvature-conductor compatibility for those comparisons. It is not a
criterion for Koszulness without \textup{(1)}--\textup{(3)} and
\textup{(5)}.

For the canonical families under their stated hypotheses,
\[
K^\kappa_{\mathrm{fam}}=
0\quad\text{for Heisenberg, affine Kac--Moody, and free fields,}
\]
\[
K^\kappa_{\mathrm{fam}}=13\quad\text{for Virasoro},\qquad
K^\kappa_{\mathrm{fam}}=250/3\quad\text{for }\mathcal W_3,
\]
\[
K^\kappa_{\mathrm{fam}}=98/3\quad\text{for Bershadsky--Polyakov},
\qquad
K^\kappa_{\mathrm{fam}}=8\quad\text{for recognized finite
Mukai-enhanced K3 Hall--Borcherds windows.}
\]
\end{theorem}

\begin{proof}
The filtered bar complex has associated graded
$\mathrm{gr}\,\cC_i=\barB(\mathrm{gr}\,\cA_i)$. Condition~(1) gives the
classical Koszul comparison on the $E_1$ page. Conditions~(2) and~(5)
let this comparison pass to the completed filtered object without
creating $\varprojlim^1$ cohomology. The result at this point is a
statement about the coalgebras $\cC_i$ and $\cA_i^{\mathrm i}$.
Condition~(3) is the separate step that turns $\cA_i^{\mathrm i}$ into
the post-Verdier algebra $\cA_i^!$; without it the theorem has no
algebra-level duality conclusion, and it says nothing about
$Z^{\mathrm{der}}_{\mathrm{ch}}(\cA_i)$.

In a curved $A_\infty$ bar-cobar complex the curvature term is not
universally anti-invariant; it is shifted by the conductor of the
family. Condition~(4) is exactly the compatibility condition for the
curved twisting morphism in the Positselski completed category
\cite{Positselski11}. The displayed constants are the canonical target
values from the landscape census:
$\kappa(V_k(\fg))=\dim(\fg)(k+h^\vee)/(2h^\vee)$,
$\kappa(\mathrm{Vir}_c)=c/2$, and
$\kappa(\mathcal W_N)=c(H_N-1)$, together with the recorded
Bershadsky--Polyakov and conditional Mukai rows. The Mukai value is
available only after the finite Hall--Borcherds recognition datum
supplies the source pairing, radical quotient, differential, and
completion. Under all five hypotheses the curved comparison is the
filtered curved bar-cobar comparison of
Theorem~\ref{thm:geometric-equals-operadic-bar}.
\end{proof}

% ================================================================
% SECTION 8.7: DERIVED CHIRAL KOSZUL DUALITY
% ================================================================

\section{Derived chiral Koszul duality}

\subsection{Derived ghost differential}

The $bc$ ghost system (weights $2,-1$) and the $\beta\gamma$ system
(weights $1,0$) do not form an ordinary quadratic pair on the nose. A
derived pairing, if it exists, must live in DG chiral algebras with a
BRST differential and with the bar coalgebra kept separate from its
post-Verdier dual algebra.

\begin{definition}[DG chiral algebra]
A DG chiral algebra is a cochain complex of factorization
$\cD$-modules
\[
\mathcal{A}^{\bullet}: \cdots \to
\mathcal{A}^{-1} \xrightarrow{d} \mathcal{A}^0
\xrightarrow{d} \mathcal{A}^1 \to \cdots
\]
whose chiral products are chain maps; equivalently, $d$ is a derivation
for every factorization multiplication map.
\end{definition}

\begin{conjecture}[Derived bc-\texorpdfstring{$\beta\gamma$}{beta-gamma} Koszul duality; \ClaimStatusConjectured]\label{conj:derived-bc-betagamma}
Let $\mathcal{F} = \Lambda^{\mathrm{ch}}(\psi^{(1)}, \psi^{(2)})$ be the
free fermion chiral algebra on two generators. The direct sum
$\beta\gamma \oplus bc$ carries a natural $\mathbb{Z}$-grading
(ghost number), making it a DG chiral algebra with differential
$d_{\mathrm{BRST}} \colon bc \to \beta\gamma$ induced by the
ghost-antighost pairing. Then:
\begin{enumerate}
\item There exists a derived chiral Koszul duality framework in which
 $\mathcal{F}$ and $(\beta\gamma \oplus bc, d_{\mathrm{BRST}})$
 form a derived Koszul pair.
\item At the level of bar complexes, the conjectural comparison is a
 coalgebra comparison
 \[
 \bar{B}^{\mathrm{ch}}_{\mathrm{derived}}(\mathcal{F})
 \;\simeq\;
 (\beta\gamma \oplus bc,\, d_{\mathrm{BRST}})^{\mathrm i},
 \]
 and only after finite-type Verdier duality would this produce the
 post-Verdier algebra
 $(\beta\gamma \oplus bc,\, d_{\mathrm{BRST}})^!$. This extends the
 proved single-pair dualities
 $(\beta\gamma)^! \cong bc$ and $\mathcal{F}(\psi)^! \cong
 \mathrm{Sym}^{\mathrm{ch}}(\gamma)$
 (Theorem~\ref{thm:betagamma-bc-koszul},
 Theorem~\ref{thm:single-fermion-boson-duality}).
\end{enumerate}
\end{conjecture}

\begin{remark}[Scope]\label{rem:scope-curved-koszul}
Conjecture~\ref{conj:derived-bc-betagamma} is tagged
\ClaimStatusConjectured{} for two independent reasons:
\begin{enumerate}
\item \emph{Algebraic.} A systematic theory of \emph{derived} chiral
 Koszul duality, extending the quadratic and PBW frameworks of
 this chapter to DG chiral algebras, has not been developed.
 The individual dualities $(\beta\gamma)^! \cong bc$ and
 $\mathcal{F}(\psi)^! \cong \mathrm{Sym}^{\mathrm{ch}}(\gamma)$
 are proved, but their interaction in a derived pair requires new
 constructions (derived bar-cobar adjunction for DG chiral algebras,
 extending Section~\ref{sec:counter-examples}).
\item \emph{Physical.} The specific role of this derived pair in
 open-closed string field theory (where $d_{\mathrm{BRST}}$ is the
 string BRST operator and the ghost number grading organizes the
 BV-BRST complex) requires external physics input from the
 formalism of~\cite{Zwi93}.
\end{enumerate}

The precise proof obligations are:
\begin{enumerate}[label=\textup{(\alph*)}]
\item construct a completed DG chiral bar-cobar adjunction whose
 coalgebras are conilpotent and whose filtrations satisfy
 Lemma~\ref{lem:completion-convergence};
\item exhibit a BRST twisting morphism
 \[
 \tau_{\mathrm{BRST}}\colon
 \bar{B}^{\mathrm{ch}}_{\mathrm{derived}}(\mathcal F)
 \longrightarrow
 (\beta\gamma\oplus bc,d_{\mathrm{BRST}})
 \]
 and prove its Maurer--Cartan equation, including the signs from the two
 fermions;
\item prove diagonal concentration for the DG bar spectral sequence, or
 identify the first off-diagonal obstruction;
\item prove finite-type Verdier dualizability of the resulting
 coalgebra before naming the post-Verdier algebra
 $(\beta\gamma\oplus bc,d_{\mathrm{BRST}})^!$;
\item for the open--closed string-field interpretation, compare
 $d_{\mathrm{BRST}}$ with the string BRST operator of~\cite{Zwi93}.
\end{enumerate}
Until \textup{(a)}--\textup{(e)} are supplied, the proved content is only
the two single-pair dualities cited in the conjecture.
\end{remark}

% ================================================================
% SECTION 8.7: COUNTER-EXAMPLES - WHEN KOSZUL DUALITY FAILS
% ================================================================

\section{Counter-examples: when Koszul duality fails}
\label{sec:counter-examples}

\subsection{Virasoro algebra: beyond the quadratic setting}
\label{subsec:virasoro-non-example}

\begin{remark}[Virasoro Koszulness requires the PBW framework]
The Virasoro algebra with central charge $c$ is not controlled by an
uncurved quadratic presentation on the single generator~$T$: the scalar
term in the $T(z)T(w)$ OPE is a curvature term, and the normal product
$:TT:$ enters the bar differential outside the quadratic relation space
generated by~$T$ alone. It is nevertheless chiral Koszul in the PBW
sense of Theorem~\ref{thm:virasoro-chiral-koszul}: the PBW filtration
yields a classical Koszul associated graded, and the spectral sequence
degenerates at~$E_2$. The completed curved same-family comparison must
therefore be stated in the filtered Positselski category of
\S\ref{sec:i-adic-completion}, with
$\mathrm{Vir}_c$, $\barB_X(\mathrm{Vir}_c)$,
$\mathrm{Vir}_c^{\mathrm i}$, $\mathrm{Vir}_c^!$, and
$Z^{\mathrm{der}}_{\mathrm{ch}}(\mathrm{Vir}_c)$ kept distinct.

\emph{Why the quadratic framework fails.}
\begin{enumerate}
\item The OPE involves a quartic pole:
\begin{equation}
T(z)T(w) = \frac{c/2}{(z-w)^4} + \frac{2T(w)}{(z-w)^2} + \frac{\partial T(w)}{z-w} + \text{regular}
\end{equation}

\item The reduced bar object is the conilpotent coalgebra
$\barB_X(\mathrm{Vir}_c)=T^c(s^{-1}\overline{\mathrm{Vir}_c})$ with
differential induced by the Virasoro OPE; it is not the post-Verdier
algebra $\mathrm{Vir}_c^!$.

\item The composite field
\begin{equation}
:TT:(z) = \lim_{w \to z} \left( T(z)T(w) - \frac{c/2}{(z-w)^4} - \frac{2T(w)}{(z-w)^2} - \frac{\partial T(w)}{z-w} \right)
\end{equation}
contributes a bar differential term outside the uncurved quadratic
relation space.

\item The central charge $c$ enters non-linearly in higher bar degrees, preventing a quadratic Koszul resolution.
\end{enumerate}
One needs the PBW spectral sequence approach
(Theorem~\ref{thm:pbw-koszulness-criterion}), or a curved
Positselski comparison with nilpotent completion and the hypotheses of
Theorem~\ref{thm:curved-koszul-pairs}. The conductor identity
$\kappa+\kappa^!=13$ is compatible with this comparison; it does not
replace the PBW collapse or the completion hypotheses.

\emph{Contrast with Heisenberg/Kac--Moody.}
\begin{center}
\begin{tabular}{|l|c|c|c|}
\hline
\textbf{Algebra} & \textbf{Quadratic?} & \textbf{PBW Koszul?} & \textbf{Post-Verdier datum} \\
\hline
Heisenberg $\mathcal{H}_k$ & Yes (genus 0) & Yes & $\mathbb D_{\Ran}\mathcal H_k^{\mathrm i}\simeq\mathrm{Sym}^{\mathrm{ch}}(V^*)$; curvature $m_0=-k\omega$ \\
Kac--Moody $\widehat{\mathfrak{g}}_k$ & Yes & Yes & $\mathbb D_{\Ran}\widehat{\mathfrak{g}}_k^{\mathrm i}$; $\kappa^!=\kappa(\widehat{\mathfrak{g}}_{-k-2h^\vee})$ \\
Virasoro $\text{Vir}_c$ & No & Yes (Thm~\ref{thm:virasoro-chiral-koszul}) & Completed curved, after Verdier duality \\
W-algebras $\mathcal{W}^k(\mathfrak{g})$ (universal) & No & Yes (Prop.~\ref{prop:pbw-universality}) & Completed curved, after Verdier duality \\
W-algebras $\mathcal{W}_k(\mathfrak{g})$ (simple quotient) & No & Open at admissible $k$ & Case-by-case \\
\hline
\end{tabular}
\end{center}

\emph{Note}: The intrinsic bar-dual coalgebra is
$\widehat{\mathfrak{g}}_k^{\mathrm i}
=H^*(\barBgeom(\widehat{\mathfrak{g}}_k))$. The Koszul dual algebra is
$\widehat{\mathfrak{g}}_k^!=\mathbb D_{\Ran}
\widehat{\mathfrak{g}}_k^{\mathrm i}$ on the finite-type Verdier lane,
not the affine algebra $\widehat{\mathfrak{g}}_{-k-2h^\vee}$ itself.
What is true is that $\kappa(\widehat{\mathfrak{g}}_k^!) =
\kappa(\widehat{\mathfrak{g}}_{-k-2h^\vee})$: the two algebras share
the same modular characteristic. For
Heisenberg (abelian, $h^\vee = 0$):
$\mathcal{H}_k^! = \mathrm{Sym}^{\mathrm{ch}}(V^*)$, the
commutative chiral algebra on the dual space
(cf.\ \S\ref{sec:heisenberg-koszul}); the curvature
$m_0 = -k\omega$ (encoding the central extension) is essential.
\end{remark}

\begin{remark}[Quadratic and PBW obstructions]
The Virasoro case isolates five facts:
\begin{itemize}
\item not every vertex algebra is quadratic Koszul in the uncurved
sense;
\item PBW Koszulness is controlled by the associated graded and the
$E_2$-collapse of the bar spectral sequence;
\item the quartic pole is an obstruction to the naive quadratic model,
not by itself an obstruction to PBW Koszulness;
\item $\mathcal W$-algebras carry analogous non-quadratic OPE terms;
the universal algebras are handled by Proposition~\ref{prop:pbw-universality},
whereas simple admissible quotients require separate null-vector
analysis.
\item PBW non-quadratic algebras require the completed bar-dual
coalgebra framework of Section~\ref{sec:three-stage-construction}.
\end{itemize}
\end{remark}

\subsection{Classical vs.\ chiral Koszulness for $\mathcal{W}$-algebras}

\begin{remark}[$\mathcal{W}$-algebras: classical quadratic Koszulness fails, chiral Koszulness holds]
\label{rem:w-algebra-classical-vs-chiral-koszul}
For $\mathcal{W}^k(\mathfrak{sl}_3)$ at generic level $k \neq -3$,
the \emph{classical quadratic} Koszul complex has nontrivial cohomology
at degree~4:
\begin{equation}
H^4(\text{classical Koszul complex}) \neq 0 \quad \text{(classical quadratic Koszulity fails)}
\end{equation}
The obstruction is concentrated at the composite field~$\Lambda$
(see~\S\ref{sec:w3-bar-degree3}).

However, $\mathcal{W}^k(\mathfrak{g})$ is \emph{chirally Koszul at all
levels}~$k$ by Proposition~\ref{prop:pbw-universality}: the universal
$\mathcal{W}$-algebra is freely strongly generated (Feigin--Frenkel),
so the PBW criterion applies. The quartic obstruction~$[\Lambda]$
measures \emph{shadow depth} (the algebra lies in class~M with
$r_{\max}=\infty$), not failure of chiral Koszulness.

\emph{Caution}: The \emph{simple quotient}~$\mathcal{W}_k(\mathfrak{g})$
at admissible levels has vacuum null vectors that may obstruct PBW; its
chiral Koszulness remains \textbf{open} (the gap is between bar-Ext and
ordinary Ext).
\end{remark}

\subsection{Tensor products and Koszulity}

\begin{remark}[Tensor products in the algebraic setting]
In the classical (non-chiral) setting, tensor products of Koszul algebras are Koszul: if $A_1$ and $A_2$ are Koszul quadratic algebras, then $A_1 \otimes A_2$ is Koszul with $(A_1 \otimes A_2)^! \cong A_1^! \otimes A_2^!$ (Polishchuk--Positselski, Theorem 3.1).

In the \emph{chiral} setting, however, the tensor product $\mathcal{A}_1 \otimes \mathcal{A}_2$ as chiral algebras involves the chiral tensor product (not the ordinary tensor product of vector spaces), and the Koszul property is no longer automatic. The bar complex of a chiral tensor product involves configuration spaces with \emph{colored} points, and the interaction between colors introduces additional structure beyond the Arnold relations for each color separately. Whether chiral Koszulity is preserved under chiral tensor products remains open.
\end{remark}

\begin{proposition}[Tensor products on the split quadratic lane; \ClaimStatusConditional]
\label{prop:koszul-dual-tensor-product}
\index{Koszul dual!tensor product}
Let $(\cA_1,\cA_1^!)$ and $(\cA_2,\cA_2^!)$ be finite-type
quadratic chiral Koszul pairs on a smooth curve~$X$. Assume the
mixed-color OPE residues vanish, the colored configuration-space
K\"unneth map is a filtered quasi-isomorphism, and the mixed relation
subspace is the standard locality subspace whose residue-orthogonal
complement gives the tensor product of duals. Then the chiral tensor
product $\cA_1\otimes^{\mathrm{ch}}\cA_2$ is Koszul, and its
bar-dual coalgebra satisfies
\begin{equation}\label{eq:koszul-dual-tensor}
(\cA_1 \otimes^{\mathrm{ch}} \cA_2)^{\mathrm i}
\;\simeq\;
\cA_1^{\mathrm i}\otimes^{\mathrm{ch}}\cA_2^{\mathrm i}.
\end{equation}
When finite duals exist this dualizes to
$(\cA_1 \otimes^{\mathrm{ch}} \cA_2)^!
\simeq \cA_1^!\otimes^{\mathrm{ch}}\cA_2^!$. In particular, for
Kac--Moody algebras the level-shift acts componentwise on the modular
characteristic:
\[
\kappa\!\left((\widehat{\mathfrak{g}}_k
\otimes^{\mathrm{ch}}\widehat{\mathfrak{h}}_\ell)^!\right)
=
\kappa(\widehat{\mathfrak{g}}_{-k-2h^\vee_{\mathfrak{g}}})
+
\kappa(\widehat{\mathfrak{h}}_{-\ell-2h^\vee_{\mathfrak{h}}}).
\]
\end{proposition}

\begin{proof}
Under the mixed-vanishing hypothesis, the colored bar complex has no
extra residue differential connecting the two colors. The K\"unneth
comparison identifies it with the completed tensor product of the two
bar complexes, with deconcatenation coproduct on each color. Since both
factors are quadratic Koszul, the tensor product of their Koszul
complexes is acyclic, hence the colored bar complex is concentrated on
the diagonal. Residue orthogonality for the standard mixed locality
relations gives the coalgebra isomorphism
\eqref{eq:koszul-dual-tensor}. Finite duality gives the algebra
statement. The Kac--Moody specialization uses only the modular
characteristic equality from
Theorem~\ref{thm:universal-kac-moody-koszul}; it does not identify the
dual algebra itself with an affine algebra at shifted level.
\end{proof}

% ================================================================
% SECTION 8.8: COMPUTATIONAL METHODS
% ================================================================

\section{Computational methods and verification}

\subsection{Algorithm for checking Koszul pairs}

\begin{construction}[Verification of Koszul pairs]\label{con:verify-koszul-pair}
For quadratic chiral algebras, extract generators and relations,
check that the residue pairing is perfect, and verify orthogonality
$R_1^\perp=R_2$ together with acyclicity of the quadratic Koszul
complex. For PBW non-quadratic algebras, low-degree bar computations
are falsification tests: a nonzero off-diagonal class in
$\bar{B}^{\leq3}$ disproves Koszulness, but acyclicity through
degree~$3$ does not prove it. A proof requires PBW spectral-sequence
degeneration, Kac--Shapovalov/Li-bar non-degeneracy where relevant,
and the completion convergence of Lemma~\ref{lem:completion-convergence}.
\end{construction}

\subsection{Complexity analysis}

For $n$ generators, $m$ relations, verification to degree $k$:
\begin{itemize}
\item Quadratic case: $O(n^2 + m^2)$ for orthogonality
\item General case: $O(n^k)$ for bar complex dimension
\item Configuration integrals: $O(k! \cdot n^k)$ worst case
\end{itemize}

% ================================================================
% SECTION 8.9: SUMMARY AND OUTLOOK
% ================================================================

\section{The Koszul-pair landscape}

The PBW criterion (Theorem~\ref{thm:pbw-koszulness-criterion})
reduces chiral Koszulness to a classical spectral sequence
degeneration, applicable uniformly to all freely strongly generated
vertex algebras (Proposition~\ref{prop:pbw-universality}). The
twelve characterizations of
Theorem~\ref{thm:koszul-equivalences-meta} read the same
genus-$0$ formality condition through distinct invariants: bar
concentration, Ext diagonal vanishing, Kac--Shapovalov
determinants, and FM boundary acyclicity. Factorization homology
enters only through the conditional comparison model of
Proposition~\ref{prop:bar-fh}; its concentration is a consequence
of Koszulness under $\mathsf{FH}_\Sigma$, and a criterion only under
the detecting strengthening $\mathsf{FH}^{\mathrm{det}}_{\mathbb{P}^1}$.
Two characterizations remain partially open: the Lagrangian
criterion awaits perfectness verification in general, and D-module
purity awaits the converse direction. The $\Eone$-chiral Koszul duality
theorem (Theorem~\ref{thm:e1-chiral-koszul-duality}) extends the
adjunction to the associative chiral operad, and the module-level
duality (Theorem~\ref{thm:e1-module-koszul-duality}) provides the
derived equivalence on the proved complete/conilpotent lane.

\begin{remark}[Structural consequences of the twisting theorem]\label{rem:ckp-dependency-digest}
\index{chiral Koszul pairs!dependency map}
The fundamental twisting theorem supplies the following structural
inputs:
\begin{itemize}
\item $\mathrm{A}_0$: Fundamental theorem of chiral twisting morphisms
 (Theorem~\ref{thm:fundamental-twisting-morphisms}), used by
 Chapters~\ref{chap:higher-genus}, \ref{chap:chiral-modules}.
\item $\mathrm{A}_1$: Bar concentration
 (Theorem~\ref{thm:bar-concentration}), used by
 Theorem~$\mathrm{C}_0$
 (Theorem~\ref{thm:fiber-center-identification}).
\item $\mathrm{A}_2$: Geometric bar-cobar duality
 (Theorem~\ref{thm:bar-cobar-isomorphism-main}), used
 everywhere.
\item Theorem~B: Higher genus inversion
 (Theorem~\ref{thm:higher-genus-inversion}) is
 proved in Chapter~\ref{chap:higher-genus} using the foundations
 above.
\item PBW Koszulness criterion
 (Theorem~\ref{thm:pbw-koszulness-criterion}), used by all
 interacting-family example chapters.
\item $\Eone$-chiral Koszul duality
 (Theorem~\ref{thm:e1-chiral-koszul-duality}), used by
 Chapter~\ref{chap:yangians}.
\end{itemize}
\end{remark}

The Koszul pair structure classifies the duality at the algebra level.
The next chapter adds the representation-theoretic layer: conformal
blocks, fusion rules, and Kazhdan--Lusztig transport. In general, the
bar-cobar machinery extends from algebras to module bar constructions
and bar-coalgebra comodules; on the separate quadratic genus-$0$
$\Eone$ complete/conilpotent lane this further yields the proved
ordinary module-category duality. Broader module-category descent is
not part of the proved statement here.

\section{K3 chiral Hall--Drinfeld double as conditional Koszul pair}
\label{sec:ckp-supplement}

\begin{remark}[Finite Hall--Borcherds recognition datum]
\label{rem:ckp-hall-drinfeld-pair}
The notation $\mathbf H_{\Delta_5}$ is used below as a recognition target,
not as an unconditional compact Hall object.  Fix a finite downward
saturated window $W$ in the Gritsenko--Nikulin root monoid.  A
Hall--to--Borcherds recognition datum on $W$ consists of:
\begin{enumerate}
\item a compact source provenance: an oriented Hall, critical CoHA, or
factorisation source for the $K3\times E$ sector, separate from the
automorphic denominator;
\item a positive-half model $Y^+_{\mathrm{Hall}}(K3\times E)_W$ with
source multiplication, coproduct, and Hall grading, in the sense of the
Kontsevich--Soibelman and Schiffmann--Vasserot/Rap\v{c}\'ak--Soibelman--Yang--Zhao
CoHA constructions~\cite{KS11,SchiffmannVasserot2012,RSYZ20};
\item parity-homogeneous simple-root representatives, real-root signs,
and full even/odd root dimensions matching the Gritsenko--Nikulin
Borcherds--Kac fixture on $W$~\cite{GritsenkoNikulin1998,KacBook};
\item a non-degenerate completed Hopf pairing, its closed radical
quotient, and the continuous Drinfeld double
$D(Y^+_{\mathrm{Hall}}(K3\times E)_W)$;
\item a source differential $d_{\mathrm{Hall}}$ from the compact Hall,
BRST, or hCS complex, compatible with the coproduct, pairing, and
completion;
\item an independent coefficient check: the finite supercharacter of the
source radical quotient agrees with the corresponding truncation of the
Gritsenko--Nikulin/Borcherds product
$64^{-1}\Delta_5(2Z)$~\cite{Bor95,GritsenkoNikulin1998}.
\end{enumerate}
Under such data the finite double is compared with the finite truncation
of the denominator algebra $\mathfrak g_{\Delta_5}$, and the
Frenkel--Kac vertex closure~\cite{FrenkelKac1980} becomes a target model
for the BKM extension.  The denominator factor supplies root
multiplicities, character, and Euler product.  It does not supply
$d_{\mathrm{Hall}}$, compact primitive representatives, the Hall pairing,
or the radical quotient.
\end{remark}

\begin{remark}[Conditional conductor on the Mukai-enhanced K3 row]
\label{rem:ckp-mukai-doubling}
For a recognized finite Hall--Borcherds window, the Mukai-enhanced K3 row
has target conductor
\[
 K^{\kappa_{\mathrm{ch}}}(\mathcal{H}_{\mathrm{Muk}}(K3)) \;=\; \kappa_{\mathrm{ch}}(\mathcal{H}_{\mathrm{Muk}}) + \kappa_{\mathrm{ch}}(\mathcal{H}_{\mathrm{Muk}}^!) \;=\; 2\,c_+(\mathrm{Mukai}(K3)) \;=\; 8,
\]
with $c_+ = 4$ from the positive signature of the rank-$24$ Mukai
pairing on $\Lambda_{\mathrm{Muk}} = \mathrm{II}_{4,20}$~\cite{Mukai1987}.
This is the $\mathcal{B}$-family value in the Theorem~C bucket
$K^\kappa_{\mathrm{fam}}\in \{0, 8, 13, 250/3, 98/3\}$.  The comparison
with the $\ell=8$ Lusztig specialisation is a consequence only after the
recognized double carries the source pairing and completion above.  Bruinier's
Heegner Chern-class reciprocity fixes the target monodromy of
$\mathcal{L}^{\Delta_5}$ around $H_1\subset\overline{\mathcal A_2}$
~\cite[Proposition~5.1]{Bruinier2002}; it is not a construction of the
compact Hall source.
\end{remark}

\section{The Koszul triple for $\mathbf H_{\Delta_5}$ on
$(\Ran(E^{\mathrm{nod,sm}}_{24}),\overline{\mathcal A}_2)$}
\label{sec:ckp-deep}

\begin{remark}[The Koszul triple $(\mathbf H_{\Delta_5}, \mathbf H^!_{\Delta_5}, B_{\mathrm{ch}}(\mathbf H_{\Delta_5}))$]
\label{rem:ckp-deep-1}
On a finite recognized window $W$, the Koszul triple on
$(\mathrm{Ran}(E^{\mathrm{nod,sm}}_{24}),\bar{\mathcal A}_2)$ is
\[
\bigl(\mathbf H_{\Delta_5,W},\ \mathbf H^!_{\Delta_5,W},\
B_{\mathrm{ch}}(\mathbf H_{\Delta_5,W})\bigr),
\qquad
B_{\mathrm{ch}}(\mathbf H_{\Delta_5,W})
=T^c(s^{-1}\overline{\mathbf H_{\Delta_5,W}}),
\]
where $\mathbf H_{\Delta_5,W}$ denotes the Hall radical quotient supplied
by Remark~\ref{rem:ckp-hall-drinfeld-pair}.  The Koszul sign rule and
cyclicity use the source Hopf pairing and the Mukai pairing; the
Verdier-Koszul dual is defined only after the completed pairing and the
source differential have been installed.  This is the finite-window form
of the Vol.~III K3 chiral-bialgebra construction of
Chapter~\ref{V3-ch:k3-chiral-bialgebra-platonic}, not an assertion that
the automorphic denominator alone produces a chiral Koszul pair.
\end{remark}

\begin{remark}[Conditional self-Koszulness at $\hbar^2 = -1/8$]
\label{rem:ckp-deep-2}
At $\hbar^2=-1/\ell=-1/8$ one may assert a quasi-isomorphism
$\mathbf H^!_{\Delta_5,W}\simeq\mathbf H_{\Delta_5,W}$ only on windows
where the recognition datum supplies the double, the pairing, the source
differential, and the exact completion.  The Bruinier class
$c_1(L^{\Delta_5})\mid_{H_1}\in\mathrm{CH}^1(H_1)$ gives the target
$\mathbb Z/8$ monodromy constraint~\cite{Bruinier2002}.  It does not, by
itself, identify the Hall bar complex with its Verdier-Koszul dual.
\end{remark}

\begin{remark}[Koszul complex with source differential]
\label{rem:ckp-deep-3}
For a recognized window $W$, the K3 Koszul complex on
$(\mathrm{Ran}(E^{\mathrm{nod,sm}}_{24}),\bar{\mathcal A}_2)$ is
\[
\mathrm{Kos}_\bullet(\mathbf H_{\Delta_5,W})\ :\
\dotsb\to
\mathrm{Sym}^k(\mathbf H_{\Delta_5,W}[-1])
\otimes\Lambda^k(\mathbf H_{\Delta_5,W}^\vee)\to\dotsb ,
\]
with differential induced by $d_{\mathrm{Hall}}$, the Hall product, the
coproduct, and the completed pairing.  The Gritsenko--Nikulin denominator
factor is a target Euler-character and coefficient check; it is not the
source differential.  If the finite-window Koszul comparison passes and
the windows form an exact cofinal system, Theorem~H applies to the limit
and computes $\mathrm{ChirHoch}^\bullet(\mathbf H_{\Delta_5})$ in degrees
$\{0,1,2\}$.  The line $\mathbb C\cdot\Delta_5$ is then the recognized
degree-$2$ target class of the Etingof--Kazhdan/Gritsenko--Nikulin
comparison~\cite{EtingofKazhdan2007PartV,GritsenkoNikulin1998}, not an
unconditional cohomology class produced by the denominator product.
\end{remark}
